%% berard-helffer-2022-HoMiNa-0214.tex (14 Feb. 2022)
%%
%% IMG: 211204-IMG-BH-HoMiNa.zip
%% BIB: BH-homina-MR-220102.bib
%% input file: BH-HoMiNa-bib-220102-short.tex (cited items)
%% or BH-HoMiNa-bib-220102-long.tex (all items)
%% Hidden text after first \end{document}
%% input files
\documentclass[12pt,a4paper]{amsart}
%\documentclass[12pt,a4paper,draft]{amsart}
%%\usepackage{showkeys}
\setcounter{secnumdepth}{3}
\setcounter{tocdepth}{3}
%%& [amsart] includes the packages [amsfonts, amsmath, amsthm]
%%& not included [amssymb]
%\documentclass[psamsfonts,11pt,a4paper]{amsart}
%\usepackage{amsmath, amssymb, theorem, latexsym}
\usepackage{amssymb,latexsym,graphicx}
%\usepackage{amsfonts, amsmath, amsthm}
\usepackage{color}
\usepackage{colortbl}
\usepackage{subcaption, float}
%\usepackage[utf8]{inputenc}
\usepackage[cp1252]{inputenc} %% For Windows
\usepackage[T1]{fontenc}
\usepackage{lmodern}
\usepackage[english]{babel}
%\usepackage[english,german]{babel} %% German is active
%%& \begin{otherlanguage}{english} / \end{otherlanguage}
\usepackage{enumerate}% needed here
\usepackage[section]{placeins}
%%& used to place floats within sections with \FloatBarrier
%\usepackage{caption}
%\usepackage{epsf}
%\usepackage[normalem]{ulem}

%\usepackage[pagebackref=true]{hyperref}
%\hypersetup{pdftex,colorlinks=true,linkcolor=blue,citecolor=blue,urlcolor=blue}
%% Inserted at the end of the preamble. Voir documentation [hyperref].

\usepackage[margin=3cm]{geometry}% needed here
%\usepackage{fullpage}
\usepackage{currfile}
%\usepackage{version}
%% list in/ex-clude just before \begin{document}

%%% Dimensions
\settowidth{\parindent}{}

%%% THEOREM alii
\theoremstyle{plain} %% Comment this line if one style only
\newtheorem{theorem}{Theorem}[section]
\newtheorem{lemma}[theorem]{Lemma}
\newtheorem{proposition}[theorem]{Proposition}
\newtheorem{conjecture}[theorem]{Conjecture}
\newtheorem{conjectures}[theorem]{Conjectures}
\newtheorem{corollary}[theorem]{Corollary}
\newtheorem{property}[theorem]{Property}
\newtheorem{properties}[theorem]{Properties}
\newtheorem{assumption}[theorem]{Assumption}
\newtheorem{guess}[theorem]{Guess}
\newtheorem{claim}[theorem]{Claim}
\newtheorem{state}[theorem]{Statement}
%%
\theoremstyle{definition}
\newtheorem{definition}[theorem]{Definition}
\newtheorem{definitions}[theorem]{Definitions}
\newtheorem{example}[theorem]{Example}
\newtheorem{remark}[theorem]{Remark}
\newtheorem{remarks}[theorem]{Remarks}
\newtheorem{notation}[theorem]{Notation}
\newtheorem{quest}[theorem]{Question}
\newtheorem{modif}[theorem]{Modifications}
\newtheorem*{remark*}{Remark}

%%% Numbering
\numberwithin{equation}{section}
\numberwithin{figure}{section}
\numberwithin{table}{section}

%%% LETTERS
\newcommand{\A}{\mathbb{A}}
\newcommand{\B}{\mathbb{B}}
\newcommand{\C}{\mathbb{C}}
\newcommand{\D}{\mathbb{D}}
\newcommand{\E}{\mathbb{E}}
\newcommand{\K}{\mathbb{K}}
\newcommand{\M}{\mathbb{M}}
\newcommand{\N}{\mathbb{N}}
\newcommand{\Nb}{\mathbb{N^{\bullet}}}
\newcommand{\bP}{\mathbb{P}}
\newcommand{\R}{\mathbb{R}}
\newcommand{\bS}{\mathbb{S}}
\newcommand{\T}{\mathbb{T}}
\newcommand{\Z}{\mathbb{Z}}

\newcommand{\cA}{\mathcal{A}}
\newcommand{\cB}{\mathcal{B}}
\newcommand{\cC}{\mathcal{C}}
\newcommand{\cD}{\mathcal{D}}
\newcommand{\cE}{\mathcal{E}}
\newcommand{\cH}{\mathcal{H}}
\newcommand{\cI}{\mathcal{I}}
\newcommand{\cJ}{\mathcal{J}}
\newcommand{\cK}{\mathcal{K}}
\newcommand{\cL}{\mathcal{L}}
\newcommand{\cN}{\mathcal{N}}
\newcommand{\cO}{\mathcal{O}}
\newcommand{\cP}{\mathcal{P}}
\newcommand{\cQ}{\mathcal{Q}}
\newcommand{\cR}{\mathcal{R}}
\newcommand{\cS}{\mathcal{S}}
\newcommand{\cT}{\mathcal{T}}
\newcommand{\cU}{\mathcal{U}}
\newcommand{\cX}{\mathcal{X}}
\newcommand{\cZ}{\mathcal{Z}}

\newcommand{\mf}{\mathfrak}

%%% OPERATORS
\DeclareMathOperator{\spc}{sp}
%% \sp already defined as alias to ^ for superscripts
\DeclareMathOperator{\mult}{mult}
\DeclareMathOperator{\ima}{img}
\DeclareMathOperator{\inte}{int}
\DeclareMathOperator{\sign}{sgn}
\DeclareMathOperator{\id}{Id}
\DeclareMathOperator{\arccot}{arccot}
\DeclareMathOperator{\ord}{ord}
\DeclareMathOperator{\spanop}{span}

%%% MATH EDITING
\newcommand{\pf}{\emph{Proof.~}}
\newcommand{\qedc}{$\checkmark$}

\newcommand{\wt}[1]{\widetilde{#1}}
\newcommand{\wh}[1]{\widehat{#1}}
\newcommand{\wb}[1]{\overline{#1}}
\newcommand{\sm}{\!\setminus\!}
\newcommand{\set}[1]{\left\lbrace #1 \right\rbrace}
\newcommand{\ps}[1]{\langle #1 \rangle}
\newcommand{\spanv}[1]{\spanop\!\left\lbrace #1 \right\rbrace}

\newcommand{\mss}{\scriptstyle} %% reduces size math equ.
\newcommand{\msss}{\scriptscriptstyle} %% reduces size math equ.

\newcommand{\vsp}{\phantom{\big|}} %% used to insert vspace in table

\newcommand{\pib}{\frac{\pi}{2}}
\newcommand{\pic}{\frac{\pi}{3}}
\newcommand{\pit}{\frac{2\pi}{3}}
\newcommand{\pid}{\frac{\pi}{4}}
\newcommand{\pif}{\frac{\pi}{6}}

%%% TEXT EDITING
\newcommand{\noi}{\noindent}
\newcommand{\noib}{\noindent $\bullet$~}
\newcommand{\noid}{\noindent $\diamond$~}
\newcommand{\noic}{\noindent $\circ$~}
\newcommand{\noih}{\noindent $\checkmark$~}
\newcommand{\noit}{\noindent $\blacktriangleright$~}

\newcommand{\resp}{\emph{resp. }}
\newcommand{\ie}{\emph{i.e. }}
\newcommand{\eg}{\emph{e.g. }}

%%% COLORS
\definecolor{purple}{RGB}{127,0,255}
\definecolor{lgray}{gray}{0.7}

%% Colors for math purposes (not for editing).
%% Do not change these tags
\newcommand{\cfr}{\color{red}}
\newcommand{\cfb}{\color{blue}}
%%

%% Colors for editing purposes
\newcommand{\clb}{\color{blue}}
%% \newcommand{\clb}{\color{black}}
 \newcommand{\clr}{\color{red}}
%% \newcommand{\clr}{\color{black}}
 \newcommand{\clg}{\color{green}}
%% \newcommand{\clg}{\color{black}}
 \newcommand{\clp}{\color{purple}}
%% \newcommand{\clp}{\color{black}}
 \newcommand{\clm}{\color{magenta}}
%%  \newcommand{\clm}{\color{black}}
\newcommand{\clc}{\color{cyan}}
\newcommand{\clga}{\color{lgray}}

%% Environment type coloring
\newenvironment{red}{\color{red}}{\ignorespacesafterend}
\newenvironment{blue}{\color{blue}}{\ignorespacesafterend}
\newenvironment{cyan}{\color{cyan}}{\ignorespacesafterend}
\newenvironment{magenta}{\color{magenta}}{\ignorespacesafterend}
\newenvironment{purple}{\color{purple}}{\ignorespacesafterend}
%\newenvironment{violet}{\color{violet}}{\ignorespacesafterend}

%% To color a cell in a table
\newcommand{\cllg}{\cellcolor{lgray}}

%%% TAGS
\newcommand{\cwhrule}[3]{\noindent\textcolor{#1}{\rule{#2}{#3}}}
\newcommand{\stoprule}{\cwhrule{red}{\textwidth}{1.5pt}}
\newcommand{\highL}{\vspace{1cm}\cwhrule{magenta}{0.5\textwidth}{1pt}}

\newcommand{\edit}{\cwhrule{magenta}{0.5\textwidth}{1pt}}
\newcommand{\edits}{\vspace{2cm}\cwhrule{magenta}{0.5\textwidth}{1pt}}
\newcommand{\editss}{\vspace{1cm}\cwhrule{magenta}{0.3\textwidth}{1pt}}

%% To mark hidden text
%\newcommand{\cache}[1]{{\clga \medskip P: #1 ~(version =
% hidden).\medskip}}

\newcommand{\cache}[1]{\vspace{0.5cm}\begin{center} \cwhrule{magenta}{0.3\textwidth}{1pt}
~~#1~~\cwhrule{magenta}{0.3\textwidth}{1pt}\end{center}}

\newcommand{\bcache}{\vspace{0.5cm}\begin{center} \cwhrule{magenta}{0.3\textwidth}{1pt}
~~b/hidden~~\cwhrule{magenta}{0.3\textwidth}{1pt}\end{center}}

\newcommand{\ecache}{\begin{center} \cwhrule{magenta}{0.3\textwidth}{1pt}
~~e/hidden~~\cwhrule{magenta}{0.3\textwidth}{1pt}\end{center}\vspace{0.5cm}}


%%% LOCAL
\newcommand{\lambdah}{\hat{\lambda}}
\newcommand{\gammat}{\widetilde{\gamma}}
\newcommand{\Gammat}{\widetilde{\Gamma}}
\newcommand{\omegat}{\widetilde{\omega}}
\newcommand{\poplus}{\overset{\bot}{\oplus}}
%% \oplus with perp symbol (\bot) above it
\newcommand{\bigpoplus}{\overset{\bot}{\bigoplus}}

\newcommand{\rP}{\mathbb{P}^2}% Real projective plane
\newcommand{\bD}{\partial \cD}
\newcommand{\vo}{\ord} %% \vo to be changed to \ord
\newcommand{\floor}[1]{\left\lfloor #1 \right\rfloor}
\newcommand{\equivo}[1]{\, \overset{#1}{\equiv}\,}

%% HYPERREF
\usepackage[pdftex,pagebackref=true,colorlinks=true,linkcolor=blue,citecolor=blue,urlcolor=blue]{hyperref}


%%%

\title[Eigenvalue multiplicites]{Upper bounds on eigenvalue multiplicities
 for surfaces of genus $0$ revisited\\[5pt]Work in progress}


\author[P. B\'{e}rard]{Pierre B\'erard}
\author[B. Helffer]{Bernard Helffer}

\address[pierrehberard@gmail.com]{Pierre B\'{e}rard, Institut Fourier, Universit\'{e} Grenoble\goodbreak Alpes and CNRS, CS 40700, 38058 Grenoble Cedex 9, France}

\address[Bernard.Helffer@univ-nantes.fr]{Bernard Helffer, Laboratoire Jean Leray, Universit\'{e} de Nantes and CNRS, and LMO (Universit\'e Paris-Sud), 44322 Nantes Cedex, France}


\thanks{We consulted several colleagues during the preparation of this paper: A.~Berdnikov, L.~Friedlander, T.~Hoffmann-Ostenhof, P.~Jammes, M.~Karpukhin, J.~Kennedy, Z.~Liqun, N.~Nadirashvili, I.~Polterovich, S.~Zelditch. We thank them all for their comments.}



\keywords{Spectral theory,  M\"{o}bius strip, Laplacian, Partitions.}

%\subjclass[2010]{58C40, 49Q10.}

\date{\today~(\currfilename)}
%\date{\today}

%% version*
%\includeversion{}
%\excludeversion{}

%%
\begin{document}

\begin{abstract}
We revisit two 1999 papers: [1]~ \emph{M.~Hoffmann-Ostenhof, T.~Hoffmann-Ostenhof, and N.~Nadirashvili. On the multiplicity of eigenvalues of the Laplacian on surfaces. Ann. Global Anal. Geom. 17 (1999) 43--48} ~~and~~ [2]~ \emph{T.~Hoffmann-Ostenhof, P.~Michor, and N.~Nadirashvili. Bounds on the multiplicity of eigenvalues for fixed membranes. Geom. Funct. Anal. 9 (1999) 1169--1188}.\\
The main result of these papers is that the multiplicity $\mult(\lambda_k(M))$ of the $k$th eigenvalue of the surface $M$ is bounded from above by $(2k-3)$ provided that $k \ge 3$. Here, $M$ is [1] a closed surface with genus $0$, or [2] a planar domain with Dirichlet boundary condition (in both cases, the starting label of eigenvalues is $1$). The proofs given in [1,2] are not very detailed, and often rely on figures or special configurations of nodal sets. In the case of closed surfaces with genus $0$, we provide full details of the proof that $\mult(\lambda_k) \le (2k-3)$ by  introducing and carefully studying the \emph{combinatorial type} (defined in Subsection~\ref{SS-h2n-s2}) of the nodal sets involved.  In the case of planar domains, we consider the three boundary conditions, Dirichlet, Neumann, Robin, and we also carefully study the \emph{combinatorial types} of the nodal sets involved. We provide full details of the proof of the inequality $\mult(\lambda_k) \le (2k-2)$ ($k \ge 3$). Unfortunately, we are so far unable to complete the proof of the inequality $\mult(\lambda_k) \le (2k-3)$, even when the domain is simply connected. In the last section, we explain the difficulties we met. Our work on the subject is still ongoing.
\end{abstract}

\maketitle
\tableofcontents

\vspace{2cm}


\section{Introduction}\label{S-intro}
 In these notes, we are concerned with upper bounds for the multiplicities of the eigenvalues $\set{\lambda_k, k\ge 1}$ of an operator $-\Delta + V$ on a compact, smooth (ie. $C^{\infty}$), connected Riemannian surface (when the boundary $\partial M$ is not empty, we consider the Dirichlet, Neumann or Robin boundary conditions).  We \emph{do not} consider the Steklov eigenvalue problem for which we refer to the papers Karpukhin, Kokarev and Polterovich \cite{KaKP2014},
Fraser and Schoen \cite{FrSc2016}, Jammes \cite{Jam2016}, and their
reference lists. \medskip

Our main purpose is to revisit the papers \cite{HoHN1999} (closed surfaces with genus zero) and \cite{HoMN1999} (planar domains with smooth boundary) whose proofs are not very detailed and often rely on figures and special configurations of nodal sets. We introduce and carefully study the \emph{combinatorial type}  (defined in Subsection~\ref{SS-h2n-s2}) of the nodal sets involved. We provide a unified treatment for the three boundary conditions  (Dirichlet, Neumann, Robin). We also illustrate our proofs with many figures.\medskip

In the sequel $\Delta$ is the Laplace-Beltrami operator for some smooth metric $g$, and $V$ is a smooth real valued function. We list the eigenvalues in nondecreasing order, multiplicities accounted for. Our convention is that, in all cases, we label the eigenvalues \emph{starting from the label} $1$,
%
\begin{equation*}
\lambda_1 < \lambda_2 \le \lambda_3 \le \ldots \,,
\end{equation*}
%
and we denote the multiplicity of $\lambda_k$ by $\mult(\lambda_k)$. We refer to Section~\ref{S-evp} for more precise definitions.\medskip

In Section~\ref{S-sketch} we provide a panorama of the main results on the multiplicity problem, and the ideas behind their proofs. Section~\ref{S-evp} is devoted to definitions, notations and preliminary results. Section~\ref{S-pargef} deals with Euler type formulas in the framework of partitions. In Section~\ref{S-h2n}, we revisit \cite{HoHN1999}. Introducing the \emph{combinatorial type} of some particular nodal sets, see Definition~\ref{D-h2n-type}, we provide a complete proof of the inequality $\mult(\lambda_k) \le (2k-3)$ in the case of closed surfaces with genus $0$, see Subsection~\ref{SS-h2n-s3}. This section constitutes a preamble to the remaining sections. Sections~\ref{S-hmn} to \ref{S-hmn6} are devoted to the upper bounds on multiplicities for smooth bounded domains in $\R^2$. They form the core of this paper.\smallskip

Section~\ref{S-hmn} contains some preliminary results, in particular a proof of Nadirashvili's estimate $\mult(\lambda_k) \le (2k-1)$ (in our restricted framework).  In Section~\ref{S-hmn2}, we give a detailed proof of the upper bound $\mult(\lambda_k) \le (2k-2)$ for a smooth bounded domain $\Omega \subset \R^2$ (not assumed to be simply connected), see Proposition~\ref{P-hmn-s2}. The strategy of the proof is as follows: assume that $\mult(\lambda_k) = (2k-1)$, construct $\lambda_k$-eigenfunctions with prescribed singular points, analyze their nodal sets, and derive a contradiction.\smallskip

In \cite[Theorem~B, p.~1172]{HoMN1999}, the authors state that the previous bound can be improved to $\mult(\lambda_k) \le (2k-3)$ for all $k\ge 3$. The strategy of the proof is similar, except that we now start from the assumption $\mult(\lambda_k) = (2k-2)$. Unfortunately, we have so far not succeeded in writing down complete details for the arguments given in \cite{HoMN1999}, even when $\Omega$ is simply connected (for comments, see Section~\ref{S-sketch} or \cite[Section~4]{Berd2018}). In Sections~\ref{S-hmn3} and \ref{S-hmn4}, we provide some properties of $\lambda_k$-eigenfunctions under the assumption that $\mult(\lambda_k) = (2k-2)$. In Section~\ref{S-hmn6}, we indicate the main steps of the proof sketched in \cite[Section~3]{HoMN1999}, and we point out where we are stuck.\smallskip

In Appendix~\ref{S-mcs} we relate the problem of bounding multiplicities from above to the question of Courant-sharp eigenvalues. The other appendices are devoted to technical results.


\section{Historical sketch}\label{S-sketch}



In the case of closed surfaces, the first upper bounds on multiplicities were obtained by  Cheng \cite{Chen1976}, Besson \cite{Bess1980}, and Nadirashvili \cite{Nadi1987}.
We denote their respective upper bounds on $\mult(\lambda_k)$ by $m_k^{*}$, with $* \in \set{b,c,n}$, where $b$ stands for ``Besson'', $c$ for ``Cheng'', and $n$ for ``Nadirashvili'', and provide a summary of their results in
Table~\ref{E-scr-2T}. \smallskip

The upper bounds for the multiplicity of the second eigenvalue (ie, the least positive eigenvalue of a closed surface) given in the fifth column are sharp. For the sphere the bound is achieved for the canonical (round) metric, \cite{Chen1976}; for the projective space the bound is achieved for the metric induced by the canonical metric of the sphere, \cite{Bess1980}; for the torus the bound is achieved for the equilateral torus $\T_e$ with metric induced from $\R^2$, \cite{Bess1980}; for the Klein bottle, the bound is achieved for a nontrivial pair $(g,V)$ constructed in \cite[\S\,2]{Nadi1987}, and for smooth metrics constructed in \cite[Th\'{e}or\`{e}me~4.2]{ColV1987}. An interesting feature of $\bS^2, \mathbb{RP}^2$ and $\T^2$ is that the bounds for $\mult(\lambda_2)$ are also achieved for metrics different from the ones mentioned above, see \cite{Bess1980}.\smallskip

In \cite[Th\'{e}or\`{e}me~1.5]{ColV1987}, Colin de Verdi\`{e}re shows that for a closed surface $M$,   $$\sup\set{\mult(\lambda_2(M,-\Delta_g + V)) ~|~ (g,V)} \ge C(M) - 1\,,$$ where the supremum is taken over the Riemannian metrics and potentials on $M$, and where $C(M)$ is the \emph{chromatic number} of $M$ (the maximal $N$ such that the complete graph on $N$ vertices $K_N$ can be embedded into $M$). Table~\ref{E-scr-2T} shows that equality holds for $\bS^2, \mathbb{RP}^2, \T^2$ and $\K^2$; it also holds for surfaces with $\chi(M) \ge -3$, \cite{Seve2002}. It is conjectured that equality holds for all closed surfaces.\smallskip


\begin{table}[!htb]
\centering
\resizebox{0.9\linewidth}{!}{
\begin{tabular}[c]{|c|c|c|c|c|l|}%
\hline
&&&&&\\[-1em]
$M$ & $\chi(M)$ & $\gamma(M)$ & $\check{\gamma}(M)$& $\mult(\lambda_2) \le$
& $\mult(\lambda_k) \le$ \\[0.2em]
\hline
&&&&&\\[-1em]
$\bS^2$ &  $2$ & $0$ & -- & $3$ &
$
\left\{
\begin{array}{ll}
m_k^{c} & = \frac 12 k(k+1)\\[5pt]
m_k^{b} & = 2k-1\\[5pt]
m_k^{n} & = 2k-1
\end{array}
\right.
$
\\[0.2em]
\hline
&&&&&\\[-1em]
$\mathbb{RP}^2$ &  $1$ & -- & $0$ & $5$ &
$
\left\{
\begin{array}{ll}
m_k^{c} & = \text{not considered} \\[5pt]
m_k^{b} & = 4k-1\\[5pt]
m_k^{n} & = 2k+1
\end{array}
\right.
$ \\[0.2em]
\hline
&&&&&\\[-1em]
$\T^2$ &  $0$ & $1$ & -- & $6$ &
$
\left\{
\begin{array}{ll}
m_k^{c} & = \frac 12  (k+2)(k+3)\\[5pt]
m_k^{b} & = 2k+3\\[5pt]
m_k^{n} & = 2k+2
\end{array}
\right.
$ \\[0.2em]
\hline
&&&&&\\[-1em]
$\K^2$ &  $0$ & -- & $1$ & $5$ &
$
\left\{
\begin{array}{ll}
m_k^{c} & =  \text{not considered} \\[5pt]
m_k^{b} & =  \text{not considered} \\[5pt]
m_k^{n} & = 2k+1
\end{array}
\right.
$ \\[0.2em]
\hline
&&&&&\\[-1em]
$M^2$ &  $\chi(M) < 0$ & $\gamma$ & -- & -- &
$
\left\{
\begin{array}{ll}
m_k^{c} & = \frac 12 (k+2\gamma)(k+2\gamma+1)\\[5pt]
m_k^{b} & = 2k + 4 \gamma - 1\\[5pt]
m_k^{n} & = 2k+ 4\gamma-3\\
\end{array}
\right.
$ \\[0.2em]
\hline
&&&&&\\[-1em]
$M^2$ &  $\chi(M) < 0$ & -- & $\check{\gamma}$ & -- &
$
\left\{
\begin{array}{ll}
m_k^{c} & = \text{not considered} \\[5pt]
m_k^{b} & = 4k + 4 \check{\gamma} - 1\\[5pt]
m_k^{n} & = 2k+ 2 \check{\gamma}-1\\
\end{array}
\right.
$ \\[0.2em]
\hline
\end{tabular}
}
\medskip
\caption{Closed surfaces: multiplicity upper bounds obtained by Cheng, Besson, and Nadirashvili}\label{E-scr-2T}
\end{table}%

Concerning Table~\ref{E-scr-2T}, recall that for an orientable surface, $\chi(M) = 2$ (the sphere) or $\chi(M) = 2 - 2 \gamma(M)$, where the genus $\gamma(M)$ is the number of handles; in this case, $M$ is the connected sum of $\gamma(M)$ tori. For a non-orientable surface, $\chi(M) = 1$ (the real projective space) or  $\chi(M) = 2 - \check{\gamma}(M)$ if the surface is a connected sum of $\check{\gamma}(M)$ projective spaces.\smallskip

Better upper bounds were later obtained by M. and T.~Hoffmann-Ostenhof and Nadirashvili \cite{HoHN1999} (surfaces with genus $0$), S\'{e}vennec \cite{Seve2002} (improved bounds on the multiplicity of $\lambda_2$ when $\chi(M) < 0$), Berdnikov, Nadirashvili and Penskoi \cite{BeNP2016} (improved bounds for the multiplicities on the projective plane),  Fortier Bourque and Petri \cite{FoBP2021} (Klein quartic).\smallskip

In \cite[Theorem~2]{Nadi1987}, Nadirashvili considers smooth bounded
domains $\Omega \subset \R^2$ and proves that the multiplicity of the $k$th eigenvalue $\lambda_k$ of an operator $-\Delta + V$ with Dirichlet or Neumann boundary condition is at most $(2k-1)$. \smallskip


In the paper \cite{HoMN1999}, Hoffmann-Ostenhof, Michor and Nadirashvili, improve Nadirashvili's bound for bounded planar domains  with $C^{\infty}$ boundary with Dirichlet boundary condition. More precisely, they state that the multiplicity of $\lambda_k$ is at most $(2k-3)$. Berdnikov \cite{Berd2018} considers the case of compact surfaces with boundary, under the assumption that $\chi(M) + b_0(\partial M)$ is negative.  He points out some problem in the proof in \cite{HoMN1999} when the domain is not simply connected.\medskip


The general idea to prove upper bounds for the eigenvalue multiplicities is a combination of the following ingredients:
\begin{enumerate}[(i)]
  \item Courant's nodal domain theorem.
  \item  Local structure theorems for eigenfunctions near a singular point.
  \item Existence of eigenfunctions with prescribed  singular points, provided the dimension of the eigenspace is large enough.
  \item Euler's formula for the graph associated with the nodal set of an eigenfunction.
  \item An argument which we call the \emph{rotating function argument} (see \S~\ref{SSS-h2n-s2r}).
  \item Energy arguments and eigenvalue monotonicity.
\end{enumerate}

In one form or another, these arguments go back to Cheng \cite{Chen1976}, Besson \cite{Bess1980}, and Nadirashvili \cite{Nadi1987}.\medskip


We refer to the papers of Burger, Colbois and/or Colin de Verdi\`{e}re \cite{BuCo1985, Colb1985, ColV1986, ColV1987, CoCo1988} for results of a different flavour. \smallskip


Two other papers,  respectively \cite{HeHO1999} by Helffer, M. and T. Hoffmann-Osten\-hof and Owen, and \cite{HeHN2002} by Helffer, M. and T. Hoffmann-Ostenhof and Nadirashvili, have used the same techniques
 for related purposes (for example the Aharo\-nov-Bohm  operators). \smallskip

Finally similar techniques are used in the analysis of the properties of minimal partitions \cite{HeHT2009,BoHe2017}.\medskip


\section{Eigenvalue problems}\label{S-evp}

\subsection{Definitions, notation and preliminary results}\label{SS-evp1}

In this section, $M$ denotes a closed surface (compact, no boundary), or a compact surface with boundary. The boundary is denoted by $\partial M$, and the interior $M\sm \partial M$ is denoted by $\inte(M)$. Unless otherwise stated, the surface is assumed to be smooth and connected. We equip $M$ with a smooth Riemannian metric $g$, and we consider a (non-magnetic) Schr\"{o}dinger operator of the form $-\Delta_g + V$, where $\Delta_g$ is the Laplace-Beltrami operator for the metric $g$ and $V$ is a smooth real valued function on $M$.\medskip

The notation are as follows. The Riemannian measure is denoted by $v_g$; $h$ is a nonnegative given constant. When $\partial M \not = \emptyset$, $\sigma_g$ is the Riemannian measure of $\partial M$ for the metric induced by $g$, and $\nu$ is the unit normal to $\partial M$ pointing outward.\smallskip

We consider the  closed\footnote{By closed eigenvalue problem, we mean an eigenvalue problem on a closed manifold (no boundary condition).} eigenvalue problem
%
\begin{equation}\label{E-evp-2c}
-\Delta u + V u = \lambda \, u \text{~~~in~} M \text{~~(here, ~}\partial M = \emptyset)\,,
\end{equation}
%
associated with the quadratic form
%
\begin{equation}\label{E-evp-2cq}
\int_M \left( |du|^2_g + V u^2 \right) \, dv_g\,, \text{~~with domain~~} H^1(M)\,.
\end{equation}

When $\partial M \not = \emptyset$, we consider the boundary eigenvalue problem
%
\begin{equation}\label{E-evp-2bc}
\left\{
\begin{array}{rll}
-\Delta u + V u & = \lambda \, u &\text{~in~} \inte(M)\,,\\[5pt]
B(u) &= 0 &\text{~on~} \partial M\,,
\end{array}
\right.
\end{equation}
%
where $B(u)$ is one the following boundary conditions:
%
\begin{equation}\label{E-evp-bc}
B(u) = \left\{
\begin{array}{ll}
u &\text{(Dirichlet),}\\[5pt]
\frac{\partial u}{\partial \nu} &\text{(Neumann),}\\[5pt]
\frac{\partial u}{\partial \nu}+h\,u &\text{($h$-Robin).}
\end{array}
\right.
\end{equation}

The associated quadratic forms are
%
\begin{equation}\label{E-evp-2dq}
\int_M \left( |du|^2_g + V u^2 \right) \, dv_g\,, \text{~~with domain~~} H_0^1(M)\,.
\end{equation}
%
for the Dirichlet problem, and
%
\begin{equation}\label{E-evp-2rq}
\int_M \left( |du|^2_g + V u^2 \right) \, dv_g + h \, \int_{\partial M} (u_{\partial M})^2\, d\sigma_g\,,\text{~~with domain~~} H^1(M)\,,
\end{equation}
%
for the Neumann problem (in this case $h=0$) and for the $h$-Robin problem.\medskip


For these eigenvalue problems, the spectrum is discrete, and consists of a sequence of non-negative eigenvalues with finite multiplicities,
%
\begin{equation}\label{E-evp-2ev}
\lambda_1 < \lambda_2 \le \lambda_3 \le \cdots \nearrow \infty \,,
\end{equation}
%
which we list in nondecreasing order, multiplicities accounted for, starting from the label $1$.

\begin{remark}\label{R-evp-2}
Whenever necessary, we shall indicate the dependence on $M, g, V, h$ and the boundary condition, for example, $\lambda_k(M,g,V,\mf{d})$ for the Dirichlet eigenvalues of $-\Delta + V$ on $(M,g)$ with the Dirichlet boundary condition on $\partial M$.
\end{remark}%

In the sequel, we only consider \emph{real valued} eigenfunctions. The \emph{eigenspace} associated with the eigenvalue $\lambda_k$ will be denoted by $U(\lambda_k)$. Its dimension, the \emph{multiplicity} of $\lambda_k$, will be denoted by $\mult(\lambda_k)$ or $\dim U(\lambda_k)$.\medskip


\begin{definitions}[Terminology]\label{D-evp2-00}\phantom{}
\begin{enumerate}[(i)]
\item The \emph{nodal set} of a nontrivial eigenfunction $u$ is denoted by $\cZ(u)$, and defined by
%
\begin{equation}\label{E-evp-ns}
\cZ(u) = \overline{\set{x \in \inte(M) ~|~ u(x) = 0}}\,.
\end{equation}
%
When $\partial M \not = \emptyset$, $\cZ(u)$ is the closure in $M$ of the set of interior zeros of $u$.
\item In dimension $2$, the nodal set $\cZ(u)$ of an eigenfunction $u$ is also called the \emph{nodal line} of $u$.
\item The \emph{nodal domains} of the eigenfunction $u$ are the connected components\footnote{In the sequel, unless otherwise stated, we shall use the word \emph{component} for the expression \emph{connected component}.} of $ \inte(M) \sm \cZ(u)$. We denote by $\kappa(u)$ the number of nodal domains of $u$.
\end{enumerate}
\end{definitions}%



According to Courant's nodal domain theorem, the number of nodal domains of an eigenfunction associated with $\lambda_k$ is at most $k$. Eigenfunctions associated with $\lambda_1$ are characterized by the fact that they have precisely one nodal domain. An eigenfunction associated with $\lambda_k, k\ge 2$, has at least two nodal domains. An eigenfunction associated with $\lambda_2$ has precisely two nodal domains. It turns out that, for $k$ large enough (depending on $M, g, V)$, $\sup\set{\kappa(u) ~|~ 0 \not = u \in U(\lambda_k)} < k$ (see \cite{Plej1956}, \cite{Polt2009}, \cite{Lena2019}), and that the bound $\kappa(u) \ge 2$, for $k \ge 2$, can generally speaking not be improved (see \cite{Ster1925,Lewy1977}, \cite{BeHe2015s,BeHe2015r}). \medskip


\subsection{Local structure of eigenfunctions near a zero}\label{SS-evp2}

We say that a function $u$ \emph{vanishes at order} $n \ge 1$ at a point $x$, and we write $\ord(u,x) = n$, if (in a local coordinate system) the function and all its derivatives of order less than or equal to $(n-1)$ vanish at $x$, and at least one derivative of order $n$ does not vanish at $x$. A \emph{critical zero} of $u$ is a point at which $u$ vanishes at order at least $2$ (i.e. $u(x)=0$ and $\nabla_x u = 0$). A critical zero $x$ of $u$ is called an \emph{interior critical zero} if $x \in \inte(M)$, and a \emph{boundary critical zero} if $x \in \partial M$.


\begin{theorem} \label{T-nodinfo}
Let $u$ be a nontrivial eigenfunction of the Schr\"odinger operator $-\Delta +V$ on a smooth compact Riemannian surface $M$ (with or without boundary),  where $V$ is a smooth real valued  potential. Then, $u\in C^\infty(M)$,  and does not vanish at infinite order at any point of $M$. Furthermore, depending on the boundary condition on $\partial M$, $u$ has the following  properties.
\begin{enumerate}[(i)]
\item If $x_0\in M$ is an interior point, and if $u$ has a zero of order $\ell$ at $x_0$, then there exist local polar coordinates $(r,\omega)$ centered at $x_0$ such that the Taylor expansion of $u$ is
\begin{equation}\label{E-harmpoly}
u(x)=r^{\ell} \big( a \sin(\ell \omega) + b \cos(\ell \omega) \big) +\cO(r^{\ell+1}),
\end{equation}
%
where $a,b \in \R$, $a^2+b^2 \not = 0$.

\item If $x_0\in \partial M$, and if a Dirichlet eigenfunction $u$ has a zero of order $\ell$ at $x_0$, then there exist local polar coordinates $(r,\omega)$ centered at $x_0$, such that the Taylor expansion of $u$ is
\begin{equation}\label{E-sinexp}
u(x)=a\, r^\ell \sin(\ell \omega) +\cO(r^{\ell +1})
\end{equation}
for some $a\in\mathbb R$, $a \not = 0$. The angle $\omega$ is chosen so that the tangent to the boundary at $x_0$ is given by the equation $\omega=0$.

\item If $x_0\in \partial M$, and if a Robin eigenfunction $u$ has a zero of order $\ell$ at $x_0$, then there exist local polar coordinates $(r,\omega)$ centered at $x_0$, such that the Taylor expansion of $u$ is
\begin{equation}\label{E-cosexp}
u(x)=b\, r^\ell \cos(\ell \omega) +\cO(r^{\ell +1})
\end{equation}
for some $b\in\mathbb R$, $b \not = 0$. The angle $\omega$ is chosen so that the tangent to the boundary at $x_0$ is given by the equation $\omega=0$.
\end{enumerate}
\end{theorem}

\pf For a proof of this theorem, and for references to the literature, we refer to \cite[Appendix~A]{GiHe2019}  (which deals with the more delicate case in which the boundary of the domain is not $C^{\infty}$). The starting point is to use the unique continuation theorem, see \cite{Aron1957} when $x_0$ is an interior point, and \cite{DoFe1990a} when $x_0$ is a boundary point  (see also \cite{YaZh2021}).\hfill \qed \medskip

From a local point of view, we have the following properties.

\begin{corollary}\label{cor:nodloc}~
\begin{enumerate}[(i)]
\item Let $x_0\in M$. If $u$ has a zero of order $\ell $ at $x_0$, then exactly $\ell$ nodal curves pass through $x_0$. More precisely, in a neighborhood of $x_0 \in \inte(M)$, the nodal set $\cZ(u)$ consists of $2\ell$ semi-arcs emanating from $x_0$ tangentially to the rays $\set{\omega = \omega_j}$ where $\omega_j := j \frac{\pi}{\ell}, 0\le j < 2\ell$. The semi-tangents to these semi-arcs dissect the full unit circle in the tangent plane at $x_0$ into $2\ell $
    equal parts.
\item Let $x_0\in \partial M$. Let $u$ be a Dirichlet eigenfunction. If $u$ has a zero of order $\ell \ge 2$ at $x_0$, then exactly $(\ell-1)$ semi-arcs hit $\partial M$ at $x_0$, their  semi-tangents at $x_0$ dissect the half unit circle in the tangent plane at $x_0$ into $\ell$ equal angles given by the equation $\sin(\ell \omega) =0$.
\item Let $x_0\in \partial M$. Let $u$ be a Robin eigenfunction. If $u$ has a zero of order $\ell \ge 1$ at $x_0$, then exactly $\ell$ semi-arcs hit $\partial M$ at $x_0$, their  semi-tangents  at $x_0$ dissect the half unit circle in the tangent plane at $x_0$ into $\ell$ equal angles given by the equation $\cos(\ell \omega) =0$.
\end{enumerate}
\end{corollary}%

\pf Assertion~(i) is proved in Appendix~\ref{S-lsns}.  For Assertions~(ii) and (iii), see the references in \cite[Appendix]{GiHe2019}. \hfill\qed \medskip

Points at which nodal arcs meet in the interior $\inte(M)$, and points at which the nodal set hits the boundary $\partial M$ play an important role in the global understanding of nodal sets. The terminology in the following definition comes from the framework of partitions, see Definition~\ref{D-parti-reg}.

\begin{definition}\label{D-evp2-0}
Define the \emph{singular points} of an eigenfunction $u$ as follows.
\begin{enumerate}[(i)]
\item A point $x_0 \in \inte(M)$ is an \emph{interior singular point} of $u$ if and only if it is an interior critical zero; the set of interior singular points of $u$ is denoted by $\cS_i(u)$. The \emph{index} $\nu(u,x_0)$ of the interior singular point $x_0$ is defined as the number of nodal semi-arcs emanating from $x_0$, $\nu(u,x_0) = 2\, \ord(u,x_0)$.
  \item A point $x_0 \in \partial M$ is a \emph{boundary singular point} of $u$ if and only if the nodal set $\cZ(u)$ hits the boundary $\partial M$ at $x_0$; the set of boundary singular points of $u$ is denoted by $\cS_b(u)$. The \emph{index} $\rho(u,x_0)$ of the boundary singular point $x_0$ is defined as the number of nodal arcs hitting $\partial M$ at $x_0$.\\
      If $u$ is a Dirichlet eigenfunction, $\rho(u,x_0) = \ord(u,x_0) - 1$; if $u$ is a Robin eigenfunction, $\rho(u,x_0) = \ord(u,x_0)$.
\end{enumerate}
The set $\cS(u)$ of singular points of $u$ is the set $\cS(u) = \cS_i(u) \cup \cS_b(u)$.
\end{definition}%

\begin{remark}\label{R-evp2-2}
The order of vanishing is \emph{semi-continuous} in the following sense. Let $\set{v_n}$ be a sequence of functions which converges to some $v$ uniformly in $C^{k+2}$ and let $\set{x_n}$ be a sequence of points which converges to some $x$ in $M$. Assume that $\ord(v_n,x_n) \ge k$ for all $n$ and some fixed $k$, then $\ord(v,x) \ge k$. Since they are defined in terms of order of vanishing, the indices $\nu$ and $\rho$ inherit this property.
\end{remark}%

 From a more global point of view, the connected components of $\cZ(u) \sm \cS(u)$ are smooth $1$-dimensional submanifolds homeomorphic to either circles or open intervals whose boundaries consist of singular points.

\begin{definitions}[Terminology]\label{D-evp2-01}\phantom{}
\begin{enumerate}[(i)]
\item We call a circle-like component, \emph{a nodal circle}; we call an interval-like component, \emph{a nodal arc}.
\item Let $I_{x,y}$ be a nodal arc with boundary $\set{x,y}$. In this case, we call $I_{x,y}\cup \set{x,y}$  \emph{a nodal semi-arc} emanating from $x$ and arriving at $y$, with some semi-tangents at $x$ and $y$ given by the local structure theorem. The point $x$ (resp. $y$) might be an interior singular point, or a boundary singular point (if the nodal set $\cZ(u)$ hits the boundary $\partial M$ at $x$, resp. $y$). If $x=y$, we say that the component $I_{x,x}$ is \emph{a nodal loop} at $x$. In this case, the semi-tangents at $x$ might not be co-linear, so that the (closed) loop is not a smooth circle, but a continuous, piecewise $C^1$ circle.
\item The collection \{singular points, nodal circles and nodal arcs\} constitutes a \emph{multigraph} called the \emph{nodal graph} of $u$, see Subsection~\ref{SS-gr}.
\end{enumerate}
\end{definitions}

\begin{corollary}\label{cor:nodinfo}\phantom{}
\begin{enumerate}
  \item The nodal set of $u$ is the union of the finitely many singular points, nodal circles in the interior of $M$, and nodal arcs some of which may hit  $\partial M$.
  \item Each component of $\partial M$ is hit by an even number of nodal arcs. If $\Gamma$ is a component of $\partial M$, then
      \begin{equation*}
        \sum_{z \in \cS_b(u) \cap \Gamma\,} \rho(u,z) \in 2\N\,.
      \end{equation*}
\end{enumerate}
\end{corollary}%

\pf The first assertion is well-known. We give the proof of the second assertion for completeness.\smallskip

\emph{\noid Dirichlet case.}~ The component $\Gamma$ is topologically a circle which meets $\cS_b(u)$ at finitely many points $z_j$, $1 \le j \le k$, which are precisely the zeros of the normal derivative $\partial_{\nu}u(z)$. Choosing a parametrization $z$ of $\Gamma$ and taking  the local structure of $u$ at each $z_j$ into account, we see that each time $z$ passes some $z_j$, the sign of $\partial_{\nu}u$ is multiplied by $(-1)^{\rho(z_j)}$. \smallskip

\emph{\noid Robin case.}~ The proof is similar, actually simpler. \hfill\qed



\subsection{Eigenfunctions with prescribed singular points}\label{SS-evp3}

\phantom{}  In order to bound multiplicities, we will use eigenfunctions with   prescribed  singular points of sufficiently high index. Their existence is given by the following lemmas.  These lemmas appear in one form or another in \cite[Theorem~3.4]{Chen1976}, \cite[Theorem~2.1]{Bess1980}, \cite[Lemma~4]{Nadi1987}, \cite[Proposition~2]{HoHN1999}, \cite[Lemma~2.9]{HoMN1999}.\medskip


The first lemma prescribes an interior singular point.

\begin{lemma}\label{L-zeroi} Let $M$ be a compact surface (with or without boundary), and $x$ an interior point. Let $U$ be a linear subspace of an eigenspace of the eigenvalue problem \eqref{E-evp-2c} or \eqref{E-evp-2bc}, with $\dim U = m \ge 2$.
\begin{enumerate}[(i)]
\item There exists a function $0 \not = u\in U$ such that $x$ is a singular point of $u$ with index $\nu(u,x) \ge 2\,\floor{\frac m2}$ (the integer part of $\frac m2$), equivalently $\ord(u,x) \ge \floor{\frac m2}$.
\item Furthermore, if $m$ is odd, there exist at least two linearly independent such functions.
\end{enumerate}
\end{lemma}

\pf We use induction on $m$.  Recall that $\nu(u,x) = 2\ord(u,x)$. The assertion is clear when $m=2$. Assume $m=3$, and let $\set{u_1,u_2,u_3}$ be a basis of $U$. Then, we can find $0 \not = v_1 \in \spanv{u_1,u_2}$ such that $\ord(v_1,x) \ge 1$. The subspace $V_1$ of $U$ orthogonal to $v_1$ has dimension $2$, and hence there exists $0 \not = v_2 \in V_1$ such that $\ord(v_2,x) \ge 1$. Then $v_1$ and $v_2$ are two linearly independent functions in $U$ vanishing at order at least $1$ at $x$.\medskip

Assume that the lemma holds for $2p$ and $2p+1$ for some $p \ge 1$.\medskip

Let $U$ be linear subspace of an eigenspace with dimension $(2p+2)$, and basis $$\set{u_1,\ldots,u_{2p+2}}.$$ By the induction hypothesis, in the subspace $V_1 := \spanv{u_1,\ldots,u_{2p+1}}$, we can find two linearly independent functions $v_1, v_2$ such that $\ord(v_i,x) \ge p$. If one of them vanishes at order at least $(p+1)$, the assertion for $U$ is satisfied. If not, according to Theorem~\ref{T-nodinfo}~(i), there exist $(a_i,b_i)$, $i=1,2$, with $a_i^2+b_i^2 \not = 0$, such that
%
\begin{equation*}
v_i = r^{p} \big( a_i \sin(p\, \omega) + b_i \cos(p\, \omega) \big) +\cO(r^{p+1}).
\end{equation*}

The subspace $V_2$ of $U$ orthogonal to $v_1$ and $v_2$ has dimension $2p$ and hence, there exists $0 \not = v_3 \in V_2$ such that $\ord(v_3,x) \ge p$. If $v_3$ vanishes at order at least $(p+1)$ at $x$, we are done. Otherwise, there exist $(a_3,b_3)$, with $a_3^2+b_3^2 \not = 0$ such that
%
\begin{equation*}
v_3 = r^{p} \big( a_3 \sin(p\, \omega) + b_3 \cos(p\, \omega) \big) +\cO(r^{p+1}).
\end{equation*}

The functions $r^p \big( a \sin(p\, \omega) + b \cos(p\, \omega) \big)$ are the homogeneous harmonic polynomials of degree $p$ in $\R^2$, a vector space of dimension $2$. The three polynomials $r^{p} \big( a_i \sin(p\, \omega) + b_i \cos(p\, \omega)\big)$, $i \in \set{1,2,3}$ must be linearly dependent, and hence there exists a nontrivial linear combination of $v_1,v_2,v_3$ which vanishes at order at least $(p+1)$ at $x$.\smallskip

Let $U$ be an eigenspace with dimension $2p+3$, with basis $\set{u_1,\ldots,u_{2p+3}}$. By the previous proof, in the subspace $V_1 := \spanv{u_1,\ldots,u_{2p+2}}$, there exists $0 \not = v_1$ such that $\ord(v_1,x) \ge (p+1)$. For the same reason, in the subspace $V_2$ orthogonal to $v_1$, there exists $0 \not = v_2$ such that $\ord(v_2,x) \ge (p+1)$. The functions $v_1, v_2$ are two linearly independent functions in $U$ vanishing at order at least $(p+1)$ at $x$. \smallskip

The proof of Lemma~\ref{L-zeroi} is complete.\hfill \qed \medskip

The next lemmas prescribe respectively one or two boundary singular points.

\begin{lemma}\label{L-zero1} Let $M$ be a compact surface with
boundary, and $x \in \partial M$. Let $U$ be a linear subspace of an eigenspace of the eigenvalue problem \eqref{E-evp-2bc}, with $\dim U = m \ge 2$. Then, there
exists a function $0 \not = u\in U$ such that $x$ is a boundary singular point of $u$ with index $\rho(u,x) \ge (m-1)$.
\end{lemma}

\pf We use induction on $m$.  Recall that $\rho(u,x) = (\ord(u,x) - 1)$ for Dirichlet eigenfunctions, resp. $\rho(u,x) = \ord(u,x)$ for Robin eigenfunctions. \smallskip

\emph{\noid Dirichlet boundary condition.}~ When $m=2$, the assertion is clear. Assume it is true for some $m\ge 2$. Let $U$ be a linear subspace of an eigenspace with dimension $(m+1)$, and basis $\set{u_1,\ldots,u_{m+1}}$. Consider the subspace $V_1 = \spanv{u_1,\ldots,u_m}$. By the induction hypothesis, there exists $0 \not = v_1 \in V_1$ such that $\ord(v_1,x) \ge m$. If $v_1$ vanishes at order at least $(m+1)$, we are done. Otherwise, by Theorem~\ref{T-nodinfo}, Equation~\eqref{E-sinexp}, there exists $a_1\not = 0$ such that, in local polar coordinates at $x$,
%
\begin{equation*}
v_1(z)=a_1\,r^m \sin(m \,\omega) +\cO(r^{m +1})\,.
\end{equation*}
%
The subspace $V_2=\set{u\in U ~|~ u \perp v_1}$ orthogonal to $v_1$ has dimension $m$, and hence there exists $0 \not = v_2 \in V_2$ such that $\ord(v_2,x) \ge m$. If $v_2$ vanishes at order at least $(m+1)$, we are done. Otherwise, as above we can write
%
\begin{equation*}
v_2(z)=a_2\,r^m \sin(m\, \omega) +\cO(r^{m +1})
\end{equation*}
%
for some $a_2 \not = 0$, and hence the linear combination $v = a_2 v_1 - a_1 v_2$ vanishes at order at least $(m+1)$. \hspace{1cm}\qedc \smallskip

\emph{\noid Robin boundary condition.}~ When $m=2$, the assertion is clear. Assume it is true for some $m\ge 2$. Let $U$ be a linear subspace of an eigenspace with dimension $(m+1)$, with basis $\set{u_1,\ldots,u_{m+1}}$. Consider the subspace $U_1 = \spanv{u_1,\ldots,u_m}$. By the induction hypothesis, there exists $0 \not = v_1 \in U_1$ such that $\ord(v_1,x) \ge (m-1)$. If $v_1$ vanishes at order at least $m$, we are done. Otherwise, by Theorem~\ref{T-nodinfo},  Equation~\eqref{E-cosexp}, there exists $b_1\not = 0$ such that, in local polar coordinates at $x$,
%
\begin{equation*}
v_1(z)=b_1\,r^m \cos((m-1) \omega) +\cO(r^{m+1}).
\end{equation*}
%
We can then consider the subspace $U_2$ orthogonal to $v_1$ in $U$, and conclude by arguing as above. \hspace{1cm}\qedc \smallskip

The proof of Lemma~\ref{L-zero1} is complete. \hfill \qed


\begin{lemma}\label{L-zero2} Let $M$ be a compact surface with
boundary, and $x, y \in \partial M$, with $x \not = y$. Let $U$ be a linear subspace of an eigenspace of the eigenvalue problem \eqref{E-evp-2bc}, with $\dim U = m \ge 3$. Then, there exists a function $0 \not = u\in U$ such that $x$ and $y$ are boundary singular points of $u$ with indices $\rho(u,x) \ge (m-2)$ and $\rho(u,y) \ge 1$.
\end{lemma}

\begin{proof}\phantom{}

\emph{\noid Dirichlet boundary condition.}~ Choose $\set{u_1,\ldots, u_m}$ a basis of $U$. Looking at a general element $u=\sum \alpha_j u_j$ in $U$, the condition at $y$ reads
%
$$
\sum_{j=1}^m  \alpha_j (\partial_\nu u_j)(y) =0\,.
$$
%
There are two cases.
\begin{itemize}
\item If $\partial_\nu u_j(y) =0$ for all $j$, the condition at $y$ is satisfied for any $u\in U$;
\item If $\partial_\nu u_j(y) \neq 0$ for some $j$, then there exists a subspace $U'\subset U$ of dimension $(m-1) \ge 2$ such that the condition at $y$ is satisfied for any $u\in U'$.
\end{itemize}
We can then apply Lemma~\ref{L-zero1} with $U$ in the first case and with $U'$ in the second case. \hspace{1cm} \qedc \smallskip

\emph{\noid Robin boundary condition.}~ The condition $\rho(u,y) \ge 1$ holds if and only if $u$ vanishes at $y$. Since $m \ge 3$ there exists a linear subspace $U' \subset U$, with $\dim U' \ge (m-1) \ge 2$ such that any $u \in U'$ satisfies $u(y) = 0$. Then, Lemma~\ref{L-zero1} implies that there exists $0 \neq u \in U'$ such that $\rho(u,x) \ge (m-2)$.


The proof of Lemma~\ref{L-zero2} is complete.
\end{proof}


\begin{lemma}\label{L-zeroc}
Let $M$ be a compact surface. Let $U$ be a linear subspace of an eigen\-space of the eigenvalue problem \eqref{E-evp-2c} or \eqref{E-evp-2bc}.
\begin{enumerate}[(i)]
  \item Let $x \in \inte(M)$, and let $u_1, u_2, u_3$ be three linearly independent functions in $U$, such that $ \nu(u_1,x) = \nu(u_2,x) = \nu(u_3,x) \ge 2$. Then, there exists $0 \neq u \in \spanop\set{u_1,u_2,u_3}$ such that $\nu(u,x) \ge \nu(u_1,x) + 2$.
  \item Let $x \in \partial M$, and let $u_1, u_2$ be two linearly independent functions in $U$, such that $ \rho(u_1,x) = \rho(u_2,x) \ge 1$. Then, there exists $0 \neq u \in \spanop\set{u_1,u_2}$ such that $\rho(u,x) \ge \rho(u_1,x) + 1$.
\end{enumerate}
\end{lemma}%

\pf Since the index of a singular point can be expressed in terms of the vanishing order, the lemma follows from Theorem~\ref{T-nodinfo}. Indeed, under the assumption of Assertion~(i), we can write
%
\begin{equation*}
u_i(z) = p_i(z-x) + \cO(|z-x|^{k+1})\,,
\end{equation*}
%
in local coordinates centered at $x$,  where $p_i$ is a nonzero harmonic homogeneous polynomial of degree $k = \frac{\nu(u_1,x)}{2}$ in two variables. Since the vector space of such polynomials has dimension $2$, there exist real numbers $\alpha_1, \alpha_2$ and $\alpha_3$, not all of them equal to zero, such that $\alpha_1 \, p_1 + \alpha_2 \, p_2 + \alpha_3 \, p_3 = 0$. It follows that $\alpha_1 \, u_1 + \alpha_2 \, u_2 + \alpha_3 \, u_3$ vanishes at order at least $(k+1)$ at $x$. This proves Assertion~(i). \smallskip

The proof of Assertion~(ii) is similar, using the local forms \eqref{E-sinexp} or \eqref{E-cosexp} depending on the boundary condition, Dirichlet or Robin.
\hfill \qed \medskip

For later purposes, we introduce the following notation.


\begin{notation}\label{N-evp-dr}
Let $u$ be an eigenfunction of the eigenvalue problem \eqref{E-evp-2bc} in the compact surface $M$. Define the function $\breve{u}$ on $\partial M$ by
%
\begin{equation}\label{E-evp-dr}
\breve{u} = \left\{
\begin{array}{ll}
u|_{\partial M} & \text{in the Robin case,}\\[5pt]
\partial_{\nu}u & \text{in the Dirichlet case.}
\end{array}
\right.
\end{equation}
Then, for any $y \in \partial M$, $\rho(u,y) \ge 1$ if and only if
$\breve{u}(y) = 0$.
\end{notation}%

The following lemma will be useful later on.


\begin{lemma}\label{L-breve}
Let $u$ be an eigenfunction of the eigenvalue problem \eqref{E-evp-2bc}.
\begin{enumerate}[(i)]
  \item If $u$ is a Dirichlet eigenfunction, and $y \in \partial M$, then $u$ vanishes at order $k$ at $y$ if and only if the function $\partial_{\nu}u$ vanishes at order $(k-1)$ at $y$ along $\partial M$.
  \item If $u$ is a Robin eigenfunction, and $y \in \partial M$, then $u$ vanishes at order $k$ at $y$ if and only if the function $u|_{\partial M}$ vanishes at order $k$ at $y$ along $\partial M$.
\end{enumerate}
Therefore, the order of vanishing of the function $\breve{u}$ at some boundary point $y$ is precisely the number $\rho(u,y)$ of nodal  arcs hitting $\partial M$ at $y$.
\end{lemma}%

\pf The proof is by induction on $k$. The equation $\Delta u = (V-\lambda)u$ implies relations between the derivatives of $u$ of degree $k$, evaluated at $y$, assuming that the derivatives of order less than or equal to $(k-1)$ vanish at $y$.\smallskip

More precisely, according to \cite[Section~2]{YaZh2021}, fixing some $y \in \partial M$, we can choose local boundary isothermal coordinates at $y$ such that the equation \break $(-\Delta + V)u = \lambda u$ in a neighborhood of $y$ is transformed into the equation
%
\begin{equation*}
\Delta v= A v \leqno{(e)}
\end{equation*}
%
in some half-ball $\set{(\xi_1,\xi_2) \in \R^2 ~|~ \xi_1^2+\xi_2^2 < a^2, \xi_2 > 0}$, where $0$ is the image of $y$. Here, $\Delta$ is the ordinary Laplacian in the variables $(\xi_1,\xi_2)$, $a$ is some given positive number, $A$ and $v$ are $C^{\infty}$ up to the boundary, and correspond to $(V-\lambda)$ and $u$ respectively.\smallskip

In the proof, we use the following conventions.
\begin{enumerate}[$\diamond$]
  \item The symbol $\equivo{t}$ indicates a trivial identity.
  \item The symbol $\equivo{e}$ indicates an identity which follows from the above identity $(e)$.
  \item The symbol $\equivo{(m)}$ indicates an identity which holds up to a linear combination of derivatives of $v$ of order less than or equal to $m$.
  \item The symbol $\partial^{(p,q)}$ stands for $\frac{\partial^{p+q}}{\partial \xi_1^p\,\partial \xi_2^q}$.
\end{enumerate}

For $k \ge 2$, we have
%
\begin{equation*}
\begin{array}{lcl}
\partial^{(k-2q,2q)}v & \equivo{t} & \partial^{(k-2q,2q-2)} \partial^{(0,2)} v\\[5pt]
& \equivo{e} & \partial^{(k-2q,2q-2)} \left( -\, \partial^{(2,0)}v + A v\right)\\[5pt]
& \equivo{(k-2)} & -\, \partial^{(k-2q+2,2q-2)}v\,.
\end{array}%
\end{equation*}

Assuming that $u$ vanishes at order larger than or equal to $(k-1)$ at $(0,0)$, we obtain that
%
\begin{equation*}
\partial^{(k-2q,2q)}u(0,0) = (-1)^q \partial^{(k,0)}u(0,0), \text{~for~} q \in \set{0,1,\ldots,\floor{\frac k2}}\,. \tag{a}
\end{equation*}

Similarly,
%
\begin{equation*}
\begin{array}{lcl}
\partial^{(k-2q-1,2q+1)}v & \equivo{t} & \partial^{(k-2q-1,2q-1)} \partial^{(0,2)} v\\[5pt]
& \equivo{e} & \partial^{(k-2q-1,2q-1)} \left( -\, \partial^{(2,0)}v + A v\right)\\[5pt]
& \equivo{(k-2)} & -\, \partial^{(k-2q+1,2q-1)}v\,.
\end{array}%
\end{equation*}

Assuming that $u$ vanishes at order larger than or equal to $(k-1)$ at $(0,0)$, we obtain that
%
\begin{equation*}
\partial^{(k-2q-1,2q+1)}u(0,0) = (-1)^q \partial^{(k-1,1)}u(0,0), \text{~for~} q \in \set{0,1,\ldots,\floor{\frac{k-1}{2}}}\,. \tag{b}
\end{equation*}


\noid \emph{Dirichlet case.}~ In this case, $\breve{v}(\xi_1) := \partial^{(0,1)}v(\xi_1,0)$. Since $v(\xi_1,0) \equiv 0$, Equation~(a) implies that
$\partial^{(k-2q,2q)}v(0,0) = 0$ for all $q \in \set{0,\ldots,\floor{\frac k2}}$. If $v$ vanishes at order greater then or equal to $(k-1)$, Equation~(b) implies that
$\partial^{(k-2q-1,2q+1)}v(0,0) = (-1)^q \partial^{(k-1)}\breve{v}(0)$ for all $q \in \set{0,\ldots,\floor{\frac{k-1}{2}}}$. It follows that if $v$ vanishes at order greater than of equal to $(k-1)$, then $v$ vanishes at order greater than or equal to $k$ if and only if $\partial^{(k-1)}\breve{v}(0) = 0$.\medskip

\noid \emph{Robin case.}~ In this case, $\breve{v}(\xi_1) := v(\xi_1,0)$. Assuming that $v$ vanishes at order at least $(k-1)$, Equation~(a) implies that
$\partial^{(k-2q,2q)}v(0,0) = (-1)^q \partial^{k}\breve{v}(0)$. Since $\partial^{(0,1)}v(\xi_1,0) \equiv B(\xi_1) u(\xi_1,0)$ (Robin condition), Equation~(b) implies that $\partial^{(k-2q-1,2q+1)}u(0,0) = 0$. Therefore, if $v$ vanishes at order at least $(k-1)$ at $(0,0)$, then $v$ vanishes at order at least $k$ if and only if $\partial^{k}\breve{v}(0)=0$.\medskip

We have proved that $v$ vanishes at order at least $k$ (resp. equal to $k$) at $(0,0)$ if and only if $\breve{v}$ vanishes at order $\rho(v,(0,0))$ at $(0,0)$. \hfill\qed



\subsection{A global property of nodal sets}\label{SS-gpns}


\begin{lemma}\label{L-hdn}
Let $(M,g)$ be a compact Riemannian surface. Let $w_n, w : M \to \R$ be continuous functions with zero sets $K_n := w_n^{-1}(0)$ and $K := w^{-1}(0)$. Assume that $w_n \to w$ uniformly.
\begin{enumerate}[(i)]
  \item The limit points of the sequence $\set{K_n}$ with respect to the Hausdorff distance associated with the Riemannian distance of $(M,g)$ are compact and contained in $K$. They are connected if the sets $K_n$ are connected.
  \item If $K_n, K$ are nodal sets of eigenfunctions, then the sequence $\set{K_n}$ converges to $K$ in the Hausdorff distance.
\end{enumerate}
\end{lemma}%

\pf The properties that the sequence $\set{K_n}$ has limit points, and that they are compact (and connected if the sets $K_n$ are connected) are general. If a subsequence  $\set{K_{n,1}}$ tends to $K'$ in the Hausdorff distance, for any $z' \in K'$ there exists a sequence $z_n$ with $z_n \in K_n$, such that $z_n \to z'$. Write $w(z') = w(z') - w(z_n) + w(z_n) - w_n(z_n)$, and use the fact that $w_n \to w$ uniformly to conclude that $K' \subset K$. The second assertion uses the nodal character: assume that there exists $z \in K \sm K'$. Then $d(z,K') =: 2\eta > 0$, and for $n$ large enough, $d(z,K_n) > \eta$. Since $K$ is a nodal set, there exists a small arc through $z$, from $z_{-}$ to $z_{+}$ such that $w(z_{-}) < 0$ and $w(z_{+}) > 0$. It follows that for $n$ large enough we also have $w_n(z_{-}) < 0$ and $w_n(z_{+}) > 0$. This shows that $w_n$ must vanish on the small arc, and hence there is some $z'_n \in K_n$ close to $z$ contradicting the fact that $d(z,K_n) > \eta$. Since the only possible limit point of $\set{K_n}$ is $K$, the assertion follows. \hfill \qed


\section{Partitions, graphs and Euler type formulas}\label{S-pargef}


In this section, $M$ denotes a closed surface (compact, no boundary), or a compact surface $M$ with nonempty boundary $\partial M$, and interior $\inte(M)$. Unless otherwise stated, the surfaces are assumed to be smooth and connected. We are interested in Euler type formulas for partitions of $M$, in particular for partitions associated with the eigenfunctions of the eigenvalue problems \eqref{E-evp-2c} or \eqref{E-evp-2bc}.


\subsection{Partitions}\label{SS-parti}

We first recall (or modify) some of the definitions given in the papers \cite{BeHe2014}, \cite{BoHe2017}.\medskip

A \emph{$k$-partition} of $M$ is a collection,  $\cD = \{D_j\}_{j=1}^{k}$, of $k$ pairwise disjoint, connected, open subsets of $M$. We furthermore assume that the $D_j$'s are piecewise $C^1$, and that
%
\begin{equation}\label{E-parti-str}
\inte(\overline{\cup_j D_j} ) = \inte(M)\,.
\end{equation}
%
The set of such partitions is denoted by $\mathfrak{D}_k(M)$. \medskip

The \emph{boundary set} $\partial \mathcal{D}$ of a partition
$\mathcal{D} = \{D_j\}_{j=1}^{k} \in \mathfrak{D}_k(M)$ is the
closed set,
%
\begin{equation}
\partial \mathcal{D} = \overline{\cup_j(\partial D_j \cap M)}\,.
\end{equation}


\begin{definition}\label{D-parti-ess}
A partition $\cD = \set{D_j}_{j=1}^k$ is called \emph{essential} if, for all $j, \,1 \le j \le k$,
%
\begin{equation}\label{E-parti-ess}
\inte(\overline{D_j} ) = D_j\,.
\end{equation}
\end{definition}%

\begin{definition}\label{D-parti-reg}
A \emph{regular $k$-partition} is a $k$-partition
whose boundary set $\bD$ satisfies the following properties:
\begin{itemize}
    \item[(i)] The boundary set $\bD$ is locally a piecewise $C^1$  immersed curve in $M$,  and it is embedded
    except possibly at finitely many points $\{y_i \in \bD \cap
    M\}$ in a neighborhood of which $\bD$ is the union of
    $\nu(y_i)$ $C^1$ semi-arcs meeting at $y_i, \nu(y_i) \ge 3$.  These points are called \emph{interior singular points. The integer $\nu(y_i)$ is called the index of $y_i$.}
    \item[(ii)] The set $\bD \cap \partial M$ consists of finitely
    many points $\{z_j\}$. Near the point $z_j$, the set $\bD$ is the
    union of $\rho(z_j)\ge 1$ $C^1$ semi-arcs hitting $\partial M$ at $z_j$.  These points are called \emph{boundary singular points. The integer $\rho(z_j)$ is called the index of $z_j$.}
    \item[(iii)] The boundary set $\bD$ has the following \emph{transversality property}:
    at any interior singular point $y_i$, the semi-arcs meet transversally; at any boundary singular point $z_j$, the semi-arcs meet transversally, and they meet  the boundary $\partial M$ transversally.
\end{itemize}
\end{definition}%

The subset of regular $k$-partitions is denoted by $\mathcal{R}_k(M) \subset \mf{D}_k(M)$. When $\cD$ is a regular partition, we denote by  $\cS(\cD) = \cS_i(\cD) \cup \cS_b(\cD)$ the set of singular points of $\bD$, where $\cS_i(\cD)$ denotes the set of interior singular points, and $\cS_b(\cD)$ the set of boundary singular points.\medskip


\begin{definition}\label{D-parti-nor}
A regular $k$-partition $\cD = \set{D_j}_{j=1}^k$ is called \emph{normal}, if it satisfies the additional condition, for all $j, ~1 \le j \le k$,
%
\begin{equation}\label{E-parti-nor}
\forall x \in \partial D_j, ~\exists \, r >0 \text{~~s. t.~~} B(x,r)\cap D_j \text{~is~} connected.
\end{equation}
\end{definition}%

\begin{remark}\label{R-parti-nor}
The definition of a normal partition implies that each domain in the partition is a topological surface with boundary (actually a piecewise $C^1$ surface with boundary, possibly with corners).  A normal partition is essential.
\end{remark}%


\begin{example}\label{Ex-parti-2}
Nodal partitions are special examples of essential, regular partitions. They are not necessarily normal, see Subsection~\ref{SS-nodp}.
\end{example}%

\begin{notation}\label{N-parti-2a}
For a partition $\cD \in \mf{D}(M)$, we introduce the following numbers.
\begin{itemize}
  \item[(a)] $\beta(\cD)$ is defined as $\beta(\cD) = b_0(\bD \cup \partial M) - b_0(\partial M)$, the difference between the number of components of $\bD \cup \partial M$, and the number of components of $\partial M$;
  \item[(b)] $\kappa(\cD)$ denotes the number of domains of the partition;
  \item[(c)] $\sigma(\cD) = \sigma_i(\cD)  + \sigma_b(\cD) $ weighs the singular points of a regular partition $\cD$,
  \begin{equation*}
  \left\{
  \begin{array}{ll}
  \sigma_i(\mathcal{D}) = \frac{1}{2}\, \sum_{x \in \mathcal{S}_i(\mathcal{D})}\left( \nu(x) - 2\right),\\[5pt]
\sigma_b(\mathcal{D}) = \frac{1}{2}\, \sum_{x \in \mathcal{S}_b(\mathcal{D})} \rho(x)\,;
  \end{array}%
  \right.
  \end{equation*}
  \item[(d)] $\omega(\cD)$ denotes the \emph{orientability character} of the partition,
%
\begin{equation*}
\hspace{5mm}\left\{
\begin{array}{ll}
\omega(\cD) = 0, & \text{if all the domains of the partition are orientable},\\[5pt]
\omega(\cD) = 1, & \text{if at least one domain of the partition is non-orientable.}
\end{array}%
\right.
\end{equation*}
\end{itemize}
%
Obviously, $\omega(\cD,M) = 0$ whenever the surface $M$ is orientable.
\end{notation}%


\begin{remarks}\label{R-parti-or}~
\begin{enumerate}
  \item[(i)] We use the definition of orientability given in \cite[Chap.~5.3]{BeGo1988} (via differential forms of degree $2$), or the similar form given in \cite[Chap.~4.5]{GaXu2013} in the setting of topological manifolds (via the degree).  A topological surface is orientable if one can choose an atlas whose changes of chart are homeomorphisms with degree $1$.
  \item[(ii)] A compact surface (with boundary) is non-orientable if and only if it contains the homeomorphic image of a M\"{o}bius strip. One direction is clear since the M\"{o}bius strip is not orientable. For the other direction, one can use the classification of compact surfaces (with boundary),  see for example \cite{GaXu2013}.
\end{enumerate}
\end{remarks}%

\begin{lemma}[Normalization]\label{L-parti-norm}
Let $\mathcal D$ be an  essential, regular partition of $M$. Then, one can construct a normal partition $\widetilde{\mathcal D}$ of $M$ such that
%
\begin{equation*}
\left\{
\begin{array}{l}
\beta(\wt{\cD}) = \beta(\cD)\,,\\[5pt]
\kappa(\wt{\cD}) - \sigma(\wt{\cD}) = \kappa(\cD) - \sigma(\cD)\,,\\[5pt]
\omega(\wt{\cD}) = \omega(\cD)\,.
\end{array}%
\right.
\end{equation*}
\end{lemma}

\begin{proof}
Using condition \eqref{E-parti-ess}, we see that an  essential, regular partition $\cD = \set{D_j}_{j=1}^k$ is normal except possibly at  points $x$ in $\cS(\cD)$, with index bigger than $1$, and for which there exists some domain $D_j$ such that $B(x,\varepsilon) \cap M_j$ has at least two components  for all $\varepsilon > 0$. Here, $B(x,\varepsilon)$ denotes the disk with center $x$ and radius $\varepsilon$ in $M$ (we can for example fix a Riemannian metric on $M$ to have a distance function). Let $x$ be such a point. For $\varepsilon$ small enough, introduce the partition $\cD_x$ whose elements are the $D_j\sm \overline{B(x,\epsilon)}$, and the extra domain $B(x,\varepsilon)\cap M$. In this procedure, we have $\kappa(\cD_x) = \kappa(\cD)+1$; an interior singular point $x$, with $\nu(x) \ge 4$ is replaced by $\nu(x)$ singular points of index $3$ for which condition \eqref{E-parti-nor} is satisfied. Hence, $\sigma (\cD_x)= \sigma (\mathcal D) +1$. It follows that $\kappa (\cD_x) - \sigma(\cD_x)= \kappa(\cD) - \sigma(\cD)$. A similar procedure is applied at a boundary singular point. This is illustrated by Figure~\ref{F-gpart-02}. By recursion, we can in this way eliminate all the singular points at which condition \eqref{E-parti-nor} is not satisfied. This procedure does not change $\beta$; choosing $\varepsilon$ small enough,  it does not change $\omega$ either, since the added disks are orientable.
\end{proof}

\begin{figure}[!ht]
\centering
\begin{subfigure}[t]{.4\textwidth}
\centering
\includegraphics[width=\linewidth]{gpart-02-i}
\caption{Interior singular point}
\end{subfigure}\qquad
\begin{subfigure}[t]{.4\textwidth}
\centering
\includegraphics[width=\linewidth]{gpart-02-b}
\caption{Boundary singular point}
\end{subfigure}
%
\caption{Normalization at a singular point}\label{F-gpart-02}
\end{figure}

\begin{remarks}\label{R-parti-norm}~(i) By the same process, one could also remove all singular points with index $\nu > 3$ or $\rho > 1$. We do not need to do that for our purposes. ~(ii) We will use the previous lemma to prove an Euler-type formula for partitions of the M\"{o}bius strip, see Theorem~\ref{T-etfm-2} and
\cite{BeHe2020m}
% Appendix~\ref{S-pfetfm}.
\end{remarks}%


\begin{notation}\label{N-parti-4}
We occasionally use the notation $\kappa(\cD,M)$ for $\kappa(\cD)$, $\cS(\cD,M)$ for $\cS(\cD)$, etc. if we need stress the fact that $\cD$ is a partition of $M$.
% This is the case for example in Appendix~\ref{S-pfetfm}.
\end{notation}%


\subsection{Graphs associated with a partition}\label{SS-gr}

\begin{notation}\label{N-gr-not}
We use the definition of a graph given in \cite{Gibl2010} (i.e. ``graph = simple graph''). Given a graph $G$, we denote by $\alpha_0(G)$ the number of vertices, by $\alpha_1(G)$ the number of edges, and by $c(G)$ the number of components of $G$. For a graph $G$ embedded in a surface $M$, we denote by $r(G,M)$ the number of components of $M \sm G$.
\end{notation}%

Let $M$ be a compact surface, with boundary $\partial M$, possibly empty. Let $\cD$ be a regular partition of $M$, with boundary set $\partial \cD$, and singular set $\cS(\cD)$. The general idea is to associate a graph to the set $\partial \cD \cup \partial M$. The vertices of such a graph should comprise the singular points of $\cD$, and the edges should comprise both sub-arcs of $\partial \cD$ and sub-arcs of $\partial M$. One snag is that the components of $\left( \partial \cD \cup \partial M \right) \sm \cS(\cD)$ might not all be intervals. Indeed, some of them might be circles (components  of $\partial \cD \cup \partial M$  which do not intersect $\cS(\cD)$). For this reason, we first define a `graph-to-be' $G_0 := G_0(\cD,M)$, as follows.\medskip

\begin{itemize}
  \item  Let $e := e(\cD,M)$ be the number of components of $\left( \partial \cD \cup \partial M \right) \sm \cS(\cD)$ which are not an interval. We first choose one vertex for each such component. Call $\set{v_1,\ldots, v_{e}}$ these vertices, if any. Define the set  $V_0$ of vertices of $G_0$ as $\cS(\cD) \cup \set{v_1,\ldots, v_{e}}$.
  \item Define the set $E_0$ of edges of $G_0$ as the set of components of $\left( \partial \cD \cup \partial M\right) \sm V_0$ (they are all intervals).
\end{itemize}

\begin{lemma}\label{L-eucl-2}
The pair $G_0 = (V_0,E_0)$ is a multigraph. The number of vertices $\alpha_0(G_0)$, and the number of edges $\alpha_1(G_0)$ of $G_0$ are given by
%
\begin{equation}\label{E-eucl-2}
\left\{
\begin{array}{l}
\alpha_0(G_0) = e + |\cS_i(\cD)| + |\cS_b(\cD)|\,,\\[5pt]
\alpha_1(G_0) = e + \frac{1}{2} \, \left( \sum_{y \in \cS_i(\cD)} \nu(y)  + \sum_{z\in \cS_b(\cD)} \rho(z)  \right) + |\cS_b(\cD)|\,,
\end{array}
\right.
\end{equation}
where $e$ is defined above, where $|\cS_i(\cD)|$ (resp $|\cS_b(\cD)|$) denotes the number of interior (resp. boundary) singular points of $\cD$, and where the numbers $\nu$ and $\rho$ are as in Definition~\ref{D-parti-reg}. In particular,
%
\begin{equation}\label{E-eucl-4}
\alpha_1(G_0) - \alpha_0(G_0) = \frac{1}{2} \, \left( \sum_{y \in \cS_i(\cD)} (\nu(y)-2)  + \sum_{z\in \cS_b(\cD)} \rho(z)  \right) \,.
\end{equation}
\end{lemma}%

\pf  Clearly $G_0$ is a multigraph in the sense of  \cite[Section~1.10]{Dies2017}. The formula for $\alpha_0$ is clear. The first term in the right-hand side for $\alpha_1$ counts both the number of points $v_j$ and the number of edges contained in the components of $\left( \partial \cD \cup \partial M \right) \sm \cS(\cD)$ which are circles. The second term counts the number of edges between two singular points, each one being a simple curve contained in $\partial \cD \sm \cS(\cD)$. The third term counts the number of edges determined by the boundary singular points on the components of $\partial M$ which intersect $\cS(\cD)$. \hfill \qed \medskip


\begin{remark}\label{R-eucl-2L}
The second relation in \eqref{E-eucl-2} also follows from the relation
%
\begin{equation}\label{E-eucl-2R}
2 \, \alpha_1(\Gamma) = \sum_{x \in V(\Gamma)} \deg_{\Gamma}(x)
\end{equation}
%
which holds for any multi-graph $\Gamma$ (here, $\deg_{\Gamma}(x)$, the degree of the vertex $x$, is the number of edges of $\Gamma$ one of whose ends is $x$).
\end{remark}%

Note that the multigraph $G_0(\cD,M)$,
\begin{itemize}
\item[(a)] might contain loops at some singular point or vertex $v_j$ ;
\item[(b)] might contain pairs of distinct vertices in $V_0$ linked by more than one edge.
\end{itemize}

In order to apply the results of \cite{Gibl2010}, we transform the multigraph $G_0(\cD,M)$ into a graph $G(\cD,M)$ (in the sense of \cite[Section~1.1]{Dies2017} or \cite[p.~10]{Gibl2010}), keeping track of the number of vertices and edges. More precisely, if necessary, we introduce additional vertices and edges by performing one of the following  vertex-edge additions.

\begin{definition}\label{D-gr-2} We call \emph{vertex-edge additions} the following modifications of the graph $G_0(\cD,M)$.
\begin{enumerate}
  \item[(i)] If a component $\Gamma$ of $\left( \partial \cD\cup \partial M\right) \sm V_0$ is bounded by only one vertex $v$ (i.e., there is a loop at $v$), we add two extra vertices $v_1, v_2$ on $\Gamma$, and replace the edge $\Gamma$ by three edges, the components of $\Gamma \sm \set{v_1,v_2}$.
  \item[(ii)] If two distinct vertices of $V_0$ are the endpoints of more than one edge in $E_0$, i.e., of more than one components $\Gamma_j$ of $\left( \partial \cD\cup \partial M\right) \sm V_0$, we add one extra vertex $w_j$ to each $\Gamma_j$, and replace $\Gamma_j$ by two edges, the components of $\Gamma_j \sm \set{w_j}$.
\end{enumerate}
\end{definition}%

 Figure~\ref{F-gpart-03} illustrates the transformation of $\bD \cup \partial M$ into a graph. Lines contained in the boundary appear in black and lines contained in $\bD$ appear in red. Blue dots represent vertices $v_i$ initially attached to each circle component. Green dots represent vertices added in the vertex-edge additions. Sub-figure (A) illustrated the transformation of circle components of $\partial M$ or $\bD$, and loop in $\bD$ into graphs. Sub-figure (B) illustrates the transformation of multiple edges into graphs.

\begin{figure}[!ht]
\centering
\begin{subfigure}[t]{.5\textwidth}
\centering
\includegraphics[width=\linewidth]{gpart-03-a}
\caption{Circle components and loops}
\end{subfigure}\qquad
\begin{subfigure}[t]{.33\textwidth}
\centering
\includegraphics[width=\linewidth]{gpart-03-b}
\caption{Multiple edges}
\end{subfigure}
%
\caption{Vertex-edge additions}\label{F-gpart-03}
\end{figure}

\begin{lemma}\label{L-gr-4}
Performing finitely many  vertex-edge additions transforms the multigraph $G_0(\cD,M)$ into a graph $G(\cD,M)$.
%(with curved edges, see \cite[\S~1.34, p. 32]{Gibl2010})
\end{lemma}

The following lemma follows from the fact that the number $\alpha_0 - \alpha_1$ remains unchanged if we perform a vertex-edge addition.

\begin{lemma}\label{L-gr-6}
For the graph $G(\cD,M)$ obtained from the multigraph $G_0(\cD,M)$ by performing  vertex-edge additions, we have,
%
\begin{equation*}
\left\{
\begin{array}{rll}
\alpha_1(G) - \alpha_0(G) & = \alpha_1(G_0) - \alpha_0(G_0) & = \sigma(\cD) \,,\\[5pt]
c(G) & = c(G_0) & = b_0(\bD \cup \partial M)\,,\\[5pt]
r(G,M) & = r(G_0,M) & = \kappa(\cD)\,.
\end{array}%
\right.
\end{equation*}
\end{lemma}%



\subsection{Euler type formulas for partitions}\label{SS-etf}

Let $\cD$ be a regular partition of a compact surface $M$, with boundary $\partial M$. In this subsection, we look at three examples.

\subsubsection{The case of the sphere}\label{SS-etf0}

\begin{proposition}\label{P-etf0-2} Let $M_0$ be the sphere $\bS^2$. For $q \ge 1$, let $M_q = M_0\sm \cup_{j=1}^q U_j$, where the $U_j, 1 \le j \le q$, are simply connected, pairwise disjoint, open sets, with piecewise $C^1$ boundary. Let $\mathcal D_q$ be a regular partition of $M_q, q\ge 0$. Then,
\begin{equation}\label{E-etf0-2}
\begin{split}
\kappa(\cD_q) & = 1 + \beta(\cD_q) + \sigma(\cD_q) \\[5pt]
& = 1 + b_0(\bD_q \cup \partial M_q) - b_0(\partial M_q)  + \sigma(\cD_q)\\[5pt]
& \ge 2 - q  + \sigma(\cD_q)\,.
\end{split}%
\end{equation}
\end{proposition}

\pf  Let $G_0 := G(\cD_0,M_0)$ be a graph associated with $\cD_0$ as in Subsection~\ref{SS-gr}.  In this case, $\beta(\cD_0) = b_0(\bD_0) = c(G_0)$. By Lemma~\ref{L-gr-6}, we can write
%
\begin{equation*}
\kappa(\cD_0) - \beta(\cD_0) - \sigma(\cD_0) = r(G_0,M_0) -c(G_0) - \alpha_1(G_0) + \alpha_0(G_0) = 1\,,
\end{equation*}
%
where the equality on the right is given by Euler's formula for graphs in $\bS^2$ (planar graphs). \medskip

Let $M_q = M_0 \sm \cup_{j=1}^q U_j$. Let $\cD_q$ be a partition of $M_q$, and let $G_q = G(\cD_q,M_q)$ be a graph as constructed in Subsection~\ref{SS-gr}. Then, $\cD_q' = \cD_q \cup \set{U_1,\ldots,U_q}$ is a partition of the sphere $M_0$. The graph $G_q$ viewed as a graph in $M_0$ is also a graph of the form $G(\cD_q',M_0)$. We have $r(G_q,M_0) = r(G_q,M_q) + q$. We also have
%
\begin{equation*}
\beta(\cD_q) := b_0(\bD_q \cup \partial M_q) - b_0(\partial M_q) = c(G_q) - q\,.
\end{equation*}

Applying Lemma~\ref{L-gr-6} to the pair $(\cD_q,M_q)$, we obtain
%
\begin{equation*}
\begin{split}
\kappa(\cD_q) - \beta(\cD_q) - \sigma(\cD_q) & = r(G_q,M_q) + q - c(G_q) + \alpha_0(G_q) - \alpha_1(G_q)\\[5pt]
& = r(G_q,M_0) - c(G_q) + \alpha_0(G_q) - \alpha_1(G_q) = 1\,,
\end{split}%
\end{equation*}
 %
 where the last equality follows from Euler's formula for the graph $G_q$ viewed as a graph in the sphere. \hfill \qed


\begin{remark}\label{R-etf0-2P}
Formula \eqref{E-etf0-2} applies in particular to a bounded domain $\Omega \subset \R^2$, with piecewise $C^1$ boundary, and $q$ boundary components.
\end{remark}%


\subsubsection{The general case}\label{SS-etfc}

More generally, we consider,  for $c\in \mathbb Z$,  a closed surface $\Sigma_{c}$ with Euler characteristic $\chi(\Sigma_c)=c$.

\begin{proposition}\label{P-etfc-2}
Let $\Sigma_c$ be a closed surface with Euler-characteristic $\chi(\Sigma_c)=c$. For $q \ge 1$, let $\Sigma_{c,q} = \Sigma_c\sm \cup_{j=1}^q U_j$, where the $U_j, 1 \le j \le q$, are simply connected, pairwise disjoint, open sets, with piecewise $C^1$ boundary. Let $\mathcal D_q$ be a regular partition of $\Sigma_{c,q}, q\ge 0$. Then,
\begin{equation}\label{E-etfc-2}
\kappa(\cD_q,\Sigma_{c,q}) \ge  \chi(\Sigma_c) - q + \sigma(\cD_q,\Sigma_{c,q}) = \chi(\Sigma_{c,q}) + \sigma(\cD_q,\Sigma_{c,q})\,.
\end{equation}
\end{proposition}

\pf The proof mimics the proof of Proposition~\ref{P-etf0-2}, and make use of Euler's inequality for graphs in a surface $\Sigma_c$, \cite[Corollary~9.27]{Gibl2010}. To the partition $\cD_q$, we associate a graph $G_q = G(\cD_q,\Sigma_{c,q})$ as in Subsection~\ref{SS-gr}. We consider $\cD_q' = \cD_q \cup \set{U_1,\ldots, U_q}$. This is a partition of $\Sigma_c$. The graph $G_q$, viewed as a graph in $\Sigma_c$, is also of the form $G(\cD_q',\Sigma_c)$, and $r(G_q,\Sigma_c) = r(G_q,\Sigma_{c,q}) + q$. We have $\beta(\cD_q) = b_0(\bD_q \cup \partial \Sigma_{c,q}) - b_0(\partial \Sigma_{c,q}) = c(G_q) - q$. Then, using Lemma~\ref{L-gr-6} and applying Euler's inequality \cite[Corollary~9.27]{Gibl2010}, we can write
%
\begin{equation*}
\begin{split}
\kappa(\cD_q,\Sigma_{c,q}) & = r(G_q,\Sigma_{c,q}) = r(G_q,\Sigma_{c}) - q\\[5pt]
& \ge \chi(\Sigma_c) + \alpha_1(G_q) - \alpha_0(G_q)  - q = \chi(\Sigma_c) - q + \sigma(\cD_q,\Sigma_{c,q})\\[5pt]
& \ge \chi(\Sigma_{c,q}) + \sigma(\cD_q,\Sigma_{c,q}).
\end{split}%
\end{equation*}
\hfill \qed


\begin{remark}\label{R-etfc-2P}
Obviously, formula \eqref{E-etfc-2} is less precise than formula \eqref{E-etf0-2}. In the case of the M\"{o}bius strip, we can give a better statement.
\end{remark}%


\subsubsection{The case of the M\"{o}bius strip}\label{SS-etfm}

\begin{theorem}\label{T-etfm-2}
Let $\cD$ be an  essential, regular partition of the M\"{o}bius strip $\M$.
Then, with the Notation~\ref{N-parti-2a},
%
\begin{equation}\label{E-etfm-2}
\kappa(\cD) = \omega(\cD) + \beta(\cD) + \sigma(\cD)\,.
\end{equation}
\end{theorem}

% We give the proof of this theorem in Appendix~\ref{S-pfetfm}.
The proof of this theorem is given in \cite{BeHe2020m}. The M\"{o}bius strip is homeomorphic to a projective plane with a disk removed. We can rewrite Formula~\ref{E-etfm-2} as follows.
%
\begin{equation*}
\kappa(\cD) = \omega(\cD) + b_0(\bD \cup \partial \M) - 1 + \sigma(\cD)\,.
\end{equation*}
%
Since $\omega(\cD) \ge 0$ and $b_0(\bD \cup \partial \M) \ge 1$, we have
%
\begin{equation*}
\kappa(\cD) \ge \chi(\rP) - 1 + \sigma(\cD) = \chi(\M) + \sigma(\cD)\,.
\end{equation*}
%
Formula~\ref{E-etfm-2} is more precise than Formula~\ref{E-etfc-2} applied to the projective plane with a disk removed.


\subsection{Nodal partitions}\label{SS-nodp}


The partition $\cD_u$ of a compact surface $M$ (with or without boundary) associated with an eigenfunction $u$ of the operator $-\Delta + V$ (with some prescribed boundary condition if $\partial M \not = \emptyset$) is called a \emph{nodal partition}. The domains of the partition $\cD_u$ are the nodal domains of $u$, the boundary set $\bD_u$ is the nodal set $\cZ(u)$, the singular set $\cS(\cD_u)$ is the set $\cS(u)$ of  singular points of $u$ (Definition~\ref{D-evp2-0}),  and we write $\cS(u) = \cS_i(u) \cup \cS_b(u)$. This is an example of an essential, regular partition. \medskip

Nodal partitions are not necessarily normal as can be seen in Figure~\ref{F-gpart-04} which displays the nodal partitions of two eigenfunctions of the flat M\"{o}bius strip with fundamental domain $(0,\pi)\times (0,\pi)$; the horizontal sides are identified, $(x,0) \sim (\pi - x,\pi)$, and the Dirichlet boundary condition holds on the sides $\set{x=0}$ and $\set{x=\pi}$,
see \cite{BeHK2021m}.  Other examples can be found in \cite{BeHK2021k}.

\begin{figure}[!ht]
\centering
\includegraphics[scale=0.3]{gpart-04}
\caption{Partition associated with eigenfunctions of the M\"{o}bius strip}\label{F-gpart-04}
\end{figure}
\FloatBarrier

For nodal partitions, we write $\kappa(u), \beta(u), \ldots$ for the corresponding $\kappa(\cD_u), \beta(\cD_u), \ldots$, and we have $\nu(\cD_u,z) = \nu(u,z)$ and $\rho(\cD_u,z) = \rho(u,z)$ for the corresponding indices.\medskip

For nodal partitions we have the following additional properties.

\begin{proposition}\label{P-euler-np}
Let $u$ be an eigenfunction of $M$. For any component $\Gamma$ of $\partial M$, we have
%
\begin{equation}\label{E-2.80}
\sum_{z\in \cS_b(u) \cap \Gamma}\, \rho(z) \in 2\N\,.
\end{equation}
%
Furthermore,
%
\begin{equation}\label{E-2.8b}
b_0(\cZ(u) \cup \partial M) -b_0(\partial M) + \frac 12 \, \sum_{y \in \cS_b(u)} \rho(y) \geq 1\,,
\end{equation}
\end{proposition}


\pf The first assertion is general and contained in Corollary~\ref{cor:nodinfo}. To prove the second assertion, we divide the components of $\partial M$ into two sets: the components $\Gamma'_i, 1 \le i \le p$, which meet $\cZ(u)$, and the components $\Gamma^{''}_j, 1 \le j \le q$, which do not meet $\cZ(u)$. Let $\Gamma(u) = \cup_{i=1}^p\Gamma'_i$. Clearly, we have the relation
%
\begin{equation*}
b_0(\cZ(u) \cup \partial M) -b_0(\partial M) = b_0(\cZ(u) \cup \Gamma(u)) -b_0(\Gamma(u))\,.
\end{equation*}

On the other-hand, according to \eqref{E-2.80}, for each $1 \le i \le p$, we have
%
\begin{equation*}
\sum_{z \in \cS_b(u) \cap \Gamma'_i} \rho(z) \ge 2\,.
\end{equation*}

Relation \eqref{E-2.8b} follows from the fact that $b_0(\cZ(u) \cup \Gamma(u)) \ge 1$. \smallskip
\hfill \qed


\begin{corollary}\label{C-euler-npm}
If $u$ is an eigenfunction of the M\"{o}bius strip $\M$, with $\cZ(u) \neq \emptyset$,
\begin{equation}\label{E-2.8c}
\kappa(u) \geq  \omega(\cD_u) + 1 + \frac 12 \, \sum_{x\in \cS_i(u)} (\nu(x) -2)\,.
\end{equation}
\end{corollary}%

\pf Relation \eqref{E-2.8c} is a consequence of \eqref{E-etfm-2} applied to the nodal partition associated with $u$, and \eqref{E-2.8b}. \hfill \qed


\section{Revisiting the multiplicity bounds for closed surfaces of genus $0$}\label{S-h2n}

In this section, we revisit the paper \cite{HoHN1999} by M. and T.~Hoffmann-Ostenhof and N.~Nadirashvili.

\subsection{Introduction}\label{SS-h2n1}

Let $(M,g)$ be a closed connected smooth Riemannian surface with genus $0$, i.e. $M$ is topologically a sphere (in particular it is simply connected with Euler characteristic $2$), and $g$ is a smooth metric. We are also given a smooth, real valued function $V$ on $M$, and we consider the eigenvalue problem \eqref{E-evp-2c} for the Schr\"{o}dinger operator $-\Delta + V$ on $M$, with eigenvalues $\set{\lambda_k}_{k=1}^{\infty}$ (with first label $1$). \medskip

According to  Cheng, \cite{Chen1976}, $\mult(\lambda_2) \le 3$ and this bound is sharp, achieved for the round metric on $\bS^2$ and $V=0$. Nadirashvili, \cite{Nadi1987}, proved that, for $ k \ge 1$, $\mult(\lambda_k) \le (2k-1)$. In \cite{HoHN1999}, M. and T. Hoffmann-Ostenhof and Nadirashvili, prove the following result.


\begin{theorem}\label{T-h2n}
For $(M,g,V)$ as above, and for $k \ge 3$, $\mult(\lambda_k) \le (2k-3)$.
\end{theorem}%

Here is the underlying idea of the proof of Theorem~\ref{T-h2n}.
Fix some $k \ge 3$, and some $x \in M$. Consider the eigenspace $U(\lambda_k)$, and recall that $M$ is topologically a sphere. \smallskip
Assume that $\mult(\lambda_k) = \dim U(\lambda_k) = m$ for some integer $m \ge 2$. Then, according to Lemma~\ref{L-zeroi}, there exists a function $0 \not = u \in U(\lambda_k)$ such that $\nu(u,x) \ge 2 \floor{\frac m2}$ (where $\floor{\frac m2}$ denotes the integer part of $\frac m2$), and actually two linearly independent such functions when $m$ is odd. By Courant's theorem, the number of nodal domains of $u$ satisfies $\kappa(u) \le k$. Since  $M$ is  topologically   a sphere, we can apply Euler's formula \eqref{E-etf0-2} to the nodal partition $\cD_{u}$ of the function $u$,
%
\begin{equation}\label{E-h2n-2}
\kappa(u) = 1 + b_0(\cZ(u)) + {\frac 12} \, \sum_{z \in \cS(u)} \left( \nu(u,z) - 2 \right).
\end{equation}

Summing up the above information, we have
%
\begin{equation}\label{E-h2n-4}
\begin{split}
0 \ge \kappa(u) - k & = \left\lbrace b_0(\cZ(u)) - 1\right\rbrace + \sum_{\substack{z \in \cS(u) \\z \not = x}} \frac{\nu(u,z) - 2}{2}
+ \lbrace \frac{\nu(u,x)}{2} - k + 1 \rbrace\\[5pt]
& \ge \left\lbrace b_0(\cZ(u)) - 1\right\rbrace + \sum_{\substack{z \in \cS(u) \\z \not = x}} \frac{\nu(u,z) - 2}{2}
+ \left\lbrace \floor{\frac m2} - k + 1 \right\rbrace\,.
\end{split}
\end{equation}

Notice that the first two terms in the right-hand side of both lines of \eqref{E-h2n-4} are nonnegative. It follows that
%
\begin{equation}\label{E-h2n-4a}
\floor{\frac m2} \le k - 1 \quad \text{i.e.~~} \dim U(\lambda_k) \le (2k-1).
\end{equation}

 The above inequality is sharp for $k=1$ and $2$, see Table~\ref{E-scr-2T}. \smallskip

In view of \eqref{E-h2n-4a}, to prove Theorem~\ref{T-h2n}, it suffices to show that the cases $\dim U(\lambda_k) = (2k-1)$, and $\dim U(\lambda_k) = (2k-2)$ cannot occur. This is the purpose of the following two subsections, revisiting the arguments of \cite{HoHN1999}.



\subsection{Surfaces with genus zero, proof that $\dim U(\lambda_k) \le (2k-2)$ for $k\ge 3$}\label{SS-h2n-s2}\phantom{}

The proof is by contradiction. Taking \eqref{E-h2n-4a} into account, we assume that
%
\begin{equation}\label{E-h2n-s2}
k \ge 3 \text{~~and~~} \dim U(\lambda_k) = 2k-1\,.
\end{equation}
%
We fix some $x \in M$, and choose an orientation in $T_xM$.\smallskip

Introduce the linear subspace
%
\begin{equation}\label{E-h2n-s2a}
U_x := \set{u \in U(\lambda_k) ~|~ \nu(u,x) \ge 2(k-1)}\,.
\end{equation}

\subsubsection{Structure and combinatorial type of $\cZ(u)$, for $0 \neq u \in U_x$}\label{SSS-h2n-s2s}

\begin{properties}\label{propertyD2}
%Assume that  $k \ge 3$, and  $\dim U(\lambda_k) = (2k-1)$.\\
Assume that \eqref{E-h2n-s2} holds. Then, for any $0 \not = u \in U_x$,
\begin{enumerate}[(i)]
 \item $\nu(u,x) = 2(k-1)$, and $x$ is the only singular point of the function $u$;
 \item $\cZ(u)$ is connected;
 \item $\kappa(u) = k$.
\end{enumerate}
\noindent Furthermore, $\dim U_x = 2$, and there exists a basis $\set{v_1,v_2}$ of $U_x$ such that, in some local polar coordinates $(r,\omega)$ centered at $x$,
%
\begin{equation}\label{E-h2n-5}
\left\{
\begin{array}{l}
v_1 = r^{k-1} \sin ((k-1)\omega) +\mathcal O (r^k)\,,\\
v_2 = r^{k-1} \cos ((k-1)\omega) +\mathcal O (r^k)\,.
\end{array}
\right.
\end{equation}
%
\end{properties}


\pf  According to Lemma~\ref{L-zeroi}~(ii), there exist two linearly independent functions $u_{1}, u_{2} \in U(\lambda_k)$, with $\nu(u_{i},x) \ge 2\floor{\frac m2} = 2(k-1)$, for $i = 1, 2$.  Given any $0 \not = u \in U_x$, we can apply Inequality~\eqref{E-h2n-4} to $u$, with $m = (2k-1)$, so that
%
\begin{equation}\label{E-h2n-6}
\begin{split}
0 \ge \kappa(u) - k &= \left\lbrace b_0(\cZ(u)) - 1\right\rbrace + \sum_{\substack{z \in \cS(u)\\ z \not = x}} \frac{\nu(u,z) - 2}{2} + \lbrace \frac{\nu(u,x)}{2} - k + 1 \rbrace \ge\\[5pt]
&\ge \left\lbrace b_0(\cZ(u)) - 1\right\rbrace + \sum_{\substack{z \in \cS(u)\\ z \not = x}} \frac{\nu(u,z) - 2}{2} \ge 0\,.
\end{split}%
\end{equation}
%
The terms in the right-hand sides of both lines of \eqref{E-h2n-6} are nonnegative, and their sum is zero. This proves Assertions~(i)--(iii).\smallskip

From the previous paragraph, we know that $U_x$ has dimension at least $2$. Assume that it has dimension at least $3$, and let $v_1, v_2, v_3$ be three linearly independent functions in $U_x$. Since $\nu(v_i,x) = 2(k-1)$,  Lemma~\ref{L-zeroc} implies the existence of a nontrivial linear combination $v$ of these functions with $\nu(v,x) \ge 2k$, contradicting Assertion~(i). Hence, $\dim U_x = 2$.\smallskip

In view of \eqref{E-harmpoly}, choosing the local polar coordinates $(r,\omega)$ appropriately, and scaling the functions if necessary, we can start from a basis $\set{v_1,w_2}$ of $U_x$ such that
%
\begin{equation*}
\left\{
\begin{array}{l}
v_1 = r^{k-1} \sin ((k-1)\omega) +\mathcal O (r^k)\,,\\
w_2 = r^{k-1} \sin ((k-1)\omega + \beta) +\mathcal O (r^k)\,,
\end{array}
\right.
\end{equation*}
 %
for some angle $\beta \in (0,\pi)$. For $\theta \in (0,\pi)$, consider the functions
%
\begin{equation*}
\begin{split}
\cos\theta \, v_1 + \sin\theta \, w_2 & = r^{k-1} \, \left( \cos\theta  + \sin\theta \cos\beta \right)\, \sin((k-1)\omega) \\[5pt]
& \hspace{10mm}+ r^{k-1}\, \sin\theta\sin\beta\cos((k-1)\omega) + \cO(r^k).
\end{split}%
\end{equation*}

Choose $\theta \in (0,\pi)$ such that $\cos\theta + \sin\theta \, \cos\beta = 0$, then $\sin\theta \, \sin\beta \not = 0$, and the function
%
\begin{equation*}
v_2 = \left( \cos\theta \, v_1 + \sin\theta \, w_2 \right)/ \left( \sin\theta \, \sin\beta \right)
\end{equation*}
%
satisfies the second condition in \eqref{E-h2n-5}. This proves the last assertion, and the proof of Properties~\ref{propertyD2} is now complete. \hfill \qedc \medskip

 Let $0 \neq u \in U_x$. The local structure theorem -- Corollary~\ref{cor:nodloc}~(i) -- implies that, in a neighborhood of $x$, the nodal set $\cZ(u)$ consists of  $2(k-1)$ nodal semi-arcs which emanate from $x$, tangentially to $2(k-1)$ rays dividing the unit circle in $T_xM$ into equal parts. Since $\cS(u) = \set{x}$, if we follow a semi-nodal arc emanating from $x$, we will eventually come back to $x$. Using the fact that $\cS(u) = \set{x}$ and the connectedness of $\cZ(u)$, we conclude that $\cZ(u)$ consists of $(k-1)$ simple loops at $x$, and that these loops only intersect each other at $x$.

\begin{definition}\label{D-h2n-rose}
 Define a \emph{$p$-bouquet of loops} at $x$ as a collection of $p$ piecewise $C^1$ loops at $x$, whose semi-tangents at $x$ are pairwise transverse in $T_xM$, and which do not intersect away from $x$.
\end{definition}%

Therefore, for any $0 \neq u \in U_x$, the nodal set $\cZ(u)$ is a $(k-1)$-bouquet of loops at the point $x$.\medskip

Choosing a direct frame $\set{\vec{e}_1, \vec{e}_2}$ in $T_xM$, we label the rays tangent to $\cZ(u)$ at $x$ counter-clockwise, according to their angle with respect to $\vec{e}_1$. We obtain an ordered list $\set{\vartheta_1, \ldots, \vartheta_{2(k-1)}}$ in which two consecutive rays make an angle $\frac{\pi}{k-1}$. The loops in the $(k-1)$-bouquet of loops $\cZ(u)$ can now be described by a map
%
\begin{equation}\label{E-h2n-tau}
\tau_{x,u} : \set{1,\ldots,2(k-1)} \to \set{1,\ldots,2(k-1)}\,,
\end{equation}
%
which is defined  in the following way. For $j \in   \set{1,\ldots,2(k-1)}$,  we consider the loop which emanates from $x$ tangentially to the ray $\vartheta_j$, and define  $\tau_{x,u}(j)$ as the labelling of the ray  tangent to the loop when it arrives back at $x$. We then denote this loop by $\gamma^{x,u}_{j,\tau_{x,u}(j)}$. We shall write $\tau$ instead of $\tau_{x,u}$ when there is no ambiguity.\smallskip

The properties of nodal sets imply that
%
\begin{equation}\label{E-h2n-taup}
\left\{
\begin{array}{l}
\tau_{x,u}(j) \neq j ~~ \text{~for all~} j \in \set{1,\ldots,
2(k-1)}\,,\\[5pt]
\tau_{x,u}^2 = \id\,.
\end{array}%
\right.
\end{equation}

Note that changing the frame $\set{\vec{e}_1,\vec{e}_2}$ at $x$ amounts to conjugating the map $\tau_{x,u}$ by a circular permutation of the set  $\set{1,\ldots,2(k-1)}$.

\begin{definition}\label{D-h2n-type}
We call the map $\tau_{x,u}$ the \emph{combinatorial type} of the nodal set $\cZ(u)$ at the point $x$.
\end{definition}%

\subsubsection{The rotating function argument}\label{SSS-h2n-s2r}

Fix the basis $\set{v_1,v_2}$ of $U_x$ provided by Properties~\ref{propertyD2}.  Fix the direct frame $\set{\vec{e}_1, \vec{e}_2}$ in $T_xM$ corresponding to the local polar coordinates $(r,\omega)$ at $x$ such that \eqref{E-h2n-5} holds.\smallskip

We now analyze the nodal sets of the one-parameter family,
%
\begin{equation}\label{E-h2n-6a}
w_\theta = \cos((k-1) \theta) \,v_1 - \sin((k-1) \theta) \, v_2\,,
\end{equation}
%
for $\theta \in \left[ 0, \frac{\pi}{k-1}\right]$. In particular, we have
%
\begin{equation}\label{E-h2n-6b}
v_1 = w_0 = - w_{\frac{\pi}{k-1}} \text{~~and~~} v_2 = - w_{\frac{\pi}{2(k-1)}}\,,
\end{equation}
%
and
%
\begin{equation}\label{E-h2n-6c}
w_\theta = r^{k-1}\, \sin((k-1)(\omega - \theta)) \, +\mathcal O (r^k)\,.
\end{equation}

Properties~\ref{propertyD2} state that $x$ is the sole singular point of the eigenfunction $w_{\theta}$, and that the nodal set $\cZ(w_{\theta})$ is connected. With respect to the frame $\set{\vec{e}_1, \vec{e}_2}$ in $T_xM$, there are $2(k-1)$ nodal semi-arcs which emanate from $x$, tangentially to the $2(k-1)$ rays $\set{\omega = \omega_j(\theta) := \omega_j + \theta}$, where $\omega_j :=j \frac{\pi}{k-1}, j \in \set{0,\ldots,2k-3}$. We call these rays $\omega_j(\theta)$ for short, and we view $j$ as defined modulo $2(k-1)$.\medskip

The nodal set $\cZ(w_{\theta})$ is a $(k-1)$-bouquet of loops described by the map $\tau_{x,w_{\theta}}$ associated with the rays $\omega_j(\theta)$. Call this map $\tau_{\theta}$ for short,  and call $\gamma_{j,\tau_{\theta}(j)}^{\theta}$ the corresponding loops at $x$. \smallskip

\begin{property}\label{propertyD3}
 Assume that \eqref{E-h2n-s2} holds.
Considering $\tau_{\theta}(j)$ instead of $j$ if necessary, we can assume that $0 \le j < \tau_{\theta}(j) \le 2k-3$. The loop $\gamma_{j,\tau_{\theta}(j)}^{\theta}$ separates  (the topological sphere) $M$ into two components. The rays $\omega_k(\theta)$ such that $j < k < \tau_{\theta}(j)$ point inside one of the two components ; the rays $\omega_k(\theta)$ with $k < j$ or $k > \tau_{\theta}(j)$ point inside the other component. In particular, $\tau_{\theta}(j) - j$ is an odd integer.
\end{property}%

The proof of this property is clear.\hfill \qedc

\begin{property}\label{propertyD4}
 Assume that \eqref{E-h2n-s2} holds.
The map $\tau_{\theta}$  (i.e. the combinatorial type of $w_{\theta}$) does actually not depend on $\theta$. More precisely, for any $j \in \set{0,\ldots,2k-3}$, and for any $\theta \in \left[ 0,\frac{\pi}{k-1}\right]$, the loop which emanates from $x$ tangentially to the ray $\omega_j(\theta)$ arrives at $x$ tangentially to the ray $\omega_{\tau_0(j)}(\theta)$.  We shall henceforth denote this map by $\tau$.
\end{property}%

\pf Since all the functions $w_{\theta}$ share the same properties, it suffices to show that $\tau_{\theta} = \tau_0$ for $\theta$ small enough. Assume the contrary. Then, there exists a sequence $\theta_n$ tending to zero, and a sequence $\set{j_n} \subset \set{0,\ldots,2k-3}$ such that $\tau_{\theta_n}(j_n) \not = \tau_0(j_n)$. Since the sequence $\{j_n\}$ takes finitely many values, we can find a constant subsequence $\set{j_{n,1}} \subset \set{j_n}$. Similarly, we can take a subsequence $\set{j_{n,2}} \subset \set{j_{n,1}}$  such that $\tau_{\theta_{n,2}}(j_{n,2})$ is constant. Hence, there exists some $\ell \in \set{0,\ldots,2k-3}$, and a sequence $\theta_n$ such that $\tau_{\theta_n}(\ell) \not = \tau_0(\ell)$. Without loss of generality, we may assume that $\ell = 0$, so that there exists a sequence $\theta_n$ tending to zero such that $\tau_{\theta_n}(0) \equiv \ell_0 \not = \tau(0)$.\smallskip

We now use a more precise version of the local structure theorem, see Appendix~\ref{S-lsns}. For any $\alpha > 0$ small enough, there exists $r_0 > 0$ such that, for all $\theta$, $ \cZ(w_{\theta}) \cap B(x,2r_0)$ consists of $2(k-1)$ nodal semi-arcs
%
\begin{equation}\label{E-h2n-8a}
A_j(r,\theta) : (0,2r_0) \ni r \mapsto \exp_x\left(r\, \omegat_j(r,\theta) \right) \in B(x,2r_0)\,,
\end{equation}
%
for $j \in \set{0,\ldots,2k-3}$. Here, we assume that an orthonormal frame $\set{e_1,e_2}$ has been chosen in $T_xM$, is such a way that the vector $e_1$ directs the ray $\omega_0$. In the polar coordinates $(r,\omega)$ associated with this frame, we identify the angle $\omega$ with a point on the unit circle. Furthermore, the functions $\omegat_j$ are smooth in $(r,\theta) \in (0,2r_0) \times [0,2\pi]$, and they satisfy,
%
\begin{equation}\label{E-h2n-8b}
\left\{
\begin{array}{l}
\omegat_j(r,\theta) \in \left( \omega_j+\theta - \alpha , \omega_j + \theta +\alpha \right)\,,\\[5pt]
\lim_{r \to 0} \omegat_j(r,\theta) = \omega_j + \theta\,,
\end{array}
\right.
\end{equation}
%
for all $j, 0 \le j \le 2k-3$. The semi-arc $A_j(r,\theta)$ is semi-tangent to the ray $\omega_j+\theta$ at the point $x$.\smallskip

We now reason as in the proof of Lemma~\ref{L-hdn}. In the closed ball $\bar{B}(x,r_0)$, the nodal set $\cZ(\omega_{\theta_n})$ consists of $2(k-1)$ nodal semi-arcs $A_j(\cdot,\theta_n)$, with end points $x$ and $\exp_x(r_0\, \omegat_j(r_0,\theta_n))$ which converge to the corresponding semi-arcs $A_j(\cdot,0)$ with end points $x$ and $\exp_x(r_0\, \omegat_j(r_0,0))$.\smallskip

In the compact set $M\sm B(x,r_0)$, the nodal set $\cZ(\omega_{\theta_n})$ consists of $(k-1)$ disjoint connected nodal arcs $C_j(r_0,\theta_n)$ with two end points $\exp_x(r_0\, \omegat_j(r_0,\theta_n))$ and $\exp_x(r_0\, \omegat_{\tau(j)}(r_0,\theta_n))$, which correspond to the intersections of the loops in $\cZ(w_{\theta_n})$ with $M\sm B(x,r_0)$. We look more precisely at the arcs $C_0(r_0,\theta_n)$. From this sequence of compact connected subsets of $M\sm B(x,r_0)$, we can extract a subsequence which converges in the Hausdorff distance to some compact connected set $C_0$. Since any $z \in C_0$ is the limit of a sequence $z_n \in C_0(r_0,\theta_n)$, and since $w_{\theta_n}$ tends to $w_0$ uniformly on $M$, we conclude that $w_0(z) = 0$, i.e., that $C_0 \subset \cZ(w_0)$. The set $C_0$ contains the points $\exp_x(r_0\, \omegat_0(r_0,0))$ and $\exp_x(r_0\, \omegat_{\ell_0}(r_0,0))$. Since $C_0$ is connected and contained in $\cZ(w_0)$, and in view of the structure of $\cZ(w_0)$, we must have $\ell_0 = \tau(0)$, and we reach a contradiction. The proof of Property~\ref{propertyD4} is complete.\hfill \qedc \bigskip

\emph{Conclusion of the rotating function argument.}~  Assuming that \eqref{E-h2n-s2} holds, since $w_{\frac{\pi}{k-1}} = - v_1$, we infer from Property~\ref{propertyD4} that $\gamma^{\frac{\pi}{k-1}}_{0,\tau(0)}= \gamma_{1,\tau(0)+1}$. Since there is only one nodal semi-arc tangent to a given ray at $x$, we conclude that $\tau(0) \not = 1$ and $\tau(0) \not = 2k-3$. It follows that $0 < 1 < \tau(0)$, and that $\tau(0) < \tau(0)+1 = \tau(1) \le 2k-3$  (here we have used the assumption $k\ge 3$). This contradicts Property~\ref{propertyD3}, and proves that $\dim U(\lambda_k) = (2k-1)$ cannot occur. \hfill \qedc \medskip

We have proved the inequality
%
\begin{equation}\label{E-h2n-b2}
\mult(\lambda_k) \le (2k-2)\,, \text{~for any~} k\ge 3\,.
\end{equation}
\hfill \qed

\begin{remark}\label{R-h2n-bess}
As far as we know, the idea to consider the family of functions $w_{\theta}$ was introduced by Besson \cite{Bess1980}, proof of Theorem~3.C.1 in which he improves the upper bound for the multiplicity of the second eigenvalue of a torus from $7$ to $6$. A similar idea was used by Nadirashvili  \cite{Nadi1987}, p.~231 lines 1--8, for higher eigenvalues as well. It is used in \cite{HoHN1999, HoMN1999} also, and will appear several times, in one form or another, in the present paper.
\end{remark}%

\begin{remark}\label{R-h2n-f}
In the next section, we will introduce the \emph{combinatorial type} for different kinds of nodal sets on a compact surface $M$ with boundary, and we will repeatedly use the ``rotating function argument''.
\end{remark}%




\subsection{Surfaces with genus zero, proof that $\dim U(\lambda_k) \le (2k-3)$ for $k\ge 3$}\label{SS-h2n-s3}\phantom{}

 The proof is by contradiction. Taking Subsection~\ref{SS-h2n-s2} into account, we assume that $\dim U(\lambda_k) = (2k-2)$.

\begin{properties} \label{propertyD6} Assume that $\dim U(\lambda_k) = (2k-2)$.  Then, the linear subspace $U_x = \set{u \in U(\lambda_k) ~|~ \nu(u,x) \ge 2(k-1)}$ has the following properties.
\begin{enumerate}[(i)]
 \item $\dim U_x = 1$, and for any $0\not = u \in U_x$, we have
 \item $\nu(u,x) = 2(k-1)$, and $x$ is the only singular point of the eigenfunction $u$;
 \item $\cZ(u)$ is connected;
 \item $\kappa(u) = k$.
 \end{enumerate}
\end{properties}

\pf By Lemma~\ref{L-zeroi}, $\dim U_x \ge 1$. Given any $0 \not = u \in U_x$, Euler's formula gives
%
\begin{equation}\label{E-h2n-10}
0 \ge \kappa(u) - k = \left\lbrace b_0(\cZ(u)) - 1\right\rbrace + \sum_{\substack{z \in \cS(u)\\ z \not = x}} \frac{\nu(u,z) - 2}{2} + \lbrace \frac{\nu(u,x)}{2} - k + 1 \rbrace\,,
\end{equation}
%
and we conclude that (ii)--(iv) hold. To prove (i), assuming that $\dim U_x \ge 2$, we can repeat the arguments of Subsection~\ref{SS-h2n-s2}, and reach a contradiction. \hfill \qed \medskip

Given any $x \in M$, there exists a function $0 \neq u_x \in U$, uniquely defined up to multiplication by a nonzero scalar, such that $\nu(u_x,x) = 2(k-1)$. Choosing a basis $\set{u_1,\dots,u_{2k-2}}$ for the eigenspace $U(\lambda_k)$, we can write $u_x = \sum_{j=1}^{2k-2} \alpha_j(x) u_j$, and the condition $\ord(u_x,x)=(k-1)$ is equivalent to a linear system in the $(2k-2)$ unknowns $\alpha_j(x)$. This linear system as constant rank $1$. It follows that,   in a neighborhood of any point $x_0$, we can find a $C^{\infty}$ function $x \to (\alpha_1(x),\ldots,\alpha_{2k-2}(x))$, with values in $\R^{2k-2}\sm \set{0}$, and defined up to multiplication by a nonzero scalar $C^{\infty}$ function. These local solutions yield a unique global $C^{\infty}$ function with values in $\bP(U(\lambda_k))$, the projective space of the eigenspace $U(\lambda_k)$. Since $M$ is simply connected, this map can be lifted to a smooth map $x \mapsto u_x$, from $M$ to $\bS(U(\lambda_k))$, the unit sphere of $U(\lambda_k)$ (for example with respect to the $L^2$ norm). \medskip

 To each $x \in M$, we associate the homogeneous polynomial $p_x$ on $T_xM$, of degree $(k-1)$, defined by $p_x : T_xM \ni v \mapsto \frac{d^{k-1}}{dt^{k-1}} u_x(\exp_x(tv))$. It is harmonic with respect to the Riemannian metric $g_x$ in $T_xM$. The map $x \mapsto p_x$ is smooth. The restriction of this polynomial to the unit circle $S_xM$ in $(T_xM,g_x)$ has simple zeros. Choose some $x_0 \in M$, and some root $e_{x_0}$ of $p_{x_0}$ in $S_{x_0}M$. Given any $x_1 \in M$, and any curve $c$ from $x_0$ to $x_1$, we can follow this root by continuity along the curve $c$ so that $e_{c(t)}$ is a root of $p_{c(t)}$. Since $M$ is simply connected, the root $e_{c(1)}$ at $x_1$ does not depend on the choice of the curve $c$. It follows that we have defined a continuous unit vector-field $x \mapsto e_x$ on $M$, contradicting the Poincar\'e-Hopf theorem \cite[p.~35]{Miln1997} which asserts that the sum of the indices of the zeros of a vector-field on $M$ equals the Euler characteristic of $M$ (here, we have a vector-field without zero, and $\chi(M)=2$). We have proved that the assumption $\dim U(\lambda_k) = (2k-2)$ leads to a contradiction, and hence that $\dim U(\lambda_k) \le (2k-3)$. \hfill \qed

\section{Plane domains, bounding $\mult(\lambda_k)$ from above, preliminaries}\label{S-hmn}


\subsection{Introduction}\label{SS-hmn-i}

Let $\Omega$ be a regular bounded domain\footnote{By ``domain'', we mean a connected open subset.} in $\R^2$. We are interested in the eigenvalue problem \eqref{E-evp-2bc} for the Laplacian or for a Schr\"{o}dinger operator of the form $-\Delta + V$ in $\Omega$.  In \cite{HoMN1999}, T.~Hoffmann-Ostenhof, P.~Michor and N.~Nadirashvili state the following theorem, and mention that the result carries over to the Neumann boundary condition (see Theorem~A and comments, p. 1170).

\begin{theorem}\label{T-hmn}
For $k \ge 3$, the Dirichlet eigenvalues of  the operator $-\Delta + V$ in a bounded domain $\Omega \subset \R^2$ satisfy $\mult(\lambda_k) \le (2k-3)$.
\end{theorem}%

More precisely, they give a $3$-step proof. They first give an easy proof of the upper bound $(2k-1)$ (which is due to Nadirashvili \cite{Nadi1987} in the case of surfaces with genus $0$). This implies in particular that the second eigenvalue has multiplicity at most $3$.  In \cite[Section~2]{HoMN1999}, they improve this bound to $(2k-2)$ when $k\geq 3$. In \cite[Section~3]{HoMN1999}, they prove the bound $(2k-3)$ (same bound as in the case of closed surfaces with genus $0$, \cite{HoHN1999}).  However, in the non simply connected case, the validity of the proof of the bound $(2k-3)$  is questioned  in \cite{Berd2018}, (Introduction, line 5; Section~2, page 544, line 10; Section~4). Therefore,  the non simply connected case seems to remain open. \smallskip

In this section, we introduce some notation and  basic inequalities to be used in Sections~\ref{S-hmn}--\ref{S-hmn4}, and we prove that $\mult(\lambda_k) \le (2k-1)$  for any $k\ge 1$. In Section~\ref{S-hmn2}, we  will prove  that  $\mult(\lambda_k) \le (2k-2)$  for any $k\ge 3$.


\subsection{Notation}\label{SS-hmn-0N}

Let us fix some notation for  Sections~\ref{S-hmn}-- \ref{S-hmn4}.\smallskip

Let $U$ denote a linear subspace of an eigenspace of the eigenvalue problem \eqref{E-evp-2bc} in $\Omega$. We assume that $U$ satisfies the inequality
%
\begin{equation}\label{E-hmn0-2}
 \set{\kappa(u) ~|~ 0 \neq u \in U} \le \ell \,,
\text{~~for some integer~} \ell \ge 2\,.
\end{equation}
%
\medskip

We write $\partial \Omega$ as the union of its components, $\partial \Omega = \bigcup_{j=1}^q \Gamma_j$, with $q \ge 1$.\smallskip

Given $0 \neq u \in U$, define the sets
%
\begin{equation}\label{E-hmn0-4}
\left\{
\begin{array}{ll}
J(u) & := \set{j ~|~ \Gamma_j \cap \cZ(u) \neq \emptyset}\,,\\[5pt]
\Gamma(u) & := \cup_{j \in J(u)} \Gamma_j\,.
\end{array}
\right.
\end{equation}

Given a function $0 \neq u \in U$, we denote by
%
\begin{equation}\label{E-hmn0-5}
[u] := \set{a\, u ~|~ a \in \R\sm \set{0}}
\end{equation}
%
the corresponding line in the projective space $\bP(U)$. We will say that $u$ is a generator of the line $[u]$. If a function $u$ is uniquely determined by some condition, up to multiplication by a nonzero scalar, we will say that $u$ is uniquely determined \emph{up to scaling} or, equivalently, that $[u]$ is uniquely determined.


\subsection{The initial inequalities}\label{SS-hmn-0B}

We shall make an extensive use of Euler's formula for the nodal partition $\cD_u$ associated with an eigenfunction $u$, see Paragraph~\ref{SS-etf0}, taking into account the assumption \eqref{E-hmn0-2} on $U$,
%
\begin{equation}\label{E-euler-or}
\ell \ge \kappa(u) = 1 + \beta(u) + \sigma_i(u) + \sigma_b(u)\,,
\end{equation}
%
where,
%
\begin{equation}\label{E-hmn0-8}
\left\{
\begin{array}{ll}
\beta(u) & := b_0(\cZ(u) \cup \partial \Omega) - b_0(\partial \Omega) \\[5pt]
& ~= b_0\left(\cZ(u) \bigcup \Gamma(u)\right) - b_0(\Gamma(u))\,,
\end{array}
\right.
\end{equation}
%
\begin{equation}\label{E-hmn0-10}
\left\{
\begin{array}{ll}
\sigma_i(u) &= \frac 12\, \sum_{z \in \cS_i(u)} \left( \nu(u,z) - 2 \right)\,,\\[5pt]
\sigma_b(u) &= \frac 12\, \sum_{z \in \cS_b(u)} \rho(u,z) = \sum_{j \in J(u)} \, \frac 12 \, \sum_{z\in \cS_b(u) \cap \Gamma_j\,} \rho(u,z)\,.\\[5pt]
\end{array}
\right.
\end{equation}

From the nodal character of the partition $\cD_u$, we also have, using Proposition~\ref{P-euler-np} and the definition of $J(u)$,
%
\begin{equation}\label{E-hmn0-12}
\forall j \in J(u), \quad \sum_{z \in \cS_b(u)\cap \Gamma_j}\rho(u,z) \text{~~is even and~} \ge 2\,.
\end{equation}

We now rewrite  the Euler inequality \eqref{E-euler-or} in the form,
%
\begin{equation}\label{E-euler-or2}
\left\{
\begin{array}{ll}
0 \geq  \kappa(u) - \ell  \,  = &\left( b_0(\cZ(u) \cup \Gamma(u))-1\right)+ \frac 12 \sum_{z \in \cS_i(u)} (\nu(u,z)-2)\\[5pt]
& + \sum_{j \in J(u)\,} \frac 12 \, \left( \sum_{z\in \cS_b(u) \cap
\Gamma_j\,} \rho(u,z) \, - 2 \right) - (\ell - 2)\,.
\end{array}
\right.
\end{equation}

We first observe that the first three terms in the right-hand side of the equality in \eqref{E-euler-or2} are nonnegative. We apply this inequality to  eigenfunctions with prescribed singular points.\medskip

Fix some $x \in \Gamma_1$, and let $m := \dim U$. By Lemma~\ref{L-zero1}, there exists $0 \neq u \in U$ such that $\rho(u,x) \ge (m-1)$. Rewrite \eqref{E-euler-or2} as,
%
\begin{equation}\label{E-euler-or2a}
\left\{
\begin{array}{ll}
0 \geq  \kappa(u) - \ell  \,  & = \left( b_0(\cZ(u) \cup \Gamma(u))-1\right)+ \frac 12 \sum_{z \in \cS_i(u)} (\nu(z)-2)\\[5pt]
& ~~ + \sum_{j \in J(u), j\not = 1\,}  \frac 12 \, \left( \sum_{z\in \cS_b(u) \cap
\Gamma_j\,} \rho(u,z) \, - 2 \right)\\[5pt]
& ~~ + \frac 12 \, \sum_{z\in \cS_b(u) \cap \Gamma_1\,} \rho(u,z) \, - \ell + 1. \end{array}
\right.
\end{equation}
%
The first three terms in the right-hand side of the equality are nonnegative. It follows that
%
\begin{equation*}
2\ell - 2 \ge \sum_{z\in \cS_b(u) \cap \Gamma_1\,} \rho(u,z) \ge m-1\,,
\end{equation*}
%
so that
%
\begin{equation}\label{E-hmn0-14}
\dim U \le (2\ell - 1).
\end{equation}

Choosing $U = U(\lambda_k)$, inequality \eqref{E-hmn0-14} yields the following result.


\begin{proposition}[\cite{Nadi1987}, Theorem~2]\label{P-hmn-s1}
Let $\Omega \subset \R^2$ be a bounded domain with smooth boundary. Let $\set{\lambda_k, k\ge 1}$ be the eigenvalues of the operator $-\Delta + V$ in $\Omega$, with Dirichlet or Robin condition. Then, for any $k \ge 1$,
\begin{equation*}
\mult(\lambda_k) \le (2k-1)\,.
\end{equation*}
\end{proposition}%

In view of Proposition~\ref{P-hmn-s1}, in order to prove Theorem~\ref{T-hmn}, it suffices to show that the equalities $\dim U(\lambda_k) = (2k-1)$ and $\dim U(\lambda_k) = (2k-2)$ cannot occur for $k\ge 3$. This is the purpose of Sections~\ref{S-hmn2} and \ref{S-hmn4} in which we revisit the arguments of \cite{HoMN1999} for the three boundary conditions \eqref{E-evp-bc}.\smallskip


Before proceeding, we examine the case of the second eigenvalue ($k=2$) which is of special interest. This is the purpose of the next subsection.


\subsection{Multiplicity of the second eigenvalue and the nodal line conjecture}\label{SS-hmn-0C}%~\phantom{}

When $k=2$, Proposition~\ref{P-hmn-s1} gives $\mult(\lambda_2) \le 3$.  A natural question, in view of Table~\ref{E-scr-2T}, is whether this bound is sharp (depending on the boundary condition). As we shall see, this question is related to the so-called  ``nodal line conjecture''.

\subsubsection{Nodal sets of second eigenfunctions}\label{SSS-hmn-0CA}

 We use the notation of Subsection~\ref{SS-hmn-0N}, and write $\partial \Omega = \bigcup_{j=1}^q \Gamma_j$, with $q \ge 1$. Let $u \in U(\lambda_2)$ be any second eigenfunction. By Courant's theorem, $u$ has exactly two nodal domains. {Euler's formula \eqref{E-euler-or2}  yields
%
\begin{equation}\label{E-hmn0-22}
\left\{
\begin{array}{ll}
0 = \kappa(u) - 2 = & \left[ b_0(\cZ(u) \cup \Gamma(u)) - 1 \right]  + \frac 12
\sum_{z \in \cS_i(u)} (\nu(u,z)-2)\\[5pt]
& + \sum_{j \in J(u)}\frac 12 \, \left( \sum_{z\in \cS_b(u) \cap \Gamma_j\,} \rho(u,z) - 2\right)\,.
\end{array}%
\right.
\end{equation}

 Since all the terms is the right hand side are nonnegative (use Corollary~\ref{cor:nodinfo} and the definition of $J(u)$), we immediately deduce that
%
\begin{equation}\label{E-hmn0-24}
\left\{
\begin{array}{l}
\cZ(u) \cup \Gamma(u) \text{~is connected}\,,\\[5pt]
\sum_{z \in \cS_i(u)} (\nu(u,z)-2) = 0 \text{~~i.e,~} \cS_i(u) = \emptyset\,,\\[5pt]
\sum_{z\in \cS_b(u) \cap \Gamma_j\,} \rho(u,z) = 2 ~~\forall j \in J(u)\,.
\end{array}%
\right.
\end{equation}

 The structure of $\cZ(u)$ depends on whether $J(u) = \emptyset$ or $J(u) \neq \emptyset$. We now consider the two simplest situations. The proofs of the following properties are clear.

\begin{property}\label{P-hmn-nl0}
Assume that $\Omega$ is \emph{simply connected}. Let $u$ be a second eigenfunction. Then, either $\cZ(u)$ does not hit $\partial \Omega$ and $\cS(u) = \emptyset$, or $\cZ(u)$ hits $\partial \Omega$, $\cS_i(u) = \emptyset$, and $\sum_{z\in \cS_b(u)} \rho(u,z) = 2$. More precisely, there are three distinct possibilities.
\begin{enumerate}[(a)]
  \item If $J(u) = \emptyset$, then $\cZ(u)$ is a nodal circle, ie a simple closed regular  connected curve contained in $\Omega$, not touching $\partial \Omega$. This case is characterized by the fact that the function $\breve{u}$  defined in \eqref{E-evp-dr} does not vanish on $\partial \Omega$.
  \item If $J(u) = \set{1}$ and $\cS_b(u) = \set{y}$ for some $y \in  \partial \Omega$ with $\rho(u,y)=2$, then $\cZ(u)$ is a nodal loop at $y$, ie $\cZ(u) \sm \set{y}$ a simple regular connected curve contained in $\Omega$.  This case is characterized by the fact that the function $\breve{u}$ vanishes only at $y_1$ on $\partial \Omega$, and does not change sign.
  \item If $J(u) = \set{1}$ and $\cS_b(u) = \set{y_1,y_2}$ for some $y_1 \neq y_2 \in  \partial \Omega$ with $\rho(u,y_1)=1$, $\rho(u,y_2)=1$, then $\cZ(u)$ is a nodal arc from $y_1$ to $y_2$, ie $\cZ(u) \sm \set{y_1,y_2}$ is a simple regular connected arc contained in $\Omega$.  This case is characterized by the fact that the function $\breve{u}$ vanishes precisely at $y_1$ and $y_2$ on $\partial \Omega$, and changes sign at these points.
\end{enumerate}
\end{property}%


\begin{property}\label{P-hmn-nl1}
Assume that $\Omega$ has \emph{one hole}, and write $\partial \Omega = \Gamma_1 \cup \Gamma_2$. Let $u$ be a second eigenfunction. Then, either $\cZ(u)$ does not hit $\partial \Omega$ and $\cS(u) = \emptyset$, or $\cZ(u)$ hits $\partial \Omega$, $\cS_i(u) = \emptyset$, and $\sum_{z\in \cS_b(u)\cap \Gamma_j} \rho(u,z) = 2$ for all $j \in J(u)$. More precisely, there are three distinct cases (up to relabeling the two components of $\partial \Omega$).
\begin{enumerate}[(1)]
  \item If $J(u) = \emptyset$, then $\cZ(u)$ is a nodal circle contained in $\Omega$, not touching $\partial \Omega$. This case is characterized by the fact that the function $\breve{u}$  defined in \eqref{E-evp-dr} does not vanish on $\partial \Omega$.
  \item If $J(u) = \set{1}$, then either $\cS_b(u) = \set{y}$ for some $y \in  \Gamma_1$ with $\rho(u,y_1)=2$, or $\cS_b(u) = \set{y_1,y_2}$ for some $y_1 \neq y_2 \in  \Gamma_1$ with $\rho(u,y_1)=\rho(u,y_2)=1$. We have either a nodal loop at $y$, or a nodal arc from $y_1$ to $y_2$.  This case is characterized by the fact that the function $\breve{u}$ vanishes only at $y$ (without changing sign along $\Gamma_1$), or vanishes at $y_1$ and $y_2$ (and changes sign along $\Gamma_1$). In both subcases, $\breve{u}$ does not vanish on $\Gamma_2$.
  \item If $J(u) = \set{1,2}$,  then $\cZ(u)$ hits both component $\Gamma_1$ and $\Gamma_2$ at one point with index $2$, or at two distinct points of index $1$. Furthermore, the components $\Gamma_1$ and $\Gamma_2$ are linked by two nodal arcs (possibly with one or two common boundary points).
\end{enumerate}
\end{property}%

\begin{remark}\label{R-hmn0-24}\phantom{}
It is not clear  a priori whether the possible nodal patterns described in Property~\ref{P-hmn-nl0} or \ref{P-hmn-nl1} are actually realized for some choice of domain $\Omega$ and potential $V$ (when $\Omega$ is convex and $V\equiv 0$, see \cite{Ales1994}). Applying Lemmas~\ref{L-zero1} or \ref{L-zero2}, one can at least prescribe one or two boundary singular points.
\begin{enumerate}[(i)]
\item Assume that $\dim U(\lambda_2(-\Delta+V)) \ge 3$. If $\Omega$ is simply connected, then there exists an eigenfunction whose nodal set satisfies (b), resp. (c), in Properties~\ref{P-hmn-nl0}. If $\Omega$ has one hole, then there exists an eigenfunction whose nodal set hits both $\Gamma_1$ and $\Gamma_2$.
  \item Assume that $\dim U(\lambda_2(-\Delta+V)) = 2$. Then, there exists an eigenfunction whose nodal domain hits $\partial \Omega$.
\end{enumerate}
\end{remark}%

Nodal sets and $\mult(\lambda_2)$ are known precisely in few circumstances only, either in very specific cases, or under additional assumptions on the domain (some convexity or symmetry conditions, see \cite{Shen1988}, \cite{Putt1990}, \cite{Putt1991} in the simply connected case, and \cite{Kiwa2018} for a convex domain with a convex sub-domain removed).



The following figures display some particular cases. The second (Dirichlet or Neumann) eigenvalue of an equilateral triangle with rounded corners has multiplicity two, with one symmetric and one antisymmetric eigenfunction.

The nodal domains of the symmetric eigenfunction appear in Figure~\ref{F-hmn1-nl-teeh} (left), see \cite{BeHe2021t}. The second (Dirichlet or Neumann) eigenvalue of an ellipse is simple with nodal domains as in Figure~\ref{F-hmn1-nl-teeh} (center); this is a particular case of the domains described in \cite{Shen1988}, \cite{Putt1990} and \cite{Putt1991}. The nodal set of a second Dirichlet eigenvalue of $D\sm B$, where $D, B$ are convex symmetric domains, has been studied in \cite{Kiwa2018}, see Figure~\ref{F-hmn1-nl-teeh} (right).


\begin{figure}[!ht]
\centering
\includegraphics[width=0.9\linewidth]{hmn1-nl-teeh}
%
\caption{Nodal patterns of second eigenfunctions}
\label{F-hmn1-nl-teeh}
\end{figure}

Numerical computations, playing with the position of the holes, give rise to
some other patterns, see Figure~\ref{F-hmn1-nl-eh}.

\begin{figure}[!ht]
\centering
\includegraphics[width=0.99\linewidth]{hmn1-nl-eh}
%
\caption{Nodal patterns of second eigenfunctions}
\label{F-hmn1-nl-eh}
\end{figure}

In view of the above remarks, the following questions are natural.

\begin{description}
  \item[Question~1] Does there exist a second eigenfunction of $-\Delta+V$ whose nodal set is a nodal circle ?
  \item[Question~2] Does there exist a second eigenfunction of $-\Delta+V$ whose nodal set is a nodal loop at some $y \in \partial \Omega$ ?
\end{description}\medskip


\subsubsection{The nodal line conjecture}\label{SSS-hmn-0CB}

For a simply connected domain $\Omega$, Pleijel \cite[p.~546]{Plej1956} observed that a second Neumann eigenfunction of $-\Delta$  cannot have a closed nodal line. The idea is to use the inequality $\lambda_2(D,-\Delta,\mf{n}) \le \lambda_1(D,-\Delta,\mf{d})$ between the second Neumann eigenvalue and the first Dirichlet eigenvalue of $-\Delta$ in a domain $D \subset \R^2$, and the monotonicity property of the first Dirichlet eigenvalue. Indeed, if there exists a second Neumann eigenfunction $u_2$ one of whose two nodal domains (say $D_1$) has the property that $\partial D_1 \cap \partial \Omega$ consists of isolated points  (which occurs in particular when $\partial D_1 \subset \Omega$), then the restriction of $u_2$ to $D_1$ is the ground state of the Dirichlet problem in $D_1$, $\lambda_1(D_1,\mf{d}) = \lambda_2(\Omega,\mf{n})$. By the strict domain monotonicity of the Dirichlet eigenvalues, we get
%
$$
\lambda_1(\Omega,\mf{d}) <  \lambda_1(D_1,\mf{d})=\lambda_2(\Omega,\mf{n})\,,
$$
hence a contradiction. This argument may fail when $\Omega$ as a hole because both nodal domains of $u$ may touch the boundary.\smallskip

The inequality $\lambda_2(D,-\Delta,\mf{n}) \le \lambda_1(D,-\Delta,\mf{d})$ is due to Szeg\"{o} (1954) when $\Omega \subset \R^2$ is simply connected and smooth.  The strict inequality is proved in an earlier paper of P\'{o}lya (1952) which does apparently not use the assumption that the domain is simply connected. It was later generalized to higher dimensions, and smooth enough domains (not necessarily simply connected) by Weinberger (1956), see \cite{Payn1967}, Theorem~3, p.~463. It has been extended to domains with $C^1$ boundary, see \cite{Maz1991} in which Mazzeo revisits the earlier paper of L.~Friedlander \cite{Frie1991}.  When $\partial D_1$ meets $\partial \Omega$ we actually need the inequality to hold for domains with piecewise $C^1$ boundary   (see \cite{ArMa2007,ArMa2012} for the Lipschitz case).
\smallskip

In view of Pleijel's observation, Payne conjectured that a second Dirichlet eigenfunction of $-\Delta$ cannot have a closed nodal line, see \cite{Payn1967}, Conjecture~5, p.~467. One can make a similar conjecture for second Robin eigenfunctions,  and also consider domains in $\R^n$, $n\ge 3$, see \cite{Four2001}, \cite{Ken2013} and their bibliographies.

\begin{remark}\label{R-hmn0-26}
As observed in \cite{Liq1995} (end of Section~2, p.~277), if a simply connected domain satisfies the \emph{nodal line conjecture}, then $\dim U(\lambda_2) \le 2$. This is an immediate consequence  of Remark~\ref{R-hmn0-24}(ii).
\end{remark}

We now consider the three boundary conditions, Dirichlet, Neumann and Robin separately.

\subsubsection{Dirichlet boundary condition}\label{SSS-hmn-0CD}

The following results are due to Lin and Ni.
\begin{enumerate}[$\diamond$]
\item \cite[Theorem~3.6]{LiNi1988}: For all $n \ge 2$, there exists a radius $R_n$ and a nonzero smooth radial potential $V_n$, such that
      %
\begin{equation*}
  \mult(\lambda_2; B(R_n), -\Delta + V_n,\mf{d}) = 1\,,
\end{equation*}
and the corresponding eigenfunction is radial, with nodal set a sphere in $B(R_n)$. Here, $B(R_n)$ is the ball of radius $R_n$ in $\R^n$.
\item \cite[Theorem~3.8]{LiNi1988}: For all $n \ge 2$, there exists a radius $R_n$ and a nonzero smooth radial potential $V_n$, such that
      %
\begin{equation*}
  \mult(\lambda_2; B(R_n), -\Delta + V_n,\mf{d}) = (n+1)\,,
\end{equation*}
and there exists a radial second eigenfunction.
\end{enumerate}

In dimension $2$, the second assertion implies that the bound of the multiplicity $\mult(\lambda_2;\Omega,-\Delta+V,\mf{d}) \le 3$ is sharp.\medskip


The following results are due to M. and T. Hoffmann-Ostenhof and Nadirashvili \cite{HoHN1997, HoHN1998} (see  for the second statement   \cite{HeHJ2020} for corrections and complements).
\begin{enumerate}[$\diamond$]
\item \cite[Theorem~2.1]{HoHN1998} There exists $N_0$ and domains $D_{N,\varepsilon} \subset \R^2$ such that for all $N \ge N_0$, and $\varepsilon$ small enough, $\lambda_2(D_{N,\varepsilon};-\Delta,\mf{d})$ is simple, with a closed nodal set contained in $D_{N,\varepsilon}$. The domain $D_{N,\varepsilon}$ is homeomorphic to a disk minus $N$ points.
  \item \cite[Theorem~2.2]{HoHN1998}: For all $N \ge 3$, and $\varepsilon$ small enough, the domains $D_{N,\varepsilon} \subset \R^2$  satisfy
      %
\begin{equation*}
  \mult(\lambda_2; D_{N,\varepsilon}, -\Delta,\mf{d}) = 3\,.
\end{equation*}
\end{enumerate}

The second assertion implies the the upper bound $3$ for the second Dirichlet eigenvalue of $-\Delta$ is sharp for non simply connected domains.\medskip

See \cite{DaGH2021} for a counterexample to the nodal line conjecture (for $-\Delta$) with six holes. However, counter-examples are still missing for domains with one or two holes.\smallskip

 The nodal line conjecture for $(\Omega,-\Delta)$ is known to be true when $\Omega$ is a bounded \emph{convex domain} in $\R^2$: \cite{Payn1973} and \cite{Lin1987} (under additional symmetry assumptions), \cite{Mela1992} (smooth convex domains) and \cite{Ales1994} (general convex domains).  This is also the case for domains which are convex in one direction only\footnote{See \cite[Corollary~2.7]{Liq1995}. Note however that the other results on the multiplicity presented in this paper are true under strong additional conditions only.  We thank the author for clarifying this point.}\medskip

As a by-product of these results, we have the upper bound
%
\begin{equation*}
\mult(\lambda_2;\Omega,-\Delta,\mf{d}) \le 2 \text{~~for any convex bounded domain~} \Omega\,.
\end{equation*}

This bound holds for domains which are convex in one direction and for domains which satisfy the nodal line conjecture (see Remark~\ref{R-hmn0-26}). This result supports the following conjecture.


\begin{conjecture}\label{C-aaa}
\begin{equation*}
\mult(\lambda_2;\Omega,-\Delta,\mf{d}) \le 2 \text{~~for any \emph{simply connected} bounded domain~} \Omega\,.
\end{equation*}
\end{conjecture}%



\subsubsection{Neumann boundary condition}\label{SSS-hmn-0CN}
The discussion in Paragraph~\ref{SSS-hmn-0CB} shows that the nodal line conjecture holds for the Neumann Laplacian in any simply connected  bounded regular domain in $\R^2$. Nadirashvili proved that the multiplicity of the second eigenvalue of a simply connected domain with nonpositive curvature is at most $2$ and that this estimate is sharp, see \cite{Nadi1987}, Theorem~2 and Corollary~1. As a matter of fact, his proof also shows that the nodal line conjecture is true in such domains (for the Neumann condition).

\subsubsection{Robin boundary condition}\label{SSS-hmn-0CR}
As observed by J.B.~Kennedy\footnote{We thank J.B. Kennedy for
useful discussions around this problem.} in \cite{Ken2011}, the proof of the nodal line conjecture for the $h$-Robin boundary condition works
in the same way as in the case of Neumann conditions provided that the following
inequality holds,
%
$$
\lambda_2 (h,\Omega) \leq  \lambda_1^D(\Omega)\,.
$$
%
Observing the monotonicity of the Robin problem with respect to $h$,
we obtain the existence of some $h_\Omega >0$ such that this inequality
holds for $h \leq  h_\Omega$.\smallskip

J.B. Kennedy also shows that, as in the Dirichlet case (which corresponds to $h=+\infty$), one can find examples of multiply connected domains for which counter examples to the nodal line conjecture can be constructed. One can also expect to construct examples for which the multiplicity is $3$
(as in  \cite{HoHN1998} and \cite{HeHJ2020})  but this is still open
at the moment.\smallskip

On the positive side, it is natural to ask if convexity is enough to ensure multiplicity at most $2$ for every Robin parameter $h>0$, following Lin's approach in \cite{Lin1987}. This is still open at the moment. What we do know is that a sufficient  condition on $\Omega$  is that nodal line conjecture holds (Remark~\ref{R-hmn0-26}).


\section{Plane domains, the estimate $\mult(\lambda_k) \le (2k-2)$ for $k\geq 3$}\label{S-hmn2}

The ultimate goal of this section is to prove that  the inequality $\mult(\lambda_k) \le (2k-2)$ holds for $k \ge 3$, see Proposition~\ref{P-hmn-s2}. Note that the inequality is not true for $k=1$ and $k=2$. For this purpose, we provide detailed proofs of the statements in \cite[Section~2]{HoMN1999}. The general idea is to prove that an a priori upper bound on the number of nodal domains of functions in a given subspace $U$ of eigenfunctions implies an upper bound on $\dim U$. Indeed, the bigger $\dim U$, the easier to construct eigenfunctions with prescribed high order singular points  and, by Euler's formula, with more nodal domains.\smallskip


We first introduce an abstract setting,  to be used when $\Omega$ is multiply connected, and some notation to describe the nodal sets in a neighborhood of a boundary singular point. The following  subsections are devoted to studying eigenfunctions with prescribed boundary singular points, and drawing consequences.


\subsection{An abstract setting}\label{SS-hmn2-10}

When the domain $\Omega \subset \R^2$ is not simply connected, its boundary $\partial \Omega$ has  $(q+1)$ components, with $q\geq 1$.  One of them $\Gamma_1$, the ``outer boundary'', bounds the unbounded component of $\R^2\sm \partial \Omega$. The other components, $\Gamma_j, ~j\neq 1$, are contained in the bounded component of $\R^2\sm \Gamma_1$. We consider the following equivalence relation in $\wb{\Omega}$:
%
\begin{equation}\label{E-hmn210-equiv}
x \sim y \text{~~if and only if~~} x, y \in \Gamma_j
\text{~for some~} j \in \set{2,\ldots,(q+1)}\,.
\end{equation}


\begin{notation}\label{N-hmn210-2}
Let $\check{\Omega}$ denote the quotient space $\wb{\Omega}/\!\!\sim$, where each $\Gamma_j, ~j\neq 1,$ is identified to one point $\xi_j$ in $\check{\Omega}$ (see \cite{Bona2009}). Define
$$\Xi := \set{\xi_2,\ldots,\xi_{q+1}}\,.$$
Generally speaking, $\check{A}$ will denote the image of the set $A \subset \wb{\Omega}$ under the projection map from $\wb{\Omega}$ to $\check{\Omega}$.\\
We also introduce $\bS_{\Omega}$, the quotient space in which each $\Gamma_j, j\ge 1$, is identified to a point.
\end{notation}%



\subsection{Describing nodal sets near a boundary singular point}\label{SS-hmn2-20}

In the plane $\R^2$ with coordinates $(\xi_1,\xi_2)$, define the sets
%
\begin{equation}\label{E-hmn-21b4}
\begin{array}{ll}
D_{+}(a) &= \set{(\xi_1,\xi_2) \in \R^2 ~|~ \xi_1^2 + \xi_2^2 < a \text{~and~} \xi_2 > 0}\,,\\[5pt]
\wb{D}_{+}(a) &= \set{(\xi_1,\xi_2) \in \R^2 ~|~ \xi_1^2 + \xi_2^2 \le a \text{~and~} \xi_2 \ge 0}\,,\\[5pt]
S_{+}(a) &= \set{(\xi_1,\xi_2) \in \R^2 ~|~ \xi_1^2 + \xi_2^2 = a \text{~and~} \xi_2 \ge 0}\,.
\end{array}%
\end{equation}


Fix $x \in \partial\Omega$, and choose a direct orthonormal frame $\set{\vec{e}_1,\vec{e}_2}$ at $x$ such that $\vec{e}_1$ is tangent to $\partial \Omega$ at $x$, and $\vec{e}_2$ is normal to $\partial \Omega$, pointing inwards. The vector $\vec{e}_1$ determines the orientation of $\partial \Omega$.\smallskip

Let $u$ be an eigenfunction with boundary singular point $x$, and whose index at the point $x$ is $\rho(u,x)$ (the number of nodal arcs hitting the boundary at this point). In order to describe the nodal set of the eigenfunction $u$ near the point $x$, we choose conformal coordinates $(\cU, E)$, where $\cU$ is a neighborhood of $x$ in $\wb{\Omega}$, and $E$  a $C^{\infty}$ conformal map from $\wb{D}_{+}(1)$ onto $\cU$, such that $E(0) = x$ (see \cite{YaZh2021}, Section~2).\smallskip

The equation $(-\Delta + V) u = \lambda\, u$ in $\cU$ is transformed into an equation of the form $(-\Delta + \tilde{V}) \tilde{u} = \lambda\, \tilde{A}\, \tilde{u}$, where $\tilde{u} = u \circ E$, $\tilde{V}$ is a real valued $C^{\infty}$ potential, and $\tilde{A}$ a positive $C^{\infty}$ function. Furthermore, $\tilde{u}$ satisfies the Dirichlet or a Robin type boundary condition on $\wb{D}_{+}(1)\cap \set{\xi_2=0}$.  The function $u$ does not vanish at infinite order at any point, see \cite{DoFe1990a}.\smallskip

Denote by $(r,\omega)$ the polar coordinates associated with $(\xi_1,\xi_2)$. If $\ord(u,x) = p$, then $\ord(\tilde{u},0) = p$ as well. Applying the local structure theorem to $\tilde{u}$ at the point $0$, and for $ r \le r_0$ with $r_0 < 1$ small enough, we can write
%
\begin{equation}\label{E-hmn-21b5}
\left\{
\begin{array}{ll}
\tilde{u}(r,\omega) = a_{u,x} r^p \, \sin(p\omega) + \cO(r^{p+1}) &\text{(Dirichlet case)},\\[5pt]
\tilde{u}(r,\omega) = a_{u,x} r^p \, \cos(p\omega) + \cO(r^{p+1}) &\text{(Robin case)},
\end{array}
\right.
\end{equation}
%
where $a_{u,x}$ is a nonzero scalar. The important point is that near the point $0$, the nodal set of $\tilde{u}$ consists of $\rho(u_1,0) = (p-1)$ (Dirichlet case), resp. $\rho(u_1,0) = p$ (Robin case), nodal semi-arcs emanating from $0$, tangentially to rays $\tilde{\omega}_j, 1 \le j \le \rho(u_1,0)$ at $0$. These rays divide the half-circle $S_{+}(r_0)$ into equal arcs. More details are given in Appendix~\ref{S-lsc}. \smallskip

Define the sets
%
\begin{equation}\label{E-hmn-21b6}
\begin{array}{l}
B_{+}(x,r) := E(D_{+}(r))\,,\\[5pt]
\wb{B}_{+}(x,r) := E(\wb{D}_{+}(r))\,,\\[5pt]
C_{+}(x,r) := E(S_{+}(r))\,,\\[5pt]
L_m := \set{1, \ldots, m} \text{~for~} m \ge 1\,.
\end{array}
\end{equation}
%
Since $\cZ(u) \cap \wb{B}_{+}(x,r_0) = E\left( \cZ(u_1) \cap
\wb{D}_{+}(r_0)\right)$, $\cZ(u) \cap \wb{B}_{+}(x,r_0)$ consists of $\rho(u,x)$ semi-arcs emanating from $x$, tangentially to the rays $\omega_j, j \in L_{\rho(u,x)}$. The rays are labeled  counter-clockwise, according to their angle with respect to the vector $\vec{e}_1$. The nodal semi-arc emanating from $x$ tangentially to the ray $\omega_j$ is denoted by $\delta_j$.\smallskip


\subsection{Analysis of eigenfunctions with two prescribed  boundary singular points}\label{SS-hmn-21}

We use the notation of Subsection~\ref{SS-hmn-0N}.  In this subsection, we assume that $U$ is a linear subspace of an eigenspace $U(\lambda)$ of \eqref{E-evp-2bc}, and that for some $\ell \ge 2$,
%
\begin{equation*}
\left\{
\begin{array}{l}
\sup\set{\kappa(u) ~|~ 0 \neq u \in U} \le \ell \text{~~and}\\[5pt]
\dim U = (2\ell - 1)\,.
\end{array}%
\right.
\end{equation*}

For $x \neq y \in \Gamma_1$, we introduce the subspace
%
\begin{equation*}%\label{E-hmn-210}
U_{x,y} := \set{u \in U ~|~ \rho(u,x) \ge (2\ell-3) \text{~and~} \rho(u,y) \ge 1}\,.
\end{equation*}

According to Lemma~\ref{L-zero2},  $U_{x,y} \neq \set{0}$. The purpose of this subsection is to investigate the properties of the functions $u \in U_{x,y}$ --precise order of vanishing, structure of their nodal sets-- under the above assumptions on $U$.

\subsubsection{ $\Omega$ simply connected. Properties of $U_{x,y}$}\label{SSS-hmn-21a}

\begin{lemma}\label{L-Uxy0}
Assume that $\Omega$ is simply connected. Let $U$ be a linear subspace of an eigenspace of \eqref{E-evp-2bc} in $\Omega$,  such that $\sup\set{\kappa(u) ~|~ 0 \not = u \in U}\le \ell$ for some $\ell \ge 2$, and $\dim U = (2\ell-1)$. Let $x \neq y \in \Gamma_1$, and define
$$U_{x,y} := \set{u \in U ~|~ \rho(u,x) \ge (2\ell - 3) \text{~and~} \rho(u,y) \ge 1}.$$ Then,
\begin{enumerate}[(i)]
  \item $\dim U_{x,y} = 1$ and, for any $0 \neq u \in U_{x,y}$\,;
  \item %$u$ does not have any other singular point, whether in the
           %interior of $\Omega$ or on the boundary $\partial\Omega$,
            $\cS_i(u) = \emptyset$ and $\cS_b(u) = \set{x,y}$\,;
  \item $\rho(u,x)=(2\ell-3)$ and $\rho(u,y) = 1$;
  \item $\kappa(u) = \ell$\,;
  \item the set $\cZ(u) \cup \partial \Omega$ is connected.
\end{enumerate}
A generator of $U_{x,y}$ will be denoted by $u_{x,y}$ (defined up to scaling).
\end{lemma}%

\pf For simplicity, in the proof, we write $\nu(z)$ for $\nu(u,z)$, \ldots.\smallskip

The fact that $\dim U_{x,y} \ge 1$ follows from Lemma~\ref{L-zero2}.
In view of our assumptions, for any $0 \neq u \in U$, Euler's formula \eqref{E-euler-or2a} gives,
%
\begin{equation}\label{L-A5}
\begin{array}{ll}
0 \geq  \kappa(u) - \ell  \,  = &\left( b_0(\cZ(u) \cup \partial \Omega)-1\right)+ \frac 12 \sum_{z \in \cS_i(u)} (\nu(z)-2)\\[5pt] &\qquad +
\frac 12 \sum_{z\in \cS_b(u),\, z\neq x,y\,}\rho (z)  + \frac 12 \left( \rho(x) + \rho(y) -2\ell +2 \right)\,.
 \end{array}
\end{equation}
%
Each term in the right-hand side of the equality being nonnegative, the inequality implies that each term is zero, thus proving Assertions~(ii)--(v). \smallskip

To prove the first assertion, assume that there exist two linearly independent functions $u_1$ and $u_2$ in $U$. By Assertion~(ii) they both satisfy $\rho(u_i,x) = (2\ell - 3)$ and $\rho(u_i,y) = 1$. Apply Lemma~\ref{L-zeroc} at the point $y$ to find a nontrivial linear combination $\tilde u$ of $u_1$ and $u_2$ such that $\rho(\tilde u,y) \ge 2$ contradicting Assertion~(ii).\hfill \qed

\begin{figure}[!ht]
\centering
\begin{subfigure}[t]{.30\textwidth}
\centering
\includegraphics[width=\linewidth]{hmn2-Uxy0-a}
\caption{}
\end{subfigure}
\begin{subfigure}[t]{.30\textwidth}
\centering
\includegraphics[width=\linewidth]{hmn2-Uxy0-b}
\caption{}
\end{subfigure}
\begin{subfigure}[t]{.30\textwidth}
\centering
\includegraphics[width=\linewidth]{hmn2-Uxy0-c}
\caption{}
\end{subfigure}
%
\caption{$\Omega$ simply connected: some possible nodal patterns for $u_{x,y}$}\label{F-hmn2-Uxy0}
\end{figure}


\begin{remark}\label{R-Fgene}
The \emph{nodal patterns} displayed in Figure~\ref{F-hmn2-Uxy0} are valid for both the Dirichlet and Robin boundary conditions. Unless otherwise stated this remark applies to all figures of this section.
\end{remark}

\begin{definition}\label{D-hmn-np}
By a \emph{nodal pattern}, we mean a nodal set, up to  continuous deformations under which singular points may move, but neither appear nor disappear (singular points  occur when two nodal arcs cross, or when a nodal arc hits the boundary).
\end{definition}%

\subsubsection{$\Omega$ simply connected. Structure and combinatorial type of nodal sets in $U_{x,y}$}\label{SSS-hmn-21as}

Fix $x \neq y \in \partial \Omega$.  Under the assumptions of Lemma~\ref{L-Uxy0}, the nodal set of an eigenfunction $u \in U_{x,y}$ can be described as follows, using the notation of Subsection~\ref{SS-hmn2-20}.\smallskip

For $r_0$ small enough, in the neighborhood $B_{+}(x,r_0)$ of $x$ the nodal set $\cZ(u)$  consists of $(2\ell - 3)$ nodal semi-arcs $\delta_j$ emanating from $x$, tangentially to the rays $\omega_j, 1 \le j \le (2\ell - 3)$. Choosing any $j$, we follow the nodal semi-arc $\delta_j$ along $Z(u)$, until we reach a singular point of $u$. Otherwise stated, we consider the connected component of $\cZ(u)) \sm \cS(u)$ which contains the semi-arc $\delta_j$ (by abuse of notation we also denote this connected component by $\delta_j$). This is a nodal arc one of whose end points is $x$. Since $\cS(u) = \cS_b(u) = \set{x,y}$, the other end point is either $x$ or $y$.
More precisely, in view of the general properties of nodal sets, we define a map
%
\begin{equation}\label{E-hmn-nt}
\tau^{u}_{x,y} : \set{\downarrow} \cup L_{(2k-3)} \to \set{\downarrow} \cup L_{(2k-3)}
\end{equation}
%
as follows.
\begin{enumerate}[(i)]
  \item There exists a unique element $a \in L_{(2k-3)}$ such that starting from $x$ along $\delta_{a}$, we reach the boundary $\partial \Omega$ at $y$. We let $\tau^{u}_{x,y}(\downarrow) = a$ and $\tau^{u}_{x,y}(a) = \downarrow$.
  \item For $j \in L_{(2k-3)}\sm \set{a}$, following $\delta_j$, we arrive back at $x$, along another nodal semi-arc, which we denote by $\delta_{\tau_{x,y}^{u}(j)}$; this semi-arc emanates from $x$ tangentially to the ray $\omega_{\tau_{x,y}^{u}(j)}$. This defines $\tau^{u}_{x,y}$ on $L_{(2k-3)}\sm \set{a}$. The local structure theorem implies that for, $j \in L_{(2k-3)}\sm \set{a}$, $\tau_{x,y}^{u}(j) \in L_{(2k-3)}\sm \set{a}$, and $\tau_{x,y}^{u}(j)\neq j$.
\end{enumerate}

Doing so, we have uniquely defined the map $\tau^{u}_{x,y}$ from $\set{\downarrow} \cup L_{(2\ell-3)}$ to itself, such that $(\tau_{x,y}^{u})^2 = \id$, and  $(\tau_{x,y}^{u})(j) \neq j$.\smallskip

The pair $\set{\downarrow,a}$ corresponds to the nodal arc $\delta_a$ from $x$ to $y$. For $j \in L_{(2k-3)}\sm \set{a}$, the pair $\set{j,\tau_{x,y}^{u}(j)}$ corresponds to a loop $\gamma_{j,\tau_{x,y}^{u}(j)}$ at $x$.  There are $(\ell - 2)$ such loops.  Since $\cS(u) = \set{x,y}$ and $\rho(u,x) = 1$, these loops and arc do not intersect away from $x$. Since $\cZ(u) \cup \partial \Omega$ is connected, the nodal set $\cZ(u)$ is actually the union of these $(\ell - 2)$ loops and arc. Otherwise stated, $\cZ(u)$ is the wedge sum $\cB_{x,(\ell-2)}^{y}$ of the simple arc $\delta_{\tau^{u}_{x,y}(\downarrow)}$ from $x$ to $y$ with an $(\ell-2)$-bouquet of loops at $x$.


\begin{definition}\label{D-hmn-nt}
By analogy with Definition~\ref{D-h2n-type}, we call the map $\tau^{u}_{x,y}$ the \emph{combinatorial type} of the eigenfunction $u$ (or of the nodal set $\cZ(u)$) with respect to the points $x$ and $y$.
\end{definition}%

We describe the map $\tau^{u}_{x,y}$ in matrix form as
%
\begin{equation}\label{E-hmn-21b2}
\tau_{x,y}^{u} = \begin{pmatrix}
                   \downarrow & 1 & \ldots & (a-1) & a & (a+1) & \ldots & (2\ell -3) \\
                   a & \tau_{x,y}^{u}(1) & \ldots & \tau_{x,y}^{u}(a-1) & \downarrow & \tau_{x,y}^{u}(a+1) & \ldots & \tau_{x,y}^{u}(2\ell - 3) \\
                 \end{pmatrix}\,.
\end{equation}
%
\smallskip


Figure~\ref{F-hmn2-Uxy0} displays some possible nodal patterns (for $\ell = 5$, $\rho(x)=7$, and $\rho(y) =1$).  The corresponding combinatorial types are given by
%
\begin{equation*}
\resizebox{.97 \textwidth}{!}
{$
\tau_A = \begin{pmatrix}
           \downarrow & 1 & 2 & 3 & 4 & 5 & 6 & 7 \\
           1 & \downarrow & 3 & 2 & 5 & 4 & 7 & 6 \\
               \end{pmatrix}~
\tau_B = \begin{pmatrix}
           \downarrow & 1 & 2 & 3 & 4 & 5 & 6 & 7 \\
           3 & 2 & 1 & \downarrow & 7 & 6 & 5 & 4 \\
               \end{pmatrix}~
\tau_C = \begin{pmatrix}
           \downarrow & 1 & 2 & 3 & 4 & 5 & 6 & 7 \\
           7 & 6 & 5 & 4 & 3 & 2 & 1 & \downarrow \\
              \end{pmatrix}\,.
$}
\end{equation*}


%\subsubsection{Case~2. $\Omega$  has one hole.}\label{SSS-hmn-21c}
\subsubsection{$\Omega$ with one hole. Properties of $U_{x,y}$}\label{SSS-hmn-21c}

\begin{lemma}\label{L-Uxy1}
Assume that $\Omega$ has one hole, with $\partial \Omega =\Gamma_1 \cup \Gamma_2$,  where $\Gamma_1$ is  the  outer boundary.  Let $U$ be a linear subspace of an eigenspace of \eqref{E-evp-2bc} in $\Omega$,  such that $\sup\set{\kappa(u) ~|~ 0 \not = u \in U}\le \ell$ for some $\ell \ge 2$, and $\dim U = (2\ell-1)$. Let $x \neq y \in \Gamma_1$, and define
$$U_{x,y} := \set{u \in U ~|~ \rho(u,x) \ge (2\ell - 3) \text{~and~} \rho(u,y) \ge 1}.$$ Then,
\begin{enumerate}[(i)]
  \item $\dim U_{x,y} = 1$.
  \item  For $0 \neq u \in U_{x,y}$, the following alternative holds,
  \begin{itemize}
    \item[$\diamond$] either $b_0(\cZ(u) \cup \partial\Omega) = 1$, in which case $\cZ(u)$ hits $\Gamma_2$, $u$ has precisely two singular points on $\Gamma_2$ (counting multiplicities), $\sum_{z \in \cS_b(u) \cap \Gamma_2} \rho (u,z) = 2$,
$\cS_b(u) \cap \Gamma_1 = \set{x,y}$ and $\cS_i(u) = \emptyset$;
    \item[$\diamond$] or $b_0(\cZ(u) \cup \partial\Omega) = 2$, in which case $\cZ(u) \cap \Gamma_2 = \emptyset$,
$\cS_b(u) \cap \Gamma_1 = \set{x,y}$, and $\cS_i(u) = \emptyset$;
\end{itemize}

\item $\rho(u,x)=(2\ell-3)$ and $\rho(u,y) = 1$.
\item $\kappa(u) = \ell$.
\end{enumerate}
A generator of $U_{x,y}$ will be denoted by  $u_{x,y}$ (defined up to scaling).
\end{lemma}%

\pf The fact that $\dim U_{x,y}\ge 1$ follows from Lemma~\ref{L-zero2}. Since we now have $b_0(\partial \Omega)=2$, Euler's formula \eqref{E-euler-or2a} applied to $u$ gives
%
\begin{equation}\label{E-A6}
\begin{array}{ll}
0 \geq \kappa(u) - \ell  \,  = & \left( b_0(\cZ(u)\cup \partial \Omega) -2\right) + \frac 12 \sum_{z\in \cS_i(u)} (\nu(z)-2)\\[5pt]
& \hspace{3mm}+ \frac 12 \sum_{z \in \cS_b(u) \cap \Gamma_2} \rho (z)\\[5pt]
& \hspace{3mm}+ \frac 12 \sum_{\substack{y\in \cS_b(u) \cap \Gamma_1,\\ z \neq x,y}\,}
\rho (z) + \frac 12 \left( \rho(x) + \rho(y) -2 \ell +2\right).
 \end{array}
\end{equation}

Except for the term $\left( b_0(\cZ(u) \cup \partial \Omega) -2\right)$, all the terms in the right-hand side of the equality are nonnegative. This implies that
%
\begin{equation*}
2 \ge b_0(\cZ(u) \cup \partial \Omega) \ge 1\,,
\end{equation*}
and we have to examine two cases.\smallskip

\noid If $b_0(\cZ(u) \cup \partial \Omega) = 1$, then the nodal set $\cZ(u)$ must hit $\Gamma_2$. According to Proposition~\ref{P-euler-np}, the sum $\sum_{z \in \cS_b(u)\cap \Gamma_2} \rho (z)$ is an even integer, and we deduce from \eqref{E-A6} that $\sum_{z \in \cS_b(u)\cap \Gamma_2} \rho (z) = 2$. This equality now implies that the other terms are zero, and hence that $\kappa(u) = \ell$.\smallskip

\noid If $b_0(\cZ(u) \cup \partial \Omega) = 2$, then all the terms in the right-hand side of \eqref{E-A6} vanish, and $\kappa(u) = \ell$. \smallskip

This proves Assertions~(ii)--(iv).  As in the proof of Lemma~\ref{L-Uxy0}, assuming that there are at least two linearly independent functions $u_1$ and $u_2$ in $U_{x,y}$, the first assertion follows from Assertion~(iii)  and Lemma~\ref{L-zeroc}. \hfill \qed\medskip

\subsubsection{$\Omega$ with one hole. Structure and combinatorial type of nodal sets in $U_{x,y}$}\label{SSS-hmn-21cs}

From a geometric point of view, once we have fixed $x \not = y \in \Gamma_1$ and under the assumptions of Lemma~\ref{L-Uxy1},  either $\cZ(u) \cap \Gamma_2 = \emptyset$ or $\cZ(u) \cap \Gamma_2 = \set{y_1,y_2}$, possibly with $y_1=y_2$.  In the first case, we simply reproduce the description given in Paragraph~\ref{SSS-hmn-21as}. In the second case, the connectivity of $\cZ(u) \cup \partial \Omega$ implies that one of the nodal arcs hitting $\Gamma_2$ also contains $x$.  The other one contains either $x$ or $y$. More precisely, we choose any $j \in L_{(2\ell - 3)}$ and follow the nodal semi-arc $\delta_j$ emanating from $x$ along $\cZ(u)$, until we meet a singular point $z$ as described in Paragraph~\ref{SSS-hmn-21as}. Since $z \in \cS(u) = \set{x,y,y_1,y_2}$, there are two possibilities. If $z \in \set{x,y}$ the description is similar to the one in Paragraph~\ref{SSS-hmn-21as}. If $z \in \set{y_1,y_2}$, say $y_1$, we continue our path from $y_1$ to $y_2$ along $\Gamma_2$, and leave $\Gamma_2$ along the second nodal arc hitting $\Gamma_2$ at $y_2$ until we meet a singular point. We then either reach the point $x$ again or the point $y$. It then follows that the nodal set of $u \in U_{x,y}$ consists of $ (\ell - 2)$ ``generalized'' nodal loops at $x$ (one of the loops may comprise some part of $\Gamma_2$), and a ``generalized'' simple arc from $x$ to $y$ (this arc may comprise a sub-arc from $y_1$ to $y_2$ on $\Gamma_2$).  In the preceding description, the points $y_1$ and $y_2$ may coincide. These ``generalized'' loops and arc do not intersect away from $x$. Then, $\cZ(u)$ is the wedge sum $\cB^{y}_{x,(\ell-2)}$ of an $(\ell - 2)$-bouquet of ``generalized'' loops at $x$, with a simple ``generalized'' arc from $x$ to $y$. We can then define the \emph{combinatorial type} $\tau_{x,y}^{u}$ of $u$ with respect to the points $x$ and $y$ as we did in Paragraph~\ref{SSS-hmn-21as}, somehow ignoring $\Gamma_2$.\smallskip

Projecting $\cZ(u)$ to $\check{\Omega}$, we obtain a set $\check{\cZ}(u) \subset \check{\Omega}$ which is the wedge sum $\cB_{\check{x}, (\ell-2)}^{\check{y}}$ of an $(\ell-2)$-bouquet of loops at $\check{x}$ with a simple arc from $\check{x}$ to $\check{y}$. One of the loops or the arc may contain the point $\xi_2$, the image of $\Gamma_2$ in $\check{\Omega}$.   Since there are only two semi-arcs at $\xi_2$, this point is a regular point of the projected nodal partition $\check{\cD}_{u}$. The general  picture is then similar to the picture in the simply connected case.


Figures~\ref{F-hmn2-Uxy1-1} and \ref{F-hmn2-Uxy1-2} display some possible nodal patterns for $0 \neq u \in U_{x,y}$ when $\Omega$ has one hole ($\ell=5$, $\rho(x)=7$, and $\rho(y)=1$).

\begin{figure}[!ht]
\centering
\begin{subfigure}[t]{.30\textwidth}
\centering
\includegraphics[width=\linewidth]{hmn2-Uxy1-a}
\caption{}
\end{subfigure}
\begin{subfigure}[t]{.30\textwidth}
\centering
\includegraphics[width=\linewidth]{hmn2-Uxy1-b}
\caption{}
\end{subfigure}
\begin{subfigure}[t]{.30\textwidth}
\centering
\includegraphics[width=\linewidth]{hmn2-Uxy1-c}
\caption{}
\end{subfigure}
%
\caption{$\Omega$ with one hole: some possible nodal patterns for $u_{x,y}$}\label{F-hmn2-Uxy1-1}
\end{figure}

\begin{figure}[!ht]
\centering
\begin{subfigure}[t]{.30\textwidth}
\centering
\includegraphics[width=\linewidth]{hmn2-Uxy1-d}
\caption{}
\end{subfigure}
\begin{subfigure}[t]{.30\textwidth}
\centering
\includegraphics[width=\linewidth]{hmn2-Uxy1-e}
\caption{}
\end{subfigure}
\begin{subfigure}[t]{.30\textwidth}
\centering
\includegraphics[width=\linewidth]{hmn2-Uxy1-f}
\caption{}
\end{subfigure}
%
\caption{$\Omega$ with one hole: some possible nodal patterns for $u_{x,y}$}\label{F-hmn2-Uxy1-2}
\end{figure}
\FloatBarrier


For the nodal patterns in Figures~\ref{F-hmn2-Uxy1-1} and  \ref{F-hmn2-Uxy1-2}, we have the  combinatorial types
%
\begin{equation*}
\begin{array}{ll}
\tau^{\ref{F-hmn2-Uxy1-1}}_A = \tau^{\ref{F-hmn2-Uxy1-2}}_B = \tau^{\ref{F-hmn2-Uxy1-2}}_C &= \begin{pmatrix}
                         \downarrow & 1 & 2 & 3 & 4 & 5 & 6 & 7 \\
                          3 & 2 & 1 & \downarrow & 7 & 6 & 5 & 4 \\
                        \end{pmatrix}\,,\\[12pt]
\tau^{\ref{F-hmn2-Uxy1-1}}_C = \tau^{\ref{F-hmn2-Uxy1-2}}_A &= \begin{pmatrix}
                         \downarrow & 1 & 2 & 3 & 4 & 5 & 6 & 7 \\
                          1 & \downarrow  & 3 & 2 & 5 & 4 & 7 & 6 \\
\end{pmatrix}\,,\\[12pt]
                        \tau^{\ref{F-hmn2-Uxy1-1}}_B &= \begin{pmatrix}
                        \downarrow & 1 & 2 & 3 & 4 & 5 & 6 & 7 \\
                        5 & 2 & 1 & 4 & 3 & \downarrow & 7 & 6 \\
                        \end{pmatrix}\,.
\end{array}%
\end{equation*}


\subsubsection{ $\Omega$ with $k$ holes. Properties of $U_{x,y}$}\label{SSS-hmn-21d}


In this case, $\partial \Omega$ has $(k+1)$ components,  $\partial \Omega = \bigcup_{j=1}^{k+1} \Gamma_j$. Fix $x \not = y \in \Gamma_1$.
With the notation of Subsection~\ref{SS-hmn-0N}, we have the following lemma.

\begin{lemma}\label{L-Uxyk} Assume that $\Omega$ has $k$ holes, with $\partial \Omega = \bigcup_{j=1}^{k+1} \Gamma_j$.  Let $U$ be an eigenspace of \eqref{E-evp-2bc} in $\Omega$, such that, for some $\ell \ge 2$,  $\sup\set{\kappa(u) ~|~ 0 \not = u \in U}\le \ell$, and  $\dim U = (2\ell-1)$. Let $x \neq y \in \Gamma_1$, and define the subspace $$U_{x,y} := \set{u \in U ~|~ \rho(u,x) \ge (2\ell - 3) \text{~and~} \rho(u,y) \ge 1}.$$  Then, $\dim U_{x,y} = 1$. Furthermore, for all $0 \neq u \in U_{x,y}$:
\begin{enumerate}[(i)]
  \item The set $\cZ(u) \cup \Gamma(u)$ is connected.
  \item If $J(u)=\set{1}$, the only singular points of $u$ are the points $x$ and $y$, with $\rho(u,x) = (2\ell - 3)$ and $\rho(u,y)=1$.
  \item If $J(u) \not = \set{1}$, each component $\Gamma_j, j \in J(u)$, is hit by exactly two nodal arcs, the function $u$ has no interior singular point, and its only singular points on $\Gamma_1$ are $x$ and $y$, with $\rho(u,x) = (2\ell - 3)$ and $\rho(u,y)=1$.
  \item In all cases, $\kappa(u) = \ell$.
  \item In all cases, the nodal set of $u$ consists of $(\ell - 2)$ simple non-intersecting ``generalized'' nodal loops at $x$ (loops comprising nodal arcs, and possibly arcs contained in some boundary components $\Gamma_j, j \in J(u)\sm \set{1}$), a simple nodal arc from $y$ to either $x$ (when $J(u)=\set{1}$) or to some inner component of $\partial \Omega$, a simple nodal arc from $x$ to some component $\Gamma_j, j\in J(u)\sm \set{1}$, and possibly some nodal arcs joining components which meet $\cZ(u)$. These nodal arcs can only intersect at $x$ or possibly on the components $\Gamma_j, j \in J(u)\sm \set{1}$. In all cases, the point $x$ is joined to the point $y$ by a simple arc comprising nodal arcs and sub-arcs of the $\Gamma_j, j \in J(u)\sm \set{1}$.
\end{enumerate}
A generator of $U_{x,y}$ will be denoted by $u_{x,y}$ (defined up to scaling).
\end{lemma}%

\emph{Proof of Lemma~\ref{L-Uxyk}.}~  With the assumptions of the lemma, Euler's formula \eqref{E-euler-or2a}
can be rewritten as,
%
\begin{equation}\label{E-bha-k10}
\begin{array}{ll}
0 \ge \kappa(u) - \ell =& \left( b_0(\cZ(u) \cup \Gamma(u)) - 1) \right)
+ \frac 12\, \sum_{z \in \cS_i(u)} (\nu(z)-2)\\[5pt]
& +\, \sum_{j \in J(u), j\neq 1\,} \frac 12  \, \left( \sum_{z\in \cS _b(u) \cap \Gamma_j}\rho (z) - 2 \right)\\[5pt]
& +\, \frac 12 \sum_{z\in \cS _b(u) \cap \Gamma_1,\, z \not = x,y\,}\rho (z) + \frac 12 \left( \rho(x) + \rho(y) - 2\ell +2\right)\,.
\end{array}%
\end{equation}

In view of our assumptions, and Proposition~\ref{P-euler-np}, the terms in the right-hand side of \eqref{E-bha-k10} are all non-negative. In view of the left hand side of the inequality, they must all be zero. We now examine two cases.\smallskip

\noid If $J(u)=\set{1}$, the second line is the right hand side disappears, the nodal set $\cZ(u)$ only meets $\Gamma_1$, $b_0(\cZ(u) \cup \Gamma_1) =1$, the only  singular points  of the function $u$ are the points $x$ and $y$, and $ \rho(x) = (2\ell-3)$, $\rho(y) = 1$.\smallskip

\noid If $J(u) \not = \set{1}$, all the terms in the right hand side must be zero: $b_0(\cZ(u) \cup \Gamma(u)) =1$, each component $\Gamma_j, j\in J(u)\sm \set{1}$, is hit by precisely two nodal arcs of $\cZ(u)$, $ \rho(x) = (2\ell-3)$, and $ \rho(y) = 1$, and the function $u$ has no other  singular point whether in the interior of $\Omega$ or on $\partial\Omega$.  Furthermore, there is a simple nodal arc from $x$ to one of the components $\Gamma_j, j \in J(u)$, a simple nodal arc from $y$ to one of the components $\Gamma_j, j \in J(u)$, and there is a simple nodal arc, possibly comprising arcs contained in $\Gamma(u)$ joining $x$ to $y$. Finally, $\kappa(u) = \ell$. This proves assertions (i)--(v).\medskip

To prove the first assertion, assuming there are at least two linearly independent functions $u_1$ and $u_2$ in $U_{x,y}$, we can apply Lemma~\ref{L-zeroc} as in the previous proofs, and construct yet another function $0 \not = \tilde{u}$ such that $\rho(\tilde{u},x) \ge (2\ell -3)$ and $\rho(\tilde{u},y)\ge 2\,$, contradicting Assertions~(ii). \hfill \qed \medskip

\begin{figure}[!ht]
\centering
\begin{subfigure}[t]{.30\textwidth}
\centering
\includegraphics[width=\linewidth]{hmn2-Uxy2-a}
\caption{}
\end{subfigure}
\begin{subfigure}[t]{.30\textwidth}
\centering
\includegraphics[width=\linewidth]{hmn2-Uxy2-b}
\caption{}
\end{subfigure}
\begin{subfigure}[t]{.30\textwidth}
\centering
\includegraphics[width=\linewidth]{hmn2-Uxy2-c}
\caption{}
\end{subfigure}
%
\caption{$\Omega$ with two holes: some possible nodal patterns for $u_{x,y}$}\label{F-hmn2-Uxy2-1}
\end{figure}

\begin{figure}[!ht]
\centering
\begin{subfigure}[t]{.30\textwidth}
\centering
\includegraphics[width=\linewidth]{hmn2-Uxy2-d}
\caption{}
\end{subfigure}
\begin{subfigure}[t]{.30\textwidth}
\centering
\includegraphics[width=\linewidth]{hmn2-Uxy2-e}
\caption{}
\end{subfigure}
\begin{subfigure}[t]{.30\textwidth}
\centering
\includegraphics[width=\linewidth]{hmn2-Uxy2-f}
\caption{}
\end{subfigure}
%
\caption{$\Omega$ with two holes: some possible nodal patterns for $u_{x,y}$}\label{F-hmn2-Uxy2-2}
\end{figure}


\subsubsection{$\Omega$ with $k$ holes. Structure and combinatorial type of nodal sets in $U_{x,y}$}\label{SSS-hmn-21ds}

 We can adapt the description of the nodal set $\cZ(u)$, $u \in U_{x,y}$ given in  Paragraph~\ref{SSS-hmn-21cs} to the present case (multiple components of $\partial \Omega$. The ``generalized'' loops or arc will then hit one or several components $\Gamma_j, j \in J(u)\sm \set{1}$. We can also define the \emph{combinatorial type}  $\tau_{x,y}^{u}$ of $u$ with respect to the points $x$ and $y$. \smallskip

Figures~\ref{F-hmn2-Uxy2-1} and \ref{F-hmn2-Uxy2-2} display possible nodal patterns for $ 0 \neq u \in U_{x,y}$  (in these examples, $\ell = 5$, $\rho(x) = 7$, and there are $3$ loops). For these nodal patterns, we have  the combinatorial types

\begin{equation*}
\begin{array}{ll}
\tau^{\ref{F-hmn2-Uxy2-1}}_A  &= \begin{pmatrix}
                       \downarrow &  1 & 2 & 3 & 4 & 5 & 6 & 7 \\
                        1 &  \downarrow & 3 & 2 & 7 & 6 & 5 & 4 \\
                        \end{pmatrix}\,,\\[12pt]
\tau^{\ref{F-hmn2-Uxy2-1}}_B  &= \begin{pmatrix}
                         \downarrow & 1 & 2 & 3 & 4 & 5 & 6 & 7 \\
                          3 & 2 & 1 & \downarrow & 7 & 6 & 5 & 4 \\
                        \end{pmatrix}\,,\\[12pt]
\tau^{\ref{F-hmn2-Uxy2-1}}_C  &= \begin{pmatrix}
                         \downarrow & 1 & 2 & 3 & 4 & 5 & 6 & 7 \\
                          5 & 4 & 3 & 2 & 1 & \downarrow & 7 & 6 \\
                        \end{pmatrix}\,,\\[12pt]
\end{array}%
\end{equation*}
%
and
%
\begin{equation*}
\begin{array}{ll}
\tau^{\ref{F-hmn2-Uxy2-2}}_A  &= \begin{pmatrix}
                        \downarrow &  1 & 2 & 3 & 4 & 5 & 6 & 7 \\
                         1 & \downarrow & 5 & 4 & 3 & 2 & 7 & 6 \\
                        \end{pmatrix}\,,\\[12pt]
\tau^{\ref{F-hmn2-Uxy2-2}}_B = \tau^{\ref{F-hmn2-Uxy2-2}}_C &= \begin{pmatrix}
                         \downarrow & 1 & 2 & 3 & 4 & 5 & 6 & 7 \\
                         3 & 2  & 1 & \downarrow & 5 & 4 & 7 & 6 \\
                        \end{pmatrix}\,.
\end{array}%
\end{equation*}

Lemma~\ref{L-Uxyk} can be reformulated in the abstract setting of Subsection~\ref{SS-hmn2-10} as follows.
For any $0 \neq u \in U_{x,y}$, the projection $\check{\cZ}(u)$ of the nodal set $\cZ(u)$ consists of $(\ell-2)$ continuous simple loops at $\check{x}$ and a continuous simple curve from $\check{x}$ to $\check{y}$. The loops and curve only intersect at $\check{x}$ and may contain points in $\Xi$. If $\xi_j \in \check{\cZ}(u)$, there are exactly two projected nodal semi-arcs at this point, and the point $\xi_j$ is a regular point of $\check{\cD}_u$.  The set $\check{\cZ}(u) \subset \check{\Omega}$ is therefore the wedge sum $\cB_{\check{x},(\ell-2)}^{\check{y}}$ of an $(\ell-2)$-bouquet of loops at $\check{x}$, with a simple arc from $\check{x}$ to $\check{y}$. The loops or the arc may contain points in $\Xi$.  This is illustrated in Figure~\ref{F-hmn2-Uxyk-p}.\smallskip

\begin{figure}[!ht]
\centering
\includegraphics[width=0.6\linewidth]{hmn2-Uxyk-p}
%
\caption{Nodal patterns in $\Omega$ and their projections in $\check{\Omega}$}
\label{F-hmn2-Uxyk-p}
\end{figure}
\FloatBarrier



\begin{remark}\label{R-hmn23-2}
From the point of view of Definition~\ref{D-parti-reg}, the points in $\Xi$ are not singular points of $\check{\cD}_u$, the projection of the nodal partition $\cD_u$ of $u$.
\end{remark}



\subsection{Analysis of eigenfunctions with one prescribed   boundary singular point}\label{SS-hmn-22}

We use the notation of Subsection~\ref{SS-hmn-0N}.  In this subsection, we assume that $U$ is a linear subspace of an eigenspace $U(\lambda)$ of \eqref{E-evp-2bc}, and that for some $\ell \ge 2$,
%
\begin{equation*}
\left\{
\begin{array}{l}
\sup\set{\kappa(u) ~|~ 0 \neq u \in U} \le \ell \text{~~and}\\[5pt]
\dim U = (2\ell - 1)\,.
\end{array}%
\right.
\end{equation*}

For $x \in \Gamma_1$, we introduce the subspaces
%
\begin{equation}\label{E-bha4-2}
\left\{
\begin{array}{l}
U^1_x = \set{u \in U ~|~ \rho(u,x) \ge (2\ell -2)}\,,\\[5pt]
U^2_x = \set{u \in U ~|~ \rho(u,x) \ge (2\ell -3)}\,.\\[5pt]
\end{array}%
\right.
\end{equation}

According to Lemma~\ref{L-zero2},  $U^1_x \neq \set{0}$. The purpose of this subsection is to investigate the properties of the functions $u \in U^1_x$ or $U^2_x$ --precise order of vanishing, structure of their nodal sets-- under the above assumptions on $U$.

\subsubsection{Properties of $U^1_x$ and $U^2_x$}\label{SSS-hmn-22a}

\begin{lemma}\label{L-Ux1}
Let $U$ be a linear subspace of an eigenspace of \eqref{E-evp-2bc} in $\Omega$, with $$ \sup\set{\kappa(u) ~|~ 0 \not = u \in U} \le \ell\,,  \text{~~and~~} \dim U = (2\ell - 1).$$ Fix some $ x \in \Gamma_1$.  For the spaces $U^1_x$ and $U^2_x$ defined in \eqref{E-bha4-2}, we have
\begin{enumerate}[(i)]
  \item $\dim U^1_x = 1$, \, $\dim U^2_x = 2$ and,
  \item  for any $0 \neq u \in U^2_x$,
  %
  \begin{equation}\label{E-bha4-2a}
\left\{
\begin{array}{l}
\kappa(u) = \ell \text{~~and~~} \cZ(u) \cup \Gamma(u) \text{~is connected},\\[5pt]
\cS_i(u) = \emptyset\,,\\[5pt]
\sum_{z \in \cS_b(u) \cap \Gamma_j} \rho(u,z) = 2 \text{~for all~} j \in J(u)\sm \set{1}\,,\\[5pt]
\sum_{z \in \cS_b(u) \cap \Gamma_1} \rho(u,z) = (2\ell - 2) \text{~and, more precisely,}\\[5pt]
\begin{array}{ll}
\text{(i) either} & \rho(u,x) = (2\ell - 2) \text{~and~} \cS_b(u) \cap \Gamma_1 = \set{x}\,,\\[5pt]
\text{(ii) or} & \rho(u,x)  = (2\ell-3),\, \exists\, y_u \in \Gamma_1 \sm \set{x} \text{~with~} \rho(u,y_u) = 1\,,\\[5pt]
& \quad \text{and~}  \cS_b(u) \cap \Gamma_1 = \set{x,y_u}\,.
\end{array}%
\end{array}%
\right.
\end{equation}

\end{enumerate}
\end{lemma}%

\pf Assume that $\partial \Omega$ has $(q+1)$ components, $\Gamma_1, \ldots, \Gamma_{q+1}$, with $x \in \Gamma_1$.\smallskip

Clearly, $\set{0} \neq U_x^1 \subset U_x^2$. Take any $0 \not = u \in U^2_x$. Euler's formula \eqref{E-euler-or2a} can be rewritten as
%
\begin{equation}\label{E-bha4-2e}
\begin{array}{ll}
0 \ge \kappa(u) - \ell =& \left( b_0(\cZ(u) \cup \Gamma(u)) - 1 \right)
+ \frac 12\, \sum_{z \in \cS_i(u)} (\nu(z)-2)\\[5pt]
& +\, \sum_{i \in J(u), i\neq 1\,} \frac 12 \, \left( \sum_{z\in \cS _b(u) \cap \Gamma_i}\rho (z) - 2 \right)\\[5pt]
& +\, \frac 12\, \left( \sum_{z\in \cS _b(u) \cap \Gamma_1\,}\rho (z)  - 2\ell +2\right)\,.
\end{array}%
\end{equation}
%
The first $ |J(u)|+2$  terms in the right-hand side of the equality are nonnegative. Since $\sum_{z\in \cS _b(u) \cap \Gamma_1\,}\rho (z)$ is even, and larger than or equal to $(2\ell -3)$, the last term is nonnegative also. In view of the first inequality, the four terms must vanish. This proves the relations \eqref{E-bha4-2a}.
\medskip

\noid \emph{Proof that $\dim U^1_x = 1$.~}  Assume that this is not the case. Then, there exist two linearly independent functions $u_1, u_2$ in $U^1_x$ such that $\rho(u_i,x) = (2\ell - 2)$. By Lemma~\ref{L-zeroc}, there would exist a nontrivial linear combination $u$ such that $u \in U^1_x$ and $\rho(u,x) \ge (2\ell - 1)$, a contradiction with \eqref{E-bha4-2a}. \medskip


\noid \emph{Proof that $\dim U^2_x = 2$.~} Choose some $0 \not = v_1 \in U^1_x$. Clearly $v_1 \in U^2_x$. On the other-hand, given any $y \in \Gamma_1\sm\set{x}$, Lemma~\ref{L-Uxy0} provides a function $u_{x,y}$ belonging to $U^2_x$, not to $U^1_x$, and hence $\dim U^2_x \ge 2$. Choose $0 \not = v_2 \in U^2_x$ orthogonal to $v_1$. Since $\dim U^1_x = 1$, the function $v_2$ satisfies $\rho(v_2,x) = (2\ell - 3)$, and by Proposition~\ref{P-euler-np}, there must exist some $y_2 \in \Gamma$ such that $\rho(v_2,y_2) \ge 1$.  By Lemma~\ref{L-Uxyk}, $\rho(v_2,y_2) = 1$ and $v_2 \in U_{x,y_2}$. The subspace $U^{1,\perp}_x := \set{u \in U^2_x ~|~ u \perp u_1}$ has dimension at least one. Assume that $\dim U^2_x \ge 3$. Then $\dim U^{1,\perp}_x \ge 2$, and we can find two linearly independent functions $u_1, u_2 \in U^{1,\perp}_x$ such that $\rho(u_i,x) = (2\ell - 3)$. By Lemma~\ref{L-zeroc}, there exists a linear combination $u \in U^{1,\perp}_x$ such that $\rho(u,x) \ge (2\ell - 2)$, a contradiction.  \hfill \qed \medskip


\begin{remark}\label{R-Ux1-2}
Up to scaling, there is a uniquely defined orthogonal basis $\set{v_1,v_2}$ of $U^2_x$, with $v_1 \in U^1_x$, $v_2 \in  U^{1,\perp}_x$, and a uniquely defined $y_2 \in \Gamma_1 \sm \set{x}$ such that $\rho(v_2,y_2) = 1$. In view of Lemma~\ref{L-breve}, we can choose $v_1$ such that $\breve{v}_1 > 0$ on $\Gamma_1 \sm \set{x}$, and $v_2$ such that $\breve{v_2} > 0$ on the arc from $x$ to $y_2$ moving counter-clockwise on $\Gamma_1$.
\end{remark}%

\begin{figure}[!ht]
\centering
\begin{subfigure}[t]{.30\textwidth}
\centering
\includegraphics[width=\linewidth]{hmn2-Ux1-a}
\caption{}
\end{subfigure}
\begin{subfigure}[t]{.30\textwidth}
\centering
\includegraphics[width=\linewidth]{hmn2-Ux1-c}
\caption{}
\end{subfigure}
\begin{subfigure}[t]{.30\textwidth}
\centering
\includegraphics[width=\linewidth]{hmn2-Ux1-f}
\caption{}
\end{subfigure}
%
\caption{Some possible nodal patterns for $0 \not = u \in U^1_x$}\label{F-hmn2-Ux1}
\end{figure}
\FloatBarrier

\subsubsection{Structure and combinatorial type of nodal sets in $U^1_x$ and $U^2_x$}\label{SSS-hmn-22as}\phantom{}~

\noid Relations~\eqref{E-bha4-2a} and an analysis as in Subsection~\ref{SS-hmn-21}, show that the nodal set of any $0\neq u \in U^1_x$ consists of $(\ell -1)$ ``generalized'' nodal loops at the point $x$, and that these loops do not intersect away from $x$.  In the abstract setting of Subsection~\ref{SS-hmn2-10},  for any $0 \neq u \in U^1_x$, the projection $\check{\cZ}(u)$ of the nodal set $\cZ(u)$ consists of $(\ell-1)$ continuous loops at $\check{x}$. The loops only intersect at $\check{x}$, and may contain points in $\Xi$. If $\xi_j \in \check{\cZ}(u)$, there are exactly two projected nodal semi-arcs at this point. It follows that $\xi_j$ is a regular point of $\check{\cD}_u$.  The set $\check{\cZ}(u)$ is an $(\ell - 1)$-bouquet of loops $\cB_{\check{x},(2\ell -2)}$ at $\check{x}$. \smallskip

Adapting the description given in Paragraph~\ref{SSS-hmn-21as},  for $0 \neq u \in U_x^1$, we define the \emph{combinatorial type} $\tau_{x}^{u}$ of the nodal set $\cZ(u)$ with respect to $x$.   This is a map from $L_{(2\ell - 2)}$ to itself.
\smallskip

Some possible nodal patterns  for $u \in U_x^1$ are displayed in Figure~\ref{F-hmn2-Ux1},  where $\ell = 5$, and $\rho(x) = 8$.  The corresponding combinatorial types  are
%
\begin{equation*}
\begin{array}{l}
 \tau^{\ref{F-hmn2-Ux1}}_A  = \begin{pmatrix}
                          1 & 2 & 3 & 4 & 5 & 6 & 7 & 8\\
                          2 & 1 & 4 & 3 & 8 & 7 & 6 & 5\\
                        \end{pmatrix}\,,\\[12pt]
\tau^{\ref{F-hmn2-Ux1}}_B = \begin{pmatrix}
                          1 & 2 & 3 & 4 & 5 & 6 & 7 & 8\\
                          4 & 3 & 2 & 1 & 8 & 7 & 6 & 5 \\
                        \end{pmatrix}\,,\\[12pt]
\tau^{\ref{F-hmn2-Ux1}}_C = \begin{pmatrix}
                          1 & 2 & 3 & 4 & 5 & 6 & 7 & 8\\
                          8 & 3 & 2 & 5 & 4 & 7 & 6 & 1 \\
                        \end{pmatrix}\,.
\end{array}%
\end{equation*}

\noid  If $u \in U^2_x$ and $u \not \in U^1_x$, there exists a unique $y_u \in \Gamma_1$, such that $y_u \neq x$ and $\cS_b(u) \cap \Gamma_1 = \set{x,y_u}$, with $\rho(u,x) = (2\ell - 3)$, $\rho(u,y_u)=1$. Furthermore, $U_{x,u_u}=[u]$. The nodal set $\cZ(u)$  and its \emph{combinatorial type} $\tau_{x,y_u}^{u}$  are described in Paragraph~\ref{SSS-hmn-21d}. Projecting  $\cZ(u)$ to $\check{\Omega}$, $\check{\cZ}(u)$ is the wedge sum $\cB^{\check{y}_u}_{\check{x},(\ell -2)}$ of a simple arc from $\check{x}$ to $\check{y}_u$ with an $(\ell-2)$-bouquet of loops at $\check{x}$.





\subsection{Application of the previous analysis}\label{SS-hmn-24}
Fix some $x \in \Gamma_1$. We now apply the analysis of Subsections~\ref{SS-hmn-21} and \ref{SS-hmn-22} to investigate the limits $u_{x,y} \in U^2_x\sm U^1_x$, when $y$ tends to $x$ on $\Gamma_1$, clockwise or anti-clockwise. The notation are the same as in  Subsection~\ref{SS-hmn-21}.\smallskip

We choose a basis $\set{v_1,v_2}$ of $U^2_x$ as described in Remark~\ref{R-Ux1-2}. In particular, $\rho(v_1,x)=(2\ell -2)$,  $v_1 \perp v_2$ in $L^2(\Omega)$,  $\rho(v_2,x)=(2\ell-3)$, there exists $y_2 \in \Gamma_1\sm \set{x}$ such that $\rho(v_2,y_2)= 1$, and $\cS(v_2) \cap \Gamma_1 = \set{x,y_2}$.
Recall the definition of the functions $\breve{v}_i$ on $\partial \Omega$,
%
\begin{equation}\label{E-bha4-3d}
\breve{v}_i := \left\{
\begin{array}{ll}
\partial_{\nu}v_i &\quad \text{in the Dirichlet case,}\\[5pt]
v_i|_{\partial \Omega} &\quad \text{in the Robin case.}\\[5pt]
\end{array}
\right.
\end{equation}


According to Lemma~\ref{L-breve}, the function $\breve{v}_1$ vanishes only at $x$ and does not change sign on $\Gamma_1$. The function $\breve{v}_2$ does not vanish on $\Gamma_1 \sm \set{x,y_2}$, and changes sign when crossing $y_2$  and $x$ along $\Gamma_1$.\smallskip


Let $\gamma : [0,2\pi] \to \Gamma_1$ be a parametrization such that $\gamma(0)=\gamma(2\pi) = x$.  Given any $y\in \Gamma_1 \sm \set{x}$, there exists a function $u_{x,y}$ which satisfies \eqref{E-bha4-2a}(ii), and this function is uniquely defined up to multiplication by a nonzero scalar. In the Dirichlet case, this function is characterized by the fact that $\breve{u}_{x,y} = \partial_{\nu}u_{x,y}|_{\Gamma_1}$ only vanishes at $x$ and $y$. In the Robin case, it is characterized by the fact that $\breve{u}_{x,y} =u_{x,y}|_{\Gamma_1}$ only vanishes at $x$ and $y$. Up to a constant factor, we may choose
%
\begin{equation}\label{E-bha4-4ab}
u_{x,y} = a(y) \, v_1 + b(y) \, v_2 \,,
\end{equation}
%
with
%
\begin{equation}\label{E-bha4-4c}
\left\{
\begin{array}{l}
a(y) = -\, \breve{v}_2(y)\, \big( \breve{v}_1^2(y) + \breve{v}_2^2(y) \big)^{- \frac 12} \,,\\[5pt]
b(y) = \breve{v}_1(y) \, \big( \breve{v}_1^2(y) + \breve{v}_2^2(y) \big)^{- \frac 12}\,,\\[5pt]
\end{array}
\right.
\end{equation}
%
where $\breve{v}_1, \breve{v}_2$ are defined in \eqref{E-bha4-3d}. \\
Then, there exists a unique $\theta(y) \in (0,\pi)$ such that $\cos(\theta(y)) = a(y)$ and $\sin(\theta(y)) = b(y)$ (this is because $\breve{v}_1$ is positive on $\Gamma_1 \sm \set{x}$). Defining
%
\begin{equation}\label{E-bha4-4d}
w_{\theta} = \cos\theta \, v_1 + \sin\theta \, v_2\,,
\end{equation}
%
we have $u_{x,y} = w_{\theta(y)}$. Conversely, according to the proof of Lemma~\ref{L-Ux1}, any function $w_{\theta}$ has exactly two singular points on $\Gamma_1$, the point $x$ and some other point $y_{\theta} \not = x$. Note that the point $y$ determines the eigenfunction $u_{x,y}$ uniquely (up to scaling) and vice versa. It follows that we have a continuous, bijective map $(0,2\pi) \ni t \mapsto \theta(\gamma(t)) \in (0,\pi)$. This map is strictly monotone (we can assume that it is increasing), and provides a diffeomorphism from $(0,2\pi)$ to $(0,\pi)$, with limits $0$ and $\pi$ respectively. Otherwise stated, the function $u_{x,\gamma(t)}$ defined in \eqref{E-bha4-4ab} tends to $v_1$ when $t$ tends to $0$ and to $-v_1$ when $t$ tends to $2\pi$. There exists $t_2$ such that $\gamma(t_2) = y_2$, and hence $\theta(y_2) = \frac \pi 2$.  We have proved the following property.

\begin{property}\label{P-hmn}
The function $u_{x,y}$ defined in \eqref{E-bha4-4ab} tends to $v_1$ when $y \neq x$ tends to $x$ clockwise,  and to $-v_1$ when when $y \neq x$ tends to $x$ counter-clockwise.
\end{property}%



\subsection{Proof that $\mult(\lambda_k) \le (2k -2)$  for all $k\ge 3$, assuming $\Omega$ to be simply connected}\label{SS-hmn-25}

In this subsection, we work with the family of functions $\set{w_{\theta} ~|~ \theta \in [0,\pi]}$ introduced in \eqref{E-bha4-4d}, and we assume that $\partial \Omega = \Gamma_1$.

\subsubsection{Preparation}\label{SSS-hmn-25A}  In view of Proposition~\ref{P-hmn-s1}, and reasoning by contradiction, we assume that $\dim U(\lambda_k) = (2k -1)$. By Courant's theorem, we have $$\sup\set{\kappa(u) ~|~ 0 \not = u \in U} \le k.$$ We can apply Lemma~\ref{L-Ux1} with $\ell = k$ and $U := U(\lambda_k)$. \smallskip


In the arguments below we keep the notation of Lemma~\ref{L-Ux1} and its proof (with $\ell = k$). We fix a basis $\set{v_1,v_2}$ of $U^2_x$ as described at the beginning of  Subsection~\ref{SS-hmn-24}, and the direct frame $\set{\vec{e_1},\vec{e}_2}$ such that $\vec{e_1}$ is tangent to $\partial \Omega$ at $x$, and $\vec{e}_2$ is normal to $\partial \Omega$, pointing inwards.


\subsubsection{Structure and combinatorial types for $v_1$ and $v_2$}\label{SSS-hmn-25B}

We apply Subsection~\ref{SS-hmn2-20} with some $r_0$ small enough. In the neighborhood $B_{+}(x,r_0)$ of $x$, the nodal set $\cZ(v_1)$ consists of $(2\ell-2)$ nodal semi-arcs $\delta_{1,j}$ emanating from $x$ tangentially to rays $\omega_{1,j}, j \in L_{(2k-2)}$; the nodal set $\cZ(v_2)$ consists of $(2\ell-3)$ nodal semi-arcs $\delta_{2,j}$ emanating from $x$ tangentially to rays $\omega_{2,j}, j \in L_{(2k-3)}$.\smallskip


The combinatorial type of the function $v_1 \in U^1_x$ with respect to $x$ is defined in Subsection~\ref{SS-hmn-22}. This is a map
%
\begin{equation}\label{E-hmn-206a}
\begin{array}{l}
\tau_{x}^{v_1} : L_{(2k-2)} \to L_{(2k-2)} \text{~such that}\\[5pt]
\tau_{x}^{v_1}(j) \not = j \text{~and~} (\tau_{x}^{v_1})^2(j) = j\,,
\text{~~for all~} j \in L_{(2k-2)}\,.
\end{array}%
\end{equation}
%
The nodal set $\cZ(v_1)$ is a $(k-1)$-bouquet of loops at $x$ described by the map $\tau_{x}^{v_1}$.\medskip

The combinatorial type $\tau_{x,y_2}^{v_2}$ of the function $v_2 \in U^{1,\perp}_x$ with respect to $x$ and $y_2$ is defined in Subsection~\ref{SS-hmn-22}. Recall that it is described as a map
%
\begin{equation}\label{E-hmn- 206b}
\left\{
\begin{array}{l}
\tau_{x,y_2}^{v_2} : \set{\downarrow} \cup L_{(2k-3)} \to \set{\downarrow} \cup  L_{(2k-3)} \text{~such that}\\[5pt]
\tau_{x,y_2}^{v_2}(\downarrow) =: a \in L_{(2k-3)} \text{~and~}
\tau_{x,y_2}^{v_2}(a) = \, \downarrow\,,\\[5pt]
\tau_{x,y_2}^{v_2}(j) \not = j \text{~and~} (\tau_{x,y_2}^{v_2})^2(j) = j\,,
\text{~~for all~} j \in L_{(2k-3)}\sm \set{a}\,.
\end{array}%
\right.
\end{equation}
%
Here, $\tau_{x,y_2}^{v_2}(\downarrow)$ is the element $a \in L_{(2k-3)}$ such that the semi-arc $\delta_a$ of $\cZ(v_2)$ which emanates from $x$ tangentially to $\omega_{2,a}$ eventually hits $\Gamma_1$ at the point $y_2$. For $a \neq j \in L_{(2k-3)}$,  the pairs $\left( j, \tau_{x,y_2}^{v_2}(j)\right)$ describe the loops of $\cZ(v_2)$ at the point $x$, so that $\cZ(v_2)$ is the wedge sum of the nodal arc $\delta_a$ with an $(k-2)$-bouquet of loops at $x$, described by the map $\tau_{x,y_2}^{v_2}$. \smallskip

Since $\Omega$ is simply connected, the arc $\delta_a$ separates $\Omega$ into
two connected components $\Omega_{a,+}$ (on the right side of $\delta_a$),
and $\Omega_{a,-}$ (on the left side of $\delta_a$). There are three cases to consider, $a = 1$, $1 < a < (2k-3)$, and $a = (2k-3)$. The following properties follow easily from looking at the local structure of $\cZ(v_2)$ at $x$.

\begin{properties}\label{P-hmn-v2}\phantom{}
\begin{enumerate}[(i)]
  \item If $a = 1$, the component $\Omega_{a,+}$ does not contain any nodal arc, and the rays $\omega_{2,j}$, $2 \le j \le (2k-3)$ point inside $\Omega_{a,-}$.
  \item If $1 < a < (2k-3)$, the rays $\omega_{2,j}$, $1 \le j \le (a-1)$ point inside $\Omega_{a,+}$ ; the rays $\omega_{2,j}$, $(a+1) \le j \le (2k-3)$ point inside $\Omega_{a,-}$.
  \item If $a = (2k-3)$, the rays $\omega_{2,j}$, $1 \le j \le (2k-4)$ point inside $\Omega_{a,+}$ ; the component $\Omega_{a,-}$ does not contain any nodal arc.
\end{enumerate}
\end{properties}%


If the ray $\omega_{2,j}$ points inside $\Omega_{a,+}$, the whole nodal arc $\delta_{j}$ of $\cZ(v_2)$ is contained in $\Omega_{a,+}$, and so does the corresponding loop $\gamma_{j,\tau_{x,y_2}^{v_2}(j)}$. There is an analogous statement for $\Omega_{a,-}$. \smallskip

This means that
%
\begin{equation}\label{E-hmn-206c}
\left\{
\begin{array}{l}
a = \tau_{x,y_2}^{v_2}(\downarrow) \in L_{(2k-3)} \text{~~ is odd,}\\[5pt]
\tau_{x,y_2}^{v_2}\left(\set{1,\ldots,(a-1)}\right) \subset \set{1,\ldots,(a-1)}\,, \\[5pt]
\tau_{x,y_2}^{v_2}\left(\set{(a+1),\ldots,(2k-3)}\right) \subset \set{(a+1),\ldots,(2k-3)}\,.
\end{array}
\right.
\end{equation}
%


More precisely, define
%
\begin{equation}\label{E-hmn-208}
\left\{
\begin{array}{l}
K_{a,+} := \set{1,\ldots,(a-1)}\,, \text{~with~~} K_{a,+} = \emptyset \text{~~if~~} a=1\,,\\[5pt]
K_{a,-} := \set{(a+1),\ldots,(2k-3)}\,, \text{~with~~} K_{a,-} = \emptyset \text{~~if~~} a=(2k-3)\,,\\[5pt]
a_{(+)} := \frac{a-1}{2}\,.
\end{array}
\right.
\end{equation}

Then, the set $K_{a,+}$ corresponds to $a_{(+)}$ loops in the component $\Omega_{a,+}$ of $\Omega\sm \delta_a$, and these loops divide $\Omega_{a,+}$ into $(a_{(+)}+1)$ nodal domains of $v_2$. The set $K_{a,-}$ corresponds to $a_{(-)} := (k-2-a_{(+)})$ loops in the component $\Omega_{a,-}$ of $\Omega\sm \delta_a$, and these loops divide $\Omega_{a,-}$ into $a_{(-)}+ 1 = (k-1-a_{(+)})$ nodal domains of $v_2$, so that we recover the fact that $v_2$ has $k$ nodal domains.\smallskip

Otherwise stated the nodal set $\cZ(v_2)$ consists of the wedge sum of the nodal arc $\delta_a$ with two bouquets of loops, one contained in $\Omega_{a,+}$ and corresponding to $ \tau_{x,y_2}^{v_2}|_{K_{a,+}}$,  another contained in $\Omega_{a,-}$ and corresponding to $ \tau_{x,y_2}^{v_2}|_{K_{a,-}}$. One of these bouquets may be empty (when $a=1$ or $a=(2k-3)$).\medskip



For $\theta \in (0,\pi)$,  let $w_{\theta} := \cos\theta \, v_1 + \sin\theta \, v_2$.  Then $\cS_b(w_{\theta}) = \set{x,y_{\theta}}$. The nodal set $\cZ(w_{\theta})$ has a structure similar to the structure of $\cZ(v_2)$:
\begin{itemize}
\item[$\diamond$] one nodal arc $\delta_{a_{\theta},\theta}$ emanating from $x$ tangentially to some ray $\omega_{2,a_{\theta}}$, and hitting the boundary $\Gamma_1$ at the point $y_{\theta} \not = x$; $a_{\theta}$ is odd, and $a_{\theta} = \tau_{x,y_{\theta}}^{w_{\theta}}(\downarrow)$;
\item[$\diamond$]  loops at $x$, on either sides of $\delta_{a_{\theta},\theta}$, described by the restriction of the combinatorial type $\tau_{x,y_{\theta}}^{w_{\theta}}$ to $L_{(2k-3)} \sm \set{a_{\theta}}$.
\end{itemize}

Note:  In the above description, the nodal arc is denoted by $\delta_{a_{\theta},\theta}$ because it does not only depend on $a_{\theta}$. This point will be needed later on.


\begin{lemma}\label{L-hmn-202} Recall the notation $a = \tau_{x,y_{2}}^{v_2}(\downarrow)$ and $a_{\theta} = \tau_{x,y_{\theta}}^{w_{\theta}}(\downarrow)$.
For all $\theta \in (0,\pi)$, we have $a_{\theta} = a$, and $\tau_{x,y_{\theta}}^{w_{\theta}} = \tau_{x,y_2}^{v_2}$, i.e., the  combinatorial type of the nodal set $\cZ(w_{\theta})$ is the same as the combinatorial type of $\cZ(v_2)$.
\end{lemma}%

\emph{Proof of Lemma \ref{L-hmn-202}.} ~ We consider local conformal coordinates as in Subsection~\ref{SS-hmn2-20}. The proof of the local structure theorem shows that one can choose a uniform $r_0$ with respect to $\theta$. We use polar coordinates $(r,\omega)$ associated with $(\xi_1,\xi_2)$ in $\R^2$, and we write $E(r,\omega)$ for $E(r\cos\omega,r\sin\omega)$.\smallskip

To prove that $a_{\theta}$ is constant, by connectedness, it suffices to prove that it is locally constant:\\ {\it For all $\theta \in (0,\pi)$ there exists $\varepsilon_{\theta} > 0$ such that $a_{\theta'} = a_{\theta}$ for $|\theta' - \theta| < \varepsilon_{\theta}$.} \\
Assume, by contradiction, that this is not the case. Then, there exists $\theta_0 \in (0,\pi)$ such that for all $n \ge 1$, there exists $\theta_n$ with $|\theta_n - \theta_0| < \frac 1n$ and $a_{\theta_n} \not = a_{\theta_0} =: a_0$. Since $a_{\theta}$ can only take finitely many values in $ L_{(2k-3)}$, taking a subsequence if necessary, we can assume that $a_{\theta_n} \equiv a_1 \not = a_0$. By the local structure theorem, there exists a uniform $r_0 > 0$ (depending on $\theta_0$) such that the nodal arc $ \delta_{a_1,\theta_n}$ intersects the   set $C_{+}(x,r_0)$ at the point $z_n := E\left( r_0, \omegat_{a_1}(r_0,\theta_n) \right)$, where the function $\omegat_{a_1}(r,\theta)$ is smooth in a neighborhood of $(r_0,\theta_0)$ (with the notation of Appendix~\ref{S-lsns}). The arcs $\delta_{a_1,\theta_n} \cap \Omega\sm B(x,r_0)$ are compact and connected, and we can find a subsequence which converges in the Hausdorff distance to some compact connected set $\bar{\delta}$ which contains the point $z_0 = E\left( r_0, \omegat_{a_1}(r_0,\theta_0) \right)$ and the point $y_{\theta_0} = \lim y_{\theta_n}$ at the boundary. The set $\bar{\delta}$ is also contained in $\cZ(w_{\theta_0})$ because $w_{\theta_n}$ tends to $w_{\theta_0}$ uniformly. Since $a_1 \not = a_0$, we have a contradiction.\smallskip

Since $a_{\theta} \equiv a$ in $(0,\pi)$, in order to prove that $\tau_{\theta} := \tau_{x,y_{\theta}}^{w_{\theta}}$ does not depend on $\theta$, it suffices to show that its restrictions to the sets $K_{a,+}$ and $K_{a,-}$ are locally constant in $\theta$. We give the proof for $K_{a,+}$ in the case $a > 1$. The other cases are similar. Reasoning by contradiction, we assume that there exists $\theta_0 \in (0,\pi)$ such that, for all $n \ge 1$, there exist $\theta_n, |\theta_n-\theta_0| < \frac 1n$, and $j_n \in K_{a,+}$, such that $\tau_{\theta_n}(j_n) \not = \tau_{\theta_0}(j_n)$. Since $K_{a,+}$ is finite, taking subsequences if necessary, we may assume that $j_n \equiv b$ and $\tau_{\theta_n}(b) \equiv c$ for some $b,c \in K_{a,+}$ with $c\not = \tau_{\theta_0}(b)$. Since $\theta_n$ is close to $\theta_0 \in (0,\pi)$ we have a uniform structure theorem, and we can reason as in the proof of Property~\ref{propertyD4} to conclude. \hfill \qed\medskip

Once we are given the basis $ \set{v_1,v_2}$, we have the associated odd integer  $a := \tau^{v_2}(\downarrow) \in L_{(2k-3)}$, where $\tau^{v_2} := \tau_{x,y_2}^{v_2}$ is the combinatorial type of $v_2$. For all $\theta \in (0,\pi)$, the nodal set $\cZ(w_{\theta})$ has the same combinatorial type $\tau^{v_2}$. In particular, it contains a single simple nodal arc $\delta_{a,\theta}$, emanating from $x$ tangentially to the ray $\omega_{2,a}$, and hitting the boundary $\Gamma_1 = \partial \Omega$ at the point $y_{\theta}$. \smallskip
%We also have a map $\tau_{x,a}:= \tau_{x,a,v_2} : L_{(2k-3)}\sm
%\set{a} \to L_{(2k-3)}\sm \set{a}$ which gives the combinatorial type
%of the nodal set $\cZ(w_{\theta})$ at $x$.

Call $\Omega_{\theta,+}$ the component of $\Omega \sm \delta_{a,\theta}$ with semi-tangent at $x$ the vector $\vec{e}_1$, and  $\Omega_{\theta,-}$ the other component, with semi-tangent at $x$ the vector $-\vec{e}_1$. The component $\Omega_{\theta,+}$ contains the $a_{(+)}$ loops corresponding to the set $K_{a,+}$. These loops bound $(a_{(+)}+1)$ nodal domains of $w_{\theta}$ which can be labeled from $1$ to $(a_{(+)}+1)$. The component $\Omega_{\theta,-}$ contains the $a_{(-)} = (k - a_{(+)} -2)$ loops corresponding to the set $K_{a,-}$. These loops bound $(k - a_{(+)} -1)$ nodal domains of $w_{\theta}$ which can be labeled from $(a_{(+)}+2)$ to $k$.\smallskip


\subsubsection{Labeling and comparing nodal domains}\label{SSS-hmn-25C}\phantom{}

 The eigenfunctions we consider in this subsection all have $k$ nodal domains, and a boundary singular point $x$ with index $(2k-2)$ or $(2k-3)$. Let $w$ be such a function, and $n_w := \rho(w,x) \in \set{(2k-3),(2k-2)}$ its index at $x$. For $r_0$ small enough, the local structure theorem applied to $w$ at $x$, implies that $\cZ(w) \cap C_{+}(x,r_0)$ consists of $n_w$ points which determine  $(n_w+1)$ intervals on $C_{+}(x,r_0)$, each contained in one nodal domain.  Note that every nodal domain actually hits at least one of these intervals. This is indeed the case because  any nodal domain is bounded by nodal arcs, and the nodal set $\cZ(w)$ consists of loops at $x$, and possibly a nodal arc from $x$ to the boundary. These intervals are labeled $I^w_1, \ldots , I^w_{n_w+1}$, in a counter-clockwise manner along $C_{+}(x,r_0)$.
\smallskip

A labeling of the nodal domains of $w$ is a set $\cJ := \set{j_1, \ldots , j_k}$  in bijection with the set of nodal domains of $w$. To any labelling $\cJ$ of the nodal domains of $w$, we associate a word of length $(n_w+1)$, $m_{\cJ,w} =\ell_{\cJ,1} \cdots \ell_{\cJ,(n_w+1)}$. The letter $\ell_{\cJ,j}$ of $m_{\cJ,w}$ is the label, in the labeling $\cJ$, of the nodal domain of $w$ which contains the interval $I_j$. Since eigenfunctions change sign across a nodal line, two consecutive letters in a word $m_{\cJ,w}$ are distinct elements of $\cJ$. Different labelings $\cJ_1$ and $\cJ_2$ of the nodal domains of $w$ might give rise to different words. To compare nodal sets, we introduce the following invariant which does actually not depend of the labeling
%
\begin{equation}\label{E-invlab}
\min \set{j ~|~ j \ge 2 \text{~and~} \ell_{\cJ,j} = \ell_{\cJ,1}}\,.
\end{equation}
%
Note that the $\inf$ does exist because the number $k$ of nodal domains is less that the number $n_w$ of intervals.\smallskip


To label the nodal domains  of the family of eigenfunctions $w_{\theta}$, we use the following convention. For a given $\theta \in (0,\pi)$, the local structure theorem for the nodal set $\cZ(w_{\theta})$ at $x$ implies that for $r_0$ small enough, the  $(2k-3)$ distinct points in $ C_{+}(x,r_0) \cap \Omega \cap \cZ(w_{\theta})$ determine $(2k-2)$ intervals, each contained in a nodal domain.
The first $a_{(+)}$ intervals are contained in $\Omega_{\theta,+}$, the remaining ones in $\Omega_{\theta,-}$. Working anti-clockwise, the nodal domain which contains the first interval is labeled $1$. The nodal domain which contains the second interval is labeled $2$. If the nodal domain which contains the 3rd interval has not yet been labeled, we label it $3$, if it has already been labeled, we keep its label.We call this labeling $\cJ_{\theta}$. Doing so, the nodal domains contained in $\Omega_{a,+}$ meet one of the first $a$ intervals, and they are labeled from $1$ to $(a_{(+)}+1)$. The remaining nodal domains are contained in $\Omega_{a,-}$, meet the next $(2k-2-a)$ intervals, and are labeled from $(a_{(+)}+2)$ to $k$.

\begin{lemma}\label{L-hmn-210}
The labeling $\cJ_{\theta}$ of nodal domains is completely determined by the map $\tau^{v_2}$, and conversely. It does not depend on $\theta \in (0,\pi)$, and we denote it by $\cJ = \set{1,\ldots, k}$.
\end{lemma}%

\pf See Appendix~\ref{S-llnd}. \hfill \qed\medskip

To the function $w_{\theta}$, we associate the word $m_{\theta} := m_{\cJ_{\theta},w_{\theta}}$. This word has length $(2k-2)$. Each letter of the word $m_{\theta}$ is an element of the set $\cJ = \set{1,\ldots,k}$. More precisely, the word $m_{\theta}$ has the form
%
$$m_{\theta} = 1p_{+}1(a_{(+)}+2)p_{-}(a_{(+)}+2),$$
%
where $p_{+}$ is a word of length $|p_{+}|=(a-2)$ with letters in $\set{1,\ldots,(a_{(+)}+1)} \subset \cJ$, and $p_{-}$ a word of length $|p_{-}|=(2k-a-4)$ with letters in $\set{(a_{(+)}+2),\ldots,k} \subset \cJ$.} The word $1p_{+}1$ corresponds to the labeling of the nodal domains contained in $\Omega_{\theta,+}$; the word $(a_{(+)}+2)p_{-}(a_{(+)}+2)$ corresponds to the labeling of the nodal domains contained in $\Omega_{\theta,-}$.


\subsubsection{The rotating function argument}\label{SSS-hmn-25D}

Letting $\theta$ tend to zero, the nodal arc $\delta_{a,\theta}$ tends to a loop in the nodal set of $w_0 = \lim_{\theta \to 0}w_{\theta}$. The  location of the nodal domains of $w_0$  with respect to $C_{+}(x,r_0)$ is  described by the word of length $(2k-1)$, $m^{(0)}:=(a_{(+)}+2)m_{\theta}$. When $\theta$ tends to $\pi$, the nodal arc $\delta_{a,\theta}$ tends to a loop in the nodal set of $w_{\pi} = \lim_{\theta \to \pi}w_{\theta}$, and   the location of the nodal domains of $w_{\pi}$  with respect to $C_{+}(x,r_0)$ is described by the word of length $(2k-1)$, $m^{(\pi)}:= m_{\theta}1$.\smallskip

To compare these words, we look  at the invariant given in \eqref{E-invlab}, ie, at the first position at which the first letter of the word is repeated. For the word $m^{(0)}$ this position must be $4 + |p_{+}|$. For the word $m^{(1)}$ this position is $2 + |p_{+}|$.  These positions being different the nodal patterns must be different, and this contradicts the fact that $w_0 = v_1 = - w_{\pi}$.\smallskip

This contradiction completes the proof that $\dim U(\lambda_k) = (2k-1)$ cannot occur  when $\Omega$ is simply connected. \medskip



The above proof is illustrated by the following figures. In Figure~\ref{F-hmn2-Uxyw-L1}, $m_{\theta} = 121343$, $m^{(\pi)} = 1213431$, $m^{(0)} = 3121343$, $p_{+} = 2$, and $p_{-} = 5$. In Figure~\ref{F-hmn2-Uxyw-L2}, $m_{\theta} = 121314$, $m^{(\pi)} = 1213141$, $m^{(0)} = 4121314$, $p_{+}= 213$, and $p_{-}$ is empty.


\begin{figure}[!ht]
\centering
\includegraphics[width=0.98\linewidth]{hmn2-Uxyw-L1}
\caption{}\label{F-hmn2-Uxyw-L1}
\end{figure}

\begin{figure}[!ht]
\centering
\includegraphics[width=0.98\linewidth]{hmn2-Uxyw-L2}
\caption{}\label{F-hmn2-Uxyw-L2}
\end{figure}
\FloatBarrier

\begin{remark}Instead of looking at the labeling of nodal domains to reach a contradiction, we can look directly at the maps $\tau$. The map $\tau^{v_2} := \tau_{x,y_2}^{v_2}$ is given by
%
\begin{equation*}
\tau^{v_2} = \begin{pmatrix}
        \downarrow &  1 & \ldots & (a-1) & a & (a+1) & \ldots & (2k-3) \\
         a &  \tau^{v_2}(1) & \ldots & \tau^{v_2}(a-1) & \downarrow & \tau^{v_2}(a+1) & \ldots & \tau^{v_2}(2k-3) \\
         \end{pmatrix}\,,
\end{equation*}
%
where the first row represents the set $ \set{\downarrow} \cup L_{(2k-3)}$, the second row its image under the map $\tau^{v_2}$, where $a \in L_{(2k-3)}$ is the label of the nodal arc $\delta_a^{w_{\theta}}$ from $x$ to $y_{\theta}$. \smallskip

When $\theta$ tends to $0$, the map $\tau^{v_2}$ yields a map $\tau^{(0)}$ associated with the nodal set of $\lim_{\theta \to 0}w_{\theta}$. This map given by
%
\begin{equation*}
\tau^{(0)} = \begin{pmatrix}
0 &1 & \ldots & (a-1) & a & (a+1) & \ldots & (2k-3) \\
a &  \tau^{v_2}(1) & \ldots & \tau^{v_2}(a-1) & 0 & \tau^{v_2}(a+1) & \ldots & \tau^{v_2}(2k-3) \\
\end{pmatrix}\,,
\end{equation*}
%
where the first row represents a set with $(2k-2)$ elements  labeled from $0$ to $(2k-3)$, and the second row its image under $\tau^{(0)}$.\smallskip

When $\theta$ tends to $\pi$, the map $\tau^{v_2}$ yields a map $\tau^{(\pi)}$ associated with the nodal set of $\lim_{\theta \to \pi}w_{\theta}$. This map is given by
%
\begin{equation*}
\tau^{(\pi)} = \begin{pmatrix}
1 & \ldots & (a-1) & a & (a+1) & \ldots & (2k-3) & (2k-2)\\
\tau^{v_2}(1) & \ldots & \tau^{v_2}(a-1) & (2k-2) & \tau^{v_2}(a+1) & \ldots & \tau^{v_2}(2k-3) & a \\
\end{pmatrix}
\end{equation*}
%
 where the first row is $L_{(2k-3)}$.\medskip

 To compare these maps without relabeling the source set, we look at an invariant, namely $\max_j\set{|\tau(j) - j|}$, for both $\tau^{(0)}$ and $\tau^{(\pi)}$, so that we have to compare the numbers
%
\begin{equation*}
\max\set{a,(2k-4-a)} \text{~~and~~} \max\set{ (a-2),(2k-2-a)}\,.
\end{equation*}
%
These numbers are different unless $a=(k-1)$. In this case they are equal to
%
\begin{equation*}
(k-1) = \max\set{(k-1),(k-3)} = \max\set{(k-3),(k-1)}\,.
\end{equation*}
%
When $a=(k-1)$, we relabel the source set of $\tau^{(0)}$ to $\set{1,\ldots,(2k-2)}$ and we find that $\tau^{(0)}(1) = k$ and $1 < \tau^{(\pi)}(1) \le (k-2)$. This shows that the maps are different.\medskip

In any case the combinatorial patterns of the limit nodal sets are different which contradicts the fact that $w_0 = v_1 = - w_{\pi}$.
\end{remark}

\begin{remarks}\label{R-hmn-202}\phantom{}
\begin{enumerate}[(i)]
  \item Assume that $1 \le a \le (2k-3)$. Then, $a\, \frac{\pi}{m} < a\, \frac{\pi}{m-1} < (a+1)\, \frac{\pi}{m}$. Since $a$ is odd, there exists some $q$, $1 \le q \le a$ such that $\tau_1(q) > a$. This implies that the nodal arc $\delta_a$ of $\cZ(v_1)$ intersects the loop $\gamma^{(1)}_{q,\tau_1(q)}$ at at least one point $z_a$.
  \item Although $a_{\theta}$ does not depend on $\theta$, the nodal semi-arc of $\cZ(w_{\theta})$ emanating from $x$ tangentially to $\omega_{2,a}$ does depend on $\theta$, and hence denoted by  $\delta_{a,\theta}$.
  \item When $\theta$ varies, the arcs $\delta_{a,\theta}$ all pass through any point $z_a$ given in (i). Such points are \emph{common zeros} of the family of functions $w_{\theta}$.
\end{enumerate}
\end{remarks}%



\subsubsection{Proof that $\mult(\lambda_k) \le (2k -2)$ for $k \ge 3$, general case}\label{SS-hmn-26}\phantom{}


Under the assumption that $\dim U(\lambda_k) = (2k-1)$, the arguments in the simply connected case only use Euler's formula applied to nodal partitions $\cD_u$, Jordan's separation theorem, the structure of the nodal sets $\cZ(v_1)$ (a bouquet of loops at $x$) and $\cZ(v_2)$ (the wedge sum of an arc from $x$ to the boundary $\Gamma_1$ with one or two bouquets  of loops at $x$), and the fact that the combinatorial type  $\tau_{x,a_{\theta},w_{\theta}}$ of the nodal sets $\cZ(w_{\theta})$ is constant for $\theta \in (0,\pi)$.\smallskip

In the general case, Euler's formula leads to a similar structure for the nodal sets $\cZ(v_1)$, $\cZ(v_2)$, and $\cZ(w_{\theta})$, with ``generalized'' loops and arcs. We can now look at the projection of these sets to $\check{\Omega}$ as in Section~\ref{SS-hmn2-10}. As observed in Remark~\ref{R-hmn23-2}, the only singular points of the projected sets $\check{\cZ}(u)$ (or partitions $\check{\cD}_u$), $u \in \set{v_2, w_{\theta}}$, are the points $\check{x}$ and $ \check{y}_{\theta}$, and the combinatorial type of $\check{\cZ}(w_{\theta})$ is constant for $\theta \in (0,\pi)$. Since $\check{\Omega}$ is simply connected, we can now apply the same arguments as in the simply connected case.\medskip


We summarize Section~\ref{S-hmn2} in the following proposition.

\begin{proposition}[\cite{HoMN1999}]\label{P-hmn-s2}
Let $\Omega \subset \R^2$ be a bounded domain with smooth boundary (no assumption on the topology of $\Omega$). Let $\set{\lambda_k, k\ge 1}$ be the eigenvalues of the operator $-\Delta + V$ in $\Omega$, with Dirichlet or Robin boundary condition  (where $V$ is a real valued $C^{\infty}$ function). Then,
\begin{equation*}
\mult(\lambda_k) \le (2k-2)~~  \text{for~}k \ge 3  \,.
\end{equation*}
\end{proposition}%


In \cite[Theorem~B, p.~1172]{HoMN1999}, the authors state that the preceding bound can be improved to $\mult(\lambda_k) \le (2k-3)$ for all $k\ge 3$. For the proof, they propose the same strategy as for the proof of Proposition~\ref{P-hmn-s2}: assume that $\mult(\lambda_k) = (2k-2)$, construct $\lambda_k$-eigenfunctions with prescribed singular points, analyze their nodal sets, and derive a contradiction. However, we have so far not succeeded in writing down complete details for the arguments in \cite{HoMN1999}, even when $\Omega$ is simply connected (see \cite[Section~4]{Berd2018}). In Sections~\ref{S-hmn3} and \ref{S-hmn4}, we provide some properties of $\lambda_k$-eigenfunctions under the assumption that $\mult(\lambda_k) = (2k-2)$. In Section~\ref{S-hmn6}, we indicate the main steps of the proof sketched in \cite[Section~3]{HoMN1999}, and we point out where we are stuck.

\section{Plane domains, properties of $\lambda_k$-eigenfunctions assuming that $\mult(\lambda_k) = (2k-2)$}\label{S-hmn3}

 The purpose of this section is to derive properties of $\lambda_k$-eigenfunctions assuming that. the upper bound in Proposition~\ref{P-hmn-s2} is achieved. Throughout this section, we assume that
%
\begin{equation}\label{E-hmn-b}
k \ge 3 \text{~~and~~} \dim U(\lambda_k) = (2k-2)\,.
\end{equation}
%
We do not yet make any assumption on the topology of $\Omega$. For simplicity, we denote $U(\lambda_k)$ by $U$ and we use the notation of Subsection~\ref{SS-hmn-0N}.  According to Courant's nodal domain theorem, $\sup\set{\kappa(u) ~|~ u \in U} \le k$, so that $\ell = k$. \smallskip

For any $0 \neq u \in U$,  \eqref{E-euler-or2} becomes,
%
\begin{equation}\label{E-hmn-3a12}
\left\{
\begin{array}{ll}
0 \ge \kappa(u) - k & = \left[ b_0\left( \cZ(u) \cup \Gamma(u)\right) -1 \right]  + \frac 12 \, \sum_{z\in \cS_i(u)} \left(  \nu(u,z) \, -2\right)\\[5pt]
& \quad +\, \sum_{\substack{j\in J(u)}\,} \frac 12 \, \left(  \sum_{z\in \cS_b(u) \cap \Gamma_j}\rho(u,z) \, -2\right) - (k -2)\,.
\end{array}
\right.
\end{equation}

In this section, we analyze eigenfunctions with prescribed singular points, under the assumption \eqref{E-hmn-b}


\subsection{Eigenfunctions with one prescribed interior singular point}\label{SS-hmn-31i}

For $x \in \Omega$, define the subspace
%
\begin{equation}\label{E-hmn3-n6a}
U_x := \set{u \in U ~|~ \nu(u,x) \ge 2k - 2}\,.
\end{equation}

In view of the assumption \eqref{E-hmn-b}, Lemma~\ref{L-zeroi} implies that $U_x \neq \set{0}$. \smallskip

The following lemma, appears as  Lemma~3.7 in \cite[p.~1183]{HoMN1999}.

\subsubsection{Properties of $U_x$}\label{SSS-hmn-31ia}

\begin{lemma}\label{L-L37}
Let $x \in \Omega$ and $U := U(\lambda_k)$. Assume that $\dim U = (2k-2)$.  The subspace $U_x$ defined in \eqref{E-hmn3-n6a} has the following properties.
\begin{enumerate}[(i)]
  \item  The dimension of $U_x$ is $1$.
  \item For all $0 \neq u \in U_x$,
  \begin{equation}\label{E-hmn3-n6b}
  \left\{
\begin{array}{l}
\kappa(u) = k\,,\\[5pt]
\cZ(u) \cup \Gamma(u) \text{~~is connected,}\\[5pt]
\cS_i(u) = \set{x} \quad\text{and~~} \nu(u,x) = 2(k -1)\,,\\[5pt]
\sum_{z \in \cS_b(u)\cap \Gamma_j} \rho(u,z) = 2 \quad \text{for all~~} j \in J(u)\,.\\[5pt]
\end{array}
\right.
  \end{equation}
\item If $u_x$ is a generator of $U_x$, the map $\Omega \ni x \mapsto [u_x] \in \bP(U)$ is $C^{\infty}$.
\end{enumerate}%
\end{lemma}


\pf  We already know that that $\dim U_x \ge 1$. \smallskip

\emph{Proof of Assertion~(ii).}

\noih We first assume that $\Omega$ is simply connected.\smallskip

The assumptions of the lemma and \eqref{E-hmn-3a12} imply that
%
\begin{equation}\label{E-hmn-312}
\begin{split}
0 \geq  \kappa(u) - k & = \left( b_0(\cZ(u) \cup \partial \Omega) -2\right) + \frac 12\, \sum_{z \in \cS_i(u), z\neq x } (\nu(u,z)-2) \\
&\hspace{5mm} + \frac 12 \, \left( \nu(u,x) - 2k +2\right)  + \frac 12 \sum_{z\in \cS_b(u) }\rho (u,z) \,.
\end{split}
\end{equation}
%
The terms in the right-hand side are nonnegative, except possibly the first one. The inequality implies that $b_0(\cZ(u) \cup \partial\Omega) \le 2$. We now consider two cases.\smallskip

\noid If $b_0(\cZ(u)\cup \partial\Omega) = 2$, the terms in the right-hand side are nonnegative, with a nonpositive sum. They must all vanish: $\kappa(u) = k$, $\cS_i(u) = \set{x}$, $\cS_b(u) = \emptyset$, and $\nu(u,x) = 2(k - 1)$. In this case, $\cZ(u) \cap \partial \Omega = \emptyset$ and $\cZ(u)$ is connected. \smallskip

\noid If $b_0(\cZ(u)\cup \partial\Omega) = 1$, the nodal set $\cZ(u)$ must hit $\partial \Omega$, which implies that $\sum_{z\in \cS_b(u) }\rho (z) \ge 2$ (use Proposition~\ref{P-euler-np}). Re-arranging the inequality, we conclude that $\kappa(u) = k$, $\cS_i(u) = \set{x}$, $\nu(u,x) = 2(k - 1)$, and $\sum_{z\in \cS_b(u) }\rho (z) = 2$. \medskip

\noih Assume now that $\partial \Omega = \cup_{j=1}^q \Gamma_j$. Let $u$ be a function such that $\nu(u,x) \ge 2(k - 1)$. With the notation of  Subsection~\ref{SS-hmn-0N}, the inequality \eqref{E-hmn-3a12} becomes
%
\begin{equation}\label{E-hmn-312b}
\begin{split}
0 \ge \kappa(u) - k &= \left( b_0(\cZ(u) \cup \Gamma(u)) - 1\right) + \frac 12\, \sum_{z \in \cS_i(u), z\not = x}(\nu(u,z) - 2) \\[5pt]
&\hspace{5mm} + \sum_{j \in J(u)\,}  \frac 12\,  \left(\sum_{z\in \cS_b(u) \cap \Gamma_j}\rho(u,z) \, - 2\right) + \frac 12\, \left( \nu(u,x) - 2k +2\right)\,.
\end{split}
\end{equation}
%
The terms in the right-hand side of \eqref{E-hmn-312b} are nonnegative and their sum is nonpositive. They must all vanish:
%
\begin{equation}\label{E-hmn-312c}
\left\{
\begin{array}{l}
\kappa(u) = k \text{~~and~~} b_0(\cZ(u) \cup \Gamma(u)) = 1\,,\\[5pt]
\cS_i(u) = \set{x} \text{~~and~~} \nu(u,x) = 2(k - 1)\,,\\[5pt]
\sum_{z\in \cS_b(u) \cap \Gamma_j}\rho(u,z) = 2\,, \text{~for all~} j \in J(u)\,.
\end{array}
\right.
\end{equation}
%
We have proved Assertion~(ii).\medskip

\emph{Proof of Assertion~(i).~} Lemma~\ref{L-zeroc} implies that $\dim U_x \le 2$. Assume by contradiction that $\dim U_x = 2\,$.  We again use a \emph{rotating function argument} similar to the one used in  Subsection~\ref{SS-h2n-s2}.\smallskip

 As in Proposition~\ref{propertyD2}, one can choose a basis $\set{v_1,v_2}$ of $U_x$ such that, in local polar coordinates centered at $x$,
%
\begin{equation*}
\left\{
\begin{array}{l}
v_1 = r^{k -1} \sin ((k -1)\omega) +\mathcal O (r^{k})\,,\\
v_2 = r^{k -1} \cos ((k -1)\omega) +\mathcal O (r^{k})\,.
\end{array}
\right.
\end{equation*}

Introducing the family of functions
%
\begin{equation*}
w_\theta = \cos((k -1) \theta) \,v_1 - \sin((k -1) \theta) \, v_2\,,
\end{equation*}
%
and letting $\theta$ tend to $0$ or $\frac{\pi}{(k - 1)}$, one can follow the arguments given in the proofs of  Properties~\ref{propertyD3} and \ref{propertyD4} to reach a contradiction.\medskip

\emph{Proof of Assertion~(iii).~} The function $u_x$ is characterized by the vanishing of the derivatives of order less than or equal to $(k -2)$. Taking a basis for $U$, this is equivalent to writing a linear system. Since the function $u_x$ is uniquely determined up to scaling, this linear system has constant rank, and we have a locally defined smooth map $x \mapsto u_x$ in the neighborhood of any given point $x_0$. This means that we have a smooth map $x \mapsto [u_x]$ from $\Omega$ into the projective space of $U$. \hfill \qed


\subsubsection{Structure and combinatorial type of nodal sets in $U_x$}\label{SSS-hmn-31ias}

%\begin{remark}\label{R-hmn3-n2}

Using Assertion~(ii), one can describe the possible nodal patterns for a generator $u_x$ of $U_x$.  If $\partial \Omega = \Gamma_1$ (i.e., $\Omega$ is simply connected), there are two cases.
\begin{enumerate}
   \item Either $\cZ(u_x)$ consists of $(k - 1)$ loops at $x$ which do not intersect away from $x$, and do not hit $\Gamma_1$.
   \item Or $\cZ(u_x)$ it consists of
   \begin{itemize}
     \item[$\diamond$] $(k - 2)$ loops at $x$ which do not hit the boundary, and
     \item[$\diamond$] two arcs emanating from $x$ and hitting $\Gamma_1$ at points $y_1 \neq y_2$, such that $\rho(u_x,y_i) = 1$ or, possibly, at one point $y$, with $\rho(u_x,y) = 2$.
   \end{itemize}
Furthermore, the loops at $x$ and the arcs from $x$ to the boundary are pairwise disjoint away from $x$.
In this case, we have a ``generalized'' nodal loop at $x$ which consists of the two arcs, and a portion of the boundary.
\end{enumerate}


When $\partial \Omega = \Gamma_1 \cup \ldots \cup \Gamma_q$ ($q \ge 2$, i.e. $\Omega$ not simply connected) the nodal set $\cZ(u_x)$ consists of arcs and loops which can be ``generalized arcs or loops'' in the sense that they may include some arcs contained in one or several of the components $\Gamma_j, j \in J(u)$. Such arcs and loops become ``ordinary'' arcs and loops once projected to $\check{\Omega}$ as in Subsection~\ref{SS-hmn2-10},  see Remark~\ref{R-hmn23-2}. Viewing $\Omega$ as a domain on the sphere, we can also identify  each component of $\partial \Omega$ to a point, and consider the quotient space $\bS_{\Omega}$, see Notation~\ref{N-hmn210-2}. The nodal sets projected from $\Omega$ to $\bS_{\Omega}$ then consist of $(k-1)$ simple loops at $\check{x}$, intersecting only at $\check{x}$, as in the closed case, Subsection~\ref{SS-h2n-s2}.

\begin{figure}[!ht]
\centering
\begin{subfigure}[t]{0.25\textwidth}
\centering
\includegraphics[width=\linewidth]{hmn3-h0-L37-a}
\caption{}
\end{subfigure}
\begin{subfigure}[t]{.25\textwidth}
\centering
\includegraphics[width=\linewidth]{hmn3-h0-L37-b}
\caption{}
\end{subfigure}
\begin{subfigure}[t]{.25\textwidth}
\centering
\includegraphics[width=\linewidth]{hmn3-h0-L37-c}
\caption{}
\end{subfigure}
%
\caption{$\Omega$ simply connected, some nodal patterns for $u_x$}\label{F-hmn3-h0-L37}% ex : F-hmn3-i-h0
\end{figure}

\begin{figure}[!ht]
\centering
\begin{subfigure}[t]{.25\textwidth}
\centering
\includegraphics[width=\linewidth]{hmn3-h1-L37-a}
\caption{}
\end{subfigure}
\begin{subfigure}[t]{.25\textwidth}
\centering
\includegraphics[width=\linewidth]{hmn3-h1-L37-b}
\caption{}
\end{subfigure}
\begin{subfigure}[t]{.25\textwidth}
\centering
\includegraphics[width=\linewidth]{hmn3-h1-L37-c}
\caption{}
\end{subfigure}
%
\caption{$\Omega$ with one hole, some nodal patterns for $u_x$}\label{F-hmn3-h1-L37}% ex : F-hmn3-i-h1
\end{figure}
\FloatBarrier

Figures~\ref{F-hmn3-h0-L37} and \ref{F-hmn3-h1-L37} display some possible nodal patterns for $u_x$.

\begin{remark}\label{R-Fgene2}
The \emph{nodal patterns} displayed in Figure~\ref{F-hmn3-h0-L37} and \ref{F-hmn3-h1-L37} are valid for both the Dirichlet and Robin boundary conditions. Unless otherwise stated this remark applies to all figures of this section.
\end{remark}


\subsection{Eigenfunctions with one prescribed boundary singular point}\label{SS-hmn-31b}\phantom{}

For $y \in \Gamma_1$, define the subspace
%
\begin{equation}\label{E-hmn3-n10a}
W_y := \set{u \in U ~|~ \rho(u,y) \ge 2k - 3}\,.
\end{equation}

In view of the assumption \eqref{E-hmn-b}, Lemma~\ref{L-zero1} implies that $W_y \neq \set{0}$. \smallskip

Lemmas~\ref{L-L32} and~\ref{L-L33} below expand and explain the lemma stated in \cite[\S3.2]{HoMN1999} and their Lemma~3.3. In particular, Lemma~\ref{L-L33}, Assertion~(iv), explains and proves the statement ``the nodal pattern of $w_y$ is constant along each connected components of $\Gamma_{1,(2\ell -3)}$'', and Assertion~(iii) states the continuity of the map $y \mapsto z(y)$ (see \cite[p.~1179]{HoMN1999}, last sentence in Lemma~3.3, and second sentence after the same lemma).

\subsubsection{Properties of $W_y$}\label{SSS-hmn-31ba}

\begin{lemma}\label{L-L32}
Let $y \in \Gamma_1$ and $U:= U(\lambda_k)$.  Assume that $\dim U = 2(k-1)$.
 The subspace $W_y$ defined in \eqref{E-hmn3-n10a} has the following properties.
\begin{enumerate}[(i)]
  \item  The dimension of $W_y$ is $1$.
  \item For all $0 \neq u \in W_y$,
%
  \begin{equation}\label{E-hmn3-n10b}
  \left\{
\begin{array}{l}
\kappa(u) = k \text{~~and~~} \cZ(u) \cup \Gamma(u) \text{~is connected,}\\[5pt]
\cS_i(u) = \emptyset\,,\\[5pt]
\sum_{z \in \cS_b(u)\cap \Gamma_j} \rho(u,z) = 2 \quad \text{for all~~} j \in J(u), ~~ j \neq 1\,,\\[5pt]
 \sum_{z \in \cS_b(u)\cap \Gamma_1} \rho(u,z) = 2k - 2\,.\\[5pt]
\end{array}
\right.
  \end{equation}
%
Furthermore,
%
  \begin{equation}\label{E-hmn3-n10d}
\left\{
\begin{array}{ll}
\text{either\quad} & \rho(u,y) = 2k -2 \text{~~and~~} \cS_b(u)\cap \Gamma_1 = \set{y}\,,\\[5pt]
\text{or\quad} & \rho(u,y) = 2k -3 \text{~~and~~} \cS_b(u)\cap \Gamma_1 = \set{y,z(y)},\\[5pt]
&\qquad \text{~~for some~} z(y) \in \Gamma_1, z(y) \neq y, \text{~with~} \rho(u,z(y)) = 1.
\end{array}
\right.
\end{equation}
  \item If $w_y$ denotes a generator of $W_y$, then the map $\Gamma_1 \ni y \mapsto [w_y] \in \bP(U)$ is $C^{\infty}$.
\end{enumerate}
\end{lemma}


\pf We already know that $\dim W_y \ge 1$. \smallskip

\noi \emph{Proof of Assertion~(ii)}.\smallskip

\noih Case~1:~ $\Omega$ is simply connected.\smallskip

Choose a function $0 \neq u \in W_y$, and apply the inequality \eqref{E-hmn-3a12} to obtain,
%
\begin{equation}\label{E-hmn-314}
\begin{split}
0 \geq  \kappa(u) - k & = \left( b_0(\cZ(u) \cup \partial \Omega) -1 \right) + \frac 12\, \sum_{z \in \cS_i(u)} (\nu(u,z)-2) \\
&\hspace{5mm} + \frac 12 \left( \sum_{z\in \cS_b(u)}\rho (u,z) - 2k +2\right) \,.
\end{split}
\end{equation}
%
Since $\rho(u,y) \ge 2k -3$, Proposition~\ref{P-euler-np} implies that the last term in \eqref{E-hmn-314} is nonnegative; all the terms in the right-hand side are nonnegative, with nonpositive sum, and hence they must all vanish. This proves \eqref{E-hmn3-n10b} in the special case $\partial \Omega = \Gamma_1$. Looking at the two possible cases, $\rho(u,y) = (2k - 2)$ or $(2k -3)$, we obtain \eqref{E-hmn3-n10d}. \smallskip

\noih Case~2:~ $\partial \Omega = \cup_{j=1}^q \Gamma_j$, with $y \in \Gamma_1$. Let $0 \neq u \in W_y$. With the notation of   Subsection~\ref{SS-hmn-0N}, we can write
%
\begin{equation}\label{E-hmn-318a}
\begin{split}
0 \ge \kappa(u) - k &= \left( b_0(\cZ(u) \cup \Gamma(u)) - 1\right) + \frac 12\, \sum_{z \in \cS_i(u)} (\nu(u,z)-2)  \\[5pt]
&\hspace{3mm}+ \sum_{j \in J(u), j\not = 1} \frac 12\, \left( \sum_{z\in \cS_b(u) \cap \Gamma_j}\rho(u,z) \, - 2\right)\\[5pt]
&\hspace{3mm} + \frac 12\, \left( \sum_{z\in \cS_b(u) \cap \Gamma_1}\rho(u,z) - 2k +2\right)\,.
\end{split}
\end{equation}

The definition of $J(u)$, and Proposition~\ref{P-euler-np}, imply that the terms in the right-hand side of \eqref{E-hmn-318a} are nonnegative, with a nonpositive sum. They must all be zero and hence,  Equations~\eqref{E-hmn3-n10b} and \eqref{E-hmn3-n10d} hold. We have proved Assertion~(ii).\medskip

\noi \emph{Proof of Assertion~(i).~} We already know that $\dim W_y  \ge 1$. Assume that there exist at least two linearly independent functions $w_1$, $w_2$ in $W_y$. By \eqref{E-hmn3-n10b}, we have $2k- 3 \le \rho(w_i,y) \le 2k - 2$.  If $\rho(w_1,y) = \rho(w_2,y) = 2k - 3$, by Lemma~\ref{L-zeroc} there exists some linear combination $w$ of $w_1$ and $w_2$ such that $\rho(w,y) = 2k -2$. This function $w$ must satisfy \eqref{E-hmn3-n10d} and hence, by Lemma~\ref{L-zeroc}, is uniquely defined (up to scaling). If $\rho(w_1,y) = 2k - 2$, then we must have $\rho(w_2,y) = 2k -3$ since $w_1$ is uniquely defined. Any other function in $W_y$ must be a linear combination of $w_1$ and $w_2$. It follows that $\dim W_y \le 2$.\smallskip

 Assume that $\dim W_y = 2$, and choose a basis $\set{w_1,w_2}$ of $W_y$, with $\rho(w_1,y) = 2k -2$,  $\rho(w_2,y) = 2k -3$, and let $y_2$ be the unique other critical zero of $w_2$ on $\partial \Omega$. We have to consider two cases, depending on whether $U$ is a Dirichlet or Robin eigenspace, and we can use a \emph{rotating function argument} as in  Subsections~\ref{SS-hmn-24} and \ref{SS-hmn-25} to conclude that $\dim W_y = 1$. The claim is proved. This completes the proof of Assertion~(i). \hfill \qedc \medskip

\emph{Proof of Assertion~(iii).}~ Using a basis for the eigenspace $U$, the condition that a generator $w_y$ of $W_y$ satisfies $\rho(w,y) \ge 2k -3$ is equivalent to a linear system. Since $\dim W_y = 1$, this linear system has constant rank when $y$ varies and we can deduce the existence of a smooth map $y \mapsto [w_y]$ from $\partial \Omega$ into the projective space $\mathbb{P}(U)$ of the eigenspace $U$, where the generator $w_y$ of $W_y$ satisfies \eqref{E-hmn3-n10b} and \eqref{E-hmn3-n10d}. \hfill \qed


\subsubsection{Structure and combinatorial type of nodal sets in $W_y$}\label{SSS-hmn-31bas}

Using Assertion~(ii) of the lemma, one can describe the possible nodal patterns of a generator $w_y$ of $W_y$, as we did in Paragraph~\ref{SSS-hmn-31ias}.
For example, in the simply connected case, if $\rho(w_y,y) = (2k-3)$, the nodal set $\cZ(w_y)$ consists of $(k-2)$ simple loops at $y$, and a simple arc from $y$ to $z(y) \in \Gamma_1$;  if $\rho(w_y,y) = (2k-2)$, the nodal set $\cZ(w_y)$ consists of $(k-1)$ simple loops at $y$. The loops and the arc do not intersect away from $y$. Figures~\ref{F-hmn3-h0-L33} and \ref{F-hmn3-h1-L33} display some possible nodal patterns. Figure~\ref{F-hmn3-h1-L33-p} displays the nodal pattern in Figure~\ref{F-hmn3-h1-L33}(A) projected to $\check{\Omega}$, or viewed on $\bS^2$.


\begin{figure}[!ht]
\centering
\begin{subfigure}[t]{.30\textwidth}
\centering
\includegraphics[width=\linewidth]{hmn3-h0-L33-a}
\caption{}
\end{subfigure}
\begin{subfigure}[t]{.30\textwidth}
\centering
\includegraphics[width=\linewidth]{hmn3-h0-L33-b}
\caption{}
\end{subfigure}
%
\caption{$\Omega$ simply connected, some nodal patterns for $w_y$}\label{F-hmn3-h0-L33}% ex : F-hmn3-b-h0
\end{figure}

\begin{figure}[!ht]
\centering
\begin{subfigure}[t]{.30\textwidth}
\centering
\includegraphics[width=\linewidth]{hmn3-h1-L33-a}
\caption{}
\end{subfigure}
\begin{subfigure}[t]{.30\textwidth}
\centering
\includegraphics[width=\linewidth]{hmn3-h1-L33-b}
\caption{}
\end{subfigure}
\begin{subfigure}[t]{.30\textwidth}
\centering
\includegraphics[width=\linewidth]{hmn3-h1-L33-c}
\caption{}
\end{subfigure}
%
\caption{$\Omega$ with one hole, some nodal patterns for $w_y$}\label{F-hmn3-h1-L33}% ex : F-hmn3-b-h1
\end{figure}

\begin{figure}[!ht]
\centering
\begin{subfigure}[t]{.35\textwidth}
\centering
\includegraphics[width=\linewidth]{hmn3-h1-L33-ac}
%\caption{Figure~\ref{F-hmn3-h1-L33} viewed in $\check{\Omega}$}
\end{subfigure}
\begin{subfigure}[t]{.35\textwidth}
\centering
\includegraphics[width=\linewidth]{hmn3-h1-L33-as}
%\caption{Figure~\ref{F-hmn3-h1-L33} viewed in $\bS^2$}
\end{subfigure}
%
\caption{Figure~\ref{F-hmn3-h1-L33} viewed in $\check{\Omega}$ (left) or in $\bS^2$ (right)}\label{F-hmn3-h1-L33-p}% ex : F-hmn3-b-p
\end{figure}



In view of Lemma~\ref{L-L32}, define the following subsets of $\Gamma_1$\,:
%
\begin{equation}\label{E-E-hmn3-n12}
\left\{
\begin{array}{l}
\Gamma_{1,(2k -3)} := \set{y \in \Gamma_1 ~|~ \rho(w_y,y) = 2k - 3}\,,\\[5pt]
\Gamma_{1,(2k -2)} := \set{y \in \Gamma_1 ~|~ \rho(w_y,y) = 2k - 2}\,.
\end{array}
\right.
\end{equation}



Let $p \in \set{(2k-3),(2k-2)}$. For a given $y \in \Gamma_1$, we apply Subsection~\ref{SS-hmn2-20} to a generator $w_y$ of $W_y$.

For both Dirichlet or Robin boundary condition, in a neighborhood of $y$,
the nodal set $\cZ(w_y)$ consists of $p$ nodal semi-arcs. As in Paragraph~\ref{SSS-hmn-25B}, for a given $j \in L_p$, we follow the
nodal semi-arc emanating from $y$ tangentially to the rays $\omega_j$
along $\cZ(w_y)$.  There are two cases.
\begin{enumerate}[$\diamond$]
\item Case $p=(2k-2)$, i.e. $y \in \Gamma_{1,(2k -2)}$. According to Lemma~\ref{L-L32}, we eventually arrive back at $y$ along a nodal semi-arc emanating from another ray which we denote by $\omega_{\tau_y^{W}(j)}$. This uniquely defines a map $\tau_y^{W} : L_{(2k-2)} \to L_{(2k-2)}$, such that $\tau_y^{W}(j) \neq j$ and $(\tau_y^{W})^2 = \id$. In this case, the nodal set $\cZ(w_y)$ consists of a $(k-1)$-bouquet of loops at $y$, denoted by $\gamma^{w_y}_{j,\tau_y^{W}(j)},\, j \in L_{(2k-2)}$.
  \item Case $p = (2k-3)$, i.e. $y \in \Gamma_{1,(2k -3)}$. We define a map $\tau_y^{W}$ from $\set{\downarrow} \cup L_{(2k-3)}$ to itself as follows. According to Lemma~\ref{L-L32}, there exists a unique $a \in L_{(2k-3)}$ (depending on $y$) such that the nodal semi-arc emanating from $y$ tangentially to the ray $\omega_{a}$ eventually hits $\Gamma_1$ at some $z(y) \neq y$. We let $\tau_y^{W}(a) = \,\downarrow$\, and $\tau_y^{W}(\downarrow)=a$. For $j \neq a$, following the nodal semi-arc, denoted by $\delta_a$, emanating from $y$ tangentially to $\omega_j$ along $\cZ(w_y)$, we will eventually arrive back at $y$ along another ray denoted by $\omega_{\tau_{y}^{W}(j)}$. This uniquely defines a map $\tau_{y}^{W}$ from $\set{\downarrow} \cup L_{(2k-3)}$ to itself such that $\tau_{y}^{W}(j) \neq j$ and $(\tau_{y}^{W})^2 = \id$. In this case, the nodal set $\cZ(w_y)$ is the wedge sum of the arc $\delta_{a}$ from $y$ to $z(y)$ and a $(k-2)$-bouquet of loops denoted by $\gamma{j,\tau_{y}^{W}(j)}$.
\end{enumerate}
In both cases, the loops and arc might be ``generalized'' in the sense that they may contain sub-arcs of $\partial \Omega \sm \Gamma_1$ (when $\Omega$ is not simply connected); they can then be viewed as loops and arc when projected to $\check{\Omega}$. When $\Omega$ is simply connected and $y \in \Gamma_{1,(2k -3)}$, the $(k-2)$-bouquet of loops actually consists of one or two bouquets of loops, one in each connected component of $\Omega \sm \delta_{\alpha(y)}$.\smallskip

 As in Paragraph~\ref{SSS-hmn-25B}, we define the map $\tau_y^{W}$ as the \emph{combinatorial type} of the nodal set $\cZ(w_y)$ at $y$, when $y \in \Gamma_{1,(2k-2)}$, resp. $y \in \Gamma_{1,(2k-3)}$. The source-set of $\tau_{y}^{W}$ is either $L_{(2k-2)}$ or $\set{\downarrow} \cup L_{(2k-3)}$.


\subsubsection{Properties of $\Gamma_{1,(2k-3)}$ and $\Gamma_{1,(2k-2)}$}\label{SSS-hmn-31bc}

\begin{lemma}\label{L-L33}
Let $\Omega \subset \R^2$ be any smooth bounded domain, and $U := U(\lambda_k)$. Assume that $\dim U = (2k-2)$. Then, the following properties hold.
\begin{enumerate}[(i)]
  \item The sets $\Gamma_{1,(2k-3)}$ and $\Gamma_{1,(2k-2)}$ are disjoint and
      $$\Gamma_1 = \Gamma_{1,(2k-3)} \bigsqcup \Gamma_{1,(2k-2)}\,.$$
  \item The set $\Gamma_{1,(2k-3)}$ is open in $\Gamma_1$; the set $\Gamma_{1,(2k-2)}$ is finite.
  \item The map $\Gamma_{1,(2k -3)} \ni y \mapsto z(y) \in \Gamma_1$, where $z(y)$ is defined in \eqref{E-hmn3-n10d} is continuous in $\Gamma_{1,(2k -3)}$, and  for any $\eta \in \Gamma_{1,(2k -2)}$ $\lim_{y \to \eta, y \in \Gamma_{1,(2k-3)}}z(y) = \eta$.
  \item   The  combinatorial type of $w_y$ is constant in any component of $\Gamma_{1,(2k-3)}$ in the sense that the maps $y \mapsto \tau_{y}^{W}(\downarrow)$ and $y \mapsto \tau_{y}^{W}$ are constant in each connected component of $\Gamma_{1,(2k-3)}$.
\end{enumerate}
\end{lemma}%


\pf Assertion~(i) follows from Lemma~\ref{L-L32}. \medskip

\emph{Proof of Assertion~(ii).~} The first part follows from the fact that $\rho(w_y,y) = (2k -3)$ is an open condition on the vanishing of derivatives,  see Lemma~\ref{L-breve}. \smallskip

For the second part, let $y_0 \in \Gamma_{1,(2k -2)}$. By a conformal change of coordinates (see for example Section~2 in \cite{YaZh2021} and Paragraph~\ref{SS-hmn2-20}), we can assume that a neighborhood of the point $y_0$ in $\overline{\Omega}$ is sent to the half-disk $\set{\xi_1^2+\xi_2^2 < 1, \xi_2\ge 0}$, with $y_0$ sent to $(0,0)$.  By abuse of notation, we identify $\Gamma_1$ to $(-1,1)\times\set{0}$ in this neighborhood of $y_0$. Let $t \in (-1,1)$ be a point corresponding  to some $y \in \Gamma_{1}$ close to $y_0$.\smallskip

According to Lemma~\ref{L-L32}, in a neighborhood  of $ (0,0)$ in $ \Gamma_1$, we have a $C^{\infty}$ family of functions $ w_t := w_{(t,0) }\in W_{(t,0)}$ in the form
%
\begin{equation}\label{E-L33-2}
w_t = \sum_{j=1}^{2k-2} a_j(t) \, \phi_j\,,
\end{equation}
%
where $\set{\phi_j\,, 1 \le j \le 2k-2}$ is a basis of $U$, and $\rho(w_t,(t,0)) = (2k-3)$ if $(t,0) \in \Gamma_{1,(2k-3)}$, or $\rho(w_t,(t,0)) = (2k-2)$ if $(t,0) \in \Gamma_{1,(2k-2)}$. Let
%
\begin{equation*}%\label{E-L33-4}
q := \left\{
        \begin{array}{l}
          (2k-2) \text{~in the Dirichlet case,}\\[5pt]
          (2k-3) \text{~in the Robin case,}
        \end{array}
      \right.
\end{equation*}

We write Taylor's formula for $w_t$ at the point $(t,0)$ in the coordinates $(\xi_1,\xi_2)$,
%
\begin{equation}\label{E-L33-6}
\left\{
\begin{array}{ll}
w_t(\xi_1,\xi_2) & = \sum_{|\alpha| = q} \frac{1}{\alpha!} \, \partial^{\alpha}w_t(t,0) \, \left( \xi_1 - t, \xi_2\right)^{\alpha}\\[5pt]
&\quad  + \sum_{|\alpha| = q+1} \frac{1}{\alpha!} \, \partial^{\alpha}w_t(t,0) \, \left( \xi_1 - t, \xi_2\right)^{\alpha}\\[5pt]
&\quad  + \sum_{|\beta| = q+2} R_{\beta}(t;\xi_1,\xi_2)\, \left( \xi_1 - t, \xi_2\right)^{\beta}\,,
\end{array}
\right.
\end{equation}
%
where
%
\begin{equation*}
R_{\beta}(t;\xi_1,\xi_2) = \frac{q+2}{\beta!} \,\int_0^1 (1-s)^{q+1} \partial^{\beta}w_t\left( t + s(\xi_1-t),\xi_2\right) \, ds\,.
\end{equation*}

It follows from \eqref{E-L33-2} that the functions $\partial^{\alpha}w_t(t,0)$ are $C^{\infty}$ in $t \in (-1,1)$, and that the functions $R_{\beta}(t;\xi_1,\xi_2)$ are $C^{\infty}$ in $(-1,1)\times \set{\xi_1^2+\xi_2^2 < 1, \xi_2\ge 0}$. Note that the first sum vanishes identically if $(t,0) \in \Gamma_{1,(2k-2)}$.\medskip

We can rewrite \eqref{E-L33-6} as
%
\begin{equation}\label{E-L33-8}
\left\{
\begin{array}{ll}
w_t(\xi_1,\xi_2) & = \sum_{|\alpha| = q} A_{0,\alpha}(t) \, \left( \xi_1 - t, \xi_2\right)^{\alpha}\\[5pt]
 &\quad  + \sum_{|\alpha| = q+1} A_{1,\alpha}(t) \, \left( \xi_1 - t, \xi_2\right)^{\alpha}\\[5pt]
& \quad  + \sum_{|\beta| = q+2} R_{\beta}(t;\xi_1,\xi_2)\, \left( \xi_1 - t, \xi_2\right)^{\beta}\,,
\end{array}
\right.
\end{equation}
%
where the functions $A_{0,\alpha}$ and $A_{1,\alpha}$ are $C^{\infty}$. Call $p_0(t,\xi_1,\xi_2)$ (resp. $p_1(t,\xi_1,\xi_2)$) the polynomial of degree $q$ (resp. $(q+1)$), with coefficients $A_{0,\alpha}$ (resp. $A_{1,\alpha}$).\smallskip

To prove that the set $\Gamma_{1,(2k -2)}$ is finite, it suffices to prove that it is discrete. Assume that the point $y_0$ is not isolated. Then, there exists a sequence $\set{t_n}$ tending to zero, such that $w_{t_n}$ satisfies $\rho(w_{t_n},(t_n,0)) = (2k-2)$ for all $n$, so that $p_0(t_n,\xi_1,\xi_2)$ is the zero polynomial.  According to Remark~\ref{R-lsns-2}, the homogeneous polynomials $p_0$ and $p_1$ are harmonic. Define
%
\begin{equation*}%\label{E-L33-10}
\left\{
\begin{array}{ll}
\varpi_n(\xi_1,\xi_2) & := t_n^{-1}\, \left( w_{t_n}(\xi_1,\xi_2) - w_{0}(\xi_1,\xi_2)\right)\\[5pt]
 & =  t_n^{-1}\, \left( p_1(t_n;\xi_1-t,\xi_2) - p_1(0;\xi_1,\xi_2)\right) \\[5pt]
& \hspace{-8mm}+ \sum_{|\beta|=q+2}t_n^{-1}\, \left( R_{\beta}(t_n;\xi_1-t,\xi_2)\, (\xi_1-t,\xi_2)^{\beta} - R_{\beta}(0;\xi_1,\xi_2)\,(\xi_1,\xi_2)^{\beta}\right).
\end{array}
\right.
\end{equation*}

The functions $\varpi_n$ belong to $U$. When $t_n$ tends to $0$, they converge to the function $ \varpi \in U$ given by
%
\begin{equation}\label{E-L33-12}
\left\{
\begin{array}{ll}
\varpi(\xi_1,\xi_2) = & \sum_{|\alpha|=q+1} A'_{1,\alpha}(0) \, (\xi_1,\xi_2)^{\alpha}\\[5pt]
 &- \sum_{|\alpha|=q+1} \alpha_1 A_{1,\alpha}(0) \, (\xi_1,\xi_2)^{\alpha -(1,0)}\\[5pt]
&+ \sum_{|\beta|=q+2} \partial_t B_{\beta}(0;\xi_1,\xi_2)\, (\xi_1,\xi_2)^{\beta}\\[5pt]
&- \sum_{|\beta|=q+2} \beta_1 B_{\beta}(0;\xi_1,\xi_2)\,
(\xi_1,\xi_2)^{\beta - (1,0)}\\[5pt]
\end{array}
\right.
\end{equation}

 The harmonic polynomial $p_1(0;\xi_1,\xi_2)$ is nonzero because it corresponds to the leading term in $w_0$. Hence, its derivative $\partial_{\xi_1}p_1$ is not zero  (otherwise, being harmonic, $p_1$ would be a polynomial in $\xi_2$, of degree one, contradicting our assumption that $k \ge 3$).  Looking at the order of vanishing of the terms in the right-hand side, the second line in \eqref{E-L33-12} shows that $w \neq 0$, and that $\rho(w,y_0) = (2k-2)$, while $\rho(w_{y_0},y_0) = (2k-3)$. This contradicts the fact that $\dim W_{y_0} = 1$. This proves that the point $y_0$ is isolated in $\Gamma_{1,(2k-2)}$, and since this is true for any point in $\Gamma_{1,(2k-2)}$, this set is discrete, hence finite. This proves Assertion~(ii). \hfill \qedc \smallskip

\emph{Proof of Assertion~(iii).~} Consider a component $C$ of $\Gamma_{1,(2k -3)}$. Define
%
\begin{equation}\label{E-bha4-3e}
\breve{w}_y := \left\{
\begin{array}{ll}
\partial_{\nu}w_y &\quad \text{in the Dirichlet case,}\\[5pt]
w_y|_{\partial \Omega} &\quad \text{in the Robin case.}\\[5pt]
\end{array}
\right.
\end{equation}

Let $y \in C$, and let $\set{y_n} \subset C$ be a sequence such that $y_n$ converges to $y$, so that $w_n := w_{y_n}$ converges to $w:=w_y$ (uniformly in the $C^p$ topology for $p$ large enough, Lemma~\ref{L-L32}).
Since $\breve{w}_n$ converges uniformly to $\breve{w}$, and since $\breve{w}$ changes sign at $z(y)$, it follows that $z(y_n)$ belongs to some neighborhood of $z(y)$, and that $z(y_n)$ tends to $z(y)$.
 This proves that $y \to z(y)$ is continuous in $C$. We now investigate the behaviour of $z(y)$ when $y$ tends to $\partial C$ (when $C \neq \Gamma_1$). Write $C$ as an open arc $(\eta_{-},\eta_{+}) \subset \Gamma_1$,
 and assume that $\set{y_n} \subset C$, with $y_n$ tending to some $\eta \in \set{\eta_{-},\eta_{+}} \subset \Gamma_{1,(2k-2)}$.
  Choose a subsequence of $\set{z(y_n)}$ which converges to some $z$. Since $w_n$ tends to $w_{\eta}$, with $\eta \in \Gamma_{1,(2k -2)}$,
   we conclude that $\breve{w}_{\eta}(z) = 0$, and hence that $z=\eta$ since $\eta$ is the unique zero of $\breve{w}_{\eta}$ in $\Gamma_1$.\hfill \qedc\medskip

\emph{Proof of Assertion~(iv).~} Let $C$ be a component of $\partial \Gamma_{1,(2k-3)}$.  Define $\alpha(y) := \tau_{y}^{W}(\downarrow) \in L_{(2k-3)}$.  Assume that the map $y \mapsto \alpha(y)$ is not locally constant. There exists $y \in C$ and a sequence $y_n$ tending to $y$ in $C$ such that $\alpha(y_n) \neq \alpha(y):=a$. Since the map $\alpha$ takes finitely many values, after taking a subsequence if necessary, we may assume that $\alpha(y_n) = b \neq a$. Let $w_n := w_{y_n}$. The nodal arc of $\cZ(w_n)$ which emanates from $y_n$ tangentially to the ray $\omega_b$ hits the boundary at the point $z_n := z(y_n)$. Since the sequence $\set{w_n}$ converges to $w_y$ in the $C^p$ topology for any fixed $p$, taking subsequences if necessary, we may assume that $\cZ(w_n)$ converges to $\cZ(w_y)$ in the Hausdorff distance, and that $\set{z_n}$ converges to $z(y)$. On the other-hand, we can apply the local structure theorem to the functions $w_n$ in a neighborhood of $y$: the arcs emanating from $y_n$ intersect a circle of radius $\varepsilon$
 (with $\varepsilon$ independent of $n$), at points $x_{n,j}, 1 \le j \le 2k -3$ and these points converge to the corresponding points $x_j, 1 \le j \le 2k -3$, for the function $w_y$.
  To prove that $\varepsilon$ can be taken independent of $n$ we use the fact that, for any fixed $m$, the derivatives of $w_n$ of order less than or equal to $m$ converge uniformly to the corresponding derivatives of $w_y$ so that the remainder term in Taylor's formula can be controlled independently of $n$, see the proof of the local structure theorem in  Appendix~\ref{S-lsc}, Remark~\ref{R-Cb}. The arc in $\cZ(w_n)$ between $x_{n,b}$ and $z_n$ must tend in the Hausdorff distance to the arc in $\cZ(w_y)$ between $x_b$ and $y$, and we get a contradiction since $b\neq a$. It follows that the map $\alpha$ is locally constant, hence constant, on the component $C$.  Since the map $y \mapsto \tau_{y}^{W}(\downarrow)$ is constant on $C$, the set $L_{(2k-3)} \sm \set{\tau_{y}^{W}(\downarrow)}$ is constant, and it suffices to look at the restriction of $\tau_{y}^{W}$ to this set. To prove that the map $y \mapsto \tau_{y,\alpha}$ is locally constant on $C$ we reproduce the arguments in the proof of  Property~\ref{propertyD4} or Lemma~\ref{L-hmn-202}. \hfill \qed


\begin{remark}\label{R-L33r}
Whether  $\Gamma_{1,(2k-2)}$ can be empty will be discussed in Section~\ref{S-hmn6}.
\end{remark}%


\section{Simply connected planar domains, properties of $\lambda_k$-eigenfunctions assuming that $\mult(\lambda_k) = (2k-2)$}\label{S-hmn4}


The purpose of this section is to derive further properties of $\lambda_k$-eigenfunctions assuming the upper bound in Proposition~\ref{P-hmn-s2} is achieved. Through out this section, we assume that

\begin{equation}\label{E-hmn-bs}
\left\{
\begin{array}{l}
\Omega \text{~is simply connected,}\\[5pt]
k \ge 3 \text{~and~} \dim U(\lambda_k) = (2k-2)\,.
\end{array}%
\right.
\end{equation}
%
In this section, to simplify the notation, we write $\Gamma$ instead of $\Gamma_1$, $\Gamma_{(2k-3)}$, instead of $\Gamma_{1,(2k-3)}$,  and $\Gamma_{(2k-2)}$ instead of $\Gamma_{1,(2k-2)}$.

\begin{remark}\label{R-simply}
 The assumption that $\Omega$ is simply connected is motivated by Berdnikov's argument (\cite{Berd2018}, Section~4)  showing why the last step in the proof of the upper bound $\mult(\lambda_k) \le (2k-3)$ given in \cite{HoMN1999} is incomplete in the non simply connected case. This assumption also makes the proofs of the following lemmas simpler. It would be worthwhile determining when it is actually necessary.
\end{remark}%


\subsection{Boundary behaviour of the map $x \mapsto u_x$}\label{SS-hmn31}

The following lemma expands Lemma~3.8 in \cite[p.~1184]{HoMN1999}.


\begin{lemma}\label{L-L38}
Let $\Omega \subset \R^2$ be a smooth bounded domain and $U := U(\lambda_k)$ for some $k \ge 3$.  Assume that $\Omega$ is simply connected, and $\dim U = 2k-2\,$. Let $\set{x_n} \subset \Omega$ be a sequence converging to some $y \in \Gamma $.  Let $\set{u_n}$ be a corresponding sequence of eigenfunctions, with $u_n  := u_{x_n} \in U_{x_n}\cap \bS(U)$.
\begin{enumerate}[(i)]
   \item If $w$ is a limit point of $\set{u_n}$, then $w \in W_y$. In particular, the continuous maps
%
\begin{equation*}
\left\{
\begin{array}{l}
\Omega \ni x \mapsto [u_x] \text{~of Lemma~\ref{L-L37}, and}\\[5pt]
\Gamma  \ni y \mapsto [w_y] \text{~of Lemma~\ref{L-L32}}
\end{array}
\right.
\end{equation*}
%
give rise to a continuous map $x\mapsto [\bar{u}_x]$ from $\overline{\Omega}$ into $\bP(U)$.
  \item  The point $y$ belongs to $\Gamma_{(2k -3)}$ if and only if,
% there exists a subsequence $\set{u_{s(n)}}$ such that  $\cS_b(u_{s(n)})
% = \set{y_{s(n)},z(y_{s(n)})}$ with $y_{s(n)} \to y$, and $z(y_{s(n)})
% \to z \not = y$
for $n$ large enough, $\cS_b(u_{n}) = \set{y_{n},z(y_{n})}$ with $y_{n} \to y$, and $z(y_{n}) \to z \not = y$.
  \item  The point $y$ belongs to $\Gamma_{(2k -2)}$ if and only if there exists an infinite subsequence $\set{u_{s(n)}}$  such that $\cS_b(u_{s(n)})= \emptyset$,  or  an infinite subsequence  $\set{u_{s(n)}}$  such that $\cS_b(u_{s(n)}) \neq \emptyset$, and  the points in $\cS_b(u_{s(n)})$ converge to $y$.
\end{enumerate}
\end{lemma}%



\begin{remark}\label{R-L38-2}
Since $\Omega$ is simply connected, we can lift the $C^{\infty}$ map $\Omega \ni x \to [u_x] \in \bP(U)$ to a $C^{\infty}$ map $x \to u_x \in \bS(U)$. Let $h_{x,(k-1)}(u_x)$ be the first nonzero term in the Taylor expansion of $u_x$ at the point $x$ (this is a harmonic polynomial of degree $(k-1)$). Then, the map $x \to h_{x,(k-1)}(u_x)$ is smooth and, by Assertion~(i), extends continuously to $\overline{\Omega}$. Unfortunately, this extension is not so interesting because $\lim_{x \to y \in \Gamma } h_{x,(k-1)}(u_x) = 0$ since $u_x$ tends to $w_y$ and $h_{y,(k-1)}(w_y) = 0$.  See also the final comment in \cite{BeNP2016}. We will mainly use Assertions~(ii) and (iii).
\end{remark}%

\emph{Proof of Lemma~\ref{L-L38}.}~ We divide the proof into several steps labeled {\bf (A), (B)}, \ldots for later cross reference.\smallskip

\noi \textbf{(A)} To the sequence of interior points, $\set{x_n} \subset \Omega$, we associate a sequence $\set{u_n = u_{x_n}}$ in the sphere $\bS(U)$ (Lemma~\ref{L-L37}). Taking a subsequence if necessary, we may assume that $\set{u_n}$ converges to some $w \in \bS(U)$. Then, the convergence is uniform in $C^m$ for any fixed $m \ge 0$. Since $\nu(u_n,x_n) = 2(k-1)$, or equivalently $\ord(u_n,x_n) = (k-1)$, with $k \ge 3$, and since the convergence is uniform, we have $\ord(w,y) \ge (k-1) \ge 2$. Define $p := \ord(w,y)$, $q:=\rho(w,y)$.\smallskip

By Lemma~\ref{L-hdn}, the (sub)sequence $\set{\cZ(u_n)}$ converges to $\cZ(w)$ in the Hausdorff distance. This in particular implies that the set $\cZ(w)$ is connected.\medskip

\noi \textbf{(B)} According to Section~\ref{SS-hmn2-20} or \cite[Section~2]{YaZh2021}, we can make a conformal change of coordinates so that, in a neighborhood of $y$, the boundary $\partial \Omega$ is a segment centered at $y$. A neighborhood of $y$ then contains a set of the form $D^{+}(r_2) := D(y,r_2) \cap \Omega$ (for these local coordinates), and we also take some $r_0 < r_1 < r_2$ to be chosen later on.\medskip

Taking $r_1 < r_2$ small enough, we may assume that the following properties hold.
\begin{enumerate}[(B-a)]
  \item The set $D^{+}(r_1)$ only contains one singular point of $w$, namely the point $y$, i.e., $D^{+}(r_1) \cap \cS_i(w) = \emptyset$ and $D^{+}(r_1) \cap \cS_b(w) = \set{y}$, so that the function $\breve{w}$ defined in \eqref{E-bha4-3d} only vanishes at $y$ on the set $\bar{D}^{+}(r_1) \cap \Gamma $.
  \item  Let $\mu_1(D^{+}(r_1))$ denote the least eigenvalue of $-\Delta + V$ in the domain $D^{+}(r_1)$ with mixed boundary condition, Dirichlet on $\partial D^{+}(r_1) \cap \Omega$, and the current boundary condition \eqref{E-evp-bc} on $\partial D^{+}(r_1) \cap \Gamma $ (Dirichlet or Robin). Choose $r_1$ small enough so that $\mu_1(D^{+}(r_1)) > \lambda_k$ ($\lambda_k$ is the eigenvalue associated with $U$). This is possible because $\mu_1(D^{+}(r))$ tends to infinity when $r$ tends to $0$ (use an extended monotonicity property for $\mu_1$ in such domains\footnote{The monotonicity follows from the min-max principle and the fact that the piece of boundary possibly with the Robin condition is always contained in $\Gamma $.},  and comparison with the eigenvalue of a half disk with Neumann condition on the diameter, and Dirichlet condition on the half circle).
  \item In the set $D^{+}(r_1)$, the local structure theorem applies to the function $w$. More precisely, let $(r,\omega)$ be local polar coordinates centered at $y$,  with respect to a direct frame $\set{\vec{e}_1,\vec{e}_2}$ such
  that $\vec{e}_1$ is tangent to $\Gamma $ at $y$, and $\vec{e}_2$ is normal to $\Gamma $, pointing inwards.
      \begin{itemize}
        \item[$\diamond$] In the Dirichlet case, $w(r,\omega) = c_w \, r^p \, \sin(p\omega) + \cO(r^{p+1})$. Define the rays $\set{\omega = \omega_j~|~ 1 \le j \le p-1}$, where $\omega_j := j \frac{\pi}{p}$. Then, in $D^{+}(r_1)$, the nodal set $\cZ(w)$ consists of $(p-1)$ arcs $r \mapsto \delta_j(r) := \big( r, \omegat_j(r) \big)$, where the functions $\omegat_j(r)$ are smooth for $0 < r < r_1$, with $\lim_{r\to 0}\omegat_j(r) = \omega_j$. These arcs are transverse to the circles $\set{r=c}$ with $c < r_1$. Note that $\rho(w,y) = q = (p-1)$.
        \item[$\diamond$] In the Robin case, $w(r,\omega) = c_w \, r^p \, \cos(p\omega) + \cO(r^{p+1})$. Define the rays
            %
            $$\set{\omega = \omega_j~|~ 1 \le j \le p},$$
            %
            where $\omega_j := (j-\frac 12) \frac{\pi}{p}$. Then, in $D^{+}(r_1)$, the nodal set $\cZ(w)$ consists of $p$ arcs $r \mapsto \delta_j(r) := \big( r, \omegat_j(r) \big)$, where the functions $\omegat_j(r)$ are smooth for $0 < r < r_1$, with $\lim_{r\to 0}\omegat_j(r) = \omega_j$. These arcs are transverse to the circles $\set{r=c}$ with $c < r_1$. Note that $\rho(w,y) = q = p$.
      \end{itemize}
  \item Let $\delta(\omega) := (r_0,\omega)$ be a parametrization of $C^{+}(r_0) := \partial D(r_0) \cap \Omega$. Then, $C^{+}(r_0) \cap \cZ(w) = \set{A_1(r_0),\ldots,A_{q}(r_0)}$. Choose $0 < \varepsilon_1 < 1$, small enough so that,  for $1 \le j \le q$, the open arcs $\left( A_j(r_0) - \varepsilon_1, A_j(r_0) + \varepsilon_1 \right)$ of $C^{+}(r_0)$, with length $2\varepsilon_1$, centered at $A_j(r_0)$ are pairwise disjoint. Define the set $C^{+}(r_0,w,\varepsilon_1)$  by
%
\begin{equation*}
 C^{+}(r_0,w,\varepsilon_1) := \bigcup_{j=1}^q \left( A_j(r_0) - \varepsilon_1, A_j(r_0) + \varepsilon_1 \right)\,,
\end{equation*}
%
see Figure~\ref{F-hmn3-L38-i}. Call $A_0(r_0)$ and $A_{q+1}(r_0)$ the extremities of $C^{+}(r_0) \cap \Gamma $. Since the (sub)sequence $\cZ(u_n)$ tends to $\cZ(w)$ in the Hausdorff distance, for $n$ large enough, we have $\cZ(u_n) \cap C^{+}(r_0) \subset C^{+}(r_0,w,\varepsilon_1)$.
\end{enumerate}\smallskip %\vspace{-3mm}

\begin{figure}[t!h]
\centering
\includegraphics[scale=0.7]{hmn3-L38-i}
\caption{Proof of Lemma~\ref{L-L38}: the set $C^{+}(r_0,w,\varepsilon_1)$} \label{F-hmn3-L38-i}
\end{figure}

\noi \textbf{(C)} From now on, we assume that $\set{x_n} \subset D^{+}(r_0/2)$, and that $n$ is large enough so that the above Property~(B-d) is satisfied. Taking Lemma~\ref{L-L37} into account,  there are two possible cases.\medskip

\noi \textbf{Case~C1}.~ \emph{There exists an infinite subsequence $\set{u_{s(n)}}$ of the sequence $\set{u_n}$ such that, for all $n$, $\cS_b (u_{s(n)}) = \emptyset$.}\smallskip

In this case, according to the proof of Lemma~\ref{L-L37}, the nodal set $\cZ(u_{s(n)})$ consists of $(k -1)$ simple loops at $x_{s(n)}$; these loops do not intersect away from $x_{s(n)}$, and do not hit $\Gamma $. Choose $\gamma$, any of these loops. Since $x_{s(n)} \in D^{+}( r_0/2)$, either the loop $\gamma$ crosses $C^{+}(r_0)$ at (at least) two distinct points $z_{\gamma,1}$ and $z_{\gamma,2}$, or its is entirely contained in $ \overline{D}^{+}(r_0)$. In the latter case, the function $u_{s(n)}$ would have a nodal domain contained in $D^{+}(r_0)$, and this would contradict the above Property~(B-b) by the monotonicity principle for Dirichlet eigenvalues, see
%Figures~\ref{F-hmn3-L38p}(A) and~\ref{F-hmn3-L38}(A).
 Figures~\ref{F-hmn3-L38-C1} (right) and \ref{F-hmn3-L38-FC1} (forbidden configurations).  For each $n$, the set $\cZ(u_{s(n)})$ consists of $(k-1)$ loops which do not intersect away from $x_{s(n)}$. It follows that we have at least $2(k-1)$ distinct points $z_{s(n),j}$ in $C^{+}(r_0) \cap \cZ(u_{s(n)})$, for $1 \le j \le (2k-2)$. Since $\set{\cZ(u_{s(n)})}$ converges to $\cZ(w)$ in the Hausdorff distance, the sequences $\set{z_{s(n),j}}, 1\le j \le (2k-2)$, converge to points $ \zeta_j \in C^{+}(r_0,w,\varepsilon_1)\cap \cZ(w) = \set{A_j(r_0), 1 \le j \le q}$. \smallskip

\begin{figure}[t!h]
\centering
\includegraphics[width=0.6\linewidth,angle=0.3]{hmn3-L38-C1}
\caption{Proof of Lemma~\ref{L-L38}: Case C1, nodal pattern.} \label{F-hmn3-L38-C1}
\end{figure}

\begin{figure}[t!h]
\centering
\includegraphics[width=0.6\linewidth,angle=0.3]{hmn3-L38-FC1}
\caption{Proof of Lemma~\ref{L-L38}: Case C1, forbidden configurations.} \label{F-hmn3-L38-FC1}
\end{figure}

\begin{claim}\label{Cl-C}
The points $\zeta_j, 1 \le j \le (2k-2)$, are pairwise distinct.
\end{claim}%

\pf Indeed, assume this is not the case, say $\zeta_1 = \zeta_2$. Then, the two sequences $\set{z_{s(n),1}}$ and $\set{z_{s(n),2}}$, with $z_{s(n),1} \not = z_{s(n),1}$, would converge to the same $A_j(r_0)$ for some $j$, $1 \le j \le q$. Writing $z_{s(n),i} = \delta(\theta_{s(n),i})$, we would find some $\theta_{s(n),3}$ in the smaller arc between $\theta_{s(n),1}$ and $\theta_{s(n),2}$ such that $(u_{s(n)} \circ \delta)'(\theta_{s(n),3}) = 0$. Since $u_n$ tends to $u$ uniformly in $C^1$, and since $z_{s(n),i}$ tends to $A_j(r_0)$, $A_j(r_0)$ would be a singular point of the function $w$. Since the nodal rays of $w$ are transverse to $C^{+}(r_0)$, this  would contradict the above Property~(B-a). \qedc \smallskip

Claim~\ref{Cl-C} implies that $q \ge (2k-2)$, i.e., $\rho(w,y) \ge (2k-2)$. Lemma~\ref{L-L32} implies that $w \in W_y$, and $\rho(w,y) = (2k-2)$, so that $y \in \Gamma_{(2k-2)}$.


\begin{remarks}\label{R-L38-4}\phantom{}
\begin{enumerate}[(i)]
  \item Note that Claim~\ref{Cl-C} also implies that, for $n$ large enough, each loop in $\cZ(u_{s(n)})$ meets $C^{+}(r_0)$ at exactly two points,  see Figure~\ref{F-hmn3-L38-C1}, right.
  \item Note that Case~C1 can only occur if $y \in \Gamma_{(2k-2)}$.
\end{enumerate}
\end{remarks}%
\medskip

\noi \textbf{Case~C2}.~ \emph{There exists an infinite subsequence $\set{u_{s(n)}}$ such that, for each $n$, we have $\cS_b (u_{s(n)}) \neq \emptyset$.}\smallskip

In this case, according to the proof of Lemma~\ref{L-L32},
 the nodal set $\cZ(u_{s(n)})$ consists of $(k -2)$ simple loops at $x_{s(n)}$, and two simple nodal arcs from $x_{s(n)}$ to the boundary.
 The loops do not intersect away from $x_{s(n)}$, and do not hit $\Gamma $; the arcs do not intersect each other (except at $x_{s(n)}$, and possibly on $\Gamma $ if they hit $\Gamma $  at the same point) and do not intersect the loops. The energy argument used in Case~C1, shows that each loop intersects $C^{+}(r_0)$ at (at least) two distinct points.
  A similar energy argument for mixed boundary conditions (Dirichlet on the nodal arcs, and the ambient boundary condition, Dirichlet or Robin, on $\Gamma $) shows that the nodal arcs cannot both be contained in $D^{+}(r_0)$,
  and at least one of them crosses $C^{+}(r_0)$, see %Figures~\ref{F-hmn3-L38p}~(B) and \ref{F-hmn3-L38}~(B).
 Figure~\ref{F-hmn3-L38-C2}, and Figures~\ref{F-hmn3-L38-FC2a} and \ref{F-hmn3-L38-FC2b} (forbidden configurations).  Finally, for each $n$, we have at least $(2k-3)$   distinct points $z_{s(n),j}$ in $C^{+}(r_0) \cap \cZ(u_{s(n)})$, for $1 \le j \le (2k-3)$.   As in Case~C1, these sequences converge to points $\zeta_j \in \cZ(w) \cap C^{+}(r_0)$. Applying Claim~\ref{Cl-C}, we conclude that $q \ge (2k-3)$, i.e. that $\rho(w,y) \ge (2k-3)$ so that $w \in W_y$, and we have two possible cases, either $\rho(w,y) = (2k-3)$ and $y \in \Gamma_{(2k-3)}$, or
$\rho(w,y) = (2k-2)$ and $y \in \Gamma_{(2k-2)}$ \smallskip

 At this stage we have proved that the only possible limit points of a sequence $\set{u_n}$ are in $W_y$,  see Figure~\ref{F-hmn3-L38-C1} (left).  Since $\dim W_y = 1$, this proves Assertion~(i) of the lemma. Claim~\ref{Cl-C} also implies that, for $n$ large enough, $Z(u_{s(n)})$ meets $C^{+}(r_0)$ at precisely $(2k-3)$ or $(2k-2)$ points. \medskip

According to  Lemma~\ref{L-L32}, we have $\cS_b(u_{s(n)}) = \set{y_{s(n),1},y_{s(n),2}}$ possibly with $y_{s(n),1} = y_{s(n),2}$. The only possible limit points of these sequences are
%
\begin{equation*}
\left\{
  \begin{array}{ll}
    y & \text{~if~} \rho(w,y) = (2k-2)\,,\\[5pt]
    y\text{~and~} z(y) & \text{~if~} \rho(w,y) = (2k-3)\,.
  \end{array}
\right.
\end{equation*}


When $\rho(w,y) = (2k-3)$, $\rho(w,z(y))=1$, and the function $\breve{w}$ vanishes and changes sign at $z(y)$. Since $\set{u_{s(n)}}$ converges to $w$ $C^{1}$-uniformly, this implies that, for $n$ large enough, the function $\breve{u}_{s(n)}$ changes sign near $z(y)$, and hence that one sequence, say $\set{y_{s(n),2}}$ tends to $z(y)$, and the other $\set{y_{s(n),1}}$ tends to $y$. Note that they cannot both tend to $z(y)$ since $\rho (w,z(y))=1$.\smallskip

When $\rho(w,y) = (2k-2)$, the sequences $\set{y_{s(n),1}}$ and $\set{y_{s(n),2}}$ must both converge to $y$.\smallskip

Applying Claim~\ref{Cl-C}, we find that there are three sub-cases.
\begin{description}
  \item[C2(i)] There exists a subsequence $\set{u_{s(n)}}$ such that $y_{s(n),1} = y_{s(n),2}$ tending to $y$. For energy reasons, the arcs from $x_{s(n)}$ to $y_{s(n),1}$ cannot both be contained in $D^{+}(r_0)$. One of these arcs, and actually only one for $n$ large enough, has to meet $C^{+}(r_0)$ at two distinct points,  see Remark~\ref{R-L38-4}(i)  and Figure~\ref{F-hmn3-L38-C2} (left).
  \item[C2(ii)] There exists a subsequence $\set{u_{s(n)}}$ such that $y_{s(n),1} \not = y_{s(n),2}$ both tending to $y$. For energy reasons, the arcs from $x_{s(n)}$ to $y_{s(n),1}$ and $y_{s(n),2}$ cannot both be contained in $D^{+}(r_0)$. One of these arcs, and actually only one for $n$ large enough, has to meet $C^{+}(r_0)$ at two distinct points.  See Figure~\ref{F-hmn3-L38-C2} (center).
  \item[C2(iii)] There exists a subsequence $\set{u_{s(n)}}$ such that $y_{s(n),1} \not = y_{s(n),2}$, with  $y_{s(n),1}$ tending to $y$ and $y_{s(n),2}$ tending to some $z \not = y$. For $n$ large enough, the arc from $x_{s(n)}$ to $y_{s(n),2}$ intersects  $C^{+}(r_0)$ at one point, and the arc from $x_{s(n)}$ to $y_{s(n),1}$ stays inside $D^{+}(r_0)$.  See Figure~\ref{F-hmn3-L38-C2} (right).
\end{description}

In subcases C2(i) and C2(ii), we have $\rho(w,y) = (2k-2)$, so that $y \in \Gamma_{(2k-2)}$. In subcase C2(iii), we have $\rho(w,y) = (2k-3)$, so that $y \in \Gamma_{(2k-3)}$ with $z(y) = z$, the limit of $\set{y_{s(n),2}}$.

\begin{figure}[t!h]
\centering
\includegraphics[width=0.9\linewidth,angle=0.3]{hmn3-L38-C2}
\caption{Proof of Lemma~\ref{L-L38}: Case C2, nodal patterns}\label{F-hmn3-L38-C2}
\end{figure}


\begin{figure}[t!h]
\centering
\includegraphics[width=0.6\linewidth,angle=0.3]{hmn3-L38-FC2a}
\caption{Proof of Lemma~\ref{L-L38}: Case C2i, forbidden configurations}\label{F-hmn3-L38-FC2a}
\end{figure}

\begin{figure}[t!h]
\centering
\includegraphics[width=0.9\linewidth,angle=0.3]{hmn3-L38-FC2b}
\caption{Proof of Lemma~\ref{L-L38}: Case C2, forbidden configurations} \label{F-hmn3-L38-FC2b}
\end{figure}


This proves Assertions~(ii) and (iii). The proof of Lemma~\ref{L-L38} is complete. \hfill \qed \medskip

Figures~\ref{F-hmn3-L38-FC2a} and \ref{F-hmn3-L38-FC2b} display forbidden configurations for $\cZ(u_{n})$ when $r_1$ is small enough and $n$ large: a loop cannot be contained in  $D^{+}(r_0)$; the arcs cannot both be contained in $D^{+}(r_0)$;  the arcs cannot both meet $C^{+}(r_0)$.

 \begin{remark}\label{R-Fgene3}
The \emph{nodal patterns} displayed in the above figures hold for both the Dirichlet and Robin boundary conditions. Unless otherwise stated this remark applies to all figures of this section.
\end{remark}



Figures~\ref{F-hmn3-L38-C1} and \ref{F-hmn3-L38-C2} display typical configurations for $\cZ(u_{n})$ when $r_1$ is small enough and $n$ large: the loops intersect $C^{+}(r_0)$ at two distinct points; one arc exits $D^{+}(r_0)$.
\medskip

 As a byproduct of the proof of Lemma~\ref{L-L38}, we obtain the configurations of the nodal sets $\cZ(u_x)$ when $x \in \Omega$ tends to some $y \in \Gamma $. When $y \in \Gamma_{(2k-3)}$, we are in Case C2(iii). When $y \in \Gamma_{(2k-2)}$, one of the cases C1, C2(i) or C2(ii) occurs.

\FloatBarrier

%
\subsection{Eigenfunctions with two prescribed boundary singular points}\label{SS-hmn32b}  Given $(y,s) \in \Gamma_{(2k -3)} \times \Gamma $, with $ y \neq s$, define the subspace
%
\begin{equation}\label{E-L34-2}
V_{y,s} := \set{u \in U ~|~ \rho(u,y) \ge 2k-4 \text{~and~} \rho(u,s) \ge 1}\,.
\end{equation}

 In view of the assumption \eqref{E-hmn-bs}, Lemma~\ref{L-zero2} implies that $V_{y,s} \neq \set{0}$. \smallskip

In this subsection, we revisit Lemmas~3.4, 3.5 and 3.6 of \cite[pp.~1180-1183]{HoMN1999}. We retain the notation of the previous subsections and introduce the following one.

\begin{notation}\label{N-L34-2}
Given two points $y_1 \neq y_2 \in \Gamma $, we denote by $\cA(y_1,y_2)$ the open arc from $y_1$ to $y_2$, moving counter-clockwise. Given $y \in \Gamma $ and $a$ smaller than half the length of $\Gamma $, we denote by $\cA_c(y;a)$ the arc centered at $y$, with length $2a$, taken counter-clockwise. In both cases, we use  the mathematical  symbols  $[$ and $]$ to denote the closed or semi-closed arcs.
\end{notation}%

\subsubsection{First properties of $V_{y,s}$}\label{SSS-hmn32ba}

\begin{lemma}\label{L-L34a}
 Let $U:=U(\lambda_k)$ with $k \ge 3$. Assume that $\Omega$ is simply connected, and $\dim U = (2k-2)$. Given $(y,s) \in \Gamma_{(2k -3)} \times \Gamma $, with $ y \neq s$,  the following properties of the subspace $V_{y,s}$ defined in \eqref{E-L34-2} hold.
\begin{enumerate}[(i)]
  \item The subspace $V_{y,s}$ has dimension $1$. Let $v_{y,s}$ denote a generator of $V_{y,s}$.
  \item Any $0 \neq u \in V_{y,s}$ satisfies
  \begin{equation}\label{E-L34-4}
  \left\{
  \begin{array}{l}
  \kappa(u) = k\,,\\[5pt]
  \cZ(u) \cup \Gamma  \text{~is connected,}\\[5pt]
  \cS_i(u) = \emptyset\,,\\[5pt]
  \sum_{z \in \cS_b(u)} \rho(u,z) = 2k-2\,, \text{~and}\\[5pt]
  \quad 2k-4 \le \rho(u,y) \le 2k-3\,,\\[5pt]
  \quad 1 \le \rho(u,s) \le 2.
  \end{array}
  \right.
  \end{equation}
  %
  More precisely, there are three distinct possibilities.
  \begin{description}
    \item[Case~(1)] $\rho(u,y) = (2k-3)$ and $\rho(u,s) = 1$.
    \item[Case~(2)] $\rho(u,y) = (2k-4)$ and $\rho(u,s) = 2$.
    \item[Case~(3)] $\rho(u,y) = (2k-4)$, $\rho(u,s) = 1$, and there exists $s' \in \Gamma  \sm \set{y,s}$ such that $\cS_b(u) = \set{y,s,s'}$, $\rho(u,s') = 1$.
  \end{description}
  In Case~(1), $u \in W_y$,  and $s = z(y)$ (with the notation of Lemma~\ref{L-L32}).
  \item If $s = z(y)$, then $V_{y,z(y)} = W_y$.
  \item The map $\set{(y,s) ~|~ (y,s) \in \Gamma_{1}\times \Gamma , s\neq y} \ni (y,s) \mapsto [V_{y,s}] \in \bP(U)$ is $C^{\infty}$.
  \end{enumerate}
\end{lemma}%

\pf  We already know that $\dim V_{y,s} \ge 1$. \smallskip

We retain the notation of Lemma~\ref{L-L32}. In particular, for $y \in \Gamma_{(2k -3)}$, we have $W_y = [w_y]$ with $0 \neq w_y \in U$ satisfying $\cS_b(w_y) = \set{y,z(y)}$ with $z(y) \neq y$, and $\rho(w_y,y) = (2k-3)$, $\rho(w_y,z(y)) = 1$. \smallskip

\noid \emph{Proof of  Assertion~(ii).~} From Euler's formula \eqref{E-hmn-3a12} we obtain,
%
\begin{equation}\label{E-L34-6}
\begin{split}
0 \geq  \kappa(u) - k & = \left( b_0(\cZ(u) \cup \Gamma ) -1 \right) + \frac 12\, \sum_{z \in \cS_i(u)} (\nu(u,z)-2) \\
&\hspace{5mm} + \frac 12 \left( \sum_{z\in \cS_b(u)}\rho (u,z) - 2k +2\right) \,.
\end{split}
\end{equation}
%
If $0 \neq u \in V_{y,s}$, we have $\sum_{z\in \cS_b(u)}\rho (u,z) \ge 2k-3$, and hence $\sum_{z\in \cS_b(u)}\rho (u,z) \ge 2k-2$ since the sum is an even integer,  by Corollary~\ref{cor:nodinfo}. All the terms in the right hand side of \eqref{E-L34-6} must vanish; this proves \eqref{E-L34-4}.  Assertion~(ii) then follows from \eqref{E-L34-4} and the assumptions that $\rho(u,y) \ge (2k-4)$ and $\rho(u,s) \ge 1$.\medskip

\noid \emph{Proof of  Assertion~(i).}~

\noi \textbf{(A)}~ We first assume that $s \neq z(y)$. Assume that there are at least two linearly independent functions $u_1,u_2 \in V_{y,s}$, then $\rho(u_i,y) = (2k-4)$ and  $\rho(u_i,s) \ge 1$. According to Lemma~\ref{L-zeroc}, there exists a nontrivial linear combination $u$ of $u_1$ and $u_2$ such that $\rho(u,y) \ge 2k-3$ and $\rho(u,s) \ge 1$. Euler's formula implies that $u$ pertains to Assertion~(ii), Case~(1), contradicting the fact that $s \neq z(y)$.\smallskip

\noi \textbf{(B)}~ We now assume that $s = z(y)$. In this case, a generator $w_y$ of $W_y$ belongs to $V_{y,z(y)}$. Assume that $\dim V_{y,z(y)} \ge 2$. Define $V'_{y,z(y)} = V_{y,z(y)} \ominus W_y$, which has dimension at least $1$. If $\dim V_{y,z(y)} \ge 3$, we can find two linearly independent $u_1,u_2 \in V'_{y,z(y)}$, such that $\rho(u_i,y) = 2k-4$, and $\rho(u_i,s) \ge 1$. By Lemma~\ref{L-zeroc}, there exists a nontrivial linear combination $u \in V'_{y,z(y)}$ such that $\rho(u,y) \ge 2k-3$ and $\rho(u,s) \ge 1$.  Hence, $u \in W_y$, a contradiction. Assuming that $\dim V_{y,z(y)} = 2$, we can choose a basis $\set{w_y,v_y}$ such that $v_y \not \in W_y$. Then, $\rho(v_y,y) = 2k-4$, and there are two cases,\smallskip
\begin{description}
  \item[Case~(a)] $\rho(v_y,z(y))=2$,\vspace{3mm}
  \item[Case~(b)] $\rho(v_y,z(y)) = 1$, and there exists some $z_1(y) \in \Gamma $, $z_1(y) \neq z(y)$,  such that $\cS_b(v_y) = \set{y,z(y),z_1(y)}$ and $\rho(v_y,z_1(y)) = 1$.
\end{description}\smallskip

Without loss of generality,  making use of Lemma~\ref{L-breve}, we may choose the functions $w_y$ and $v_y$ as follows  (we consider open arcs).  First we choose $w_y$ so that $\breve{w}_y > 0$ on the arc $\cA(y,z(y))$, and $\breve{w}_y < 0$ on the arc $\cA(z(y),y)$. \smallskip

\noic In Case (a), we choose $v_y$ such that $\breve{v}_y > 0$ on $\cA(y,z(y)) \cup \cA(z(y),y)$.\smallskip

\noic In Case (b), assuming that $z_1(y) \in \cA(y,z(y))$, we choose $v_y$ such $\breve{v}_y > 0$ on $\cA(z(y),y) \cup \cA(y,z_1(y))$, and $\breve{v}_y < 0$ on $\cA(z_1(y),z(y))$.\smallskip

Figure~\ref{F-hmn3-L34-4} displays the signs of $\breve{v}_y$ in both cases.

\begin{figure}[!ht]
  \centering
  \includegraphics[scale=0.8]{hmn3-L34-4}
  \caption{Signs of $\breve{v}_y$, Cases (a) and (b)}\label{F-hmn3-L34-4}
\end{figure}


\emph{Claim~1.~ Under the assumption that $\dim V_{y,z(y)} = 2$, there exists some function $v \in V_{y,z(y)}$ such that $\rho(v,y)=2k-4$ and $\rho(v,z(y))=2$.} \smallskip

\emph{Proof of Claim~1.}~ If $v_y$ satisfies Claim~1, there is nothing to prove. If not, $v_y$ falls into Case (b) above.\smallskip

Given $t \in \Gamma  \sm \set{y, z(y)}$, $\breve{w}_y(t) \neq 0$, and we can define the function
%
\begin{equation}\label{E-L34-10}
\xi_t := a(t)\, w_y - b(t)\, v_y \in V_{y,z(y)}\,,
\end{equation}
%
where
%
\begin{equation}\label{E-L34-10a}
\left\{
\begin{array}{ll}
a(t) & =\breve{v}_y(t)\, \left( \breve{v}^2_y(t) + \breve{w}^2_y(t) \right)^{- \frac 12}\,,\\[5pt]
b(t) & =\breve{w}_y(t) \, \left( \breve{v}^2_y(t) + \breve{w}^2_y(t) \right)^{- \frac 12}\,.
\end{array}
\right.
\end{equation}

For $t \not \in \set{y, z(y)}$, $b(t) \neq 0$, and hence $\rho(\xi_t,y) = 2k-4$, $\rho(\xi_t,z(y)) \ge 1$, $\rho(\xi_t,t) \ge 1$. Euler's formula applied to $\xi_t$ implies that  $\rho(\xi_t,z(y)) = \rho(\xi_t,t) = 1$,  and $\cS_b(\xi_t) = \set{y,z(y),t}$. According to Lemma~\ref{L-breve}, the function $\breve{\xi}_t$ has precisely three zeros at $y, z(y)$ and $t$, changes sign at $z(y)$ and $t$, and does not change sign at $y$.  For $t \in \cA(z_1(y),z(y))$, $\breve{\xi}_t(z_1(y)) < 0$, and we conclude that
%
\begin{equation}\label{E-L34-10b}
\text{for~} t \in \cA(z_1(y),z(y)), \qquad\left\{
\begin{array}{l}
\breve{\xi}_t > 0 \text{~~in~~} \cA(t,z(y)), \text{~and}\\[5pt]
\breve{\xi}_t < 0 \text{~~in~~} \cA(z(y),y) \bigcup \cA(y,t)\,.
\end{array}%
\right.
\end{equation}
%
where $y, z(y)$ and $z_1(y)$ are as in Figure~\ref{F-hmn3-L34-4}(b).

Choose a sequence $\set{t_n} \subset \cA(z_1(y),z(y))$, with $t_n \to z(y)$. Taking a subsequence if necessary, we may assume that the sequence $\set{\left( a(t_n),b(t_n)\right)}$ converges to some $(a,b) \in \bS^1$, so that the sequence $\set{\xi_{t_n}}$ converges to the function $\xi := a\, w_y - b\, v_y$. From \eqref{E-L34-10b}, we conclude that $\breve{\xi} \le 0$ on $\Gamma $. Since $\xi \in V_{y,z(y)}$ we have three possibilities,
\begin{enumerate}[(i)]
  \item $\rho(\xi,y) = 2k-3$ and $\rho(\xi,z(y)) =1$,
  \item $\rho(\xi,y) = 2k-4$, $\rho(\xi,z(y)) =1$, and $\rho(\xi,z_2)$ for some $z_2 \neq y, z(y)$,
  \item $\rho(\xi,y) = 2k-4$ and $\rho(\xi,z(y)) =2$.
\end{enumerate}

Since (i) and (ii) are incompatible with $\breve{\xi} \le 0$ on $\Gamma $, we conclude that $\rho(\xi,y) = 2k-4$ and $\rho(\xi,z(y)) =2$. This proves Claim~1.\hfill \qedc \medskip


We now continue  with part (B) in  the proof of Assertion~(i). In view of Claim~1,  assuming that $\dim V_{y,z(y)}= 2$, we may choose a basis $\set{w_y,v_y}$ of $V_{y,z(y)}$ such that
%
\begin{equation}\label{E-L34-10d}
\begin{array}{ll}
\left\{
\begin{array}{l}
\rho(w_y,y) = 2k-3\,,\\[5pt]
\rho(w_y,z(y)) = 1\,,\\[5pt]
\breve{w}_y|_{\cA(y,z(y))} > 0 \text{~and~} \breve{w}_y|_{\cA(z(y),y)} < 0\,,\\[5pt]
\end{array}
\right.
\text{~~and~~} & ~~\left\{
\begin{array}{l}
\rho(v_y,y) = 2k-4\,,\\[5pt]
\rho(v_y,z(y)) = 2\,,\\[5pt]
\breve{v}_y|_{\Gamma  \sm \set{y,z(y)}} > 0\,.
\end{array}
\right.
\end{array}
\end{equation}


Examples of nodal sets of these functions are displayed in Figure~\ref{F-hmn3-L34-8}: on the left $\cZ(w_y)$, on the right $\cZ(v_y)$, with two possible cases.\smallskip

\begin{figure}[!ht]
  \centering
  \includegraphics[scale=0.38]{hmn3-L34-8}\\
  \caption{Nodal sets of $w_y$ (left) and $v_y$ (right), with $k=4$.}\label{F-hmn3-L34-8}
\end{figure}

From now on, to simplify the notation in the proof,  we denote the arc $\mathcal A(y,z(y))$  by $A_1$, and the arc $\mathcal A(z(y),y)$ by $A_2$.\smallskip

For $s \not \in \set{y,z(y)}$, we consider the function
%
\begin{equation}\label{E-L34-12}
\xi_s = a(s) \, w_y - b(s) \, v_y\,,
\end{equation}
%
where $w_y$ and $v_y$ satisfy \eqref{E-L34-10d}, and $a(s)$, $b(s)$ are given by
\begin{equation}\label{E-L34-14}
\left\{
\begin{array}{l}
a(s) =\breve{v}_y(s)\, \left( \breve{v}^2_y(s) + \breve{w}^2_y(s) \right)^{- \frac 12}\,,\\[5pt]
b(s) = \breve{w}_y(s) \, \left( \breve{v}^2_y(s) + \breve{w}^2_y(s) \right)^{- \frac 12}\,.
\end{array}
\right.
\end{equation}
%
In particular, $a(s) > 0$ in $A_1 \cup A_2$, $b(s) > 0$ in $A_1$ and $b(s) < 0$ on $A_2$. \smallskip
Since $a(s)$ and $b(s)$ are different from $0$, $\rho(\xi_s,y) = (2k-4)$ and $\rho(\xi_s,z(y))=1$. Since $\breve{\xi}_s(s) =0$, $\rho(\xi_s,s) \ge 1$. Since $\xi_s \in V_{y,z(y)}$,  Equation~\eqref{E-L34-4} implies that $\rho(\xi_s,s) = 1$, $\cS_b(\xi_s) = \set{y,z(y),s}$, and $\breve{\xi_s}$ changes sign at $z(y)$ and $s$  (use Lemma~\ref{L-breve} again).\smallskip

Taking $s_2 \in A_2$ and $s \in A_1$, we find that $\breve{\xi}_s(s_2) < 0$ and hence,
%
\begin{equation}\label{E-L34-14a}
\text{for~} s \in A_1, \quad \left\{
\begin{array}{ll}
\breve{\xi_s} > 0 & \text{in~~} \cA(s, z(y))\,,\\[5pt]
\breve{\xi_s} < 0 & \text{in~~} A_2 \bigcup \cA(y,s)\,.
\end{array}
\right.
\end{equation}

Similarly, taking $s_1 \in A_1$ and $s \in A_2$, we find that $\breve{\xi}_s(s_1) > 0$ and hence,
%
\begin{equation}\label{E-L34-14b}
\text{for~} s \in A_2, \quad \left\{
\begin{array}{ll}
\breve{\xi_s} < 0 & \text{in~~} \cA(z(y), s)\,,\\[5pt]
\breve{\xi_s} > 0 & \text{in~~} \cA(s,y) \bigcup A_1 \,.
\end{array}
\right.
\end{equation}

\begin{figure}[!ht]
\centering
\begin{subfigure}[t]{.30\textwidth}
\centering
\includegraphics[width=\linewidth]{hmn3-L34-20a}
\caption{$s \in A_1$}
\end{subfigure}
\begin{subfigure}[t]{.30\textwidth}
\centering
\includegraphics[width=\linewidth]{hmn3-L34-20b}
\caption{$s\in A_2$}
\end{subfigure}
%
\caption{ Signs of the function $\breve{\xi}_s$}\label{F-hmn3-L34-20}
\end{figure}

There exists a unique $\theta(s) \in (-\pib , \pib) $ such that $a(s) = \cos(\theta(s))$ and $b(s) = \sin(\theta(s))$, so that $\xi_s = \cos(\theta(s)) \, w_y - \sin(\theta(s))\, v_y$. Then, $\theta(s) \in (0,\pib)$ when $b(s) > 0$ or equivalently when $s \in A_1$;  $\theta(s) \in (-\pib,0)$ when $b(s) < 0$ or equivalently when $s \in A_2$. Furthermore, the map $s \mapsto \theta(s)$ is injective because $\cS_b(\xi_s) = \set{y,z(y),s}$. \smallskip

For $s_1, s_2 \in A_1$ or $A_2$, we have
%
\begin{equation}\label{E-L34-14c}
\left\{
\begin{array}{ll}
\xi_{s_1} - \xi_{s_2} & = \big( \cos(\theta(s_1) - \cos(\theta(s_2)\big) w_y - \big( \sin(\theta(s_1) - \sin(\theta(s_2)\big) v_y \\[5pt]
& = - 2 \sin\frac{\theta(s_1) - \theta(s_2)}{2}\big( \sin\frac{\theta(s_1) + \theta(s_2)}{2}\, w_y + \cos\frac{\theta(s_1) + \theta(s_2)}{2}\, v_y\big)\,.
\end{array}
\right.
\end{equation}



If $s_1, s_2 \in A_1$, $\theta(s_1), \theta(s_2) \in (0,\pib)$, and the second factor in the second line of Equation~\eqref{E-L34-14c} is positive in $A_1$. If $s_1, s_2 \in A_2$, $\theta(s_1), \theta(s_2) \in (-\pib,0)$, and the second factor in the second line of Equation~\eqref{E-L34-14c} is positive in $A_2$. \\  Since $\breve{\xi}_{s_1}(s_2) - \breve{\xi}_{s_2}(s_2) = \breve{\xi}_{s_1}(s_2) $, using Equations~\eqref{E-L34-14a} and \eqref{E-L34-14b} we conclude that,
%
\begin{equation}\label{E-L34-14e}
\left\{
\begin{array}{l}
s_1 \in A_1 \text{~and~} s_2 \in \cA(s_1,z(y)) \Rightarrow \breve{\xi}_{s_1}(s_2) > 0\\[5pt]
 \quad \Rightarrow  \sin\frac{\theta(s_1)-\theta(s_2)}{2}
< 0  \quad \Rightarrow \theta(s_2) > \theta(s_1) \text{~in~} (0,\pib)\,,\\[5pt]
s_1 \in A_2  \text{~and~} s_2 \in \cA(s_1,y) \Rightarrow \breve{\xi}_{s_1}(s_2) > 0\\[5pt]
 \quad \Rightarrow  \sin\frac{\theta(s_1)-\theta(s_2)}{2}
< 0  \quad \Rightarrow \theta(s_2) > \theta(s_1)  \text{~in~} (-\pib,0)\,,\\[5pt]
\end{array}%
\right.
\end{equation}
%
otherwise stated, when $s$ moves counter-clockwise on $A_1$, resp. $A_2$,  the function $\theta(s)$ increases from $0$ to $\pib$, resp. from $-\pib$ to $0$, see Figure~\ref{F-hmn3-L34-22}.\medskip


\begin{figure}[!ht]
\centering
\begin{subfigure}[t]{.30\textwidth}
\centering
\includegraphics[width=\linewidth]{hmn3-L34-22a}
\caption{$s_1, s_2 \in A_1$}
\end{subfigure}
\begin{subfigure}[t]{.30\textwidth}
\centering
\includegraphics[width=\linewidth]{hmn3-L34-22b}
\caption{$s_1, s_2 \in A_2$}
\end{subfigure}
%
\caption{Sign of $\breve{\xi}_{s_1}(s_2)$}\label{F-hmn3-L34-22}
\end{figure}

 Under the assumption that $\dim V_{y,z(y)}=2$, we are in a framework similar to that of Section~\ref{SS-hmn-25}, with $V_{y,z(y)}$ replacing $U^2_x$. We use a rotating function argument similar to the one used in Paragraph~\ref{SSS-hmn-25D}. For this purpose, we first investigate the limits of $\xi_s$ and $\theta(s)$ when $s$ tends to $z(y)$ or to $y$.\smallskip

Let $\gamma_z$ (resp. $\gamma_y$) denote a local parametrization of $\Gamma $ in a neighborhood of $z(y)$ (resp. $y$), such that $\gamma_z(0) = z(y)$, $\gamma_z(-\varepsilon) \in A_2$, and $\gamma_z(\varepsilon) \in A_1$ (resp. $\gamma_y(0) = y$, $\gamma_y(-\varepsilon) \in A_2$, and $\gamma_y(\varepsilon) \in A_1$). Using our choice of sign for  $\breve{w}_y$ and $\breve{v}_y$, the vanishing properties of these functions, and  Lemma~\ref{L-breve}, we find that, in a  pointed neighborhood of $z(y)$,
%
\begin{equation*}
\left\{
\begin{array}{l}
\breve{w}_y(\gamma_z(t)) = \alpha_w\, t + o(t), {\text{~with~}}\alpha_w > 0\,,\\[5pt]
\breve{v}_y(\gamma_z(t)) = \alpha_v\, t^2 + o(t^2), {\text{~with~}}\alpha_v > 0\,,\\[5pt]
a(\gamma_z(t)) = \frac{\alpha_v}{\alpha_w} |t| + o(t)\,,\\[5pt]
b(\gamma_z(t)) = \sign(t) + o(1)\,.\\[5pt]
\end{array}%
\right.
\end{equation*}
%
Similarly, in a neighborhood of $y$,

\begin{equation*}
\left\{
\begin{array}{l}
\breve{w}_y(\gamma_y(t)) = \beta_w\, t^{2k-3} + o(t^{2k-3}), {\text{~with~}}\beta_w > 0\,,\\[5pt]
\breve{v}_y(\gamma_y(t)) = \beta_v\, t^{2k-4} + o(t^{2k-4}), {\text{~with~}}\beta_v > 0\,,\\[5pt]
a(\gamma_y(t)) =  1 + o(1)\,,\\[5pt]
b(\gamma_y(t)) = \frac{\beta_w}{\beta_v}\, t + o(t)\,.\\[5pt]
\end{array}%
\right.
\end{equation*}

This gives us the limits of $\xi_s$ and $\theta(s)$ when $s$ tends to $z(y)$ in $A_1$ or $A_2$ (resp. when $s$ tends to $y$ in $A_1$ and $A_2$),


\begin{equation}\label{E-L34-14g}
\left\{
\begin{array}{llll}
\lim_{\substack{s \to z(y)\\s \in A_1}}\xi_s &= - v_y\,,
& \qquad  \lim_{\substack{s \to z(y)\\s \in A_1}}\theta(s) &= \pib\,,\\[10pt]
\lim_{\substack{s \to y\\s \in A_1}}\xi_s &= w_y\,,
& \qquad  \lim_{\substack{s \to y\\s \in A_1}}\theta(s) &= 0\,,\\[10pt]
\lim_{\substack{s \to z(y)\\s \in A_2}}\xi_s &=  v_y\,,
& \qquad  \lim_{\substack{s \to z(y)\\s \in A_2}}\theta(s) &= 0\,,\\[10pt]
\lim_{\substack{s \to y\\s \in A_2}}\xi_s &= w_y\,,
& \qquad  \lim_{\substack{s \to y\\s \in A_2}}\theta(s) &=  -\pib \,.
\end{array}
\right.
\end{equation}
\medskip

\begin{figure}[!ht]
  \centering
  \includegraphics[scale=0.7]{hmn3-L34-22n}\\
  \caption{Lemma~\ref{L-L34a}: limits of $\xi_s$ when $s$ tends to $y$ or $z(y)$}\label{F-hmn3-L34-22n}
\end{figure}

When $s$ moves counter-clockwise from $y$ to $z(y)$ on $A_1$, $\theta(s)$ increases from $0$ to $\pib$; when $s$ moves counter-clockwise from $z(y)$ to $y$ on $A_2$, $\theta(s)$ increases from $-\pib$ to $0$.\medskip

Let $\sigma \mapsto \Gamma (\sigma)$ be a parametrization of $\Gamma $ such that $\sigma \in [-\pib,\pib]$, $\gamma_1(-\pib) = \gamma_1(\pib) = z(y)$, $\gamma_1(0)=y$,  $\gamma_1\left( (-\pib,0) \right) = A_2$, and $\gamma_1\left( (0,\pib) \right) = A_1 \,$.\smallskip

Consider the map
%
\begin{equation*}
\theta_1 : (-\pib,0) \cup (0,\pib) \ni \sigma \mapsto \theta(\gamma_1(\sigma)) \in (-\pib,\pib)\,.
\end{equation*}
%
Then, $\theta_1$ extends to a continuous, increasing map from $(-\pib,\pib)$ to $(-\pib,\pib)$ such that \goodbreak  $\lim_{\sigma \to \pm \pib}\theta_1(\sigma) = \pm \pib$ and $\lim_{\sigma \to 0}\theta_1(\sigma) = 0$. \medskip


For $t \in [-\pib,\pib]$, introduce the functions
%
\begin{equation}\label{E-L34-15a}
\zeta_{t} := \cos t \, w_y - \sin t \, v_y\,,
\end{equation}
%
$\zeta_{-\pib} = v_y$, $\zeta_0 = w_y$, and $\zeta_{\pib} = - v_y$. If $t \not \in \set{-\pib, 0, \pib}$ there exists a unique $s(t) \in \Gamma \sm\set{y,z(y)}$ such that
%
\begin{equation}\label{E-L34-15c}
\left\{
\begin{array}{l}
\rho(\zeta_{t},y) = (2k-4)\,, ~~\rho(\zeta_{t},z(y))=1\,,~~\rho(\zeta_{t},s(t))=1\,,\\[5pt]
\cS_b(\zeta_{t}) = \set{y,z(y),s(t)}\,,\\[5pt]
s(t) \in A_1 \text{~if~} t \in (0,\pib)\,, \text{~~and~~} s(t) \in A_2 \text{~if~} t \in (-\pib,0)\,.
\end{array}
\right.
\end{equation}
%
Indeed, near $z(y)$, $\cos t > 0$ implies that $\zeta_t$ has the sign of $w_y$. In a small pointed arc $J_y$ around $y$, $\zeta_t \sim - \sin t \, v_y$ which has the sign of $(-t)$. \medskip

We now apply the ``rotating function'' argument, see Paragraph~\ref{SSS-hmn-25D}, to the family of nodal sets $\cZ(\zeta_t)$. We have $\cZ(\zeta_{-\pib}) = \cZ(\zeta_{\pib}) = \cZ(v_y)$; when $t \in (0,\pib)$, $\cS_b(\zeta_t) = \set{y, z(y),s(t)}$ with $s(t) \in A_1$; when $t \in (-\pib,0)$, $s(t) \in A_2$. \smallskip


Figure~\ref{F-hmn3-L34-24n} (resp. Figure~\ref{F-hmn3-L34-26n}) illustrates the deformation of the nodal set $\cZ(v_y)$ given in Subfigure (A) in a particular case with $k=4$.\medskip

In Figure~\ref{F-hmn3-L34-24n} the nodal set $\cZ(v_y)$ is connected. When $t$ decreases from $\pib$ to $0$ (top line),
 the nodal set $\cZ(\zeta_t)$ deforms from  $\cZ(v_y)$ in Subfigure~(A) to $\cZ(w_y)$ in Subfigure~(L).  When $t$ increases from $-\pib$ to $0$ (bottom line), the nodal set $\cZ(\zeta_t)$ deforms from $\cZ(v_y)$ in Subfigure~(A) to $\cZ(w_y)$ in Subfigure~(R). In Figure~\ref{F-hmn3-L34-26n} the nodal set $\cZ(\zeta_t)$ has two connected components and deforms to (R) or (L).\medskip

\begin{figure}[!ht]
  \centering
  \includegraphics[width=0.99\textwidth]{hmn3-L34-24n}
  \caption{The point $s$  tends to $y$ clockwise (top) or counter-clockwise (bottom), here $k=4$}\label{F-hmn3-L34-24n}
\end{figure}

\begin{figure}[!ht]
  \centering
  \includegraphics[width=0.99\textwidth]{hmn3-L34-26n}
  \caption{The point $s$  tends to $y$ clockwise (top) or counter-clockwise (bottom), here $k=4$}\label{F-hmn3-L34-26n}
\end{figure}


 The nodal patterns $(L)$ and $(R)$ belong to a function $w_y$. We claim that they are different. For this purpose, we label the loops as in Paragraph~\ref{SSS-hmn-25B}, and we use the combinatorial type of the function $w_y$, see Paragraph~\ref{SSS-hmn-31bas}.

\begin{figure}[!ht]
  \centering
  \includegraphics[scale=0.25]{hmn3-L34-26LR}
  \caption{The nodal patterns $(L)$ and $(R)$ are different}\label{F-hmn3-L34-26LR}
\end{figure}

\FloatBarrier

The maps $\tau$ describing the combinatorial types of the nodal patterns $(L)$ and $(R)$ of Figure~\ref{F-hmn3-L34-26LR} are given by
%
\begin{equation*}
\tau_L = \begin{pmatrix}
           \downarrow & 1 & 2 & 3 & 4 & 5 \\
          5 & 2 & 1 & 4 & 3 & \downarrow \\
         \end{pmatrix}
\text{~~and~~}
\tau_R = \begin{pmatrix}
           \downarrow & 1 & 2 & 3 & 4 & 5 \\
           1 & \downarrow & 3 & 2 & 5 & 4 \\
              \end{pmatrix},
\end{equation*}
%
where $\downarrow$ corresponds to the arc hitting the boundary. Correspondingly, we label the nodal domains as in Paragraph~\ref{SSS-hmn-25C}, and we find the words,
%
\begin{equation*}
m_L = 121314 \text{~~and~~} m_R = 123242.
\end{equation*}

 The criterion  \eqref{E-invlab} given in Paragraph~\ref{SSS-hmn-25C} shows that the nodal patterns  $(L)$ and $(R)$ are different, although they should both be the nodal pattern of $w_y$, a contradiction. Recall that we already proved that $\dim V_{y,s} \le 2$. Since the assumption $\dim V_{y,s} = 2$ leads to a contradiction, at least in the example at hand, we conclude that  $\dim V_{y,s} = 1$. The proof in the general case follows the same lines, as in  Subsection~\ref{SS-hmn-25}. This proves Assertion~(i). \medskip

\begin{remark}\label{R-hmn3-L34-4}
When $t$ varies in $(-\pib,\pib)$, the nodal sets $\cZ(\zeta_{t})$
vary continuously with respect to the Hausdorff distance, see Lemma \ref{L-hdn}.
It follows that they are either all connected, or that they all have two connected components, so that their type in Figure~\ref{F-hmn3-L34-2} is either (b) \& (c)
or (d) \& (e).
\end{remark}%

\noid \emph{Proof of Assertion~(iii).~} This is a consequence of  Assertion~(i) and its proof.\medskip

\noid \emph{Proof of Assertion~(iv).~} The fact that $\dim V_{y,s} = 1$ implies that the linear system which defines $v_{y,s}$ up to scaling has constant rank, so that it has a solution which depends smoothly on the parameters $y, s$ locally. \medskip

Lemma~\ref{L-L34a} is proved.\hfill \qed


\subsubsection{Structure and combinatorial type of nodal sets in $V_{y,s}$}\label{SSS-hmn32bas}

The last two lines in \eqref{E-L34-4} give rise to three cases.\smallskip

\noi \textbf{Case~1}.~ $\rho(u,y) = 2k-3$, $\rho(u,s) = 1$,  and $\cS_b(u) = \set{y,s}$.  This means that $u \in W_y$, and this case only occurs when $s = z(y)$, see Figure~\ref{F-hmn3-L34-2}(a).\smallskip

\noi \textbf{Case~2}.~ $\rho(u,y) = 2k-4$, $\rho(u,s) = 2$, and $\cS_b(u) = \set{y,s}$, with two possibilities for $\cZ(u)$,
\begin{itemize}
  \item[$\diamond$] \emph{either} $\cZ(u)$ consists of $(k-2)$ loops at $y$ which do not intersect nor meet $\Gamma $ away from $y$, and one loop at $s$ which does not hit $\Gamma $ away from $s$, and does not meet the loops at $y$,
  \item[$\diamond$] \emph{or} $\cZ(u)$ consists of $(k-3)$ loops at $y$  which do not intersect nor meet $\Gamma $ away from $y$, and two simple arcs from $y$ to $s$ which do not meet except at $y$ and $s$, and do not meet the loops except at $y$; in this case we have a ``generalized loop'' which hits $\Gamma $ at $s$,
\end{itemize}
see Figures~\ref{F-hmn3-L34-2}(b) and \ref{F-hmn3-L34-2}(d).\smallskip

 \noi \textbf{Case~3}.~ $\rho(u,y) = 2k-4$, $\rho(u,s) = 1$, and there exists another point $s_1 \in \Gamma $, $s_1 \neq s,y$ such that  $\cS_b(u) = \set{y,s,s_1}$ and $\rho(u,s_1)=1$. In this case there are two possibilities for $\cZ(u)$,
\begin{itemize}
  \item[$\diamond$] \emph{either} $\cZ(u)$ consists of $(k-2)$ loops at $y$ which do not intersect nor meet $\Gamma $ away from $y$, and one arc from $s$ to $s_1$ which does not hit $\Gamma $ away from $s,s_1$, and does not meet the loops at $y$,
  \item[$\diamond$] \emph{or} $\cZ(u)$ consists of $(k-3)$ loops at $y$  which do not intersect nor meet $\Gamma $ away from $y$, and two simple arcs, one from $y$ to $s$ and one from $y$ to $s_1$ which do not meet except at $y$, and do not meet the loops except at $y$; in this case we have a ``generalized loop'' which contains a sub-arc of $\Gamma $ from $s$ to $s_1$,
\end{itemize}
see Figures~\ref{F-hmn3-L34-2}(c) and  \ref{F-hmn3-L34-2}(e).

\begin{figure}[!ht]
  \centering
  \includegraphics[width=0.5\textwidth]{hmn3-L34-2}
\caption{Nodal patterns for $u \in V_{y,s}$ (k=4)}\label{F-hmn3-L34-2}
\end{figure}


\begin{remarks}\label{R-hmn3-L34-2}\phantom{}
\begin{enumerate}[(i)]
  \item The subcases in Cases~2 and 3 are distinguished by the fact that $b_0(\cZ(u)) = 1$ (as in Figures~\ref{F-hmn3-L34-2}(d) \& (e)) or $b_0(\cZ(u)) = 2$, as in Figures~\ref{F-hmn3-L34-2}(b) \& (d), if $s \neq z(y)$. The subcases cannot occur simultaneously. Indeed, by Lemma~\ref{L-zeroc} we would otherwise find a function $u$ such that $\rho(u,y) \ge (2k-3)$ and $\rho(y,s) \ge 1$, with $s \neq z(y)$, contradicting Case~1.
  \item  At this stage of the discussion, the location of $z(y)$ with respect to $y$, $s$, and $s_1$ in Sub-figures~\ref{F-hmn3-L34-2}(b)--(e) is not clear. This will be explained in Lemma~\ref{L-L35a}.
\end{enumerate}
\end{remarks}%


For a pair $(y,s) \in \Gamma_{(2k -3)} \times \Gamma $ with $y \neq s$, and  $\dim V_{y,s} = 1$, we can define the \emph{combinatorial type} of a generator $v_{y,s}$ at the point $y$, as we did for the generator $w_y$ of $W_y$ in Paragraph~\ref{SSS-hmn-31bas}, taking the above cases into account.

When $(y,s) = (y,z(y))$, the combinatorial type is that of $w_y$, and we denote it by $\tau_{w_y}$. For example, the combinatorial type $\tau_{w_y}$ of the function $w_y$ whose nodal pattern appears in Figure~\ref{F-hmn3-L34-2}(a) is given by
%
\begin{equation*}
\tau_{a}^{\ref{F-hmn3-L34-2}} = \begin{pmatrix}
             \downarrow &  1 & 2 & 3 & 4 & 5 \\
               3 & 2 & 1 & \downarrow & 5 & 4 \\
             \end{pmatrix}\,,
\end{equation*}

When $s \neq y$, the combinatorial type $\tau := \tau_{v_{y,s}}$ of a function $v_{y,s}$ is described as follows. When the nodal set is connected,  as in Figure~\ref{F-hmn3-L34-2}(d) and (e), we write $\tau(j) = s$ (resp. $s_1$) to indicate that the nodal semi-arc emanating from $y$ tangentially to the ray labeled $j$ ends up at $s$ (resp. $s_1$). When the nodal set has two connected components, as in Figure~\ref{F-hmn3-L34-2}(b) and (c), we write $\tau(s)=s$ to indicate that there is a loop at $s$, and $\tau(s) = s_1$ to indicate that there is a nodal arc from $s$ to $s_1$. We describe the maps $\tau$ by $2\times (2k-2)$ matrices. The first row enumerates the rays at $y$ and the rays at $s$ and $s_1$ (counter-clockwise). With this convention, the combinatorial type $\tau_{v_{y,s}}$ of a function $v_{y,s}$ whose nodal pattern appears in Figure~\ref{F-hmn3-L34-2}(b)--(e), is given by one of the following formulas.
%
\begin{equation*}
\tau_{2,b}^{\ref{F-hmn3-L34-2}} = \begin{pmatrix}
               1 & 2 & 3 & 4 & s & s\\
               2 & 1 & 4 & 3 & s & s\\
             \end{pmatrix}\,,~~~
\tau_{2,c}^{\ref{F-hmn3-L34-2}} = \begin{pmatrix}
               1 & 2 & 3 & 4 & s_1 & s\\
               2 & 1 & 4 & 3 & s & s_1\\
             \end{pmatrix}\,,
\end{equation*}
%
and
%
\begin{equation*}
\tau_{1,d}^{\ref{F-hmn3-L34-2}} = \begin{pmatrix}
               1 & 2 & 3 & 4 & s & s\\
               s & 3 & 2 & s  & 1 & 4\\
             \end{pmatrix}\,,~~~
\tau_{1,e}^{\ref{F-hmn3-L34-2}} = \begin{pmatrix}
               1 & 2 & 3 & 4 & s_1 & s\\
               s_1 & 3  & 2 & s & 1 & 4\\
             \end{pmatrix}\,.
\end{equation*}
%


\subsubsection{Precise description of $V_{y,s}$}\label{SSS-hmn32bf}

 In this paragraph we analyze the behavior of a generator $v_{y,s}$ of $V_{y,s}$ when $y$ and $s$ vary. More precisely, Lemma~\ref{L-L34b} describes the behavior of $v_{y,s}$ when $y \in \Gamma_{(2k -3)}$ is fixed and $s$ tends to $y$ or to $z(y)$, and the behavior when $s \in \Gamma_{(2k -3)}$ is fixed and $y$ tends to $s$. Lemma~\ref{L-L35a} describes the global behavior of $v_{y,s}$ for a given $y \in \Gamma_{(2k -3)}$.

\begin{lemma}\label{L-L34b}
Let $U:=U(\lambda_k)$ with $k \ge 3$. Assume that  $\Omega$ is simply connected and  $\dim U = (2k-2)$. Given $(y,s) \in \Gamma_{(2k -3)} \times \Gamma $, with $ y \neq s$, recall that
%
\begin{equation*}
V_{y,s} := \set{u \in U ~|~ \rho(u,y) \ge 2k-4 \text{~and~} \rho(u,s) \ge 1}\,.
\end{equation*}
The functions in the one-dimensional space $V_{y,s}$ have the following properties.
  \begin{enumerate}[(i)]
  \item  When $s=z(y)$, the function $\breve{v}_{y,z(y)}=\breve{w}_y$ vanishes on $\partial \Omega$ precisely at the points $y$ and $z(y)$, and changes sign at these points. When $s \neq z(y)$ and $\cS_b(v_{y,s}) = \set{y,s,s'}$ with $s \neq s'$, $\rho(v_{y,s},s)=\rho(v_{y,s},s')=1$, the function $\breve{v}_{y,s}$ vanishes on $\partial \Omega$ precisely at the points $y, s$ and $s'$, does not change  sign at $y$, and changes sign at $s$ and $s'$. When $s \neq z(y)$ and $\cS_b(v_{y,s}) = \set{y,s}$ with  $\rho(v_{y,s},s)=2$, the function $\breve{v}_{y,s}$ vanishes on $\partial \Omega$ precisely at the points $y$ and $s$, and does not change sign.
   \item For fixed $y \in \Gamma_{(2k -3)}$, and $s$ close enough to $z(y)$, $\cS_b(v_{y,s}) = \set{y,s,s'}$, with $s' \neq s, y$,  and $\rho(v_{y,s},s) = 1$, $\rho(v_{y,s},s') = 1$. Furthermore, when $s$ tends to $z(y)$, $[v_{y,s}]$ tends to $[w_{y}]$, and $s'$ tends to $y$.
  \item For fixed $y \in \Gamma_{(2k -3)}$, and $s$ close enough to $y$, $\cS_b(v_{y,s}) = \set{y,s,s'}$, with $s' \neq s, y$,  and $\rho(v_{y,s},s) = 1$, $\rho(v_{y,s},s') = 1$. Furthermore, when $s$ tends to $y$, $[v_{y,s}]$ tends to $[w_{y}]$, and $s'$ tends to $z(y)$.
  \item  For fixed $s \in \Gamma_{(2k -3)}$, and $y\in \Gamma_{(2k -3)}$ close enough to $s$, $\cS_b(v_{y,s}) = \set{y,s,s'}$, with $s' \neq y,s$,  and $\rho(v_{y,s},s) = 1$, $\rho(v_{y,s},s') = 1$. Furthermore, when $y$ tends to $s$, $[v_{y,s}]$ tends to $[w_s]$, and $s'$ tends to $z(s)$.
\end{enumerate}
\end{lemma}%


\pf\phantom{}\\
\noid \emph{Proof of  Assertion~(i).~} Use Lemma~\ref{L-breve}, and the description of the possible nodal patterns for $V_{y,s}$ which follows from Lemma~\ref{L-L34a}.\medskip

\noid \emph{Proof of  Assertion~(ii).~} Assume that the first statement is not true. Then, we can find a sequence $\set{s_n}$ tending to $z(y)$, and a corresponding sequence $\set{v_n := v_{y,s_n}} \subset \bS(U)$ such that $\rho(v_n,y) = (2k-4)$,  $\rho(v_n,s_n) = 2$, and $v_n$ tends to some $v \in \bS(U)$. Since the convergence is uniform in $C^m$ for fixed $m \ge 0$, it follows that $\rho(v,y) \ge (2k-4)$ and $\rho(v,z(y)) \ge 2$, see Remark~\ref{R-evp2-2}  (lower semi-continuity of $\rho$),  and hence $v \in V_{y,z(y)}$. By Lemma~\ref{L-L34a}(iii), we must have $[v]=[w_y]$, and we reach a contradiction since $\rho(w_y,z(y))=1$. This proves the first statement. \smallskip

 We now prove the second statement. Considering $\set{v_n} \subset \bS(U)$ such that $\cS_b(v_n) = \set{y,s_n,s'_n}$ with $s_n \neq s'_n$ and $s_n$ tends to $z(y)$, we may also assume that $s'_n$ tends to some $s'$, and $v_n$ tends to $v$. Then, $\rho(v,y) \ge (2k-4)$, and $\rho(v,z(y)) \ge 1$,  $\rho(v,s') \ge 1$. Lemma~\ref{L-L34a}(iii) implies that $[v] =[w_y]$, and $\breve{v}(s') = 0$ so that $s' \in \set{y,z(y)}$ (where $\breve{v}$ is defined by \eqref{E-evp-dr}). Assuming that $s' = z(y)$, we would have $\rho(v,z(y)) = 2$, leading to a contradiction. Indeed, in that case,  $s_n$ and $s'_n$ would both tend to $z(y)$, and the derivative $\partial_b \breve{v}_n$ of the function $\breve{v}_n$ along the boundary would vanish at some $t_n$ on the smallest arc between $s_n$ and $s'_n$, with $t_n$ tending to $z(y)$. Passing to the limit, we would have $\partial_b \breve{w}_y(z(y)) = 0$, implying that $\rho(w_y,z(y)) \ge 2$. Assertion~(ii) is proved.\medskip


\noid \emph{Proof of  Assertion~(iii).~} Assume that the first statement is false. Then, we can find a sequence $\set{s_n}$ tending to $y$,  and a corresponding sequence $\set{v_n := v_{y,s_n}} \subset \bS(U)$ such that $\rho(v_n,y) = (2k-4)$,  $\rho(v_n,s_n) = 2$, and $v_n$ tends to some $v \in \bS(U)$. Then, $\rho(v,y) \ge (2k-4)$. Applying the local structure theorem to $v$, and following the arguments in the proof of Lemma~\ref{L-L38},  Part~(C), we may choose $r_1 > 0$ such that the neighborhood $D^{+}(r_1)$ of $y$ satisfies Properties~(B-a)--(B-d) in the proof of Lemma~\ref{L-L38}. We may also assume that the sequence $\set{s_n}$ is contained in $D^{+}(r_1)$. The nodal set of $v_n$ consists of either $(k-2)$ loops at $y$ and a loop at $s_n$, or $(k-3)$ loops at $y$ and two arcs from $y$ to $s_n$, see Figure~\ref{F-hmn3-L34-2}(b-d).
 As in the proof of Lemma~\ref{L-L38},  Part~(C), for $r_0 < r_1$ small enough, each loop at $y$ must cross $C^{+}(r_0)$ at (at least) two distinct points; so does each loop at $s_n$, and each arc from $y$ to $s_n$. We then conclude that $\rho(v,y) \ge (2k-2)$. Euler's formula then implies that actually $\rho(v,y) = (2k-2)$. This is not possible since $y \in \Gamma_{(2k -3)}$.   This proves the first statement in Assertion~(iii).\smallskip

We now prove the second statement. If $s_n$ tends to $y$, we have $\cS_b(v_n) = \set{y,s_n,s'_n}$, with $s_n \neq s'_n$. The preceding argument also shows that no subsequence of $s'_n$ can tend to $y$. We may then assume that $s_n$ tends to $y$ and $s'_n$ tends to some $s' \neq y$, with $v_n$ tending to some $v$. Then, $\rho(v,y) \ge (2k-4)$, and the previous argument shows that $\rho(v,y) \ge (2k-3)$, and $\rho(v,s') \ge 1$. Lemma~\ref{L-L32} then shows that $s' = z(y)$, and that $[v]=[w_y]$. Assertion~(iii) is proved. \medskip

\noid \emph{Proof of  Assertion~(iv).~} See Figure~\ref{F-hmn3-L34-8i}.\smallskip

\noi Assume that the first statement in false. Then, there exists a sequence $\set{y_n} \subset \Gamma_{(2k -3)}$ which tends to $s$, with a corresponding sequence $\ \set{v_n:=v_{y_n,s}} \subset \bS(U)$ tending to some $v \in \bS(U)$, and such that $\rho(v_{y_n,s},s) = 2$. The convergence of $v_n$ to $v$ being uniform in $C^{2k}$, we have $\rho(v,s) \ge (2k-4)$. Lemma~\ref{L-hdn} and Lemma~\ref{L-L34a}(ii) applied to $v_n$  imply that $ \cZ(v) \cup \Gamma $ is connected.  Applying Euler's formula to $v$, we conclude that $(2k-4) \le \sum_{z \in \cS_b(v)} \rho(v,z) \le (2k-2)$. Since the functions $\breve{v}_n$ do not change sign on $\Gamma $, $\breve{v}$ does not change sign either, so that $\rho(v,s) \neq (2k-3)$. Applying the local structure theorem to $v$ at $s$, and using the same proof as in Lemma~\ref{L-L38},  Part~(C),  (with the disc $D^{+}(s,r_0)$ and circle $C^{+}(s,r_0)$ centered at $s$), we infer that $\rho(v,s) = (2k-2)$.   Indeed, the nodal sets $\cZ(v_n)$ consist either is $(k-2)$ loops at $y_n$ (including a special loop touching $s$), or $(k-3)$ loops at $y_n$ and a loop at $s$. Since $y_n$ tends to $s$, for $r_0$ small enough, these loops must intersect $C^{+}(s;r_0)$ at $(2k-2)$ distinct points, and we can conclude as in the proof of Lemma~\ref{L-L38}.\smallskip

\noi To prove the second statement, we can now choose a sequence $\set{y_n}$ such that $\rho(v_n,s) = 1$ for $n$ large enough, so that $\cS_b(v_n) = \set{y_n,s,s'_n}$, with $s'_n \neq s$. An argument similar to the previous one, shows that no subsequence of $\set{s'_n}$ can tend to $s$. Since $s \in \Gamma_{(2k -3)}$, the only remaining possibility is that $\rho(v,s) = (2k-3)$, and hence that $v \in W_s$. Assertion~(iv) is proved. \medskip

The proof of Lemma~\ref{L-L34b} is complete.\hfill \qed \medskip

\begin{figure}[!ht]
  \centering
  \includegraphics[width=0.8\textwidth]{hmn3-L34-8i}
  \caption{Proof of Lemma~\ref{L-L34b}(iv)}\label{F-hmn3-L34-8i}
\end{figure}

We can enhanced the previous lemma by the following properties.


\begin{properties}\label{P-hmn3-P38}
Assume that $\Omega$ is simply connected, and that $\dim U = (2k-2)$. Given $(y,s) \in \Gamma_{(2k-3)} \times \Gamma $ with $y\neq s$, there exists $\varepsilon_0$ such that
\begin{enumerate}[(i)]
  \item for all $s \in \cA_c(z(y);\varepsilon_0)\cup \cA_c(y;\varepsilon_0) \sm \set{y,z(y)}$, $\rho(v_{y,s},s)=1$;
  \item for all $\varepsilon < \varepsilon_0$, there exists $\eta > 0$ such that for all $s \in \cA_c(z(y);\eta) \sm \set{z(y)}$, $\cS_b(v_{y,s}) = \set{y,s,s'}$ with $s' \in \cA(y;\varepsilon)\sm \set{y}$;
  \item for all $\varepsilon < \varepsilon_0$, there exists $\eta > 0$ such that for all $s \in \cA_c(y;\eta) \sm \set{y}$, $\cS_b(v_{y,s}) = \set{y,s,s'}$ with $s' \in \cA_c(z(y);\varepsilon)\sm \set{z(y)}$.
\end{enumerate}%

Let $s_1 \not \in \set{y,z(y)}$. Assume that the function $v_{y,s_1}$ satisfies $\rho(v_{y,s_1},s_1)=1$, ie $\cS_b(v_{y,s_1}) = \set{y,s_1,s'_1}$, with $s'_1 \neq s_1$. Then,
\begin{enumerate}[(a)]
  \item there exists $\varepsilon_1 > 0$, such that for all $s \in \cA_c(s_1;\varepsilon)$, $\rho(v_{y,s},s) = 1$;
  \item for all $\varepsilon > 0$, there exists $\eta < \varepsilon_1$ such that for all $s \in \cA_c(s_1;\eta)$, $s' \in \cA_c(s'_1;\varepsilon)$, i.e. the map $s \mapsto s'$ is continuous.
\end{enumerate}
\end{properties}%


\begin{lemma}\label{L-L35a}
 Assume that $\Omega$ is simply connected, and that $\dim U = (2k-2)$.  Let $y \in \Gamma_{(2k-3)}$. Then, the following properties hold.
\begin{enumerate}[(i)]
  \item For any $s \in \cA(y,z(y))$, the function $v_{y,s}$ satisfies $\cS_b(v_{y,s}) = \set{y,s,s'}$ with $s' \in \cA(y,z(y))$, possibly with $s=s'$.
  \item There exists  a unique $s_1 \in \cA(y,z(y))$ such that $v_{y,s_1}$ satisfies $\rho(v_{y,s_1},s_1) = 2$.
  \item For all $s \in \cA(s_1,z(y))$, $\cS_b(v_{y,s}) = \set{y,s,s'}$ with $s' \in \cA(y,s_1)$. Furthermore when $s$ moves counter-clockwise in $\cA(s_1,z(y))$, $s'$ moves clockwise in $\cA(y,s_1)$.
  \item For all $s \in \cA(y,s_1)$, $\cS_b(v_{y,s}) = \set{y,s,s'}$ with $s' \in \cA(s_1,z(y))$. Furthermore when $s$ moves counter-clockwise in $\cA(y,s_1)$, $s'$ moves clockwise in $\cA(s_1,z(y))$.
\end{enumerate}
Similar statements hold for the arc $\cA(z(y),y)$.
\end{lemma}%

The statements in Lemma~\ref{L-L35a} are illustrated in Figure~\ref{F-hmn3-L35a-1} (for the arc $\cA(y,z(y))$).  The corresponding nodal patterns appear in Figure~\ref{F-hmn3-L35a-9a}. \medskip

\begin{figure}[!ht]
  \centering
  \includegraphics[scale=0.25]{hmn3-L35a-1}
  \caption{ Lemma~\ref{L-L35a}: Assertions~(ii) and (iii)}\label{F-hmn3-L35a-1}
\end{figure}


\pf Choose a generator $w_y$ of $W_y$ such that $\breve{w}_y$ is positive in $\cA(y,z(y))$, and negative in $\cA(z(y),y)$.\smallskip

\emph{Assertion~(i).}~ Assume that the assertion is false: there exists some $s_0 \in \cA(y,z(y))$ such that $v_0 := v_{y,s_0}$ satisfies $\cS_b(v_0) = \set{y,s_0,s'_0}$, with $s'_0 \in \cA(z(y),y)$. Since $s_0 \neq s'_0$, Lemma~\ref{L-L34b}(i) implies that $\breve{v}_0$ vanishes on $\partial \Omega$ only at the points $y$, $s_0$ and $s'_0$, does not change sign at $y$, and changes sign at $s_0$ and $s'_0$. Choose $v_0$ such that $\breve{v}_0 < 0$ in $\cA(s_0,s'_0)$. For $s \neq y,z(y)$, introduce the function
%
\begin{equation*}
\xi_s = a_0(s) w_y - b_0(s) v_0\,,
\end{equation*}
%
where
%
\begin{equation*}
\left\{
\begin{array}{l}
a_0(s) = \breve{v}_0(s) \left( \breve{v}_0^2(s) + \breve{w}_y^2(s)
\right)^{-\frac 12}\,,\\[5pt]
b_0(s) = \breve{w}_y(s) \left( \breve{v}_0^2(s) + \breve{w}_y^2(s)
\right)^{-\frac 12}\,.
\end{array}
\right.
\end{equation*}
%
Since $b_0(s) \neq 0$, $\rho(\xi_s,y) = 2k-4$ and $\rho(\xi_s,s) \ge 1$, so that $\xi_s \in V_{y,s}$.\smallskip

Choose $s \in \cA(s_0,z(y))$. Since $\breve{v}_0$ only changes sign at $s_0$ and $s'_0$, $\breve{\xi}_s(s'_0) > 0$. By Lemma~\ref{L-breve}, at $y$ along $\partial \Omega$, $\breve{w}_y$ vanishes at order $(2k-3)$, while $\breve{v}_0$ vanishes at order $(2k-4)$. In a pointed neighborhood $J\sm \set{y}$ of $y$ in $\partial \Omega$, we have that $\breve{\xi}_s \sim - b_0(s) \breve{v}_0 < 0$, and hence $\breve{\xi}_s$ vanishes at some $s' \in \cA(s'_0,y)$. It follows that $\xi_s \in V_{y,s}$ with $\cS_b(\xi_s) = \set{y,s,s'}$, $\rho(\xi_s,s) = 1$, $\rho(\xi_s,s') = 1$.\smallskip

Similarly, choosing $t \in \cA(y,s_0)$, we have $\breve{\xi}_t(s'_0) < 0$ and $\breve{\xi}_t(z(y)) > 0$, and $\breve{\xi}_t$ vanishes at some $t' \in \cA(z(y),s'_0)$. \\ Finally, we conclude as above that $\xi_t \in V_{y,t}$ with $$\cS_b(\xi_t) = \set{y,t,t'}\,, \, \rho(\xi_t,t) = 1\,, \, \rho(\xi_t,t') = 1\,.
$$
These arguments can be visualized on Figure~\ref{F-hmn3-L35a-2}.%\smallskip

\begin{figure}[!ht]
  \centering
  \includegraphics[scale=0.25]{hmn3-L35a-2}%hmn3-L35a-2
  \caption{Proof of Lemma~\ref{L-L35a}(i)\\}\label{F-hmn3-L35a-2}
\end{figure}

From the assumed existence of $s_0$, we conclude that, for all $s \in \cA(y,z(y))$, $\xi_s \in V_{y,s}$ and $\cS_b(\xi_s) = \set{y,s,s'}$, with $s' \in \cA(z(y),y)$, $\rho(\xi_s,s) = 1$, $\rho(\xi_s,s') = 1$. Because $s \in \cA(s_0,z(y))$ implies that $s' \in \cA(s'_0,y)$, with a parallel statement for $t$, the previous argument also shows that when the point $s$ moves counter-clockwise in $\cA(y,z(y))$, the point $s'$ moves counter-clockwise in $\cA(z(y),y)$.  According to Lemma~\ref{L-L34b}, Assertions~(ii) and (iii), $\xi_s$ tends to $w_y$ when $s$ tends to $y$ or to $z(y)$ in $\cA(y,z(y))$. Looking at the behaviour of the nodal sets, we reach a contradiction as Figure~\ref{F-hmn3-L35a-4} shows. Assertion~(i) is proved. \hfill \qedc \medskip

\begin{figure}[!ht]
  \centering
  \includegraphics[width=0.9\textwidth]{hmn3-L35a-4}%hmn3-L35a-4
  \caption{Proof of Lemma~\ref{L-L35a}(i)}\label{F-hmn3-L35a-4}
\end{figure}

\textbf{Remark}. The previous arguments implicitly use the fact that the combinatorial type of $v_{y,s}$ does not change when $y$ is fixed and $s$ varies in $\Gamma  \sm \set{y}$ (the proof is similar to the proof of Lemma~\ref{L-L33}(iv)), see Paragraph~\ref{SSS-hmn32bas}. Note also that the reasoning in Figure~\ref{F-hmn3-L35a-4} is actually quite general, and only depends on the position of the arc from $y$ to $z(y)$ with respect to the loops.\medskip

\emph{Proof of Assertion~(ii).}

\begin{figure}[!ht]
  \centering
  \includegraphics[scale=0.25]{hmn3-L35a-6}
  \caption{Proof of Lemma~\ref{L-L35a}(ii)\\}\label{F-hmn3-L35a-6}
\end{figure}


Let $\gamma_1 : [0,\ell_1] \to \Gamma $ be an arc-length parametrization of $\Gamma $, such that $\gamma_1(t)$ moves counter-clockwise, and $\gamma_1(0) = \gamma_1(\ell_1) = y$, $\gamma_1(\ell) = z(y)$. \smallskip

According to Lemma~\ref{L-L34b}(i), and to the above Assertion~(i), taking $s \in \cA(y,z(y))$ close to $y$, we have $\cS_b(v_{y,s}) = \set{y,s,s'}$ with $s' \in \cA(y,z(y))$ close to $z(y)$. If $s=\gamma_1(t)$ and $s'=\gamma_1(t')$, for $t$ positive small enough, we have $t < t'$. Choose any such point $s_0 = \gamma_1(t_0)$ such that $v_0 = v_{y,s_0}$ satisfies $\cS_b(v_0) = \set{y,s_0,s'_0}$, with $s'_0 = \gamma_1(t'_0)$ and $t_0 < t'_0$. By Lemma~\ref{L-L34b}~(i), the function $\breve{v}_0$ vanishes precisely at the points $s_0$ and $s'_0$ and changes sign at these points. We choose it so that it is positive on $\cA(s_0,s'_0)$. Define $\xi_s$ as in the proof of Assertion~(i) (now with $s_0' \in \cA(y,z(y))$). Take $s \in \cA(y,s_0)$, $s=\gamma_1(t)$ with $0 < t < t_0$. With our previous choice of signs for $\breve{w}_y$, we find that $\breve{\xi}_s(z(y)) > 0$ and $\breve{\xi}_s(s'_0) < 0$, so that $\breve{\xi}_s$ vanishes at some $s'$ in $\cA(s'_0,z(y))$. Since $\xi_s \in V_{y,s}$, we have $\cS_b(\xi_s) = \set{y,s,s'}$, with $s \in \cA(y,s_0)$, $s' \in \cA(s'_0,z(y))$, $\rho(\xi_s,s)=\rho(\xi_s,s')=1$. Then, $s' = \gamma_1(t')$ with $0 < t < t_0 < t'_0 < t' < \ell$, so that the map $(0,t_0) \ni t \mapsto t'$ is decreasing.\smallskip

Introduce the set
%
\begin{equation*}
J := \set{t \in (0,\ell) ~|~ \cS_b(v_{y,\gamma_1(t)}) = \set{y,\gamma_1(t),\gamma_1(t')} \text{~with~} t < t'}\,.
\end{equation*}

If $t \in J$, $\rho(v_{y,\gamma_1(t)},\gamma_1(t)) = \rho(v_{y,\gamma_1(t)},\gamma_1(t')) = 1$. We have $(0,t_0) \subset J$ so that $J \neq \emptyset$, and hence $s_1 := \sup J$ exists. Since $\rho(v_{y,s},s)=1$ is an open condition, $s_1 \not \in J$. Take a subsequence $\set{t_n} \subset J$ tending to $s_1$, and choose corresponding functions $v_{y,t_n} \in \bS(U) \cap V_{y,t_n}$. A subsequence converges to some function $v_1$ in $\bS(U)$ which satisfies $\rho(v_1,y) \ge (2k-4)$ and $\rho(v_1,s_1) \ge 1$. By Lemma~\ref{L-L34a}, this implies that  $v_1 \in V_{y,s_1}$ (use Remark~\ref{R-evp2-2}). Since $s_1 \not \in J$, we must have $\rho(v_1,s_1)= 2$, so that $v_1$ is a generator $v_{y,s_1}$ of $V_{y,s_1}$. \medskip



\emph{Proof of Assertions~(iii) and (iv).} Take $v_1 = v_{y,s_1}$ given by Assertion (ii). By Lemma~\ref{L-L34b}~(i), we may choose this function such that $\breve{v}_1 \ge 0$ and vanishes only at $y$ and $s_1$.  For $s\neq y, z(y)$, introduce the functions,
%
\begin{equation*}
\xi_s = a_1(s) w_y - b_1(s) v_1\,,
\end{equation*}
%
where
%
\begin{equation*}
\left\{
\begin{array}{l}
a_1(s) = \breve{v}_1(s) \left( \breve{v}_1^2(s) + \breve{w}_y^2(s)
\right)^{-\frac 12}\,,\\[5pt]
b_1(s) = \breve{w}_y(s) \left( \breve{v}_1^2(s) + \breve{w}_y^2(s)
\right)^{-\frac 12}\,.
\end{array}
\right.
\end{equation*}
%


\begin{figure}[!ht]
  \centering
  \includegraphics[scale=0.25]{hmn3-L35a-8}
  \caption{Proof of Lemma~\ref{L-L35a}, Assertion~(iii)}\label{F-hmn3-L35a-8}
\end{figure}


When $s = s_1$, we have $\breve{\xi}_{s_1} = - b_1(s_1) \breve{v}_1 \le 0$, and $\breve{\xi}_{s_1}$ only vanishes at $y$ and $s_1$; $\xi_{s_1} \in V_{y,s_1}$. When $s \in \cA(y,z(y))\sm \set{s_1}$, taking into account the properties of the functions involved, we find that
%
\begin{equation*}
\left\{
\begin{array}{l}
\breve{\xi}_s(z(y)) < 0\,,\\[5pt]
\breve{\xi}_s(s_1) > 0\,,\\[5pt]
\breve{\xi}_s < 0 \text{~in~} J_y\sm\set{y}\,,\\[5pt]
\end{array}%
\right.
\end{equation*}
%
where $J_y$ is a small arc on $\Gamma_1$ centered at $y$.\smallskip

It follows that, for $s \in \cA(y,z(y))\sm \set{s_1}$, $\xi_s \in V_{y,s}$, $\cS_b(\xi_s) = \set{y,s,s'}$, $\rho(\xi_s,s) = \rho(\xi_s,s') = 1$, with the point $s$ on one side of $s_1$ in $\cA(y,z(y))$, and the point $s'$ on the other side. For these points $s, s'$, we have $V_{y,s} = V_{y,s'}$, so that $\xi_s$ and $\xi_{s'}$, must be proportional (Lemma~\ref{L-L34a}~(i)), and since the functions $\breve{\xi}_s$ and $\breve{\xi}_{s'}$ both take a positive value at $s_1$, $\xi_{s'} = a \, \xi_{s}$, with $a > 0$. \smallskip

We also have $\breve{\xi}_s(t) \breve{\xi}_t(s) \le 0$, since
%
\begin{equation*}
\breve{\xi}_s(t) \breve{\xi}_t(s) = - \left( \breve{v}_1^2(s) +
\breve{w}_y^2(s) \right)^{-\frac 12}\left( \breve{v}_1^2(t) +
\breve{w}_y^2(t) \right)^{-\frac 12}\left( \breve{v}_1(s)
\breve{w}_y(t)  - \breve{w}_y(s) \breve{v}_1(t) \right)^2\,.
\end{equation*}

Choose $s \in \cA(s_1,z(y))$, and $t \in \cA(s,z(y))$. Then $s',t' \in \cA(y,s_1)$.  The functions $\breve{\xi}_s$ and $\breve{\xi}_t$ are positive at $s_1$, and hence positive respectively on the arcs $\cA(s',s)$ and $\cA(t',t)$. We have $\breve{\xi}_s(t) < 0$, and hence, using the above properties, $\breve{\xi}_{s'}(t) > 0$, and $\breve{\xi}_t(s') < 0$. This implies that $t' \in \cA(y,s')$. We have proved that when $s$ moves counter-clockwise in $\cA(s_1,z(y))$, $s'$ move clockwise in $\cA(y,s_1)$. This is coherent with Lemma~\ref{L-L34b}(ii). The proof of Assertion~(iii) is similar.
\hfill  \qed

\begin{remarks} ~(i) Lemma~\ref{L-L35a} corresponds to the first part of the proof of Lemma~3.5 in \cite{HoMN1999} (from p.~1181, line (-7), ``We consider the function'' to p.~1182, line (+5), ``the following nodal domain'').

Note however that, in Lemma~\ref{L-L35a}, we do not assume that the point $z(y)$ is in the same connected component of $\Gamma_{(2k-3)}$ as the point $y$. \smallskip

(ii)  Figure~\ref{F-hmn3-L35a-9a} displays the possible nodal patterns of $w_y$, and the corresponding nodal patterns for the function $v_{y,s_1}$ with $s_1 \in \cA(y,z(y))$ (see  Lemma~\ref{L-L35a}(ii)), and for the function $v_{y,s}$  with $s \in \cA(s_1,z(y))$ (see Lemma~\ref{L-L35a}(iii)).
\end{remarks}%

\begin{figure}[!ht]
  \centering
  \includegraphics[width=0.8\linewidth]{hmn3-L35a-9a-LCR}
  \caption{Lemma~\ref{L-L35a}~(ii) and (iii): nodal patterns for $w_y$ and $v_{y,s_1}$}\label{F-hmn3-L35a-9a}
\end{figure}
\FloatBarrier
%\clearpage

\section{ Possible steps towards a proof that $\mult(\lambda_k) \le (2k-3)$ for $k \ge 3$}\label{S-hmn6}


Let $U := U(\lambda_k)$. throughout this section, we assume that
%
\begin{equation}\label{E-hmn6-0}
\left\{
\begin{array}{l}
\Omega \text{~is simply connected,}\\[5pt]
k \ge 3 \text{~~and~~} \dim U = (2k-2)\,.
\end{array}
\right.
\end{equation}

In \cite[Theorem~B]{HoMN1999} the authors claim that under the above assumption, the bound $\mult(\lambda_k) \le (2k-3)$ holds for $k \ge 3$ (see \cite[Section~4]{Berd2018} for comments on the non simply connected case).\smallskip

The purpose of this section is to revisit their scheme of proof, and point out some difficulties. We use the notation of the preceding sections.

\subsection{Preamble}\label{SS-hmn6-1}

We fix a direct orthonormal frame $\set{\vec{e}_1, \vec{e}_2}$ in $\R^2$, and the associated coordinates $\xi = (\xi_1,\xi_2)$. Without loss of generality, we may assume that the length of $\Gamma := \partial \Omega$ is equal to $2\pi$. We parametrize $\Gamma$ by arc-length, counter-clockwise, by some $\gamma_1 : [0,2\pi] \to \Gamma_1$, with unit tangent $\tau_1(t)$, and unit normal $\nu_1(t)$, so that the frame $\set{\tau_1(t),\nu_1(t)}$ is direct and $\nu_1$ points inside $\Omega$. For $\varepsilon$ small enough, we define the curve $c(t,\varepsilon) = \gamma_1(t) +\varepsilon\, \nu_1(t)$.\smallskip

We also fix $\set{\phi_1,\ldots,\phi_{(2k-2)}}$, an orthonormal basis of the eigenspace $U$. According to Lemma~\ref{L-L37}, we have a $C^{\infty}$ map $x \mapsto [u_x]$ from (the interior of) $\Omega$ to $\bP(U)$, where the function $u_x$ is a generator of the space $U_x$ defined by \eqref{E-hmn3-n6a}. The function $u_x$ is given by
%
\begin{equation}\label{E-hmn6-2}
u_x = \sum_{j=1}^{(2k-2)} \alpha_j(x) \phi_j
\end{equation}
%
where the nonzero vector $\alpha(x) := (\alpha_1(x),\ldots,\alpha_{(2k-2)}(x))$ is well defined up to scaling, given locally by a linear system, and $C^{\infty}$. Since $\Omega$ is simply connected, we can lift the map $x \mapsto \alpha(x)$ to a $C^{\infty}$ map from $\Omega$ to the unit sphere $\bS^{(2k-3)}$ or, equivalently, lift the map $x \mapsto [u_x]$ to a map $x \mapsto u_x$ from $\Omega$ to $\bS(U)$. The main property of this map $x \mapsto u_x$ is that $\nu(u_x,x) = (2k-2)$ for all $x \in \Omega$. The least order homogeneous term in the Taylor expansion of the function $u_x$ near $x$ has degree $(k-1)$. More precisely, according to Bers's Theorem,
%
\begin{equation}\label{E-hmn6-4}
h_x(\xi) := \frac{d^{(k-1)}}{dt^{(k-1)}}u_x(x+t\xi)\big|_{t=0}
\end{equation}
%
is a nonzero, harmonic, homogeneous polynomial of degree $(k-1)$ in $\xi$. The expression \eqref{E-hmn6-2} shows that the map $x \mapsto h_x$ is $C^{\infty}$. Since $h_x \neq 0$ for all $x$, we also obtain a map $\bar{h}$
%
\begin{equation}\label{E-hmn6-6}
\bar{h} : \left\{
\begin{array}{l}
x \mapsto \bar{h}_x \\[5pt]
\Omega \to S_{\cH}
\end{array}
\right.
\end{equation}
%
where $S_{\cH}$ is the unit circle in the $2$-dimensional vector space $\cH$ of harmonic, homogenous polynomials of degree $(k-1)$ in $\R^2$.\medskip

Let $\gamma$ be any Jordan curve in $\Omega$. Restricting the map $\bar{h}$ to $\gamma$, we obtain a map $s \mapsto \bar{h}_{\gamma(s)}$ from $\gamma$ to $S_{\cH}$. This map has a degree $\deg(\bar{h},\gamma)$. Since $\gamma$ is homotopic to a point, and since $\bar{h}$ is smooth, we conclude that $\deg(\bar{h},\gamma) = 0$.\medskip

In order to show that the assumption \eqref{E-hmn6-0} is absurd, hence that $\mult(\lambda_k) \le (2k-3)$ for any $k \ge 3$, it suffices to show that $\deg(\bar{h},c(\cdot,\varepsilon))\neq 0$ for $\varepsilon$ small enough. As we already pointed out in Remark~\ref{R-L38-2}, there is a difficulty here: the map $x \mapsto u_x$ extends to the boundary (see Lemma~\ref{L-L38}), and so does the map $x \mapsto h_x$. However, for $y \in \Gamma$, $\lim_{x\to y} h_x = 0$, and we have to look more closely at what happens near $\Gamma$, using Assertions~(ii) and (iii) of Lemma~\ref{L-L38}.\smallskip

The steps to achieve the previous goal are given in the next subsection, unfortunately with yet incomplete proofs.

\subsection{Steps for the proof that $\mult(\lambda_k) \le (2k-3)$ when $k \ge 3$}\label{SS-hmn6-2}

The following two lemmas do not appear in \cite{HoMN1999} and seem to be implicitly assumed.

\begin{lemma}\label{L-L35b0}
Assume that $\Omega$ is simply connected. Assume that  $k \ge 3$ and that $\dim U = (2k-2)$. Then,  $\Gamma_{(2k-2)}$ is not empty.
\end{lemma}

\pf We are only able to give a partial proof.\smallskip

\emph{Claim~1.~Assume that $\Gamma_{(2k-3)} = \Gamma $.  Then, $\inf_{y\in \Gamma} d(y,z(y)) >0$.}

Indeed, the map $y \mapsto z(y)$ (see Lemma~\ref{L-L33}) is continuous on $\Gamma$ and $z(y) \neq y$. The result follows by compactness. Another argument is as follows. Assume that $\inf_{y\in \Gamma} d(y,z(y)) = 0$. Then, there exists a sequence $y_n$ such that $0 < d(y_n,z(y_n)) < 1/n$. To the sequence $\set{y_n}$ we associate a sequence of functions $\set{w_n:=w_{y_n}} \in \bS(U)$ (Lemma~\ref{L-L32}). Up to taking subsequences, we may assume that $y_n$ converges to some $\bar{y}$, and that $w_n$ converges to some $w \in \bS(U)$. Since $y \mapsto z(y)$ is continuous,  $z(y_n)$ converges to $z(\bar{y})$, and the assumption on the sequence implies that $z(y) = y$.  By Lemma~\ref{L-L38}, $w \in W_{\bar{y}}$, and the arguments in the proof of this lemma imply that $\rho(w,\bar{y}) = (2k-2)$ which contradicts our assumption. \hfill \qedc \medskip

\emph{Claim~2.~Assume that $\Gamma_{(2k-3)} = \Gamma$. Then, there exists $\varepsilon_0 > 0$ such that for all $\varepsilon < \varepsilon_0$, and all $t \in [0,2\pi]$, $\cS_b(u_{c(t,\varepsilon)}) \neq \emptyset$.}

Assume this is not the case. Then, there exists a sequence $\set{\varepsilon_n}$ tending to zero, and a sequence $\set{t_n} \subset [0,2\pi]$ such that $u_n := u_{c(t_n,\varepsilon_n)}$ satisfies $\cS_b(u_{c(t_n,\varepsilon_n)}) = \emptyset$. Without loss of generality, we may assume that $\set{t_n}$ tends to some $\bar{t}$. Then, $x_n := c(t_n,\varepsilon_n)$ tends to $y = c(\bar{t},0) \in \Gamma$. It follows that $[u_n]$ tends to $[w_y]$ and the proof of Lemma~\ref{L-L38} implies that $ord(w_y,y) = (2k-2)$, contradicting our assumption. \hfill \qedc \medskip

\emph{Claim~3.~Assume that $\Gamma_{(2k-3)} = \Gamma$ Then, there exists $\varepsilon_0 > 0$ and $\delta_0 > 0$ such that for all $\varepsilon < \varepsilon_0$, and all $t \in [0,2\pi]$, $\cS_b(u_{c(t,\varepsilon)})$ has two elements whose distance on $\Gamma$ is bounded from below by $\delta_0$.}

Assume this is not the case. From Claim~1, we already know that  $\cS_b(u_{c(t,\varepsilon)}) \neq \emptyset$. Then there exists sequences $\set{\varepsilon_n}$ and $\set{\delta_n}$ tending to zero, and $\set{t_n}$ tending to some $\bar{t}$ such that $u_n := u_{c(t_n,\varepsilon_n)}$ has two boundary singular points $y_{n,1}$ and $y_{n,2}$ (possibly equal) whose distance is less than $\delta_n$. Taking subsequences, we may assume that these sequences tend to some $y \in \Gamma$, and the proof of Lemma~\ref{L-L38} shows that $u_n$ tends to $w_y$ with $\rho(w_y,y) = (2k-2)$, a contradiction. \hfill \qedc \medskip

Figure~\ref{F-hmn6-L35b0-2} displays the nodal sets $\cZ(u_n)$ (blue) and $\cZ(w_y)$ (red).

\begin{figure}[!ht]
  \centering
  \includegraphics[width=0.7\textwidth]{hmn6-L35b0-2}
  \caption{$\cZ(u_n)$ (blue) and $\cZ(w_y)$ (red)}\label{F-hmn6-L35b0-2}
\end{figure}
The assumption of the lemma implies that the combinatorial type of $w_y$ is constant along $\Gamma$, see Lemma~\ref{L-L33}, Assertion~(iv). This means that the nodal arc from $y$ to $z(y)$ makes a constant angle $\theta$ with the tangent to $\Gamma$ at $y$. \smallskip

 Here is a  \emph{heuristic} argument to proceed. \emph{For $\varepsilon$ small enough, the nodal set $\cZ(u_{c(t, \varepsilon)})$ contains a nodal arc which goes from $c(t,\varepsilon)$ to some $y_{c(t,\varepsilon),2}$ close to the point $z(\gamma_1(t))$. Furthermore, this nodal arc is close to the arc from $\gamma_1(t)$ to $z(\gamma_1(t))$ (in the Hausdorff distance). As illustrated in Figure~\ref{F-hmn6-L35b0-2}, its tangent at $c(t,\varepsilon)$ should make an angle $\theta_t$ with the tangent to $\Gamma$ at $\gamma_1(t)$, and $\theta_t$ should be close to $\theta$. For $\varepsilon$ small enough, the degree of the map $\bar{h}$ restricted to the curve $c(\cdot,\varepsilon)$ should therefore be one. A contradiction.} We however do not know how to turn this argument into a rigorous proof of the lemma.\smallskip

 The (unproven) arguments in \cite{HoMN1999} are as follows.\smallskip

p.~1184, line (-2),\\ ``We want to analyze how the star of tangents at $c(t,\varepsilon)$ to the $(k-1)$ nodal lines of $u_{c(t,\varepsilon)}$ crossing at $c(t,\varepsilon)$ turns if we follow $t$ from $0$ to $2\pi$."\smallskip

p.~1185, line (+1),\\ ``To make this precise, we consider the continuous function $f:\Omega \to \bS^1$ given by $f(x) = (2k-2)\alpha(x)$ modulo $2\pi$, where
%
$$
t \mapsto t \exp\left( \frac{2\pi i k}{(2k-2)} + i \alpha(x) \right), ~~k=0,\ldots , (2k-3), ~t\ge 0
$$
%
are the tangents rays of the nodal lines through $x \in \cZ(u_x)$. We want to analyze $f(c(t,\varepsilon))$."\smallskip

p.~1185, line (+ 9),\\ ``If $x \in \Omega$ is near enough to some point $y$ in the open set $\Gamma_{(2k-3)}$ then the nodal pattern $\cZ(u_x)$ can be read off the nodal pattern $\cZ(w_y)$, since the nodal domains move continuously. The nodal line leaving $\partial \Omega$ vertically stays connected to $\partial \Omega$ near $y$." \smallskip

p.~1185, line (+ 22),\\ ``This already implies that $f(c(t))$ follows the angle of $\partial D$ along each arc $(\partial D)_{2\ell -3}$." \medskip
\hfill \qed

\begin{lemma}\label{L-L35b1}
Assume that $\Omega$ is simply connected. Assume that  $k \ge 3$ and that $\dim U= (2k-2)$. Then,  $\Gamma_{(2k-2)}$ is not reduced to one point.
\end{lemma}

\pf The proof would be clear if we knew that  when $y \in \Gamma_{(2k-3)}$ tends to $\eta \in \Gamma_{(2k-2)}$, both $y$ and $z(y)$ are close to $\eta$, one on each side, compare with Remark~\ref{R-L352b-2}. For the time being, we do not have a proof of this lemma. \hfill \qed \medskip

The next lemma appears as Lemma~3.5 in \cite{HoMN1999}. It implicitly assumes that $\Gamma_{(2k-2)}$ contains at least two points.

\begin{lemma}\label{L-L35b2}
Assume that $\Omega$ is simply connected. Assume that  $k \ge 3$ and that $\dim U =(2k-2)$. Let $C$ be a connected component of $\Gamma_{(2k-3)}$, and $y \in C$. Then, $z(y) \not \in C$.
\end{lemma}%

\pf According to Lemmas~\ref{L-L35b0} and \ref{L-L35b1}, $\Gamma_{(2k-2)}$ has at least two elements.\smallskip

Let $C:=\cA(\eta_1,\eta_2)$ be a connected component of $\Gamma_{(2k-3)}$, and $y \in \cA(\eta_1,\eta_2)$. According to Lemma~\ref{L-L35a}~(ii), there exists (a unique) $s_1 \in \cA(y,z(y))$ such that the function $v_{y,s_1}$ satisfies $\rho(v_{y,s_1},s_1)=2$. Similarly, there is a unique $s_2$ in the arc $\cA(z(y),y)$ such that $\rho(v_{y,s_1},s_1)=2$. If $z(y) \in \cA(\eta_1,\eta_2)$, one of the points $s_1$ or $s_2$ also belongs to  $\cA(\eta_1,\eta_2)$. \smallskip

Without loss of generality, we may start out with a point $y_1 \in C$, such that $z(y_1) \in C$,  and there is a unique $s_1 \in \cA(y_1,z(y_1))$ such that $\rho(v_{y_1,s_1}, s_1)=2$. According to Assertion~(iv) of Lemma~\ref{L-L34b}, for $y \in \cA(y_1,s_1)$ close enough to $s_1$, we have $\rho(v_{y,s_1},s_1)=1$. \medskip

We now  come back to the analysis of Lemma 3.5 in \cite{HoMN1999} and consider the validity of the arguments given in \cite{HoMN1999}, starting from p.~1182, line (+6). \smallskip

The first one claims (in our notation):\\
{\it ``Then we move $y \in \cA(y_1,s_1)$ towards $s_1$ and consider $\cZ(v_{y,s_1})$. One of the nodal arcs hitting at $s_1$ must move away before $y$ hits $s_1$ since there is no point of $\Gamma_{(2k -2)}$ in between,...''.}\smallskip

Assuming $y$ is close enough to $s_1$ so that $\rho(v_{y,s_1},s_1)=1$, applying  Lemma~\ref{L-L34b}, Assertion~(iv), we have $\cS_b(v_{y,s_1}) = \set{y,s_1,s'(y)}$, and there are two possibilities, either $s'(y) \in \cA(y,s_1)$ or $s'(y) \not \in \cA(y,s_1)$. (We write $s'(y)$ to mention the dependence with respect to $y$ ($s_1$ being fixed).) \emph{This claim is correct}.\smallskip

Then, \cite{HoMN1999} claims that only the first case can occur for $y$ close enough to $s_1$, see p.~1182, line (+8),\\ ``and it must move eventually towards $y$ [\ldots] since the nodal type of $w_y$ [in our notation] is constant in $C$ [the connected component]''.\smallskip

This claim is questionable but let us continue the proof (see below the discussion of the second possibility) .\smallskip

{\bf First possibility.} Assuming that  the first possibility actually holds, namely that for $y$ close enough to $s_1$, $s'(y) \in \cA(y,s_1)$. Letting $y$ tend to $s_1$, we find some $y_2 \in \cA(y_1,z(y_1))$ such that $s'(y_2)=y_2$ which implies that $\rho ( v_{y_2,s_1} , y_2)=2k-3$, hence  $z(y_2) = s_1$, see Figure~\ref{F-hmn6-L35b2-2}.  Summing up, we have $y_2, z(y_2) \in \cA(y_1,z(y_1))$ and, measuring distances on the arc $\cA(y_1,z(y_1))$, $d(y_2,z(y_2)) < d(y_1,z(y_1))$. \smallskip

To end up the argument in \cite{HoMN1999}, p.~1182, line (+10),\\
``We have then the same situation \ldots a contradiction.'',\\
we can consider the set
%
$$
K:= \set{y \in \cA [y_1,z(y_1)] ~|~ z(y) \in \cA [y_1,z(y_1)]}
$$
%
and the number $d:= \inf\set{d(y,z(y)) ~|~ y \in K}$.  Since the set $K$ is non-empty and closed, the infimum is achieved.
 The previous argument shows that $d=0$, and hence that there exists some $y_0$ in $K$ such that $z(y_0)=y_0$ which contradicts the fact that $y_0 \in \Gamma_{(2k-3)}$. So finally, if $y \in \Gamma_{(2k-3)}$, $y$ and $z(y)$ cannot belong to the same connected component of $\Gamma_{(2k-3)}$. This would end the proof of Lemma~\ref{L-L35b2}, i.e., the proof of \cite{HoMN1999}, Lemma~3.5,  assuming that at each step only the first possibility occurs. \medskip

\begin{figure}[!ht]
  \centering
  \includegraphics[width=0.9\textwidth]{hmn6-L35b2-2}
  \caption{1st possibility}\label{F-hmn6-L35b2-2}
\end{figure}


{\bf Second possibility.} \cite{HoMN1999} seems to claim that this  possibility leads to an immediate contradiction and refers to  the nodal type of $w_y$. Figure~\ref{F-hmn6-L35b2-4} seems to invalidate this claim. As far as we understand,  Lemma~\ref{L-L35b2} remains unproven. \hfill \qed \medskip

\begin{figure}[!ht]
  \centering
  \includegraphics[width=0.9\textwidth]{hmn6-L35b2-4}%hmn3-L35a-8
  \caption{2nd possibility: apparent counterexample to the claim on the nodal type}\label{F-hmn6-L35b2-4}
\end{figure}

\begin{remark}\label{R-L352b-2}
It might be sufficient to prove the following weaker statement: if $C = \cA(\eta_1,\eta_2)$, when $y$ is close enough to $\eta_i$, the point $z(y)$ is not in $C$ (so that $y$ and $z(y)$ are not on the same side of $\eta_i$).
\end{remark}%


Lemma~3.6 in \cite{HoMN1999} contains the following assertions.
\begin{enumerate}[(i)]
  \item \it If $y$ moves clockwise to $\eta_1$ then  $z(y)$ moves anticlockwise to $\eta_1$.
  \item \it If $y$ moves anti-clockwise to $\eta_2$ then  $z(y)$ moves clockwise to $\eta_2$.
  \item \it In particular the nodal pattern of $\cZ (w_{\eta_1})$ and  $\cZ (w_{\eta_2})$ are different.
  \item \it Moreover  $\Gamma_{(2k-3)}$ consists of only finitely many open arcs and $\Gamma_{(2k-2)}$ is a finite set of $\Gamma$.
\end{enumerate}

\begin{figure}[!ht]
  \centering
  \includegraphics[width=0.7\textwidth]{hmn6-L35b2-R1}
  \caption{$Z(w_{\eta_1})$ and $Z(w_{\eta_2})$ have different combinatorial types}\label{F-hmn6-L35b2-R1}
\end{figure}


Some comments.
\begin{enumerate}[$\diamond$]
  \item If \emph{clockwise} and \emph{anticlockwise} are meant to indicate some monotonic behavior, Assertions (i) and (ii) are not proved in \cite{HoMN1999}. It is not clear whether knowing this precise behavior of the map $y \mapsto z(y)$ is necessary.
  \item Assuming Lemma~\ref{L-L35b2} (their Lemma~3.5) is true, and taking Lemma~\ref{L-L33}(iii) into account, we see that when $y \in \Gamma_{(2k-3)}$ tends to some $\eta \in \Gamma_{(2k-2)}$ the points $y$ and $z(y)$ do not lie on the same side of $\eta$. According to Remark~\ref{R-L352b-2}, this property seems sufficient to prove that the eigenfunctions $w_{\eta_1}$ and $w_{\eta_2}$ have different combinatorial types (same proof as in Subsection~\ref{SS-hmn-25}), giving a precise meaning to Assertion~(iii). This is illustrated in Figure~\ref{F-hmn6-L35b2-R1}. This fact does not seem to be used in the sequel.
  \item More interestingly, assuming that when $y \in \Gamma_{(2k-3)}$ tends to $\eta \in \Gamma_{(2k-3)}$ the points $y$ and $z(y)$ do not lie on the same side of $\eta$, we see that the combinatorial type of $w_y$ changes when $y$ passes through some $\eta \in \Gamma_{(2k-2)}$. This is illustrated in Figure~\ref{F-hmn6-L35b2-R2}
  \item Assertion~(iv) is proved in our Lemma~\ref{L-L33}(ii).
\end{enumerate}


\begin{figure}[!ht]
  \centering
  \includegraphics[width=0.7\textwidth]{hmn6-L35b2-R2}
  \caption{The combinatorial types of $Z(w_y)$ are different on the right and on the left of $\eta \in \Gamma_{(2k-2)}$}\label{F-hmn6-L35b2-R2}
\end{figure}

The last argument in the proof should be as follows.

\begin{lemma}\label{L-L50}
For $\varepsilon$ small enough, the smooth mapping $t\mapsto \bar{h}_{c(t,\varepsilon})$ from the circle $c(\cdot,\varepsilon)$ to $\bS_{\cH}$ has nonzero degree.
\end{lemma}

\pf We do not know how to prove this lemma. \smallskip

 The argument given in \cite[p.~1186]{HoMN1999} is:\smallskip

p.~1186, line (-5),\\ ``So finally, the smooth mapping $t \mapsto f(c(t,\varepsilon)), \bS^1 \to \bS^1$ has mapping degree $2\pi \#(\partial D)_{2\ell - 2} + 2\pi (2\ell -2) > 0$ and cannot be null homotopic. But by construction it is continuously extended into the interior of the circle and this is null homotopic, a contradiction."\smallskip

Note that we have the same difficulty as in the proof of Lemma~\ref{L-L35b0}.
\hfill \qed

\vspace{1cm}

\textbf{Concluding remarks}. Unless we modify the strategy of proof, we are led to two main questions.\medskip

\noi \emph{Question 1.} Given $\eta \in \Gamma_{(2k-2)}$ and $y \in \Gamma \sm \set{\eta}$ close to $\eta$, $z(y)$ is close to $\eta$ as well, so that $y$ and $z(y)$ are both in a small arc $\cA$ centered at $\eta$. Is it true that $y$ and $z(y)$ are not on the same connected component of $\cA \sm \set{\eta}$ (i.e., that they are on either sides of $\eta$)?\smallskip

If so, we could prove that the combinatorial type of $w_y$ (which is constant in each connected component of $\Gamma_{(2k-3)}$) changes each time $y$ crosses some $\eta \in \Gamma_{(2k-2)}$. This would in particular prove Lemma~\ref{L-L35b1}. \medskip

One way to prove Lemma~\ref{L-L35b0} and Lemma~\ref{L-L50} would be to better understand how the \emph{star} at $x$, i.e. the zero set of the harmonic polynomial $h_x$, changes (rotates) when $x \in \Omega$ tends to $y \in \Gamma_{(2k-3)}$. If $x_n$ tends to $y$, the stars at $x_n$ have limit points because they can be identified to sets of equidistant points on the unit circle.  This leads us to ask the following question.  \smallskip

\noi \emph{Question 2.} Does the star at $x$ have a (unique) limit when $x$ tends to $y \in \Gamma_{(2k-3)}$, and can we identify one branch of the limit star as the ray at $y$ corresponding to the nodal arc from $y$ to $z(y)$ ? \smallskip

One difficulty here is that Hausdorff convergence of nodal sets is not sufficient.




\vspace{2cm}
\newpage

\appendix

\section{A remark on upper bounds for multiplicities vs Courant-sharp eigenvalues}\label{S-mcs}

For simplicity, let us only consider Dirichlet eigenvalues in a smooth bounded domain $\Omega \subset \R^2$. The upper bounds on the eigenvalue multiplicities strongly rely on Courant's nodal domain theorem. As observed in \cite{HoMN1999}, they are actually a consequence of the following 3-step result, $n \in \set{1,2,3}$,  provided that the third step ($n=3$) is correct, see the final paragraph  in Section~\ref{S-hmn2}.

\begin{proposition}\label{P-propC1} Let $U(\lambda)$ is a Dirichlet eigenspace. Assume that $$\sup\set{\kappa(u) ~|~ u \in U(\lambda)} = \ell$$ for some $\ell \ge 3$. Then $\dim(U(\lambda)) \le (2\ell -n)$.
\end{proposition}

Indeed, since $u \in U(\lambda_k)$ implies that $\kappa(u) \le k$:
\begin{enumerate}
  \item The first step, $n=1$, yields the upper bound $\mult(\lambda_k \le (2k-1)$  (Theorem~1 in \cite{Nadi1987}).
  \item The second step, $n=2$, yields the upper bound $\mult(\lambda_k) \le (2k-2)$ (Lemma~2.13 in \cite{HoMN1999} or Proposition~\ref{P-hmn-s2}).
  \item The third step, $n=3$, yields the upper bound $\mult(\lambda_k) \le (2k-3)$ (Theorem~B in \cite{HoMN1999}).
\end{enumerate}

According to Pleijel \cite{Plej1956}, the equality in Courant's theorem can only occur for finitely many eigenvalues, the so-called \emph{Courant-sharp} eigenvalues. If $\lambda_k$ is not a Courant-sharp eigenvalue, in particular if $k$ is large enough (depending on the geometry of the domain), $u \in U(\lambda_k)$ implies that $\kappa(u) \le (k-1)$, and the above proposition implies that $\mult(\lambda_k) \le (2k-2-n)$. When $\lambda_k$ is not Courant-sharp, the upper bound for the multiplicity is improved by $2$.\medskip


The above proposition can be restated as
%
\begin{equation}\label{E-mcs-2}
\mult(\lambda_k) \leq 2 \sup_{u \in U(\lambda_k)} \kappa(u)\, -1\,.
\end{equation}
%
Note that this implies that
%
\begin{equation*}\label{E-mcs-2a}
\limsup_{k\rightarrow +\infty}\frac{\mult(\lambda_k)}{k} \le 2\,.
\end{equation*}

We can continue the discussion a little further recalling Pleijel's proof. Let $u$ be an eigenfunction in $U(\lambda_k)$ with precisely $\kappa(\lambda_k) := \sup_{u \in U(\lambda_k)} \kappa(u)$ nodal domains. Applying the Faber-Krahn inequality to each nodal domain, we obtain
%
\begin{equation*}\label{E-mcs-4f}
\lambda_k |\Omega| \ge \kappa(\lambda_k) \pi j_{0,1}^2\,,
\end{equation*}
%
where $|\Omega|$ denotes the area of $\Omega$. Then,
%
\begin{equation}\label{E-mcs-4}
\mult(\lambda_k) \leq \frac{2}{\pi j_{0,1}^2} \lambda_k |\Omega|\,  -1\,.
\end{equation}

For a regular bounded domain $\Omega$ in $\R^2$, Weyl's asymptotic formula reads
%
\begin{equation}\label{E-mcs-W}
N(\lambda) := \# \set{j~|~ \lambda_j < \lambda} = \frac{|\Omega|}{4\pi} \, \lambda + O(\sqrt{\lambda})\,.
\end{equation}

\begin{remark}\label{R-mcs-2}
For this formula, we refer to Ivrii's papers \cite{Ivri1980r, Ivri1980}, and to H\"{o}rmander's book \cite[Theorem~29.3.3]{Horm2009d}. These authors actually give a much more precise formula which yields a two-term asymptotic formula under an assumption on the set of periodic billiard trajectories (see Corollary~29.3.4 in H\"{o}rmander's book).
\end{remark}%

It follows from \eqref{E-mcs-W} that
%
\begin{equation}\label{E-mcs-6}
\limsup_{k\rightarrow +\infty} \frac{\kappa(\lambda_k)}{k} \le \frac{|\Omega|}{\pi j_{0,1}^2}\, \lim_{k\rightarrow +\infty} \frac{\lambda_k}{k} = \frac{4}{j_{0,1}^2} =: \gamma < 1.
\end{equation}

As a consequence, we obtain the better estimate (see \cite{Plej1956})
%
\begin{equation}\label{E-mcs-10}
\limsup_{k\rightarrow +\infty} \frac{\mult(\lambda_k)}{k} \le 2 \, \gamma < 2\,.
\end{equation}

On the other hand, for any $\delta$ small enough, $\mult(\lambda_k) = N(\lambda_k+\delta) - N(\lambda_k - \delta)$. By Weyl's asymptotic formula, we should have
%
$$
\mult(\lambda_k) = O(\sqrt{\lambda_k})\,,
$$
%
and hence
%
$$
\limsup_{k\rightarrow +\infty} \frac{\mult(\lambda_k)}{k} = 0\,,
$$
%
which shows that the inequality \eqref{E-mcs-10} is not pertinent.\medskip

For a bounded domain $\Omega \subset \R^n$ with $C^2$ boundary, and Dirichlet boundary condition, one can show that there exists a constant $C(\Omega)$, depending on $\Omega$ and invariant under dilations, such that for $k > C(\Omega)$, the $k$th Dirichlet eigenvalue $\lambda_k$ is not Courant-sharp. In dimension $2$, the constant $C(\Omega)$ can be estimated in terms of the area $|\Omega|$, the length $|\partial \Omega|$ of the boundary, the curvature and the cut-distance of $\partial \Omega$. We refer to \cite[Theorem~1.3]{BeHe2016} for more details, and to \cite[Theorem~1]{BeGi2016} for an extension to less regular domains. The proof of this result makes use of a lower bound on the remainder term $R(\lambda) := N(\lambda) - \frac{|\Omega|}{4\pi} \, \lambda$ in Weyl's asymptotic estimate, as given for example in \cite{BeLi2001}.

\begin{remark}\label{R-mcs-4}
For the case of domains with Neumann or Robin boundary condition, we refer to \cite[Proposition~1.1]{GiLe2020}. This paper also considers the special case of convex domains.
\end{remark}%


\begin{remark}\label{R-mcs-6}~ As in B\'erard-Helffer \cite{BeHe2016}  (see also \cite{BeGi2016}) we can think of a geometrical control of the constants (see  Safarov  \cite{Safa2001} or  Van den Berg-- Lianantonakis \cite{BeLi2001})
 who give estimates in the form
 $$
|N(\lambda) - C_n |\Omega| \lambda^{n/2}| \leq C_{geom} (\Omega) \lambda^{(n-1)/2} \ln \lambda\,.
$$
Although we only use lower bounds of the remainder  in the proof of Pleijel's theorem, we will also need here upper bounds.
\end{remark}


\section{The local structure of the nodal set in the neighborhood of a critical zero}\label{S-lsns}

In this Appendix, we provide a proof of the local structure theorem for the nodal set of an eigenfunction $u$ in the neighborhood of an interior singular point $x$. Cheng \cite{Chen1976} proposes a more precise result (existence of a local diffeomorphism which sends the nodal set of $u$ in a neighborhood of $x$ onto the nodal set of the higher order term in the Taylor expansion of $u$ at $x$ (a harmonic polynomial), which consists of rays. Our proof has the advantage of being more quantitative in particular when applied to eigenfunctions depending on a parameter. This proof is probably known although we did not find any reference except the analysis in \S\,128 of the book \emph{Georges Valiron, Cours d'analyse math\'{e}matique. Th\'{e}orie des fonctions. Masson, 1966.} cited in \cite{Bess1980}.

Let $(M,g)$ be a compact Riemannian surface, and let $V : M \to \R$ be a real valued smooth potential. Let $u \not = 0$ be a real valued function, satisfying
%
\begin{equation}\label{E-lsns-1}
(-\Delta_g + V)u = \lambda u
\end{equation}
%
for some real number $\lambda$. Here $\Delta_g$ denotes the Laplace-Beltrami operator of $(M,g)$. Let $x$ be a given (interior) point of $M$ (in case $M$ has a boundary).\medskip

Choose some $r_0 > 0$ such that the exponential map $\exp_x : T_xM \to M$ is a diffeomorphism from the disk $D(2r_0)$, with center $0$ and  radius $2r_0$ in $T_xM$, onto the geodesic disk $D(x,2r_0)$, with center $x$ and radius $2r_0$ in $M$. Choose an orthonormal frame in $T_xM$, call $(\xi_1,\xi_2)$ the corresponding coordinates in $T_xM$, and $(r,\omega)$ the associated polar coordinates. In the normal coordinates $(\xi_1,\xi_2)$, the Riemannian metric is given by the $2\times2$ matrix $G = (g_{ij})$, where $g_{ij} = g(\frac{\partial}{\partial \xi_i},\frac{\partial}{\partial \xi_j})$; the Riemannian measure is given by $v_g\, d\xi_1 d\xi_2$, where $v_g = \sqrt{\det G}$. Write the matrix $G^{-1}$ as $G^{-1} = (g^{ij})$. Then (see for example \cite[\S~2.89bis, p. 87]{GaHuLa2004}),
%
\begin{equation}\label{E-lsns-2a}
\left\{
\begin{array}{l}
G(0,0) = \id\,, \text{~~i.e.~~} g_{ij}(0,0) = \delta_{ij}\,,
\quad 1 \le i,j \le 2\,, \\[5pt]
\frac{\partial G}{\partial \xi_k}(0,0) = 0\,, \text{~~i.e.~~} \frac{\partial g_{ij}}{\partial \xi_k}(0,0) = 0\,, \quad 1 \le i,j,k \le 2.
\end{array}%
\right.
\end{equation}

It follows that
%
\begin{equation}\label{E-lsns-2b}
\left\{
\begin{array}{l}
v_g(0,0) = 1\,,\\[5pt]
G^{-1}(0,0) = \id\,, \text{~~i.e.~~} g^{ij}(0,0) = \delta_{ij}\,,
\quad 1 \le i,j \le 2\,, \\[5pt]
\frac{\partial v_g}{\partial \xi_k}(0,0) = 0\,, \quad 1 \le k \le 2. \\[5pt]
\frac{\partial G^{-1}}{\partial \xi_k}(0,0) = 0\,, \text{~~i.e.~~} \frac{\partial g^{ij}}{\partial \xi_k}(0,0) = 0\,, \quad 1 \le i,j,k \le 2.
\end{array}%
\right.
\end{equation}

Given a function $u$ on $M$, let $f = u\circ \exp_x$. In the local coordinates $(\xi_1,\xi_2)$, the Laplace-Beltrami operator $\Delta_g$ is given (see \cite[\S G.{III}, p.~126]{BeGM1971}) by
%
\begin{equation}\label{E-lsns-2c}
\left\{
\begin{array}{ll}
\Delta_g f & = v_g^{-1}\, \sum_{1 \le i,j \le 2\,} \frac{\partial}{\partial \xi_i}\left( v_g\, g^{ij} \frac{\partial f}{\partial \xi_j}\right)\,,\\[5pt]
& = \sum_{1 \le i,j \le 2\,} g^{ij}\, \frac{\partial^2 f}{\partial \xi_i \partial \xi_j} + \sum_{1 \le j \le 2\,}  b_j \, \frac{\partial f}{\partial \xi_j}\,,\\[5pt]
& \text{~~~where~}b_ j = \sum_{1 \le i \le 2\,} v_g^{-1}\, \frac{\partial}{\partial \xi_i} (v_g g^{ij})\,, \quad 1 \le j \le 2\,.
\end{array}%
\right.
\end{equation}

Letting $\Delta_0 = \sum_{1 \le j \le 2\,} \frac{\partial^2}{\partial \xi_j^2}$ denote the Laplacian in the Euclidean space $(T_xM,g_x)$, and taking relations \eqref{E-lsns-2a} and \eqref{E-lsns-2b} into account, we obtain the following expression for the Laplace-Beltrami operator,
%
\begin{equation}\label{E-lsns-2d}
\left\{
\begin{array}{l}
\Delta_g = \Delta_0 + \sum_{1 \le i,j \le 2\,} a_{ij}\, \frac{\partial^2}{\partial \xi_i \partial \xi_j} + \sum_{1 \le j \le 2\,} b_{j}\, \frac{\partial}{\partial \xi_j}\,,\\[5pt]
\text{where~~} \ord\left( a_{ij},(0,0) \right) = 2 \text{~and~} \ord\left( b_j,(0,0) \right))=1\,.
\end{array}
\right.
\end{equation}

If $u$ satisfies \eqref{E-lsns-1} and $u(x) = 0$, the unique continuation theorem \cite{Aron1957,DoFe1990a} implies that $f$ does not vanish at infinite order at $0$.  If $ord(u,x)= ord(f,0)=p$, Taylor's formula at $0$, gives
%
\begin{equation}\label{E-lsns-4}
\left\{
\begin{array}{l}
f(\xi_1,\xi_2) = \sum_{|\alpha| = p} \frac{1}{\alpha!}\, D^{\alpha}\!f(0,0)\,(\xi_1,\xi_2)^{\alpha} + R_{p+1}(\xi_1,\xi_2)\,, \text{~where}\\[8pt]
R_{p+1}(\xi_1,\xi_2) = \sum_{|\alpha| = p+1} \frac{p+1}{\alpha!}\, \,(\xi_1,\xi_2)^{\alpha} \, {\displaystyle \int_0^1 (1-t)^p\, D^{\alpha}\!f(t\,\xi_1,t\,\xi_2)\,dt}\,.
\end{array}%
\right.
\end{equation}
%
Here,  as usual,
%
\begin{equation}\label{E-lsns-4a}
\left\{
\begin{array}{l}
\alpha = (\alpha_1,\alpha_2)\,, \quad|\alpha| = \alpha_1 + \alpha_2\,, \quad (\xi_1,\xi_2)^{\alpha} = \xi_1^{\alpha_1}\, \xi_2^{\alpha_2}\,, \text{and}\\[5pt]
D^{\alpha}\!f = \dfrac{\partial^{|\alpha|}\!f}{\partial \xi_1^{\alpha_1}\partial \xi_2^{\alpha_2}}\,.
\end{array}%
\right.
\end{equation}

Using relations \eqref{E-lsns-1} and \eqref{E-lsns-2d}, and identifying the terms with lowest order, we find that the polynomial $P_p(\xi_1,\xi_2) := \sum_{|\alpha| = p} \frac{1}{\alpha!}\, D^{\alpha}\!f(0,0)\,(\xi_1,\xi_2)^{\alpha}$ is homogenous of degree $p$, and harmonic with respect to $\Delta_0$.


\begin{remark}\label{R-lsns-2}
As a matter of fact, we may write a $2$-term Taylor formula for the function $f$,
\begin{equation*}
f(\xi_1,\xi_2) = P_p(\xi_1,\xi_2) + P_{p+1}(\xi_1,\xi_2) + R_{p+1}(\xi_1,\xi_2)\,,
\end{equation*}
where $P_p$ and $P_{p+1}$ are homogeneous polynomials of degrees $p$ and $(p+1)$ respectively, and where the remainder term $R_{p+2}$ vanishes at order at least $(p+2)$. Then, we actually have that $\Delta_0 P_p = 0$ and $\Delta_0 P_{p+1} = 0$.
\end{remark}%

Writing the harmonicity condition $\Delta_0 P_p = 0$ in polar coordinates $(r,\omega)$ in $T_xM$, we find that the polynomial $P_p$ has the form
%
\begin{equation}\label{E-lsns-4b}
P_p(r\, \cos \omega, r\, \sin \omega) = a \, r^p \, \sin(p\,\omega - \omega_0)\,.
\end{equation}
%
for some $0 \not = a \in \R$ and some $\omega_0 \in [0,2\pi]$.\medskip

Multiplying the function $f$ by some constant, and rotating the coordinates $(\xi_1,\xi_2)$ in $\R^2$ if necessary, we can assume that $a=1$ and $\omega_0 = 0$. It follows that $f$ can be written as
%
\begin{equation}\label{E-lsns-6}
f(r\, \cos \omega, r\, \sin \omega) = r^p \, \sin(p\,\omega) + R_{p+1}(r\, \cos\omega, r\, \sin\omega)\,,
\end{equation}
%
and we can write the second term as
%
\begin{equation}\label{E-lsns-8}
\left\{
\begin{array}{l}
R_{p+1}(r\, \cos\omega, r\, \sin\omega) = r^{p+1}T_{p+1}(r\,, \omega)\,, \text{~where~} T_{p+1}(r,\omega) \text{~is given by}\\[5pt]
\sum_{|\alpha| = p+1} \frac{p+1}{\alpha!}\, \,(\cos \omega,\sin \omega)^{\alpha} \, {\displaystyle \int_0^1 (1-t)^p\, D^{\alpha}\!f(tr\,\cos\omega,tr\,\sin\omega)\,dt}\,.
\end{array}%
\right.
\end{equation}

Define
%
\begin{equation}\label{E-lsns-10}
W(r,\omega) := \sin(p\, \omega) + r\, T_{p+1}(r,\omega)\,.
\end{equation}

The function $W(0,\omega)$ vanishes precisely for the values
%
\begin{equation}\label{E-lsns-12}
\omega_j := j\frac{\pi}{p}\,, \quad j \in \set{0,\ldots,2p-1}.
\end{equation}

Choose $\alpha_1 \in \left( 0,\frac{\pi}{8} \right)$ and define
$\alpha_p := \frac{\alpha_1}{p}$. We have the following relations.



\begin{equation}\label{E-lsns-14}
\left\{
\begin{array}{l}
\sin\left(p\,(\omega_j \pm \alpha_p)\right) = \pm (-1)^j \, \sin(\alpha_1)\,,\\[5pt]
|\sin(p\, \omega)| \ge \sin \alpha_1\,, \text{~for~} \omega \not \in \bigcup_{j=0}^{2p-1}\left( \omega_j-\alpha_p,\omega_j+\alpha_p\right)
\end{array}%
\right.
\end{equation}

Define
%
\begin{equation}\label{E-lsns-16}
r_1 := \min\set{r_0 \,,\, \frac 12 \, \sin(\alpha_1)\, \|T_{p+1}\|^{-1}_{\infty,D(r_0)}\,,\frac 12 \, \cos(\alpha_1)\, \|\partial_{\omega}T_{p+1}\|^{-1}_{\infty,D(r_0)}}\,,
\end{equation}
%
where $\|\cdot\|_{\infty,D(\frac 12)}$ denotes the $L^{\infty}$ norm of functions on the disk $D(r_0)$ of radius $r_0$ in $TxM$.


\begin{proposition}\label{P-lsns-2}
For any $r \le r_1$,
\begin{enumerate}
  \item[(i)] the function $\omega \mapsto W(r,\omega)$ does not vanish in $$
      [0,2\pi] \sm \bigcup_{j=0}^{2p-1}\left( \omega_j-\alpha_p,\omega_j+\alpha_p\right) = \bigcup_{j=0}^{2p-1}\left[ \omega_j+\alpha_p,\omega_{j+1}-\alpha_p\right]\,;
$$
  \item[(ii)]for each $j \in \set{0,\ldots,2p-1}$, the function $\omega \mapsto W(r,\omega)$ has exactly one zero
  $\omegat_j(r) \in \left( \omega_j - \alpha_p \,,\, \omega+\alpha_p \right)$;
  \item[(iii)] for each $j \in \set{0,\ldots,2p-1}$, the function $r \mapsto \omegat_j(r)$ is $C^{\infty}$ in $(0, r_1)$ and tends to $\omega_j$ as $\rho$ tends to zero;
  \item[(iv)] for each $j \in \set{0,\ldots,2p-1}$, the curve
  $$
  \left( 0 , r_1\right) \ni r \mapsto a_j(r) = \left( r\, \cos(\omegat_j(r) \,,\, r \, \sin(\omegat_j(r)  \right)
  $$
  is smooth and has semi-tangent $\omega_j$ at the origin.
\end{enumerate}
\end{proposition}%

\pf To prove (i), we observe that in each interval $\set{r} \times \left[ \omega_j+\alpha_p,\omega_{j+1}-\alpha_p\right]$, $|W(r,\omega)| \ge \frac 12 \sin(\alpha_1)$. To prove (ii), we observe that the function $W(r,\omega)$ changes sign in $\set{r}\times \left( \omega_j-\alpha_p,\omega_j+\alpha_p\right)$ and that its partial derivative with respect to $\omega$ does not vanish.  Assertion~(iii) follows from the implicit function theorem. Assertion~(iv) follows from the previous ones. \hfill \qed

\begin{remark}\label{R-lsns-4}
Later on, we will encounter the following situation. There exist two eigenfunctions $v_1$ and $v_2$ such that the functions $h_i = v_i \circ \exp_x$ satisfy the relations
%
\begin{equation}%\label{}
\left\{
\begin{array}{l}
h_1(r\, \cos\omega, r\, \sin\omega) = r^p \, \sin(p\,\omega) + R_{1,p+1}(r\, \cos \omega, r\, \sin\omega)\,,\\[5pt]
h_2(r\, \cos\omega, r\, \sin\omega) = r^p \, \cos(p\,\omega) + R_{2,p+1}(r\, \cos \omega, r\, \sin\omega)\,.
\end{array}%
\right.
\end{equation}
%
Defining the family of functions $w_{\theta} = \cos\theta \, v_1 - \sin\theta \, v_2$, the associated family of functions $h_{\theta} = w_{\theta}\circ \exp_x$, satisfies
%
\begin{equation}%\label{}
\left\{
\begin{array}{l}
h_{\theta}(r\, \cos\omega, r\, \sin\omega) = r^p \, \sin(p\,(\omega - \theta)) + R_{\theta,p+1}(r\, \cos\omega, r\, \sin\omega)\,, \text{~with}\\[5pt]
R_{\theta,p+1} = \cos\theta \, R_{1,p+1} - \sin\theta \, R_{2,p+1}.
\end{array}%
\right.
\end{equation}
%
Then, Proposition~\ref{P-lsns-2} remains valid for the family $h_{\theta}$, uniformly with respect to the variable $\theta \in [0,2\pi]$, with $\omega_j$ replaced by $(\omega_j-\theta)$, and the corresponding functions $r \mapsto \omegat_j(r,\theta)$ and $r \mapsto a_j(r,\theta)$ are smooth in $(r,\theta)$.
\end{remark}%


\begin{remark}\label{R-lsns-6}
Proposition~\ref{P-lsns-2} tells us that, in a neighborhood of the critical zero $x$ of $u$, the nodal set $\cZ(u)$ consists of $p$ smooth semi-arcs emanating from $x$ tangentially to the rays $\omega_j$.
\end{remark}%


\section{Local structure of the nodal set near a critical zero, another approach}\label{S-lsc}

In this Appendix, we use a local conformal chart to determine the local structure of the nodal set $\cZ(u)$ near an interior critical zero of $u$, instead of normal coordinates as in Appendix~\ref{S-lsns}.

\begin{remark}\label{R-Cb}
Using ``boundary isothermal coordinates'' as given in \cite[Section~2]{YaZh2021}, one can also determine the local structure of the nodal set $\cZ(u)$ near a boundary singular point of $u$ (assuming Dirichlet, Neumann or Robin boundary condition).
\end{remark}%

Let $(M,g)$ be a compact Riemannian surface, and let $V : M \to \R$ be a real valued smooth (ie, $C^\infty$) potential. Let $u \not = 0$, a real valued function, satisfying $(-\Delta + V)u = \lambda u$, for some real number $\lambda$. Let $x$ be a given (interior) point of $M$ (in case $M$ has a boundary). Assume that $u(x) = 0$.\smallskip

Choose some $r_0 > 0$ such that the geodesic disk $D(x,r_0)$ is contained in a conformal coordinate chart centered at $x$: there exists some open neighborhood $\cU_x$ of $x$, such that $D(x,r_0) \subset \cU_x \subset M$, and a conformal diffeomorphism $\phi : \cU_x \to \D$, such that $\phi(x) = 0$ and $\phi^{*}(d\xi_1^2 + d\xi_2^2) = \sigma(\psi(\xi_1,\xi_2)) \, g$. Here $\D$ is the unit disk in $\R^2$,  $(\xi_1,\xi_2)$ denote the coordinates in $\R^2$, $\psi = \phi^{-1}$, and $\sigma$ is a smooth positive function on $\cU_x$.

Then, we have
%
\begin{equation}\label{E-lsc-2}
- \Delta_0 f = \sigma \, \left( \lambda - V\circ \psi\right)\, f\,,
\end{equation}
%
where, $f = u\circ \psi$, and $\Delta_0$ is the Laplacian in $\R^2$ with the metric $d\xi_1^2 + d\xi_2^2$. To determine the local structure of the nodal set $\cZ(u)$ in a neighborhood of $x \in \inte(M)$, is sufficient to determine the structure of the zero set of $f$ in a neighborhood of $0$ in $\D$.\medskip

The unique continuation theorem \cite{Aron1957,DoFe1990a} implies that $f$ (equivalently $u$) does not vanish at infinite order at $0$.  If $ord(u,x)= ord(f,0)=p$, Taylor's formula at $0$, gives
%
\begin{equation}\label{E-lsc-4}
\left\{
\begin{array}{l}
f(\xi_1,\xi_2) = \sum_{|\alpha| = p} \frac{1}{\alpha!}\, D^{\alpha}\!f(0,0)\,(\xi_1,\xi_2)^{\alpha} + R_{p+1}(\xi_1,\xi_2)\,, \text{~where}\\[8pt]
R_p{p+1}(\xi_1,\xi_2) = \sum_{|\alpha| = p+1} \frac{p+1}{\alpha!}\, \,(\xi_1,\xi_2)^{\alpha} \, {\displaystyle \int_0^1 (1-t)^p\, D^{\alpha}\!f(t\,\xi_1,t\,\xi_2)\,dt}\,.
\end{array}%
\right.
\end{equation}
%
Here,  as usual,
%
\begin{equation}\label{E-lsc-4a}
\left\{
\begin{array}{l}
\alpha = (\alpha_1,\alpha_2)\,, \quad|\alpha| = \alpha_1 + \alpha_2\,, \quad (\xi_1,\xi_2)^{\alpha} = \xi_1^{\alpha_1}\, \xi_2^{\alpha_2}\,, \text{and}\\[5pt]
D^{\alpha}\!f = \dfrac{\partial^{|\alpha|}\!f}{\partial \xi_1^{\alpha_1}\partial \xi_2^{\alpha_2}}\,.
\end{array}%
\right.
\end{equation}

Using \eqref{E-lsc-2}, and identifying the terms with lowest order, we find that the polynomial $P_p(\xi_1,\xi_2) := \sum_{|\alpha| = p} \frac{1}{\alpha!}\, D^{\alpha}\!f(0,0)\,(\xi_1,\xi_2)^{\alpha}$ is homogenous of degree $p$, and harmonic (with respect to $\Delta_0$). Writing the harmonicity condition in polar coordinates $(\rho,\omega)$ in $\R^2$, we find that the polynomial $P_p$ has the form
%
\begin{equation}\label{E-lsc-4b}
P_p(\rho\, \cos \omega, \rho\, \sin \omega) = a \, \rho^p \, \sin(p\,\omega - \omega_0)\,.
\end{equation}
%
for some $0 \not = a \in \R$ and some $\omega_0 \in [0,2\pi]$.\medskip

Multiplying the function $f$ by some constant, and rotating the coordinates $(\xi_1,\xi_2)$ in $\R^2$ if necessary, we can assume that $a=1$ and $\omega_0 = 0$. It follows that $f$ can be written as
%
\begin{equation}\label{E-lsc-6}
f(\rho\, \cos \omega, \rho\, \sin \omega) = \rho^p \, \sin(p\,\omega) + R_{p+1}(\rho\, \cos \omega, \rho\, \sin \omega)\,,
\end{equation}
%
and we can write the second term as
%
\begin{equation}\label{E-lsc-8}
\left\{
\begin{array}{l}
R_{p+1}(\rho\, \cos \omega, \rho\, \sin \omega) = \rho^{p+1}T_{p+1}(\rho\,, \omega)\,, \text{~where~} T_{p+1}(\rho,\omega) \text{~is given by}\\[5pt]
\sum_{|\alpha| = p+1} \frac{p+1}{\alpha!}\, \,(\cos \omega,\sin \omega)^{\alpha} \, {\displaystyle \int_0^1 (1-t)^p\, D^{\alpha}\!f(t\rho\,\cos\omega,t\rho\,\sin\omega)\,dt}\,.
\end{array}%
\right.
\end{equation}

Define
%
\begin{equation}\label{E-lsc-10}
W(\rho,\omega) := \sin(p\, \omega) + \rho\, T_{p+1}(\rho,\omega)\,.
\end{equation}

The function $W(0,\omega)$ vanishes precisely for the values
%
\begin{equation}\label{E-lsc-12}
\omega_j := j\frac{\pi}{p}\,, \quad j \in \set{0,\ldots,2p-1}.
\end{equation}

Choose $\alpha_1 \in \left( 0,\frac{\pi}{8} \right)$ and define
$\alpha_p := \frac{\alpha_1}{p}$. We have the following relations.



\begin{equation}\label{E-lsc-14}
\left\{
\begin{array}{l}
\sin\left(p\,(\omega_j \pm \alpha_p)\right) = \pm (-1)^j \, \sin(\alpha_1)\,,\\[5pt]
|\sin(p\, \omega)| \ge \sin \alpha_1\,, \text{~for~} \omega \not \in \bigcup_{j=0}^{2p-1}\left( \omega_j-\alpha_p,\omega_j+\alpha_p\right)
\end{array}%
\right.
\end{equation}

Define
%
\begin{equation}\label{E-lsc-16}
\left\{
\begin{array}{l}
\rho_0 := \min\set{\frac 12 \,,\, \frac 12 \, \sin(\alpha_1)\, \|T_{p+1}\|^{-1}_{\infty,D(\frac 12)}}\,,\\[5pt]
\rho_1 := \min\set{\frac 12 \,,\, \frac 12 \, \cos(\alpha_1)\, \|\partial_{\omega}T_{p+1}\|^{-1}_{\infty,D(\frac 12)}}\,,\\[5pt]
\end{array}%
\right.
\end{equation}
%
where $\|\cdot\|_{\infty,D(\frac 12)}$ denotes the $L^{\infty}$ norm of functions on the disk $D(\frac 12)$ of radius $\frac 12$.


\begin{proposition}\label{P-lsc-2}
For any $\rho \le \min\set{\rho_0,\rho_1}$,
\begin{enumerate}
  \item[(i)] the function $\omega \mapsto W(\rho,\omega)$ does not vanish in $$
      [0,2\pi] \sm \bigcup_{j=0}^{2p-1}\left( \omega_j-\alpha_p,\omega_j+\alpha_p\right) = \bigcup_{j=0}^{2p-1}\left[ \omega_j+\alpha_p,\omega_{j+1}-\alpha_p\right]\,;
$$
  \item[(ii)]for each $j \in \set{0,\ldots,2p-1}$, the function $\omega \mapsto W(\rho,\omega)$ has exactly one zero
  $\omegat_j(\rho) \in \left( \omega_j - \alpha_p \,,\, \omega+\alpha_p \right)$;
  \item[(iii)] for each $j \in \set{0,\ldots,2p-1}$, the function $\rho \mapsto \omegat_j(\rho)$ is $C^{\infty}$ in the interval $\left( 0 \,,\, \min\set{\rho_0,\rho_1}\right)$ and tends to $\omega_j$ as $\rho$ tends to zero;
  \item[(iv)] for each $j \in \set{0,\ldots,2p-1}$, the curve
  $$
  \left( 0 \,,\, \min\set{\rho_0,\rho_1}\right) \ni \rho \mapsto a_j(\rho) = \left( \rho \, \cos(\omegat_j(\rho) \,,\, \rho \, \sin(\omegat_j(\rho)  \right)
  $$
  is smooth and has semi-tangent $\omega_j$ at the origin.
\end{enumerate}
\end{proposition}%

\pf To prove (i), we observe that in each interval $\set{\rho} \times \left[ \omega_j+\alpha_p,\omega_{j+1}-\alpha_p\right]$, $|W(\rho,\omega)| \ge \frac 12 \sin(\alpha_1)$. To prove (ii), we observe that the function $W(\rho,\omega)$ changes sign in $\set{\rho}\times \left( \omega_j-\alpha_p,\omega_j+\alpha_p\right)$ and that its partial derivative with respect to $\omega$ does not vanish.  Assertion~(iii) follows from the implicit function theorem. Assertion~(iv) follows from the previous ones. \hfill \qed

\begin{remark}\label{R-lsc-2}
Later on, we will encounter the following situation. There exist two eigenfunctions $v_1$ and $v_2$ such that the functions $h_i = v_i \circ \psi$ satisfy the relations
%
\begin{equation}%\label{}
\left\{
\begin{array}{l}
h_1(\rho\, \cos \omega, \rho\, \sin \omega) = \rho^p \, \sin(p\,\omega) + R_{1,p+1}(\rho\, \cos \omega, \rho\, \sin \omega)\,,\\[5pt]
h_2(\rho\, \cos \omega, \rho\, \sin \omega) = \rho^p \, \cos(p\,\omega) + R_{2,p+1}(\rho\, \cos \omega, \rho\, \sin \omega)\,.
\end{array}%
\right.
\end{equation}
%
Defining the family of functions $w_{\theta} = \cos\theta \, v_1 - \sin\theta \, v_2$, the associated family of functions $h_{\theta} = w_{\theta}\circ \psi$, satisfies
%
\begin{equation}%\label{}
\left\{
\begin{array}{l}
h_{\theta}(\rho\, \cos \omega, \rho\, \sin \omega) = \rho^p \, \sin(p\,(\omega - \theta)) + R_{\theta,p+1}(\rho\, \cos \omega, \rho\, \sin \omega)\,, \text{~with}\\[5pt]
R_{\theta,p+1} = \cos\theta \, R_{1,p+1} - \sin\theta \, R_{2,p+1}.
\end{array}%
\right.
\end{equation}
%
Then, Proposition~\ref{P-lsc-2} remains valid for the family $h_{\theta}$, uniformly with respect to the variable $\theta \in [0,2\pi]$, with $\omega_j$ replaced by $(\omega_j-\theta)$, and the corresponding functions $\rho \mapsto \omegat_j(\rho,\theta)$ and $\rho \mapsto a_j(\rho,\theta)$ are smooth in $(\rho,\theta)$.
\end{remark}%


\begin{remark}\label{R-lsc-4}
Proposition~\ref{P-lsc-2} tells us that, in a neighborhood of the critical zero $x$ of $u$, the nodal set $\cZ(u)$ consists of $p$ smooth semi-arcs emanating from $x$ tangentially to the rays $\omega_j$.  An alternative proof is to take exponential coordinates at $x$, and polar coordinates in $T_xM$, as in Appendix~\ref{S-lsns}. Cheng \cite{Chen1976} gives a more precise, yet less quantitative, result by showing that there exists a local diffeomorphism sending the nodal set of $u$ in a neighborhood of the singular point $x$ onto the nodal set of higher order terms of the Taylor expansion of $u$ at $x$ (a harmonic polynomial) which consists of rays.
\end{remark}%



\section{Labeling loops and nodal domains}\label{S-llnd}

In this Appendix, $(M,g)$ is a surface of genus $0$ and we consider eigenfunctions of the operator $-\Delta + V$, where $g$ is a $C^{\infty}$ metric and $V$ a $C^{\infty}$ real valued potential.

 Similar arguments can be applied to smooth bounded domains in $\R^2$.


\subsection{Labeling of loops and bouquets of loops}\label{SS-llnd-w}

Let $u$ be an eigenfunction such that $\ord(u,x) =p$, with $x \in \inte(M)$. Furthermore, assume that the point $x$  is the sole critical zero of $u$, and that the nodal set $\cZ(u)$ is connected. \smallskip

In a small neighborhood of $x$, the nodal set $\cZ(u)$ consists of $2p$ nodal semi-arcs emanating from $x$, tangentially to $2p$ distinct\footnote{The rays actually form an equiangular system, but this is not needed here.} rays.  Fix an orientated $2$-frame in $T_xM$, with first vector supported by one of the rays, denoted by $\omega_0$, and label the $2p$ rays counter-clockwise, $\omega_0, \omega_1, \ldots, \omega_{2p-1}$. Denote by $(r,\omega)$ the polar coordinates associated with this frame in $T_xM$.\smallskip

Since nodal arcs can only meet at critical zeros, the nodal semi-arc emanating from $x$ tangentially to the ray $\omega_j$ must eventually come back to $x$, arriving tangentially to some ray $\omega_{\tau(j)}$, forming a loop $\gamma_{j,\tau(j)}$ at $x$. Hence, we have a map,
%
\begin{equation}\label{E-llnd-2}
\tau : L_0 \mapsto L_0\,, \text{~where~~} L_0:= \set{0,\ldots,2p-1}\,,
\end{equation}
%
with the following properties,
%
\begin{equation}\label{E-llnd-4}
\left\{
\begin{array}{l}
\tau^2 = \id\,,\\[5pt]
\tau(j) \not = j\,, \quad \forall j \in L_0\,,\\[5pt]
\tau(j) - j \text{~is odd}\,, \quad \forall j \in L_0\,.
\end{array}%
\right.
\end{equation}

The first property is clear. The second and third ones follow from the local structure theorem. \smallskip

The global picture is that the nodal set $\cZ(u)$ is a \emph{$p$-bouquet of loops} $W(L_0)$ at $x$: $p$ simple loops at $x$ which do not intersect away from $x$, and which meet transversally at $x$.\smallskip

Take any loop, $\gamma_{j,\tau(j)}$. Exchanging, $j$ and $\tau(j)$ if necessary, we may always assume that $j < \tau(j)$. Consider the sets
%
\begin{equation}\label{E-llnd-6}
\left\{
\begin{array}{l}
L_{(j)} := \set{j,(j+1),\ldots,(\tau(j)-1),\tau(j)}\,,\\[5pt]
J'_{(j)} := L_{(j)} \sm \set{j,\tau(j)}
\end{array}
\right.
\end{equation}

Since $M$ is topologically a sphere, $M\sm \set{\gamma_{j,\tau(j)}}$ has two components. The local structure of $\cZ(u)$ at $x$ shows that the rays $\omega_k, k \in J'_{(j)}$, point inside one of the components, call it $A_{(j)}$.
The corresponding nodal arcs must stay in $A_{(j)}$, so that  $J'_{(j)}$ is invariant under $\tau$, and corresponds to a bouquet of loops $W(J'_{(j)}) \subset A_{(j)}$. Similarly, the rays $\omega_k, k \in L_0\sm L_{(j)}$, point inside the other component, call it $A'_{(j)}$. The corresponding nodal arcs stay in $A'_{(j)}$, so that $L_0\sm L_{(j)}$ is invariant under $\tau$, and corresponds to a bouquet of loops $W(L_0\sm L_{(j)}) \subset A'_{(j)}$.\medskip

In the polar coordinates $(r,\omega)$, and for $r$ small enough, the nodal semi-arcs  emanating from $x$ are given by equations $r \mapsto \exp_x(r\, \omegat_j(r))$, $0 \le j \le 2p-1$. These arcs determine $2p$ intervals
%
\begin{equation}\label{E-llnd-8}
I_j(r) := \set{ \exp_x(r \,\omega) ~|~  \omega \in \left( \omegat_j(r),\omegat_{j+1}(r) \right) }, \quad 0 \le j \le 2p-1\,,
\end{equation}
%
on the circle $S_x(r) = \set{\exp_x(r\, \omega) ~|~ \omega \in [0,2\pi]}$. The function $u$ does not vanish in these open intervals, and changes sign while crossing an end point $\exp_x(r\, \omegat_j(r))$ along the circle. In particular, each open interval is contained in a unique nodal domain, and two contiguous intervals are contained in different nodal domains. Note that the point $x$ is in the closure of each nodal domain of $u$, so that each nodal domain meets at least one interval.


\begin{lemma}\label{L-llnd-2}
The complement of a $p$-bouquet of loops $W$ in $M$ has $(p+1)$ components.
\end{lemma}%

\pf Looking at the $p$-bouquet as a connected multi-graph with $1$ vertex and $p$ edges, we find that the number of components in the complement of the bouquet $W$ is $\alpha_2(W,\bS^2) = \alpha_1(W) - \alpha_0(W) + c(W) + 1 =p+1$ (with the notation of Section~\ref{S-pargef}).  Note that a $p$-bouquet of loops is not a graph in the sense of Giblin \cite{Gibl2010}. We can modify the bouquet as in Subsection~\ref{SS-parti} to turn it into a graph, so that we can apply Euler's formula for the sphere. \hfill \qed


\subsection{Labeling of nodal domains}\label{SS-llnd-d}

We now define a labeling of nodal domains.

\begin{definition}\label{D-llnd-2}
Label the nodal domains by introducing the map
%
$$
\delta : L_0 \to \set{1,\ldots,p}
$$
%
such that the interval $I_j(r)$ is contained in the nodal domain $\Omega_{\delta(j)}$. More precisely,
\begin{enumerate}
\item let $\Omega_1$ be the nodal domain which contains the interval $I_0(r)$, i.e., $\delta(0) := 1$;
\item let $\Omega_2$ be the nodal domain which contains  the interval $I_1(r)$, i.e., $\delta(1) := 2$;
\item for $j \ge 1$, assuming that $\delta(1), \ldots, \delta(j)$ are defined, let
\begin{enumerate}
\item[$\diamond$] $\delta(j+1) = \delta(\ell)$, if the interval $I_{j+1}(r)$ is contained in the nodal domain $\Omega_{\delta(\ell)}$ for some $\ell \le j$, i.e. in a nodal domain already labeled;
  \item[$\diamond$] let $\delta(j+1) = \delta(j) + 1$, if the interval $I_{j+1}(r)$ is not contained in an already labeled nodal domain.
\end{enumerate}
\end{enumerate}
\end{definition}%

\begin{property}\label{P-llnd-2}
Once we have chosen an orientation in $T_xM$, an initial ray $\omega_0$, and labeled the other rays counter-clockwise, the map $\tau : L_0 \to L_0$ determines the map $\delta : L_0 \to \set{1,\ldots,p}$, and conversely.
\end{property}%


\subsubsection{Working out an example} Choose $p=8$, so that
%
\begin{equation*}
L_0 = \set{0,1,2,3,4,5,6,7,8,9,10,11,12,13,14,15}\,,
\end{equation*}
%
and define the map $\tau$ in matrix form ($\tau$ maps the first line to the second line) by
%
\setcounter{MaxMatrixCols}{20}
\begin{equation}%\label{}
\tau ~=~
\begin{pmatrix}
  0 & 1 & 2 & 3 & 4 & 5 & 6 & 7 & 8 & 9 & {10} & 11 & 12 & 13 & 14 & 15\\
  3 & 2 & 1 & 0 &9  & 8 & 7 & 6 & 5 & 4 & 15 & 12 & 11 & 14 & 13 & 10\\
\end{pmatrix}\,.
\end{equation}

Since $\tau(0) = 3$, $\tau(4) = 9$ and $\tau(10) = 15$, we decompose $L_0$ into three subsets $L_1 = \set{0,1,2,3}$, $L_2 = \set{4,5,6,7,8,9}$ and $L_3 = \set{10, 11, 12, 13, 14, 15}$,
%
\begin{equation}%\label{}
\tau ~=~\left(
\begin{vmatrix}
  0 & 1 & 2 & 3 \\
  3 & 2 & 1 & 0 \\
\end{vmatrix}
\begin{vmatrix}
  4 & 5 & 6 & 7 & 8 & 9 \\
  9 & 8 & 7 & 6 & 5 & 4 \\
\end{vmatrix}
\begin{vmatrix}
  10 & 11 & 12 & 13 & 14 & 15 \\
  15 & 12 & 11 & 14 & 13 & 10 \\
\end{vmatrix}
\right)\,,
\end{equation}
%
It is clear from the definition of $\tau$ that these subsets are $\tau$-invariant. The reason behind this fact is explained in Subsection~\ref{SS-llnd-w}.\smallskip

\noid Since $\tau(0)=3$, the first subset we consider is $L_1 = \set{0,1,2,3}$. Let $\ell_1 = 0$, $k_1 := |L_1|/2 = 2$. The set $L_1$ corresponds to a $k_1$-bouquet of loops, $W(L_1) := \set{\gamma_{0,3},\gamma_{1,2}}$. The set $J'_1 := L_1\sm\set{0,3} = \set{1,2}$ corresponds to a $(k_1-1)$-bouquet of loops, $W(J'_1) := \set{\gamma_{1,2}}$. \smallskip

Since $M$ is a sphere, $M\sm \gamma_{0,3}$ has two components $A_1$ and $A'_1$, and we choose $A_1$ such that $W(J'_1) \subset A_1$ (see Subsection~\ref{SS-llnd-w}). We begin labeling the nodal domains.\smallskip

We have $\delta(0) = 1$, $\delta(1) = 2$. Since $k_1 = 2$, $W(L_1)$ contains two loops, $I_0(r), I_1(r)$ and $I_2(r)$ are contained in $A_1$, and we must have $\delta(2)= \delta(0)=1$. Furthermore, since $A_1$ is simply connected, the set $A_1\sm W(J'_1)$ has two components which are nodal domains of $u$. It also follows that $I_3(r) \subset A'_1$, which implies that $\delta(3) = k_1+1=3$. We have actually determined $\delta$ on $L_1$. \smallskip

\noid Since $\tau(4) = 9$, we turn to the second set,  $L_2 = \set{4,5,6,7,8,9}$, with $\ell_2 = 4$, $k_2 = |L_2|/2 = 3$, and $J'_2 = \set{5,6,7,8}$. The set $M\sm \gamma_{4,9}$ has two components, $A_2$ and $A'_2$, and we choose $A_2$ such that $W(J'_2) \subset A_2$. Notice that the nodal domain $\Omega_3$ (corresponding to $\delta(3)$ is contained in $A'_1 \cap A'_2$. Since $A_2$ is simply connected, the set $A_2 \sm W(J'_2)$ has $k_2=3$ components which are nodal domains of $u$. These nodal domains are different from the previous labeled ones ($\Omega_1, \Omega_2$ and $\Omega_3$). We label them beginning with $k_1+2=4$; furthermore $\delta(\ell_2) = \delta(\tau(\ell_2)-1) = k_1+2 = 4$, i.e. $\delta(4) = \delta(8) = 4$. To continue, we look at $J'_2 = \set{5,6,7,8}$ which is invariant by $\tau$. The action of $\tau$ on this set is similar to the action of $\tau$ on $L_1$. Taking into account the fact that we must have $\delta(5) = k_1+ 3 = 5$, we conclude that $\delta(5) = 5, \delta(6) = 6$ and $\delta(7) = 5$. Notice that $\delta(9) = k_1+1 = 3$. We have labeled $1+k_1+k_2$ nodal domains. \smallskip

\noid Since $\tau(10) = 15$, we turn to the third set,  $L_3 = \set{10,11,12,13,14,15}$, with $\ell_3 = 10$, $k_3 = |L_3|/2 = 3$ and $J'_3 = \set{11,12,13,14}$. Notice that the actions of $\tau$ on $L_2$ and $L_3$ are different. The set $M\sm \gamma_{10,15}$ has two components, $A_3$ and $A'_3$, and we choose $A_3$ such that $W(J'_3) \subset A_3$. Notice that the nodal domain $\Omega_3$ (corresponding to $\delta(3)$ is contained in $A'_1 \cap A'_2 \cap A'_3$. Since $A_3$ is simply connected, the set $A_3 \sm W(J'_3)$ has $k_3=3$ components which are nodal domains of $u$. These nodal domains are different from the previous labeled ones ($\Omega_1, \ldots, \Omega_6$). We label them beginning with $k_2+k_1+2=7$; furthermore $\delta(\ell_3) = \delta(\tau(\ell_3)-1) = k_2+k_1+2 = 7$, i.e. $\delta(10) = \delta(14) = 7$. To continue, we look at the sets $\set{11,12}$ and $\set{13,14}$, and we easily infer that $\delta(11) = 8, \delta(12) = 7$, and $\delta(13)=9$.\smallskip

We have proved, for this explicit example, that the map $\tau$ determines the map $\delta$.\medskip


We can describe the map $\delta$ by the following matrix
%

\begin{equation}%\label{}
\delta ~=~
\begin{pmatrix}
  0 & 1 & 2 & 3 & 4 & 5 & 6 & 7 & 8 & 9 & 10 & 11 & 12 & 13 & 14 & 15 \\
  1 & 2 & 1 & 3 & 4 & 5 & 6 & 5 & 4 & 3 & 7 & 8 & 7 & 9 & 7 & 3  \\
\end{pmatrix}\,,
\end{equation}
%
where $\delta$ sends the entry $j \in L_0$ in the first line, the label of the interval $I_j(r)$, to the entry $\delta(j)$ in the second line, the label of the nodal domain $\Omega_{\delta(j)}$. \smallskip

To determine $\tau(0)$, we look at the occurrences of $1$, the label of the first nodal domain $\Omega_1$: we find $1 = \delta(0) = \delta(2)$, and notice that $\delta(3)=\delta(15)=3$. We have determined that $L_1 = \set{0,1,2,3}$. This also tells us that the nodal domain $\Omega_2$ is bounded by a single loop. It follows that $\tau$ on $L_1$ is given by
%
\begin{equation*}
\tau|_{L_1} ~=~
\begin{pmatrix}
  0 & 1 & 2 & 3 \\
  3 & 2 & 1 & 0 \\
\end{pmatrix}\,.
\end{equation*}

The next nodal domain which appears is $\Omega_4$. To determine $\tau(4)$, we look at the occurrences of $4$, $4 = \delta(4) = \delta(8)$, and we note that $\delta(3) = \delta(9)= 3$. This means that $\tau(4) = 9$, so that $L_2 = \set{4,5,6,7,8,9}$. In $A_2$, the label $4$ occurs twice, the label $5$ occurs twice as well, and the label $6$ once. We conclude that
%
\begin{equation*}
\tau|_{L_2} ~=~
\begin{pmatrix}
  4 & 5 & 6 & 7 & 8 & 9 \\
  9 & 8 & 7 & 6 & 5 & 4 \\
\end{pmatrix}\,.
\end{equation*}

The next nodal domain to appear is $\Omega_7$. To determine $\tau(10)$, we look at the occurrences of $7$: $7 = \delta(10) = \delta(12) = \delta(14)$, while $3 = \delta(9) = \delta(15)$. This determines $L_3$, and actually $\tau|_{L_(3)}$ as well,
%
\begin{equation*}
\tau|_{L_3} ~=~
\begin{pmatrix}
  10 & 11 & 12 & 13 & 14 & 15 \\
  15 & 12 & 11 & 14 & 13 & 10 \\
\end{pmatrix}\,.
\end{equation*}

We have recovered the map $\tau$ from the map $\delta$ in the example at hand.


\subsubsection{Proof of Property~\ref{P-llnd-2}}~

\noib We first show that the map $\tau$ determines the map $\delta$. \smallskip

\noid We first consider the loop $\gamma_{0,\tau(0)}$. Note that $0 < \tau(0) \le 2p-3$. Define
%
\begin{equation}\label{E-llnd-9}
\left\{
\begin{array}{l}
\ell_1 = 0\,,\\[5pt]
L_1 = \set{0,\ldots,\tau(0)} \text{~and~} L_1' = L_1\sm \set{0,\tau(0)}\,,\\[5pt]
k_1 = \frac{1}{2}|L_1| = \frac{1}{2}\left(\tau(0) +1 \right)\,.
\end{array}%
\right.
\end{equation}

The loop $\gamma_{0,\tau(0)}$ divides $M$ into two components, both homeomorphic to a disk (recall that $M$ is topologically a sphere). The rays with index in $J'_1$ point inside one of the components, call this first component $A_1$; the rays with index in $L_0\sm L_1$ point inside the second component $A'_1$. In case $J'_1=\emptyset$, there is no ray pointing into the component $A_1$, and $A_1$ is actually a nodal domain bounded by the loop $\gamma_{0,1}$. Because the loops at $x$ do not intersect away from $x$, the loops with semi-tangents in $J'_1$ must stay inside $A_1$; the loops with semi-tangents in $L_0\sm L_1$ must stay inside $A'_1$.\smallskip

In particular, this implies that the map $\tau$ leaves the sets $L_1$, $J'_1$ and $L_0\sm L_1$ globally invariant. To the set $L_1$ we can therefore associate a $k_1$-bouquet $W(L_1)$ of loops at $x$, whose complement in $M$ has $(k_1+1)$ components. Similarly, to the set $J'_1$ we can associate a $(k_1-1)$-bouquet $W(J'_1)$ which is contained in $A_1$, and whose complement in $A_1$ has $k_1$ components which are actually nodal domains of $u$. The remaining nodal domains are all contained in $A'_1$. The $k_1$ nodal domains contained in $A_1$ are labeled from $1$ to $k_1$; the labeling of the nodal domains contained in $A'_1$ begins with $k_1+1$. The intervals $I_j(r), 0 \le j \le \tau(0)-1$ are contained in $A_1$; the intervals $I_j(r), \tau(0) \le j$ are contained in $A'_1$. From these facts, we infer that
%
\begin{equation}\label{E-llnd-10}
\left\{
\begin{array}{l}
\delta(0) = \delta(\tau(0)-1) = 1\,,\\[5pt]
\delta(1) = 2\,,\\[5pt]
\delta(\tau(0)) = \delta(2p-1) = k_1+1\,,\\[5pt]
\delta(\tau(0)+1) = k_1+2\,.
\end{array}%
\right.
\end{equation}
\medskip

\noid Since the labeling of the nodal domains contained in $A_1$ and the labeling of the nodal domains contained in $A'_1$ are essentially independent matters, we can pursue labeling the nodal domains in $A_1$ as if we only worked with the bouquet $W(L_0\sm L_1)$. \smallskip

The domain $A_1$ is simply connected, with boundary $\gamma_{0,\tau(0)}$. Any loop $\gamma$ in the bouquet $W(J'_1)$ is contained in $A_1$ and hence separates $A_1$ into two components, a component $A_{\gamma}$ which only touches $\gamma_{0,\tau(0)}$ at the point $x$, and a component $A'_{\gamma}$ one of whose boundaries is $\gamma_{0,\tau(0)}$, and hence meets the nodal set $\Omega_1$. \smallskip

We now look at the bouquets of nodal loops contained in $A_1$. If $\tau(1) = 2$, the corresponding loop $\gamma_{1,2}$ bounds the nodal domain $\Omega_2$, we have $\delta(3) = 3$, and we are left with analyzing the smaller bouquet associated with the set $\set{3,\ldots,(\tau(0)-1)}$. If $\tau(1) > 2$, we apply the above reasoning with the bouquet associated with the set $ \set{1,\ldots,\tau(1)}$.\medskip


\noid We now define
%
\begin{equation}\label{E-llnd-12}
\left\{
\begin{array}{l}
\ell_2 = \tau(\ell_1)+1 = \tau(0) + 1\,,\\[5pt]
L_2 = \set{\ell_2,\ldots,\tau(\ell_2)} \text{~and~} L_2' = L_2\sm \set{\ell_2,\tau(\ell_2)}\,,\\[5pt]
k_2 = \frac{1}{2}|L_2| = \frac{1}{2}\left(\tau(\ell_2) - \ell_2+1 \right)\,.
\end{array}%
\right.
\end{equation}

We reason as above. The loop $\gamma_{\ell_2,\tau(\ell_2)}$ divides $M$ into two components. One of them $A_2$ contains the $(k_1-1)$-bouquet of loops $W(J'_2)$; the other one, denoted by $A'_2$, contains the nodal domains $\Omega_{k_0+1}$. The complement of $W(J'_2)$ in $A_2$ has $k_2$ components which are nodal domains, which are labeled $\Omega_{k_0+2}, \ldots \Omega_{k_0+k_1+1}$. In particular, we have $\delta(k_0+1) = k_0+2$ and $\delta(k_0+k_1+1) = k_0+1$.\smallskip

We then continue as previously by analyzing smaller and smaller bouquets of loops inside $A_2$.\medskip

After finitely many such steps, we exhaust the set $L_0$. This proves that the map $\tau$ determines the map $\delta$.\medskip

\noib We now show that the map $\delta$ determines the map $\tau$.\smallskip

It is not clear which maps $\delta$ correspond to maps $\tau$ (with the above properties). Assume that we are given some $\delta$ associated with some $p$-bouquet of loops as above. The $p$-bouquet actually defines some map $\tau$ which induces $\delta$. In order to recover $\tau$ from $\delta$, the idea is to recover the sets $L_i$ introduced above, and to reason by induction on the length of the size of the bouquets.\smallskip

Assume that we are given a map $\delta : L_0  := \set{0,\ldots,2p-1} \to \set{1,\ldots,p}$, induced by some map $\tau : L_0 \to L_0$ such that $\tau^2 = \id$ and, for all $j \in L_0$, $\tau(j) \not = j$, and $\tau(j)-j$ is odd. Does $\delta$ determine $\tau$?\smallskip

\noid We first look at the nodal domain $\Omega_1$, with label $1$, and at the set $\delta^{-1}(1)$. If $\delta^{-1}(1) = \set{0}$, then we must have $\tau(0)=1$, and the loop $\gamma_{0,1}$ bounds $\Omega_1$. If $|\delta^{-1}(1)| > 1$, we look at $m_1 := \max \delta^{-1}(1)$. Then, the only possibility is that $\tau(0) = m_1+1$. Because of our framework, see Subsection~\ref{SS-hmn-25}, the subsets $L_1 := \set{0,\ldots, (m_1+1)}$, $J'_1:=\set{1,\ldots,m_1}$, and $L_0\sm L_1$ must be $\tau$ invariant. Defining $A_1$ and $A'_1$ as above, we see that there are exactly $k_1$ nodal domains inside $A_1$, where $2k_1 = |L_1|$, and we have $\delta(m_1+1) = k_1+1$. We can then use induction.

\vspace{1cm}

\newpage
%% Bibliography using BibTeX
%%\bibliographystyle{plain}
%\bibliographystyle{abstract} %% Citation Keys printed
%% \nocite{AAA} %% Fait appara\^{\i}tre la r\'{e}f. non cit\'{e}e AAA
%\nocite{*} %% Fait appara\^{\i}tre toutes les r\'{e}f\'{e}rences non cit\'{e}es
%\bibliography{BH-homina-MR-220102} % name of .BIB file
%%% possibly \bibliography{file1,file2}
%\end{document}

% Calling bbl files (under TeX names)
%\begin{thebibliography}{HHOHOO99}
%% BH-HoMiNa-bib-220102-short.tex (only cited items)
%% BH-HoMiNa-bib-220102-long.tex (all items)

%\input{BH-HoMiNa-bib-220102-long.tex}
\input{BH-HoMiNa-bib-220102-short.tex}

%\end{thebibliography}

\end{document} %% Hidden text after this tag

%% Extras on berard-helffer-HoMiNa
%\newpage
%\input{}

\newpage
%% Comments and questions on the current version
\input{BH-HoMiNa-extra-00-comments-questions.tex}

\newpage
%% Attempts at understanding the end of the proof in [HoMN1999]
\input{BH-HoMiNa-extra-10-old-notes-comments-HoMiNa-Sect3.tex}

\newpage
%% An euler-type formula for the M\"{o}bius strip and applications
\input{BH-HoMiNa-extra-20-euler-moebius.tex}
\newpage
\input{BH-HoMiNa-extra-21-euler-moebius-cylinder-1.tex}
\newpage
\input{BH-HoMiNa-extra-22-euler-moebius-cylinder-2.tex}
\newpage
\input{BH-HoMiNa-extra-23-euler-moebius-projective-plane.tex}
\newpage
\input{BH-HoMiNa-extra-24-multiplicity-bounds-dirichlet-eigenval-moebius.tex}

\newpage
%% Prescribing singular points, a more genral statement with a proof
%% \`{a} la Besson
\input{BH-HoMiNa-extra-30-prescribing-singular-pts-ala-besson.tex}

\newpage
%% Tentative proof using the double surface
\input{BH-HoMiNa-extra-40-multiplicity-with-symmetry.tex}

\newpage
%% Topology questions to Fran\c{c}ois Laudenbach
\input{BH-HoMiNa-extra-90-questions-to-FLaudenbach.tex}

%\newpage
%\input{}

\newpage
%% Last \newpage to insert address
\end{document} %% Hidden text between bibliography and this tag


\begin{thebibliography}{HHOHOO99}

\bibitem[Ales1994]{Ales1994}
Giovanni Alessandrini.
\newblock Nodal lines of eigenfunctions of the fixed membrane problem in
  general convex domains.
\newblock {\em Comment. Math. Helv.}, 69(1):142--154, 1994.

\bibitem[ArMa2007]{ArMa2007}
Wolfgang Arendt and Rafe Mazzeo.
\newblock Spectral properties of the dirichlet-to-neumann operator on lipschitz
  domains.
\newblock {\em Ulmer Seminare}, 12:28--38, 2007.

\bibitem[ArMa2012]{ArMa2012}
Wolfgang Arendt and Rafe Mazzeo.
\newblock Friedlander's eigenvalue inequalities and the
  {D}irichlet-to-{N}eumann semigroup.
\newblock {\em Commun. Pure Appl. Anal.}, 11(6):2201--2212, 2012.

\bibitem[Aron1957]{Aron1957}
Nachman Aronszajn.
\newblock A unique continuation theorem for solutions of elliptic partial
  differential equations or inequalities of second order.
\newblock {\em J. Math. Pures Appl. (9)}, 36:235--249, 1957.

\bibitem[BeGM1971]{BeGM1971}
Marcel Berger, Paul Gauduchon, and Edmond Mazet.
\newblock {\em Le spectre d'une vari\'{e}t\'{e} riemannienne}.
\newblock Lecture Notes in Mathematics, Vol. 194. Springer-Verlag, Berlin-New
  York, 1971.

\bibitem[BeGi2016]{BeGi2016}
Michiel van~den Berg and Katie Gittins.
\newblock On the number of {C}ourant-sharp {D}irichlet eigenvalues.
\newblock {\em J. Spectr. Theory}, 6(4):735--745, 2016.

\bibitem[BeGo1988]{BeGo1988}
Marcel Berger and Bernard Gostiaux.
\newblock {\em G\'{e}om\'{e}trie diff\'{e}rentielle: vari\'{e}t\'{e}s, courbes
  et surfaces}.
\newblock Math\'{e}matiques. [Mathematics]. Presses Universitaires de France,
  Paris, 1987.

\bibitem[BeHK2021k]{BeHK2021k}
Pierre B\'{e}rard, Bernard Helffer, and Rola Kiwan.
\newblock Courant-sharp eigenvalues of compact flat surfaces: {K}lein bottles
  and cylinders.
\newblock {\em Proc. Amer. Math. Soc.}, 150(1):439--453, 2021.

\bibitem[BeHK2021m]{BeHK2021m}
Pierre B\'{e}rard, Bernard Helffer, and Rola Kiwan.
\newblock Courant-sharp property for {D}irichlet eigenfunctions on the
  {M}\"{o}bius strip.
\newblock {\em Port. Math.}, 78(1):1--41, 2021.

\bibitem[BeHe2014]{BeHe2014}
Pierre B\'{e}rard and Bernard Helffer.
\newblock Remarks on the boundary set of spectral equipartitions.
\newblock {\em Philos. Trans. R. Soc. Lond. Ser. A Math. Phys. Eng. Sci.},
  372(2007):20120492, 15, 2014.

\bibitem[BeHe2015r]{BeHe2015r}
Pierre {B\'erard} and Bernard {Helffer}.
\newblock {Nodal sets of eigenfunctions, Antonie Stern's results revisited}.
\newblock In {\em {Actes de S\'eminaire de Th\'eorie Spectrale et
  G\'eom\'etrie. Ann\'ee 2014--2015}}, pages 1--37. St. Martin d'H\`eres:
  Universit\'e de Grenoble I, Institut Fourier, 2015.
\newblock %zbMATH06787674 ; zbl 1379.35203.

\bibitem[BeHe2015s]{BeHe2015s}
Pierre {B\'erard} and Bernard {Helffer}.
\newblock {Edited extracts from Antonie Stern's thesis}.
\newblock In {\em {Actes de S\'eminaire de Th\'eorie Spectrale et
  G\'eom\'etrie. Ann\'ee 2014--2015}}, pages 39--72. St. Martin d'H\`eres:
  Universit\'e de Grenoble I, Institut Fourier, 2015.
\newblock %zbMATH06787675 ; zbl 1379.35204.

\bibitem[BeHe2016]{BeHe2016}
Pierre B\'{e}rard and Bernard Helffer.
\newblock The weak {P}leijel theorem with geometric control.
\newblock {\em J. Spectr. Theory}, 6(4):717--733, 2016.

\bibitem[BeHe2020m]{BeHe2020m}
Pierre Bérard and Bernard Helffer.
\newblock An {E}uler-type formula for partitions on the {M\"{o}bius} strip.
\newblock arXiv:2005.12571, 2020.

\bibitem[BeHe2021t]{BeHe2021t}
Pierre B\'{e}rard and Bernard Helffer.
\newblock Level sets of certain {N}eumann eigenfunctions under deformation of
  {L}ipschitz domains application to the extended {C}ourant property.
\newblock {\em Ann. Fac. Sci. Toulouse Math. (6)}, 30(3):429--462, 2021.

\bibitem[BeLi2001]{BeLi2001}
Michiel van~den Berg and Maria Lianantonakis.
\newblock Asymptotics for the spectrum of the {D}irichlet {L}aplacian on
  horn-shaped regions.
\newblock {\em Indiana Univ. Math. J.}, 50(1):299--333, 2001.

\bibitem[BeNP2016]{BeNP2016}
Aleksandr Berdnikov, Nikolai Nadirashvili, and Alexei Penskoi.
\newblock Bounds on multiplicities of {Laplace-Beltrami} operator eigenvalues
  on the real projective plane.
\newblock arXiv:1612.04805, 2016.

\bibitem[Berd2018]{Berd2018}
Aleksandr Berdnikov.
\newblock Bounds on multiplicities of {L}aplace operator eigenvalues on
  surfaces.
\newblock {\em J. Spectr. Theory}, 8(2):541--554, 2018.

\bibitem[Bess1980]{Bess1980}
G\'{e}rard Besson.
\newblock Sur la multiplicit\'{e} de la premi\`ere valeur propre des surfaces
  riemanniennes.
\newblock {\em Ann. Inst. Fourier (Grenoble)}, 30(1):x, 109--128, 1980.

\bibitem[BoHe2017]{BoHe2017}
Virginie Bonnaillie-No\"{e}l and Bernard Helffer.
\newblock Nodal and spectral minimal partitions---the state of the art in 2016.
\newblock In {\em Shape optimization and spectral theory}, pages 353--397. De
  Gruyter Open, Warsaw, 2017.

\bibitem[Bona2009]{Bona2009}
Francis Bonahon.
\newblock {\em Low-dimensional geometry}, volume~49 of {\em Student
  Mathematical Library}.
\newblock American Mathematical Society, Providence, RI; Institute for Advanced
  Study (IAS), Princeton, NJ, 2009.
\newblock From Euclidean surfaces to hyperbolic knots, IAS/Park City
  Mathematical Subseries.

\bibitem[BuCo1985]{BuCo1985}
Marc Burger and Bruno Colbois.
\newblock {\`A} propos de la multiplicit\'{e} de la premi\`ere valeur propre du
  {Laplacien} d'une surface de {R}iemann.
\newblock {\em C. R. Acad. Sci. Paris S\'{e}r. I Math.}, 300(8):247--249, 1985.

\bibitem[Chen1976]{Chen1976}
Shiu~Yuen Cheng.
\newblock Eigenfunctions and nodal sets.
\newblock {\em Comment. Math. Helv.}, 51(1):43--55, 1976.

\bibitem[CoCo1988]{CoCo1988}
Bruno Colbois and Yves Colin~de Verdi\`ere.
\newblock Sur la multiplicit\'{e} de la premi\`ere valeur propre d'une surface
  de {R}iemann \`a courbure constante.
\newblock {\em Comment. Math. Helv.}, 63(2):194--208, 1988.

\bibitem[ColV1986]{ColV1986}
Yves Colin~de Verdi\`ere.
\newblock Sur la multiplicit\'{e} de la premi\`ere valeur propre non nulle du
  {Laplacian}.
\newblock {\em Comment. Math. Helv.}, 61(2):254--270, 1986.

\bibitem[ColV1987]{ColV1987}
Yves Colin~de Verdi\`ere.
\newblock Construction de laplaciens dont une partie finie du spectre est
  donn\'{e}e.
\newblock {\em Ann. Sci. \'{E}cole Norm. Sup. (4)}, 20(4):599--615, 1987.

\bibitem[Colb1985]{Colb1985}
Bruno Colbois.
\newblock Petites valeurs propres du {L}aplacien sur une surface de {R}iemann
  compacte et graphes.
\newblock {\em C. R. Acad. Sci. Paris S\'{e}r. I Math.}, 301(20):927--930,
  1985.

\bibitem[DaGH2021]{DaGH2021}
Joel Dahne, Javier G\'{o}mez-Serrano, and Kimberly Hou.
\newblock A counterexample to {P}ayne's nodal line conjecture with few holes.
\newblock {\em Commun. Nonlinear Sci. Numer. Simul.}, 103:Paper No. 105957, 13,
  2021.

\bibitem[Dies2017]{Dies2017}
Reinhard Diestel.
\newblock {\em Graph theory}, volume 173 of {\em Graduate Texts in
  Mathematics}.
\newblock Springer, Berlin, fifth edition, 2017.

\bibitem[DoFe1990a]{DoFe1990a}
Harold Donnelly and Charles Fefferman.
\newblock Nodal sets of eigenfunctions: {R}iemannian manifolds with boundary.
\newblock In {\em Analysis, et cetera}, pages 251--262. Academic Press, Boston,
  MA, 1990.

\bibitem[FoBP2021]{FoBP2021}
Maxime {Fortier Bourque} and Bram Petri.
\newblock The {Klein} quartic maximizes the multiplicity of the first positive
  eigenvalue of the {Laplacian}.
\newblock arXiv: 2111.14699.

\bibitem[Four2001]{Four2001}
S{\o}ren Fournais.
\newblock The nodal surface of the second eigenfunction of the {L}aplacian in
  {${\bf R}^D$} can be closed.
\newblock {\em J. Differential Equations}, 173(1):145--159, 2001.

\bibitem[FrSc2016]{FrSc2016}
Ailana Fraser and Richard Schoen.
\newblock Sharp eigenvalue bounds and minimal surfaces in the ball.
\newblock {\em Invent. Math.}, 203(3):823--890, 2016.

\bibitem[Frie1991]{Frie1991}
Leonid Friedlander.
\newblock Some inequalities between {D}irichlet and {N}eumann eigenvalues.
\newblock {\em Arch. Rational Mech. Anal.}, 116(2):153--160, 1991.

\bibitem[GaHuLa2004]{GaHuLa2004}
Sylvestre Gallot, Dominique Hulin, and Jacques Lafontaine.
\newblock {\em Riemannian geometry}.
\newblock Universitext. Springer-Verlag, Berlin, third edition, 2004.

\bibitem[GaXu2013]{GaXu2013}
Jean Gallier and Dianna Xu.
\newblock {\em A guide to the classification theorem for compact surfaces},
  volume~9 of {\em Geometry and Computing}.
\newblock Springer, Heidelberg, 2013.

\bibitem[GiHe2019]{GiHe2019}
Katie Gittins and Bernard Helffer.
\newblock Courant-sharp {R}obin eigenvalues for the square and other planar
  domains.
\newblock {\em Port. Math.}, 76(1):57--100, 2019.

\bibitem[GiLe2020]{GiLe2020}
Katie Gittins and Corentin Léna.
\newblock Upper bounds for {C}ourant-sharp {N}eumann and {R}obin eigenvalues.
\newblock {\em Bull. Soc. Math. {F}rance}, 148(1):99--132, 2020.

\bibitem[Gibl2010]{Gibl2010}
Peter Giblin.
\newblock {\em Graphs, surfaces and homology}.
\newblock Cambridge University Press, Cambridge, third edition, 2010.

\bibitem[HeHJ2020]{HeHJ2020}
Bernard Helffer, Thomas Hoffmann-Ostenhof, François Jauberteau, and Corentin
  Léna.
\newblock On the multiplicity of the second eigenvalue of the {Laplacian} in
  non simply connected domains -- with some numerics --.
\newblock {\em Asymptotic Analysis}, 121(1):35--57, 2021.

\bibitem[HeHN2002]{HeHN2002}
Bernard Helffer, Maria Hoffmann-Ostenhof, Thomas Hoffmann-Ostenhof, and Nikolai
  Nadirashvili.
\newblock Spectral theory for the dihedral group.
\newblock {\em Geom. Funct. Anal.}, 12(5):989--1017, 2002.

\bibitem[HeHO1999]{HeHO1999}
Bernard Helffer, Maria Hoffmann-Ostenhof, Thomas Hoffmann-Ostenhof, and Mark~P.
  Owen.
\newblock Nodal sets for groundstates of {S}chr\"{o}dinger operators with zero
  magnetic field in non-simply connected domains.
\newblock {\em Comm. Math. Phys.}, 202(3):629--649, 1999.

\bibitem[HeHT2009]{HeHT2009}
Bernard Helffer, Thomas Hoffmann-Ostenhof, and Susanna Terracini.
\newblock Nodal domains and spectral minimal partitions.
\newblock {\em Ann. Inst. H. Poincar\'{e} Anal. Non Lin\'{e}aire},
  26(1):101--138, 2009.

\bibitem[HoHN1997]{HoHN1997}
Maria Hoffmann-Ostenhof, Thomas Hoffmann-Ostenhof, and Nikolai Nadirashvili.
\newblock The nodal line of the second eigenfunction of the {L}aplacian in
  {${\bf R}^2$} can be closed.
\newblock {\em Duke Math. J.}, 90(3):631--640, 1997.

\bibitem[HoHN1998]{HoHN1998}
Maria Hoffmann-Ostenhof, Thomas Hoffmann-Ostenhof, and Nikolai Nadirashvili.
\newblock On the nodal line conjecture.
\newblock In {\em Advances in differential equations and mathematical physics
  ({A}tlanta, {GA}, 1997)}, volume 217 of {\em Contemp. Math.}, pages 33--48.
  Amer. Math. Soc., Providence, RI, 1998.

\bibitem[HoHN1999]{HoHN1999}
Maria Hoffmann-Ostenhof, Thomas Hoffmann-Ostenhof, and Nikolai Nadirashvili.
\newblock On the multiplicity of eigenvalues of the {L}aplacian on surfaces.
\newblock {\em Ann. Global Anal. Geom.}, 17(1):43--48, 1999.

\bibitem[HoMN1999]{HoMN1999}
Thomas Hoffmann-Ostenhof, Peter~W. Michor, and Nikolai Nadirashvili.
\newblock Bounds on the multiplicity of eigenvalues for fixed membranes.
\newblock {\em Geom. Funct. Anal.}, 9(6):1169--1188, 1999.

\bibitem[Horm2009d]{Horm2009d}
Lars H\"{o}rmander.
\newblock {\em The analysis of linear partial differential operators. {IV}}.
\newblock Classics in Mathematics. Springer-Verlag, Berlin, 2009.
\newblock Fourier integral operators, Reprint of the 1994 edition.

\bibitem[Ivri1980]{Ivri1980}
Victor~Ja. Ivri\u{\i}.
\newblock Second term of the spectral asymptotic expansion of the
  {L}aplace-{B}eltrami operator on manifolds with boundary.
\newblock {\em Funct Anal Its Appl.}, 14(2):98--106, 1980.

\bibitem[Ivri1980r]{Ivri1980r}
Victor~Ja. Ivri\u{\i}.
\newblock The second term of the spectral asymptotics for a
  {L}aplace-{B}eltrami operator on manifolds with boundary.
\newblock {\em Funktsional. Anal. i Prilozhen.}, 14(2):25--34, 1980.

\bibitem[Jam2016]{Jam2016}
Pierre Jammes.
\newblock Multiplicit\'{e} du spectre de {S}teklov sur les surfaces et nombre
  chromatique.
\newblock {\em Pacific J. Math.}, 282(1):145--171, 2016.

\bibitem[KaKP2014]{KaKP2014}
Mikhail Karpukhin, Gerasim Kokarev, and Iosif Polterovich.
\newblock Multiplicity bounds for {S}teklov eigenvalues on {R}iemannian
  surfaces.
\newblock {\em Ann. Inst. Fourier (Grenoble)}, 64(6):2481--2502, 2014.

\bibitem[Ken2011]{Ken2011}
James~B. Kennedy.
\newblock The nodal line of the second eigenfunction of the {R}obin {L}aplacian
  in {$\Bbb R^2$} can be closed.
\newblock {\em J. Differential Equations}, 251(12):3606--3624, 2011.

\bibitem[Ken2013]{Ken2013}
James~B. Kennedy.
\newblock Closed nodal surfaces for simply connected domains in higher
  dimensions.
\newblock {\em Indiana Univ. Math. J.}, 62(3):785--798, 2013.

\bibitem[Kiwa2018]{Kiwa2018}
Rola Kiwan.
\newblock On the nodal set of a second {D}irichlet eigenfunction in a doubly
  connected domain.
\newblock {\em Ann. Fac. Sci. Toulouse Math. (6)}, 27(4):863--873, 2018.

\bibitem[Lena2019]{Lena2019}
Corentin L\'{e}na.
\newblock Pleijel's nodal domain theorem for {N}eumann and {R}obin
  eigenfunctions.
\newblock {\em Ann. Inst. Fourier (Grenoble)}, 69(1):283--301, 2019.

\bibitem[Lewy1977]{Lewy1977}
Hans Lewy.
\newblock On the minimum number of domains in which the nodal lines of
  spherical harmonics divide the sphere.
\newblock {\em Comm. Partial Differential Equations}, 2(12):1233--1244, 1977.

\bibitem[LiNi1988]{LiNi1988}
Chang~Shou Lin and Wei-Ming Ni.
\newblock A counterexample to the nodal domain conjecture and a related
  semilinear equation.
\newblock {\em Proc. Amer. Math. Soc.}, 102(2):271--277, 1988.

\bibitem[Lin1987]{Lin1987}
Chang~Shou Lin.
\newblock On the second eigenfunctions of the {L}aplacian in {${\bf R}^2$}.
\newblock {\em Comm. Math. Phys.}, 111(2):161--166, 1987.

\bibitem[Liq1995]{Liq1995}
Zhang Liqun.
\newblock On the multiplicity of the second eigenvalue of {L}aplacian in {${\bf
  R}^2$}.
\newblock {\em Comm. Anal. Geom.}, 3(1-2):273--296, 1995.

\bibitem[Maz1991]{Maz1991}
Rafe Mazzeo.
\newblock Remarks on a paper of {L}. {F}riedlander concerning inequalities
  between {N}eumann and {D}irichlet eigenvalues.
\newblock {\em Internat. Math. Res. Notices}, 1991(4):41--48, 1991.
\newblock %``{S}ome inequalities between {D}irichlet and {N}eumann
  eigenvalues'' [{A}rch. {R}ational {M}ech. {A}nal. {\bf 116} (1991), no. 2,
  153--160; {MR}1143438 (93h:35146)].

\bibitem[Mela1992]{Mela1992}
Antonios~D. Melas.
\newblock On the nodal line of the second eigenfunction of the {L}aplacian in
  {${\bf R}^2$}.
\newblock {\em J. Differential Geom.}, 35(1):255--263, 1992.

\bibitem[Miln1997]{Miln1997}
John~W. Milnor.
\newblock {\em Topology from the differentiable viewpoint}.
\newblock Princeton University Press, Princeton, NJ, 1997.
\newblock Revised reprint of the 1965 original.

\bibitem[Nadi1987]{Nadi1987}
Nikolai~S. Nadirashvili.
\newblock Multiple eigenvalues of the {L}aplace operator ({R}ussian).
\newblock {\em Mat. Sb. (N.S.)}, 133(175)(2):223--237, 272, 1987.
\newblock Translation, Math USSR-Sb. 61 225-238 1988.

\bibitem[Payn1967]{Payn1967}
Lawrence~E. Payne.
\newblock Isoperimetric inequalities and their applications.
\newblock {\em SIAM Rev.}, 9:453--488, 1967.

\bibitem[Payn1973]{Payn1973}
Lawrence~E. Payne.
\newblock On two conjectures in the fixed membrane eigenvalue problem.
\newblock {\em Z. Angew. Math. Phys.}, 24:721--729, 1973.

\bibitem[Plej1956]{Plej1956}
{\AA}ke Pleijel.
\newblock Remarks on {C}ourant's nodal line theorem.
\newblock {\em Comm. Pure Appl. Math.}, 9:543--550, 1956.

\bibitem[Polt2009]{Polt2009}
Iosif Polterovich.
\newblock Pleijel's nodal domain theorem for free membranes.
\newblock {\em Proc. Amer. Math. Soc.}, 137(3):1021--1024, 2009.

\bibitem[Putt1990]{Putt1990}
Rolf P\"{u}tter.
\newblock On the nodal lines of second eigenfunctions of the fixed membrane
  problem.
\newblock {\em Comment. Math. Helv.}, 65(1):96--103, 1990.

\bibitem[Putt1991]{Putt1991}
Rolf P\"{u}tter.
\newblock On the nodal lines of second eigenfunctions of the free membrane
  problem.
\newblock {\em Appl. Anal.}, 42(3-4):199--207, 1991.

\bibitem[Safa2001]{Safa2001}
Yuri Safarov.
\newblock Fourier {T}auberian theorems and applications.
\newblock {\em J. Funct. Anal.}, 185(1):111--128, 2001.

\bibitem[Seve2002]{Seve2002}
Bruno S\'{e}vennec.
\newblock Multiplicity of the second {S}chr\"{o}dinger eigenvalue on closed
  surfaces.
\newblock {\em Math. Ann.}, 324(1):195--211, 2002.

\bibitem[Shen1988]{Shen1988}
Chao~Liang Shen.
\newblock On the nodal sets of the eigenfunctions of the string equation.
\newblock {\em SIAM J. Math. Anal.}, 19(6):1419--1424, 1988.

\bibitem[Ster1925]{Ster1925}
Antonie Stern.
\newblock {\em {Bemerkungen \"uber asymptotisches Verhalten von Eigenwerten und
  Eigenfunktionen}}.
\newblock {Math-Naturwiss. Diss.}, Universität Göttingen, 1925.

\bibitem[YaZh2021]{YaZh2021}
Yunyan Yang and Jie Zhou.
\newblock Blow-up analysis involving isothermal coordinates on the boundary of
  compact {R}iemann surface.
\newblock {\em J. Math. Anal. Appl.}, 504(2):Paper No. 125440, 39, 2021.

\end{thebibliography}
